% !TeX program = xelatex
% 数学 LaTeX 排版 Skill v4.0.0；版式保持 v3 金标准. XeLaTeX 编译至少两次. 不分发字体.
% WZ-LOCKED-CLASS-BEGIN
\begin{filecontents*}[overwrite]{elegantbook-original-adapter.cls}
\NeedsTeXFormat{LaTeX2e}
\ProvidesClass{elegantbook-original-adapter}[2026/09/23 v3.0 fixed mathematical book environments]
\LoadClass[10pt,letterpaper,oneside]{book}
\RequirePackage[UTF8,scheme=plain,fontset=fandol]{ctex}
\RequirePackage{fontspec}
\setmainfont{cmunrm.otf}[BoldFont=cmunbx.otf,ItalicFont=cmunti.otf,BoldItalicFont=cmunbi.otf]
\setCJKfamilyfont{sourceheader}[ItalicFont=FandolKai-Regular.otf]{FandolKai-Regular.otf}
\RequirePackage{amsmath,amssymb,mathrsfs}
\RequirePackage{amsthm,etoolbox}
\RequirePackage{xcolor,geometry,setspace,fancyhdr,enumitem,needspace,xurl}
\definecolor{structurecolor}{HTML}{880E4F}
\definecolor{main}{HTML}{9C27B0}
\geometry{paperwidth=612bp,paperheight=792bp,left=20mm,right=20mm,top=95.04bp,bottom=20mm,headheight=12pt,headsep=25pt,footskip=30pt}
\setstretch{1.6}
\setlength{\parindent}{2em}
\setlength{\parskip}{0pt}
\setlist{nosep,leftmargin=2em}
\AtBeginDocument{\setlength{\abovedisplayskip}{4pt plus 1pt minus 1pt}\setlength{\belowdisplayskip}{4pt plus 1pt minus 1pt}\setlength{\abovedisplayshortskip}{2pt}\setlength{\belowdisplayshortskip}{3pt}}
\fancyhf{}
\fancyhead[L]{{\itshape\CJKfamily{sourceheader}\rightmark}}
\fancyhead[R]{{\itshape\CJKfamily{sourceheader}\leftmark}}
\fancyfoot[C]{\thepage}
\renewcommand{\headrulewidth}{0.4pt}
\renewcommand{\headrule}{{\color{structurecolor}\hrule height\headrulewidth width\headwidth\vskip-\headrulewidth}}
\pagestyle{fancy}
\fancypagestyle{cover}{\fancyhf{}\fancyfoot[C]{\thepage}\renewcommand{\headrulewidth}{0pt}\renewcommand{\headrule}{}}
\RequirePackage{hyperref}
\hypersetup{unicode=true,colorlinks=true,linkcolor=structurecolor,urlcolor=structurecolor,citecolor=structurecolor,linktoc=section,bookmarksnumbered=true,bookmarksopen=true,pdfborder={0 0 0},pdfstartview=FitH}
\setcounter{tocdepth}{1}
\renewcommand*\l@section{\@dottedtocline{1}{1.5em}{2.4em}}
\renewcommand{\thesection}{\thechapter.\arabic{section}}
\newcommand{\originalchapter}[1]{\par\addvspace{14pt}\Needspace{5\baselineskip}\refstepcounter{chapter}\setcounter{section}{0}\addcontentsline{toc}{chapter}{\protect\numberline{\thechapter}#1}\markboth{\thechapter\quad #1}{}\begingroup\centering\color{structurecolor}\bfseries\fontsize{14.4}{20}\selectfont\thechapter\quad #1\par\endgroup\nobreak\vspace{8pt}}
\newcommand{\originalsection}[1]{\par\addvspace{9pt}\Needspace{4\baselineskip}\refstepcounter{section}\addcontentsline{toc}{section}{\protect\hspace*{1.5em}\protect\numberline{\protect\textcolor{structurecolor}{\thesection}}#1}\markright{\thesection\quad #1}\noindent{\fontsize{12}{17}\selectfont\color{structurecolor}\thesection\quad}{\bfseries\fontsize{12}{17}\selectfont #1}\par\nobreak\vspace{4pt}}

% Semantic environments are registered once here, not re-designed in manuscripts.
\newtheoremstyle{wzstatement}{6pt}{5pt}
  {\normalfont\kaishu}{0pt}{\normalfont\bfseries\color{structurecolor}}
  {}{.7em}{\thmname{#1}\thmnumber{ #2}\thmnote{\normalfont\ (#3)}}
\newtheoremstyle{wzdefinition}{6pt}{5pt}
  {\normalfont\songti}{0pt}{\normalfont\bfseries\color{structurecolor}}
  {}{.7em}{\thmname{#1}\thmnumber{ #2}\thmnote{\normalfont\ (#3)}}
\newtheoremstyle{wzproblem}{7pt}{5pt}
  {\normalfont\songti}{2em}{\normalfont\bfseries\color{main}}
  {}{.7em}{\thmname{#1}\thmnumber{ #2}}
\newtheoremstyle{wzremark}{4pt}{4pt}
  {\normalfont\songti}{0pt}{\normalfont\bfseries}
  {}{.7em}{\thmname{#1}}
\theoremstyle{wzstatement}
\newtheorem{theorem}{定理}[chapter]
\newtheorem{lemma}[theorem]{引理}
\newtheorem{proposition}[theorem]{命题}
\newtheorem{corollary}[theorem]{推论}
\theoremstyle{wzdefinition}
\newtheorem{definition}[theorem]{定义}
\theoremstyle{wzproblem}
\newtheorem{example}{例题}[chapter]
\newtheorem{exercise}{习题}[chapter]
\theoremstyle{wzremark}
\newtheorem*{remark}{注}
\renewcommand{\proofname}{证明}
% Retain the native amsthm QED stack; only adjust the fixed typography.
\expandafter\patchcmd\csname\string\proof\endcsname{\normalfont \topsep6\p@\@plus6\p@\relax}
  {\normalfont\songti \topsep5\p@\@plus1\p@\relax}
  {}{\ClassError{elegantbook-original-adapter}{Cannot patch proof spacing}{Check the amsthm version.}}
\expandafter\patchcmd\csname\string\proof\endcsname{\itshape}{\normalfont\bfseries}
  {}{\ClassError{elegantbook-original-adapter}{Cannot patch proof heading}{Check the amsthm version.}}
\expandafter\patchcmd\csname\string\proof\endcsname{\ignorespaces}
  {\setlength{\parindent}{2em}\ignorespaces}
  {}{\ClassError{elegantbook-original-adapter}{Cannot patch proof paragraphs}{Check the amsthm version.}}
\AtBeginEnvironment{proof}{\par\Needspace{3\baselineskip}}
\renewcommand{\qedsymbol}{\ensuremath{\square}}
\newenvironment{solution}{\begin{proof}[解答]}{\end{proof}}
\newcommand{\wzcheckchapter}{\ifnum\value{chapter}<1
  \ClassError{elegantbook-original-adapter}{Numbered mathematics before chapter 1}{Move it after the first originalchapter.}\fi}

\newcommand{\wzregisterguard}[1]{\AtBeginEnvironment{#1}{\wzcheckchapter\par\Needspace{3\baselineskip}}}
\forcsvlist{\wzregisterguard}{theorem,lemma,proposition,corollary,definition,example,exercise}
% Front matter and bibliography are not numbered mathematical chapters.
\newcommand{\originalpreface}{\par\addvspace{10pt}\Needspace{5\baselineskip}
  \noindent{\bfseries\fontsize{12}{17}\selectfont 前言}\par\nobreak\vspace{4pt}}
\newcommand{\originalcontents}{\par\addvspace{14pt}
  \begin{center}{\color{structurecolor}\bfseries\fontsize{14.4}{20}\selectfont 目录}\end{center}
  \vspace{-10pt}\@starttoc{toc}}
\renewcommand{\bibname}{参考文献}
\renewenvironment{thebibliography}[1]{
  \par\addvspace{12pt}\Needspace{3\baselineskip}\phantomsection
  \addcontentsline{toc}{chapter}{\bibname}
  \markboth{\bibname}{}
  \noindent{\color{structurecolor}\bfseries\fontsize{14.4}{20}\selectfont\bibname}\par\nobreak\vspace{6pt}
  \list{\@biblabel{\@arabic\c@enumiv}}
    {\settowidth\labelwidth{\@biblabel{#1}}\leftmargin\labelwidth
     \advance\leftmargin\labelsep\usecounter{enumiv}
     \let\p@enumiv\@empty\renewcommand\theenumiv{\@arabic\c@enumiv}
     \setlength{\itemsep}{3pt}\setlength{\parsep}{0pt}}
  \sloppy\clubpenalty4000\widowpenalty4000\sfcode`\.=1000\relax
}{\def\@noitemerr{\@latex@warning{Empty bibliography}}\endlist}

\numberwithin{equation}{chapter}
\renewcommand{\theequation}{\thechapter.\arabic{equation}}
\allowdisplaybreaks[2]
\emergencystretch=2em
\clubpenalty=6000
\widowpenalty=6000
\raggedbottom
\endinput
\end{filecontents*}
% WZ-LOCKED-CLASS-END
\documentclass{elegantbook-original-adapter}
\usepackage{tikz,pgfplots,longtable,booktabs}
\usetikzlibrary{arrows.meta,calc,positioning,patterns,decorations.pathmorphing}
\pgfplotsset{compat=1.18}
\newcommand{\dd}{\,\mathrm d}
\newcommand{\ee}{\mathrm e}
\newcommand{\ii}{\mathrm i}
\newcommand{\Ms}{M_{\odot}}
\newcommand{\Ls}{L_{\odot}}
\newcommand{\Rs}{R_{\odot}}
\newcommand{\kb}{k_{\mathrm B}}
\newcommand{\mpl}{M_{\mathrm{Pl}}}
\newcommand{\src}[1]{\par\Needspace{5\baselineskip}\noindent{\color{structurecolor}\textbf{原题定位:} #1}\par}
\renewcommand{\figurename}{图}
\renewcommand{\tablename}{表}
\hypersetup{pdftitle={天体物理与宇宙学},pdfauthor={中文重编贡献者}}
\begin{document}
\frontmatter
\thispagestyle{cover}
\begin{center}{\color{structurecolor}\bfseries\fontsize{17}{23}\selectfont 天体物理与宇宙学}\par
\vspace{5pt}{\bfseries 从恒星结构到宇宙热史}\par
\vspace{5pt}MIT OpenCourseWare 8.901、8.902 与 8.286 中文主题重编\par
2026年10月12日\end{center}
读者应已学过大学基础力学、电磁学、热学和量子力学, 能使用多元微积分与常微分方程. 曲率、张量和爱因斯坦方程承接本仓库广义相对论卷; 此处只调用均匀宇宙的 Friedmann 方程, 不重写该卷的几何体系. 本卷先建立光与物质之间的联系, 再研究自引力系统, 最后把同一套能量、输运和统计思想用于膨胀宇宙.

正文为中文主题重编, 并非原文件逐页翻译. 原题题面翻译与解答分开, 原编号由题前定位保留. 8.901 和 8.902 的答案由 AI 独立推导; 8.286 六道测验题核对官方解答后重新展开, 两道官方仅回指作业的题为独立补解. 数值计算随源码附上, 数值检查不代替一般推导.

\begin{itemize}
\item 8.901 Astrophysics I, Spring 2006, Deepto Chakrabarty. 官方课程有 26 次课历与十份作业, 没有公开整套讲义. 本卷只从公开课程范围和入选自足作业组织恒星基础, 不宣称全部课堂内容已译出.
\item 8.902 Astrophysics II, Fall 2023, Mark Vogelsberger. 讲义署名 Mark Vogelsberger、Stephanie O'Neil、David DePalma. 官方提供 122 页合订本与 21 份分讲文件, 两种形式有重复, 不叠加为独立课程覆盖. 六组作业、期中题组和数值项目实际可取得; 作业答案与正式期末卷未公开. 期中题组按作业计分, 不是 2023 年闭卷期中考试.
\item 8.286 The Early Universe, Fall 2013, Alan Guth. 官方列 23 次课历、22 份幻灯文件、十份作业和三次测验及解答. 本卷以均匀膨胀、宇宙热史和暴胀为主, 选八道测验题完整译解, 不把复习题库或每段视频都计为已全译.
\end{itemize}

基于 MIT OCW 的改编及本项目增补采用 \href{https://creativecommons.org/licenses/by-nc-sa/4.0/}{CC BY-NC-SA 4.0}: 保留署名, 非商业使用, 相同方式共享. 翻译、重排、补充推导及自绘图不代表 MIT 或原作者认可. 第三方照片、SDSS 图像、受限教材题面及未经确认可复用的数据不收录; 原件只列来源链接和哈希. 逐件清单与未选范围见源码包中的 \texttt{sources/} 与 \texttt{coverage.md}.

各题所用历史参数沿用题面, 例如 $H_0=71\,\mathrm{km\,s^{-1}Mpc^{-1}}$, 不将其作为当前观测最佳值. $\rho$ 表示质量密度, $u=\rho c^2$ 表示能量密度; 同一公式中不可混用二者.

\clearpage
\originalcontents
\clearpage
\mainmatter
\originalchapter{星系观测如何变成可比较的物理量}

本章对应 8.902 Fall 2023 第1讲与第2讲开头的恒星种群说明. 课程由 Mark Vogelsberger 主讲, 笔记由 Stephanie O'Neil 与 David DePalma 整理. 我们从探测器实际收到的光出发, 依次建立通量、星等、面亮度、光度函数和速度的联系. 多波段照片中的颜色只是观测量的一种显示, 要讨论恒星年龄或暗物质, 还须说明通带、距离与动力学模型. 本章不复制原讲义中第三方照片; 图均由给定数学模型自绘.

\originalsection{通量密度与比强度的区别}
设一束频率在 $[\nu,\nu+\dd\nu]$ 内的辐射穿过面积 $\dd A$. 面的法线与光传播方向夹角为 $\theta$, 比强度 $I_\nu$ 按
\[
\dd E=I_\nu\cos\theta\,\dd A\,\dd t\,\dd\nu\,\dd\Omega
\]
定义. 它把方向也分辨出来; 通量密度则对方向积分:
\[
F_\nu=\int I_\nu\cos\theta\dd\Omega.
\]
单位分别为 $\mathrm{erg\,s^{-1}cm^{-2}Hz^{-1}sr^{-1}}$ 与 $\mathrm{erg\,s^{-1}cm^{-2}Hz^{-1}}$. 对一个小角度的远方天体, 探测面朝向它, 可近似 $F_\nu=\int_{\rm source}I_\nu\dd\Omega$. 球对称、无吸收且静态的源若光谱光度为 $L_\nu$, 穿过任意包围它的球面的功率相同, 因此
\[
4\pi D^2F_\nu=L_\nu,\qquad F_\nu\propto D^{-2}.
\]
原笔记先列出的 $r^2$ 与随后能量守恒推导不一致; 此处以守恒式确定负二次方. 同一小块物理发光面的张角也随 $D^{-2}$ 变化, 故静态 Euclidean 条件下其比强度不因离远而改变. 总光变暗与每单位张角的光不变并不矛盾: 远处的源在图上占据更少的面积.

这条面亮度性质需附条件. 吸收、分辨率不足、宇宙红移都能改变测量. 后文将得到光子传播中的 $I_\nu/\nu^3$ 不变量, 由它推出宇宙学总面亮度的 $(1+z)^{-4}$ 衰减; 不能把这里的静态结论无条件用于高红移星系.

\originalsection{星等以对数记录通量比}
在同一通带和零点定义下,
\[
m_1-m_2=-2.5\log_{10}\frac{F_1}{F_2}.
\]
所以低星等表示更亮, 五个星等对应一百倍通量. Vega 系统以织女星的参考光谱定零点; AB 系统使用频率通量密度恒定的参考源 $F_{\nu,0}=3631\,\mathrm{Jy}$, $1\,\mathrm{Jy}=10^{-23}\,\mathrm{erg\,s^{-1}cm^{-2}Hz^{-1}}$. 单色 AB 星等写成
\[
m_{\rm AB}=-2.5\log_{10}\!\left[\frac{F_\nu}{\mathrm{erg\,s^{-1}cm^{-2}Hz^{-1}}}\right]-48.60.
\]
这里取对数的是规定单位下的数值, 不是直接对有量纲物理量取对数. 原课程参考通量的 $3.63\times10^{-20}$ 是上述零点的近似.

实际滤镜收一段频率, 并非一个点. 若已把响应及探测方式吸收到归一化权重 $W_X(\nu)$ 中, $\int W_X\dd\nu=1$, 则能量响应约定下定义 $\langle F_\nu\rangle_X=\int F_\nu W_X\dd\nu$, 再代入同一零点比值. 光子计数器的计数率正比于 $\int F_\nu T_X(\nu)\dd\nu/(h\nu)$, 因而应对目标和参考源使用相同的光子权重. 忽略这个区别会把不同光谱误当作相同通量. Vega 与 AB 的通带偏移由参考光谱与响应积分决定, 不是所有滤镜通用的一组常数.

\begin{example}
一个 AB 星等为20的源有 $F_\nu=3631\times10^{-8}\,\mathrm{Jy}=36.31\,\mu\mathrm{Jy}$. 若另一源星等22.5, 其通量为前者的十分之一. 这只比较同一频率或同一响应约定下的通量, 尚未给出两者总光度之比.
\end{example}

绝对星等 $M_X$ 定义为将同一源放到 $10\,\mathrm{pc}$ 时的视星等. 静态无消光条件下, 通量平方反比给出距离模数
\[
\mu\equiv m_X-M_X=5\log_{10}\frac{D}{10\,\mathrm{pc}},
\qquad D=10\,\mathrm{pc}\,10^{\mu/5}.
\]
例如 $\mu=30$ 对应 $D=10\,\mathrm{Mpc}$. 消光使 $m_X$ 增加 $A_X$; 对宇宙学源还需光度距离和把红移光谱映入观测通带的 $K$ 修正. 因而实际使用的是 $m_X=M_X+\mu(D_L)+A_X+K_X$, 并应声明各项是否已校正.

颜色 $m_X-m_Y$ 是两通带通量比的对数. 同一近邻源的距离项抵消, 故它比较光谱形状而不是总光度. 消光一般不抵消, 因为 $A_X\ne A_Y$. 恒星年轻、金属丰度较低或尘埃较少都可能使某些颜色变蓝; 只凭一个颜色不能唯一恢复年龄.

\originalsection{面亮度轮廓与星系形态}
令 $R$ 为投影半径, $I(R)$ 为单位投影面积的光度. 若轴对称, 环带面积为 $2\pi R\dd R$, 故 $L=2\pi\int_0^\infty I(R)R\dd R$. 指数盘有
\[
I(R)=I_0\ee^{-R/R_s},\quad L=2\pi I_0R_s^2,
\quad \frac{L(<R)}L=1-\ee^{-x}(1+x),\quad x=R/R_s.
\]
半光半径 $R_e$ 由右式等于 $1/2$ 确定, $R_e\simeq1.678R_s$. $R_s$ 和 $R_e$ 不能混用.

更一般的 Sérsic 形式可按半光半径写为
\[
I(R)=I_e\exp\left\{-b_n\left[(R/R_e)^{1/n}-1\right]\right\}.
\]
代换 $u=b_n(R/R_e)^{1/n}$ 给出
\[
L=2\pi nI_e\ee^{b_n}R_e^2b_n^{-2n}\Gamma(2n),\qquad
\gamma(2n,b_n)=\tfrac12\Gamma(2n).
\]
第二式才使 $R_e$ 成为半光半径. 指数盘为 $n=1$, $b_1\simeq1.678$; de Vaucouleurs 轮廓为 $n=4$, $b_4\simeq7.669$. 原讲义的中心归一化形式 $I_0\exp[-(R/R_0)^{1/n}]$ 与此式等价, 但 $R_0=R_e/b_n^n$.

\begin{center}
\begin{tikzpicture}
\begin{axis}[width=8.5cm,height=5cm,xlabel={$R/R_e$},ylabel={$I/I_e$},ymode=log,xmin=0.03,xmax=4,ymin=0.002,ymax=300,grid=major,legend style={font=\small}]
\addplot[structurecolor,thick,domain=0.03:4,samples=150]{exp(-1.678*(x-1))};
\addlegendentry{指数盘 $n=1$}
\addplot[black,dashed,domain=0.03:4,samples=150]{exp(-7.669*(x^0.25-1))};
\addlegendentry{$n=4$}
\end{axis}
\end{tikzpicture}
\end{center}
图比较同一 $I_e,R_e$ 下的理想轮廓, 不代表两个模型总光度相同. $n=4$ 的中心更集中、外翼也更宽. 测得一个轮廓并不直接给出质量: 还需质光比及其径向变化.

螺旋星系通常具有恒星和气体盘, 有持续恒星形成, 旋转对支持作用重要; 椭圆星系常以老恒星为主, 随机速度弥散较重要; 透镜状星系有恒星盘而冷气体和恒星形成较弱. 这些是群体上的趋势, 不应把每个椭圆星系都视为完全不转动, 或把形态分类当作唯一的演化时间序列. 环境与并合史也是需要动力学解释的变量.

\originalsection{光度函数及它的积分条件}
按单位体积、单位光度区间定义 $\phi(L)\dd L$. Schechter 经验函数为
\[
\phi(L)\dd L=\phi_*\left(\frac{L}{L_*}\right)^\alpha
\ee^{-L/L_*}\frac{\dd L}{L_*}.
\]
亮端由指数截断, 暗端由 $\alpha$ 控制. 令 $x=L/L_*$, 总数密度与光度密度分别是
\[
n=\phi_*\Gamma(\alpha+1)\quad(\alpha>-1),\qquad
j=\phi_*L_*\Gamma(\alpha+2)\quad(\alpha>-2).
\]
两个收敛条件不同. 课程例子的 $\alpha\simeq-1.09$ 不能直接积分到零光度求有限总数, 却允许有限总光度. 实际样本有选择限, 可写 $n(>L_{\min})=\phi_*\Gamma(\alpha+1,L_{\min}/L_*)$. 此处不把经验暗端外推到零当作无限多真实星系的观测证明.

课程给出的 $\phi_*\sim0.02h^3\,\mathrm{Mpc^{-3}}$, $L_*\sim10^{10}h^{-2}\Ls$ 是相应样本与通带的数量级例子. 这里的 $h$ 由\[
H_0=100h\,\mathrm{km\,s^{-1}Mpc^{-1}}
\]
定义; 改变距离标定同时改变所推断的体积与光度, 因而这两个 $h$ 次幂应一起保留.

\originalsection{速度、引力质量与暗物质线索}
光谱线中心的偏移给出红移, 定义 $1+z=\lambda_{\rm obs}/\lambda_0$. 对低红移近邻样本, 在忽略特有速度影响时, $D\simeq cz/H_0$. 星系内不同位置的谱线还给出视线速度. 倾斜盘的测量近似 $v_{\rm los}=v_c\sin i\cos\varphi$; 倾角不明会影响从投影速度恢复真实旋转速度.

若近似球对称且圆轨道成立, 牛顿引力给出 $v_c^2=GM(<r)/r$. 中心质量外应有 $v_c\propto r^{-1/2}$; 若观测外部旋转曲线近乎平坦, 则
\[
M(<r)\simeq\frac{v_c^2r}{G},\qquad
\rho(r)\simeq\frac{v_c^2}{4\pi Gr^2}.
\]
后一式是球模型下的推断, 不可直接当作薄盘的精确反演. 平坦曲线提示可见光之外有较广的引力质量分布, 但单一观测并不确定粒子种类. 随机运动支持的星系则需速度弥散、轨道各向异性与分布函数, 后文用 Jeans 方程建立这种联系.

原课程按2023年的背景讨论冷、温、自相互作用与极轻玻色子暗物质, 以及修改动力学的不同解释路线. 这些模型对小尺度结构、散射或波动尺度的预测不同. 重子气体、昏暗恒星和已知轻中微子不能在同一组观测限制下无条件承担全部暗物质; 微引力透镜等限制依赖致密天体的质量范围和空间分布. 本章不把课程中的概述改写成任何候选已经被发现或在全部参数区间排除的结论.

\originalsection{恒星种群合成是一个逆问题}
星系光谱由不同恒星的谱叠加, 简化模型写为 $L_\nu=\sum_aN_a\ell_{\nu,a}$. 恒星类别的温度不同, 其连续谱可用黑体作初步比较, 但真实谱含吸收线、金属丰度和表面重力信息. 一个更物理的前向模型为
\[
L_\nu(t)=\int_0^t\psi(t')\,\ell_\nu^{\rm SSP}(t-t',Z(t'))\dd t',
\]
其中 $\psi$ 为恒星形成率, $\ell_\nu^{\rm SSP}$ 是按初始质量函数归一化的单次形成恒星种群光谱, $Z$ 为金属质量分数. 再加入尘埃传输与滤镜响应才得到探测器数据. 不同年龄、金属丰度和尘埃可能产生近似相同颜色, 所以从少数通带恢复完整形成史通常不唯一. 后面的化学演化章会补充独立的守恒条件, 使观测、动力学和演化相互约束.

\originalchapter{星系支持与维里定理的边界条件}

本章对应 8.902 第2讲的整体动力学. 恒星与暗物质可以在共同引力势中运动, 气体还受压力、冷却与碰撞影响. 星系看起来稳定, 并不表示每颗恒星静止. 维里定理把微观运动与整体尺度联系起来, 但它是对系统边界和时间平均作出假设后的陈述, 不能只因一个物体受引力就写 $2K+U=0$.

\originalsection{从大小的二阶导数得到恒等式}
在惯性参考系中, 对固定的一组粒子取相对质心的位置 $\boldsymbol r_i$, 定义质量二次矩 $\mathcal I=\sum_i m_ir_i^2$, 动能 $K=\frac12\sum_im_iv_i^2$. 此处 $\mathcal I$ 是整体大小的加权度量, 不是某一旋转轴的惯性矩. 两次微分给出
\[
\frac12\dot{\mathcal I}=\sum_i\boldsymbol p_i\cdot\boldsymbol r_i,
\qquad
\frac12\ddot{\mathcal I}=2K+\sum_i\boldsymbol r_i\cdot\boldsymbol F_i.
\]
这一式尚未使用引力, 也没有假设平衡. 定义 $W=\sum_i\boldsymbol r_i\cdot\boldsymbol F_i$ 为力的维里项. 对两两内力, 由作用反作用将两个有向求和合并为一个无重复配对:
\[
W=\sum_{i<j}(\boldsymbol r_i-\boldsymbol r_j)\cdot\boldsymbol F_{i\leftarrow j}.
\]
若配对势能是 $V_{ij}(r_{ij})$, 则该项为 $-r_{ij}\dd V_{ij}/\dd r_{ij}$. 对齐次势 $V_{ij}=a_{ij}r_{ij}^{n}$, 有 $W=-nU$, $U=\sum_{i<j}V_{ij}$. 因此
\[
\frac12\ddot{\mathcal I}=2K-nU.
\]
牛顿引力 $n=-1$, 势零点取无穷远, 故 $U=-\sum_{i<j}Gm_im_j/r_{ij}$ 且 $W=U$. 加一个任意常数到 $U$ 不改变力, 却会破坏 $W=U$ 的写法; 这说明势的零点约定也属于条件.

\originalsection{时间平均不是瞬时静止}
在时间窗 $\mathcal T$ 上积分,
\[
2\langle K\rangle-n\langle U\rangle
=\frac{\dot{\mathcal I}(\mathcal T)-\dot{\mathcal I}(0)}{2\mathcal T}.
\]
对保持束缚、边界项随 $\mathcal T$ 增长比 $\mathcal T$ 慢的系统, 长时间右端趋零. 牛顿引力于是给出 $2\langle K\rangle+\langle U\rangle=0$. 单个有界轨道也能满足时间平均关系, 不要求速度分布已经是 Maxwell 分布. 相反, 正在整体塌缩或自由膨胀的系统有非零 $\ddot{\mathcal I}$, 不能用同一平衡式.

真实星系的测量常选一个空间孔径, 有粒子跨过边界, 还可能受外部潮汐场. 这时固定粒子集合的推导与固定空间区域的推导不同, 后者出现动量流和表面项. 本章的质量估算先使用孤立、近似稳定的整体模型; 后面若引入截断晕或气体外压, 将显式恢复相应项.

\originalsection{动力学质量需要几何系数与速度投影}
写引力势能 $U=-aGM^2/R$, $a$ 由密度轮廓和 $R$ 的定义决定. 令 $\langle v^2\rangle$ 是相对质心的三维质量加权速度二次矩, 则
\[
2K=M\langle v^2\rangle,\qquad
M=\frac{R\langle v^2\rangle}{aG}.
\]
等向随机运动且无整体旋转时 $\langle v^2\rangle=3\sigma_{\rm los}^2$, 才能把观测视线弥散写入. 对均匀球半径 $R$, 积分组装各壳层得到
\[
U=-\int_0^R\frac{GM(r)}r\dd M(r)
=-\frac35\frac{GM^2}{R},\qquad
M=\frac{5R\sigma_{\rm los}^2}{G}.
\]
天文常用 $G\simeq4.301\times10^{-3}\,\mathrm{pc\,\Ms^{-1}(km\,s^{-1})^2}$. 例如 $R=10\,\mathrm{kpc}$、$M=10^{11}\Ms$ 给 $\sqrt{GM/R}\simeq207\,\mathrm{km\,s^{-1}}$, 与课程给出的星系速度数量级相符; 均匀球模型的三维均方根速度是它乘 $\sqrt{3/5}$. 两种数值不同的原因是一个只作尺度估算, 另一个固定了几何系数.

\begin{example}
若把三维均方根速度误当成视线弥散, 同一 $a,R$ 下会把质量低估三倍. 若实际各向异性或质光比变化显著, 即使恢复这个三倍因子, 仍不能保证质量估计精确. 需要空间分辨运动学和轨道模型来解除退化.
\end{example}

平衡下总机械能 $E=K+U=U/2=-K<0$. 若系统缓慢失去能量并在相近结构中重新达到维里平衡, 最终 $K$ 反而增大, 系统收缩. 自引力系统的负热容由这条能量关系表现出来. 这不表示恒星在每次辐射后瞬间升温; 气体输运、两体弛豫和重新平衡的时间尺度需要分别比较.

\originalsection{压力项如何给出理想气体状态方程}
考虑理想气体被体积 $V$ 的壁面限制. 壁面对气体的宏观力沿内法向, 力的维里为
\[
W_{\rm wall}=-\oint_{\partial V}P\boldsymbol r\cdot\boldsymbol n\dd A.
\]
若壁面压力均匀, 散度定理 $\nabla\cdot\boldsymbol r=3$ 给 $W_{\rm wall}=-3PV$. 气体粒子无相互作用势, 平衡维里式于是为 $2K=3PV$. 用非相对论单原子气体 $K=3N\kb T/2$ 得 $PV=N\kb T$.

此处壁面力不必被假装成一个处处满足 $V\propto r^2$ 的配对势能. 原笔记用齐次势的记法给理想气体例子, 容易让力的维里与势能混淆; 上述表面力推导直接说明了压力的来源. 对自引力气体的静水平衡, 无整体流动时相应积分式是
\[
3\int_VP\dd V+U-3P_sV=0
\]
（球形且边界压力均匀）. 若 $P_s=0$, 压力积分支持负引力能; 若外压不零, 忽略它会改错平衡关系. 后面的恒星结构与 Jeans 不稳定性将再次使用这一物理联系.

\originalchapter{两体弛豫、相混合与快速集体演化}

本章对应 8.902 第2讲末至第3讲. 要把恒星视作连续的无碰撞分布, 须比较轨道穿越时间与粒子离散引力引起的弛豫时间. “无碰撞”并非“无引力”: 恒星不断受共同引力势加速, 只是单独近距离遭遇改变轨道的速度通常较慢. 星系仍可在时变集体势中很快调整, 这种过程需要与两体随机扩散区分.

\originalsection{一次弱引力遭遇的冲量}
先取质量 $m$ 的场恒星近似静止, 另一恒星以相对速度 $v$ 沿直线 $x=vt$ 通过, 冲量参数为 $b$. 被测恒星的横向加速度大小
\[
a_\perp=\frac{Gmb}{(b^2+v^2t^2)^{3/2}}.
\]
积分奇对称的纵向分量为零, 横向分量为
\[
\delta v_\perp=\int_{-\infty}^{\infty}a_\perp\dd t
=\frac{Gm}{bv}\int_{-\infty}^{\infty}\frac{\dd u}{(1+u^2)^{3/2}}
=\frac{2Gm}{bv}.
\]
这个公式依赖 $\delta v\ll v$, 不能把它外推成任意大角散射. 两体相对轨道的精确散射角满足 $\tan(\Theta/2)=G(m_1+m_2)/(bv^2)$, 因而 $90^\circ$ 的尺度是 $b_{90}=G(m_1+m_2)/v^2$. 对等质量、维里尺度 $v^2\sim GN m/R$, 得 $b_{90}\sim2R/N$. 系数取决于这里的 $v$ 是相对速度还是单星速度, 后面估算只保留一致的数量级.

\begin{center}
\begin{tikzpicture}[>=Latex,scale=0.9]
\draw[->](-3,1.1)--(3,1.1) node[right]{$v$};
\fill(0,0)circle(2pt) node[below]{场恒星};
\draw[dashed](0,0)--(0,1.1) node[midway,right]{$b$};
\fill(-1.6,1.1)circle(2pt) node[above]{被测恒星};
\draw[->,structurecolor](-1.6,1.05)--(-1.3,0.65) node[left]{$a_\perp$};
\end{tikzpicture}
\end{center}
图用于直线弱偏转近似, 不画成精确双曲轨道. 在星系里 $b_{90}\ll R$, 强遭遇截面 $\pi b_{90}^2$ 很小. 均匀球的数密度 $n=3N/(4\pi R^3)$ 给一次穿越强遭遇的概率尺度 $n\pi b_{90}^2R\sim3/N$. 对 $N\sim10^{11}$ 极小; 对球状星团 $N\sim10^4$ 较大, 仍须累积许多次穿越才能估计整体演化.

\originalsection{许多弱遭遇产生对数积分}
在 $\dd t$ 内, 冲量参数在 $[b,b+\dd b]$ 的遭遇数近似 $\dd N_{
\rm enc}=2\pi b\dd b\,nv\dd t$. 若不同遭遇方向随机且相关可忽略, 平均冲量相消, 速度二次变化累积:
\[
\frac{\dd\langle(\Delta v)^2\rangle}{\dd t}
=\int_{b_{\min}}^{b_{\max}}2\pi bnv\left(\frac{2Gm}{bv}\right)^2\dd b
=\frac{8\pi G^2m^2n}{v}\ln\Lambda,
\quad\Lambda=\frac{b_{\max}}{b_{\min}}.
\]
每个对数区间的贡献相当, 因此大量远距离弱遭遇总效应可以超过罕见强遭遇. $b_{\min}$ 不能任意趋零: 应由大角偏转、恒星半径或其他小尺度物理截断; $b_{\max}$ 约为系统尺度或使均匀独立遭遇近似失效的尺度. $\ln\Lambda$ 被称作 Coulomb 对数, 名称来自相似积分结构, 此处相互作用仍是引力.

以速度改变达到原速度数量级定义 $t_{\rm rel}$, 得
\[
t_{\rm rel}\simeq\frac{v^3}{8\pi G^2m^2n\ln\Lambda}.
\]
用 $t_{\rm cross}=R/v$, $v^2\simeq GNm/R$ 及均匀球 $n$,
\[
\frac{t_{\rm rel}}{t_{\rm cross}}\simeq\frac{N}{6\ln\Lambda}
\sim\frac{N}{6\ln N}.
\]
系数6来自这一简化冲量与均匀球模型, 更完整的速度分布会改变系数. 物理结论由大的 $N/\ln N$ 决定. 星系 $R=10\,\mathrm{kpc}$、$v=200\,\mathrm{km\,s^{-1}}$ 的穿越时间约 $4.9\times10^7$ 年, 取 $N=10^{11}$ 给弛豫时间约 $3\times10^{16}$ 年, 远超课程宇宙年龄尺度. 无碰撞近似因而合理. 这个时间不是一条圆轨道周期, 后者另有 $2\pi$ 因子.

\begin{example}
简化球状星团取 $N=10^4$, $R=3\,\mathrm{pc}$, $v=5\,\mathrm{km\,s^{-1}}$, 有 $t_{\rm cross}\simeq0.59$ 百万年, $N/(6\ln N)\simeq181$, 故 $t_{\rm rel}$ 约一亿年. 这是一致尺度估算, 并未声称所有球状星团都有相同弛豫时间; 中心密度和质量谱会显著改变局部演化.
\end{example}

\originalsection{无碰撞分布仍可改变粗粒化外观}
若 $f(\boldsymbol x,\boldsymbol v,t)$ 是相空间数密度, 无碰撞动力学意味着沿轨道 $\dd f/\dd t=0$. 这并不表示固定位置的 $f$ 不变. 相空间团块可被轨道流拉伸和卷绕, 精细分布保持信息, 有限分辨的平均却趋于较平滑的外观. 将后者叫作粗粒化弛豫, 必须保留精细动力学的可逆性.

对可积定常势, 用作用量 $\boldsymbol J$ 与角变量 $\boldsymbol\theta$, 有 $\dot{\boldsymbol J}=0$, $\dot{\boldsymbol\theta}=\boldsymbol\Omega(\boldsymbol J)$. 不同 $J$ 的频率不同, 角向不均匀性逐渐失相. 一个 Fourier 分量的可观测平均包含
\[
\int\dd J\,a_k(J)\exp[-\ii k\Omega(J)t].
\]
若 $a_k$ 有有限宽度且 $\Omega'(J)\ne0$, 不同相位在积分中抵消, 混合时间量级 $t_{\rm mix}\sim[|k\Omega'|\Delta J]^{-1}$. 这是相混合. 若势严格谐振且所有轨道频率相同, 团块只共同转动, 不会凭这一机制变成均匀角分布. 因而“等一段时间必然相混合”需要频率展宽条件.

没有碰撞时精细 $f$ 的最大值也不会被动力学任意增加. 从光滑观测恢复一个粗粒化 $\bar f$ 常会丢失这些细丝, 但这属于分辨率与平均的选择, 不能当作已证明熵在每条精细轨道上严格增加.

\originalsection{时变引力势直接交换单星能量}
令单星的单位质量能量 $E=v^2/2+\Phi(\boldsymbol x,t)$, 运动方程 $\dot{\boldsymbol v}=-\nabla\Phi$. 用链式法则,
\[
\frac{\dd E}{\dd t}=\boldsymbol v\cdot\dot{\boldsymbol v}
+\frac{\partial\Phi}{\partial t}+\boldsymbol v\cdot\nabla\Phi
=\frac{\partial\Phi}{\partial t}.
\]
定常势保留每颗星的能量, 时变势却可交换它们的能量而无需近距离碰撞. 若势井在恒星通过时变深, 出井时可能带走较少能量; 变浅时可带走较多. 在自洽系统中恒星能量之和还需计入共享势能的配对计数, 不能简单把所有单位质量 $E$ 相加代替系统总机械能.

星系形成、塌缩和并合使集体势在动力学时间内大幅改变, 由此产生快速的暴烈弛豫（violent relaxation）. 可用
\[
t_{\rm vr}\sim\left\langle\left(\frac{\partial_t\Phi}{E}\right)^2\right\rangle^{-1/2}
\]
估计势变化引起的能量调整时间, 在非零典型 $E$ 的模型中它常与塌缩或自由落体时间同阶. 这不是普适精确公式, 也没有保证最终分布一定 Maxwell 化. 调整可能不完整, 还保留形成史与守恒量.

三种过程至此有清楚分工: 两体弛豫是离散遭遇的随机扩散; 相混合是在定常势中由频率差造成角相位失相; 暴烈弛豫依靠时变集体势改变单星能量. 后面以 Poisson 方程求共同势, 再以无碰撞 Boltzmann 方程和分布函数描述恒星, 正是建立在这些时间尺度比较上.

\originalchapter{共同引力势、薄盘与暗晕模型}

本章对应 8.902 第4讲及第5讲开头的势—密度对. 对许多恒星的运动, 先不逐个计算所有配对力, 而用平滑质量密度 $\rho$ 产生共同势 $\Phi$. 一个可用模型需同时满足力的方向、Poisson 方程和边界条件. 原笔记中数处势的正负号、测试质量和 $G$ 因子互相不一致, 以下固定 $\Phi$ 为单位质量势, 以这些条件逐项复核.

\originalsection{Poisson 方程与势零点}
加速度 $\boldsymbol a=-\nabla\Phi$, 引力 Gauss 定律为 $\nabla\cdot\boldsymbol a=-4\pi G\rho$, 故 $\nabla^2\Phi=4\pi G\rho$. 使用 $\nabla^2(1/|\boldsymbol r-\boldsymbol r'|)=-4\pi\delta^{(3)}(\boldsymbol r-\boldsymbol r')$, 对有限孤立质量分布取无穷远势零,
\[
\Phi(\boldsymbol r)=-G\int\frac{\rho(\boldsymbol r')}{|\boldsymbol r-\boldsymbol r'|}\dd^3r'.
\]
负号给吸引力. $[\Phi]=L^2T^{-2}$; 测试质量的势能是 $m\Phi$, 不应把 $m$ 写入共同势本身. 自引力系统总势能为 $U=\frac12\int\rho\Phi\dd^3r$, 其中 $1/2$ 避免配对重数.

球对称时, Gauss 定律直接给
\[
\Phi'(r)=\frac{GM(<r)}{r^2},\qquad
M(<r)=4\pi\int_0^r\rho(s)s^2\dd s.
\]
于是 $\boldsymbol a=-GM(<r)\hat{\boldsymbol r}/r^2$. 反过来给定一个球对称势, 可由
\[
\rho(r)=\frac1{4\pi Gr^2}\frac{\dd}{\dd r}\left(r^2\Phi'(r)\right)
\]
恢复密度, 再检查是否非负. 给定任意看起来平滑的 $\Phi$ 并不保证它对应物理质量分布.

\originalsection{点质量与均匀球的内外连接}
点质量给 $\Phi=-GM/r$, 在 $r>0$ 满足 Laplace 方程, 质量由原点的 delta 源承担. 均匀球半径 $R$, $\rho=3M/(4\pi R^3)$, 内部质量 $M(<r)=M(r/R)^3$. 积分径向力并使 $r=R$ 的势连续, 得
\[
\Phi(r)=\begin{cases}
-\dfrac{GM}{2R^3}(3R^2-r^2),&r\le R,\\[3pt]
-\dfrac{GM}{r},&r\ge R.
\end{cases}
\]
势和一阶导数均连续, 二阶导数跨密度突变可以不同. 球内加速度正比于 $-\boldsymbol r$, 因此粒子作谐振运动, $\omega^2=GM/R^3=4\pi G\rho/3$. 球壳外质量不贡献局部力, 却改变内部势的常数; 要算逃逸速度就不能丢掉这个常数.

\begin{example}
在有限均匀球中心, $v_{\rm esc}^2=-2\Phi(0)=3GM/R$; 表面则为 $2GM/R$. 若只把内部 $\Phi\propto r^2$ 的变动部分当完整势, 会误认为中心逃逸速度为零. 逃逸问题同时依赖整个质量分布与无穷远边界条件.
\end{example}

\originalsection{Mestel 薄盘的质量在导数跳跃中}
令柱坐标为 $(R,\varphi,z)$, $s=\sqrt{R^2+z^2}$. 课程的 Mestel 势可等价写为
\[
\Phi(R,z)=v_0^2\ln\frac{s+|z|}{2r_0}.
\]
$r_0$ 只确定势的加法常数. 在 $z\ne0$ 处直接微分得到 $\nabla^2\Phi=0$. 但不能据此说总密度为零, 因为 $|z|$ 的导数在盘面跳跃. 对 $z>0$, $\partial_z\Phi=v_0^2/s$; 对 $z<0$ 则为 $-v_0^2/s$. 因而
\[
\partial_z\Phi|_{0^+}-\partial_z\Phi|_{0^-}
=\frac{2v_0^2}{R}=4\pi G\Sigma(R),\qquad
\Sigma(R)=\frac{v_0^2}{2\pi GR}.
\]
这就是 $\rho=\Sigma(R)\delta(z)$ 的分布意义 Poisson 方程. 跳跃方向给 $z>0$ 的加速度向下, 与盘的吸引相符. 在盘面 $\partial_R\Phi=v_0^2/R$, 所以 $v_c^2=R\partial_R\Phi=v_0^2$, 旋转曲线恰为平坦.

注意这与球密度 $\rho\propto r^{-2}$ 是两种不同几何的模型. 薄盘的 $M(<R)=2\pi\int_0^R\Sigma(s)s\dd s=v_0^2R/G$ 恰有同样线性尺度, 但不能据此把球壳定理套给任意盘. 两者各自以完整势满足圆轨道条件. 无限 Mestel 盘的总质量发散, 中心也奇异; 它只能在有限半径区间近似真实盘. 盘外截断会改变整体势和旋转曲线, 不能保持每条公式都精确.

\originalsection{NFW 暗晕及集中度的物理定义}
NFW 是课程用于无碰撞冷暗物质模拟的经验球形轮廓. 取尺度半径 $a$, $x=r/a$,
\[
\rho(r)=\frac{\rho_s}{x(1+x)^2}.
\]
内侧 $r\ll a$ 为 $r^{-1}$ 尖峰, 外侧为 $r^{-3}$. 质量积分用 $x/(1+x)^2=1/(1+x)-1/(1+x)^2$ 得
\[
M(<r)=4\pi\rho_sa^3\left[\ln(1+x)-\frac{x}{1+x}\right].
\]
令 $m(x)$ 为方括号, 则圆速度 $v_c^2=4\pi G\rho_sa^2m(x)/x$. 势满足无穷远零点时
\[
\Phi(r)=-4\pi G\rho_sa^2\frac{\ln(1+x)}x.
\]
微分得到 $\Phi'=GM(<r)/r^2$, 是对常数和 $G$ 因子的独立检查. 原课式漏写的 $G$ 在此恢复; 否则量纲不是速度平方. 当 $x\to0$ 势趋有限常数, 尖峰密度不必使势本身发散. 当 $x\to\infty$ 质量对数发散但势趋零, 两者也不矛盾.

定义临界密度 $\rho_c=3H^2/(8\pi G)$, $r_{200}$ 内平均密度为 $200\rho_c$, 集中度 $c_{200}=r_{200}/a$. 此处下标避免把集中度与光速 $c$ 混淆. 由两种质量表达式相等,
\[
M_{200}=\frac{4\pi}{3}200\rho_cr_{200}^3,
\qquad
\frac{\rho_s}{\rho_c}=\frac{200}{3}\frac{c_{200}^3}{m(c_{200})}.
\]
因此 $M_{200},c_{200}$ 决定模型两个参数. 模拟中的质量—集中度趋势可减少平均意义的自由度, 但同质量晕仍有形成史和散布, 不应把原笔记“同质量完全相同”当作精确唯一性定理. 以平均宇宙物质密度而非临界密度定义的 $r_{200}$ 又是另一约定, 需声明参考密度和红移.

\begin{center}
\begin{tikzpicture}
\begin{axis}[width=8cm,height=4.8cm,xmode=log,ymode=log,xlabel={$r/a$},ylabel={$\rho/\rho_s$},xmin=0.03,xmax=30,grid=major]
\addplot[structurecolor,thick,domain=0.03:30,samples=160]{1/(x*(1+x)^2)};
\end{axis}
\end{tikzpicture}
\end{center}
图直接采样轮廓, 不是对观测数据的拟合. 要把模型作为有限系统使用, 还需截断或与环境匹配; 在截断半径内力可保留同一包围质量, 势的常数及外部逃逸条件却会改变.

\originalsection{小尺度问题应比较同一种可观测量}
课程讨论核—尖峰、旋转曲线多样性、缺失卫星和过大而难以没有恒星的子晕等问题. 比较时不能把纯暗物质子晕数直接等同于已观测发光卫星数: 气体冷却、反馈、潮汐剥离及观测选择都介入. 由恒星轨道推断中央暗物质密度又受到各向异性退化, 需把动力学误差与模型预测一起比较.

这些问题激发重子反馈以及不同暗物质模型的检验, 并非其中任意一项已经单独确定了粒子性质. 本章保留2023年课程的问题结构, 不复制含权利例外的图片, 也不把课程的概述包装为2026年最新观测结论. 下一章将从给定势建立轨道, 再解释从运动学反推势时还缺哪些条件.

\originalchapter{恒星轨道、导引中心与局部旋转测量}

本章对应 8.902 第4讲末、第5讲与第6讲开头. 已知共同势后, 每颗恒星的轨道由初值与守恒量决定. 球对称势的运动可降到一个平面; 轴对称盘则通常只保留轴向角动量. 这两种对称性不同, 正是圆轨道、径向摆动和竖直摆动能够分开近似的前提.

\originalsection{守恒量减少自由变量, 不自动关闭轨道}
定常势中单位质量能量 $E=v^2/2+\Phi$ 守恒, 因为 $\dd E/\dd t=\partial_t\Phi=0$. 球对称势的加速度沿径向, 所以单位质量角动量 $\boldsymbol\ell=\boldsymbol r\times\boldsymbol v$ 也守恒. 非零 $\boldsymbol\ell$ 把运动限制在垂直于它的平面. 取该平面的极坐标 $(r,\psi)$,
\[
\boldsymbol v=\dot r\hat{\boldsymbol e}_r+r\dot\psi\hat{\boldsymbol e}_\psi,
\quad \ell=r^2\dot\psi,
\quad E=\tfrac12\dot r^2+\Phi_{\rm eff}(r),
\quad\Phi_{\rm eff}=\Phi+\frac{\ell^2}{2r^2}.
\]
径向转折点满足 $E=\Phi_{\rm eff}$, 径向加速度 $\ddot r=-\Phi_{\rm eff}'$. 角变量仍通过 $\dot\psi=\ell/r^2$ 重建. 当 $\ell=0$ 时是径向运动, 不宜再用角变量参数化轨道.

只保留轨道平面时有 $(r,\psi,\dot r,\dot\psi)$ 四个变量, 再固定 $E,\ell$ 后可进一步降维. 原课把“有守恒量”直接说成四维相空间, 容易混淆消去方向、固定守恒量与整个系综的维数. 在一般轴对称势中, $E,\ell_z$ 也不保证所有轨道可积; 不能凭两个积分写出普遍第三积分.

\originalsection{Binet 方程与闭合条件}
令 $u=1/r$, 由于 $\dd/\dd t=\ell u^2\dd/\dd\psi$, 有 $\dot r=-\ell u'$ 和 $\ddot r=-\ell^2u^2u''$. 代入 $\ddot r-r\dot\psi^2=-\Phi'(r)$,
\[
u''+u=\frac{\Phi'(1/u)}{\ell^2u^2}.
\]
这不是“去掉 $r=0$ 奇点”的保证, 而是在非零角动量区用角度取代时间. Kepler 势 $\Phi=-GM/r$ 给常数右端 $GM/\ell^2$, 所以
\[
u=\frac{GM}{\ell^2}[1+e\cos(\psi-\psi_0)],\qquad
r=\frac{p}{1+e\cos(\psi-\psi_0)},\quad p=\ell^2/(GM).
\]
对 $0\le e<1$, 是闭合椭圆, $E=-GM/(2a)$ 与 $p=a(1-e^2)$. 一般势中的径向周期与角向周期未必可公度, 轨道常为进动玫瑰形. 近圆轨道的一径向周期积累角度 $2\pi\Omega/\kappa$; 只有合适频率比时轨道闭合.

课程用 $\Phi=-GM/r-2G^2M^2/(c^2r^2)$ 演示进动. 该玩具中心势给
\[
u''+\left(1-\frac{4G^2M^2}{\ell^2c^2}\right)u=\frac{GM}{\ell^2},
\qquad \Delta\psi\simeq\frac{4\pi GM}{pc^2}.
\]
它可以说明频率改变会产生进动, 但不等于 Schwarzschild 的完整弱场轨道方程. 后者给 $6\pi GM/(pc^2)$, 含速度依赖的相对论动力学, 见已有 GR 卷. 本卷不重复 GR 推导, 也不把上述玩具模型声称为精确水星进动公式. 原课的振子式中位移与速度导数混排, 此处以明确的 $u''+k^2u$ 恢复其示范目的.

\originalsection{轴对称势的二维有效运动}
取光滑定常势 $\Phi(R,z)$, 不依赖方位角 $\varphi$, 则 $\ell_z=R^2\dot\varphi$ 守恒. 柱坐标单位向量的时间导数给
\[
\ddot R-R\dot\varphi^2=-\partial_R\Phi,
\quad R\ddot\varphi+2\dot R\dot\varphi=0,
\quad\ddot z=-\partial_z\Phi.
\]
定义 $\Phi_{\rm eff}=\Phi+\ell_z^2/(2R^2)$ 后,
\[
\ddot R=-\partial_R\Phi_{\rm eff},\quad
\ddot z=-\partial_z\Phi_{\rm eff},\quad
E=\tfrac12(\dot R^2+\dot z^2)+\Phi_{\rm eff}.
\]
可达区域由 $E\ge\Phi_{\rm eff}$ 限制, 方位运动由瞬时 $R$ 重建, 因而通常不是角速度严格均匀的圆周运动. 若势对 $z\mapsto-z$ 对称且光滑, 盘面圆轨道 $R=R_g,z=0$ 满足
\[
\partial_R\Phi(R_g,0)=\frac{\ell_z^2}{R_g^3},
\qquad \Omega_g^2=\frac{\partial_R\Phi(R_g,0)}{R_g}.
\]
导引中心 $R_g$ 是与给定 $\ell_z$ 相应的圆轨道半径, 不是恒星每时每刻的实际位置. 圆轨道是否稳定, 还需有效势二阶导数的符号.

\originalsection{本轮近似由有效势 Hessian 得到}
令 $x=R-R_g$, 在光滑对称盘面附近展开. 一阶项由圆轨道条件消失, 混合项由反射对称消失, 因而
\[
\Phi_{\rm eff}\simeq\Phi_{\rm eff}(R_g,0)+\tfrac12\kappa^2x^2+\tfrac12\nu_z^2z^2,
\quad
\kappa^2=\Phi_{RR}+3\Omega_g^2,\quad \nu_z^2=\Phi_{zz}.
\]
$\kappa$ 为径向本轮频率, $\nu_z$ 为竖直频率, 后者加下标避免与辐射频率混淆. 比较 $\Omega^2(R)=\Phi_R(R,0)/R$, 得
\[
\kappa^2=R\frac{\dd\Omega^2}{\dd R}+4\Omega^2
=\frac1{R^3}\frac{\dd}{\dd R}(R^4\Omega^2).
\]
两频率平方都正才有两个稳定小振动. 取 $x=X\cos(\kappa t+\alpha)$, $z=Z\cos(\nu_zt+\beta)$, 又由角动量守恒
\[
\dot\varphi\simeq\Omega_g(1-2x/R_g),\quad
\varphi\simeq\Omega_gt+\varphi_0-
\frac{2\Omega_gX}{\kappa R_g}\sin(\kappa t+\alpha).
\]
随导引中心旋转的局部切向坐标 $y=R_g(\varphi-\Omega_gt-\varphi_0)$, 故 $y=-2\Omega_gX\sin(\kappa t+\alpha)/\kappa$. 平面内运动是逆向的本轮椭圆, 长短轴比由 $2\Omega_g/\kappa$ 决定.

\begin{example}
若 $v_c\propto R^q$, 则 $\Omega\propto R^{q-1}$, $\kappa^2=2(q+1)\Omega^2$. Kepler 区 $q=-1/2$ 给 $\kappa=\Omega$; 平坦曲线 $q=0$ 给 $\kappa=\sqrt2\Omega$; 刚体旋转 $q=1$ 给 $\kappa=2\Omega$. 若 $q<-1$, 径向频率平方负, 对应圆轨道不稳定, 不应仍画成稳定小椭圆.
\end{example}

\begin{center}
\begin{tikzpicture}[>=Latex]
\draw[->](-1.6,0)--(1.6,0)node[right]{$x/X$};
\draw[->](0,-1.7)--(0,1.7)node[above]{$y/X$};
\draw[structurecolor,thick,domain=0:360,samples=120,variable=\t]plot({cos(\t)},{-sqrt(2)*sin(\t)});
\draw[->,structurecolor,thick](1,0)--(0.98,-0.25);
\node[below right] at(0,0){导引中心};
\end{tikzpicture}
\end{center}
图为平坦曲线的小振幅极限, 箭头取从 $x=X$ 开始的运动方向. 无限薄 Mestel 盘的势在 $z=0$ 有 $|z|$ 尖角, $\Phi_{zz}$ 不是普通连续函数, 不满足这个光滑竖直 Taylor 近似. 真实有限厚盘或其他光滑模型才能直接使用有限 $\nu_z$.

\originalsection{Oort 常数需要区分速度与角速度}
在太阳附近作 $d\ll R_0$ 的局部展开, 假设扣除太阳特有运动后恒星随近圆轨道运动. 定义
\[
A=-\tfrac12R_0\Omega'_0=\tfrac12(v_c/R_0-v_c'),
\qquad B=-\Omega_0-\tfrac12R_0\Omega'_0.
\]
因此 $A-B=\Omega_0$, $\kappa_0^2=-4B(A-B)$. $A,B$ 都是速度除长度的量纲; 原笔记的 $A$ 缺少 $R_0$ 时不满足量纲.

若银经 $l$ 从朝银心方向起算, 银纬为 $b$, 正视线速度指远离太阳, 则局部一阶表达式为
\[
v_{\rm los}\simeq dA\sin(2l)\cos^2b,\qquad
v_l\simeq d[A\cos(2l)+B]\cos b.
\]
相应沿银经的大圆角速度 $\mu_l^*\equiv\dot l\cos b=v_l/d$; 故在盘面 $b=0$, $\mu_l^*=A\cos(2l)+B$, 而不是再乘一个距离. 原课用 $\mu$ 表示带 $d$ 的式子, 若解释为角自行会量纲错误; 此处分别记 $v_l$ 与 $\mu_l^*$. 实际拟合还要考虑距离误差、恒星速度弥散、非轴对称流和非圆轨道, 局部公式不是整个银河的通用速度场.

\originalsection{光度—速度幂律中的额外假设}
维里尺度 $v^2\sim GM/R$, 质光比 $\Upsilon=M/L$, 平均面亮度尺度 $\Sigma_L\sim L/R^2$ 合并给
\[
v^4\sim G^2\Upsilon^2\Sigma_LL,\qquad
L\sim\frac{v^4}{G^2\Upsilon^2\Sigma_L}.
\]
若同类星系 $\Upsilon,\Sigma_L$ 及结构系数近似相同, 才得到 $L\propto v^4$. 课程以此解释螺旋星系 Tully--Fisher 与椭圆星系 Faber--Jackson 关系的尺度动机. 前者比较旋转速度, 后者比较弥散, 不能互换测量定义. 经验斜率、散布和波段依赖需由样本标定, 不是维里定理单独固定. 一旦标定可由速度估计光度, 再与通量比较求距离, 这把本章动力学与第1章观测连接起来.

\originalchapter{无碰撞分布函数与 Jeans 质量推断}

本章对应 8.902 第6讲. 知道一个势里的单条轨道还不足以描述星系: 需知道每类轨道上放了多少恒星. 用相空间分布函数记录这种系综, 对其演化方程取速度矩, 才得到观测常用的密度、平均速度与弥散关系. 矩方程简化了问题, 同时留下闭合与各向异性的不确定性.

\originalsection{概率、数密度与质量密度须统一归一化}
本卷用 $f(\boldsymbol x,\boldsymbol v,t)\dd^3x\dd^3v$ 表示该相空间体元中的恒星数. 因而 $\int f\dd^3x\dd^3v=N$, $n(\boldsymbol x)=\int f\dd^3v$. 原课先以随机抽星的概率密度归一化为1, 后又把其速度积分叫数密度; 若使用该概率约定, 应将积分乘 $N$ 才有相同物理数密度.

等质量恒星给 $\rho_*=mn$. 多质量种群需分别定义 $f_a,m_a$, $\rho=\sum_am_a\int f_a\dd^3v$ 再加气体与暗物质. 观测的发光恒星可只是引力势里的示踪种群, 它的 $mn$ 不一定等于生成全部势的 $\rho$.

定义平均速度 $u_i=n^{-1}\int v_if\dd^3v$ 和弥散张量
\[
\sigma_{ij}^2=\overline{v_iv_j}-u_iu_j
=\frac1n\int(v_i-u_i)(v_j-u_j)f\dd^3v.
\]
“零阶矩、一次矩、二次矩”按速度幂次命名, 分别给 $n,nu_i,n\overline{v_iv_j}$. 这样不会把密度误叫含一个速度的第一阶矩.

\originalsection{六维连续性给无碰撞 Boltzmann 方程}
相空间速度为 $(\dot{\boldsymbol x},\dot{\boldsymbol v})=(\boldsymbol v,-\nabla\Phi)$. 粒子数守恒给六维连续方程
\[
\partial_tf+\nabla_x\cdot(f\boldsymbol v)
+\nabla_v\cdot(-f\nabla_x\Phi)=0.
\]
因 $\boldsymbol v$ 不依赖 $\boldsymbol x$, 牛顿势力不依赖 $\boldsymbol v$, 相空间流散度为零, 得
\[
\partial_tf+\boldsymbol v\cdot\nabla_xf
-\nabla_x\Phi\cdot\nabla_vf=0.
\]
这就是沿轨道 $\dd f/\dd t=0$, 与前章相混合的精细守恒相符. 自洽问题还要同求 $\nabla^2\Phi=4\pi G\rho[f]$, 称作 Vlasov--Poisson 系统. 对固定背景中的示踪种群, 则 $\Phi$ 是给定输入. 两者的自由度与稳定性问题不同.

在定常势中, 若 $I_a$ 是轨道积分, 任意非负 $f=F(I_a)$ 都满足定常无碰撞方程, 因为沿轨道每个 $I_a$ 不变. 这给建立平衡模型的充分方法. 球对称非旋转系常取 $f(E,\ell^2)$, 各向同性可取 $f(E)$. 反过来把所有定常解都表示成某个简单有限积分集, 还需轨道的规则性与混合条件; 存在混沌区时不可无条件套最强形式的 Jeans 定理. 定常也不等于稳定, 下一章将显式引入扰动.

\originalsection{逐步取矩得到连续性与动量方程}
假设 $f$ 及 $v_if$ 在速度无穷远衰减足够快, 对无碰撞方程积分. 力项为速度空间全导数, 边界项零, 得
\[
\partial_tn+\partial_i(nu_i)=0.
\]
再乘 $v_j$ 积分. 分部积分给 $\int v_j\partial_{v_i}f\dd^3v=-\delta_{ij}n$, 故
\[
\partial_t(nu_j)+\partial_i[n\overline{v_iv_j}]+n\partial_j\Phi=0.
\]
代入 $\overline{v_iv_j}=u_iu_j+\sigma_{ij}^2$, 用连续方程消去乘 $u_j$ 的项,
\[
\partial_tu_j+u_i\partial_iu_j
=-\frac1n\partial_i(n\sigma_{ij}^2)-\partial_j\Phi.
\]
左侧是平均流的物质导数, 其中对流加速度不是黏性系数; 原笔记将它标成 viscosity/shear, 此处区分几何对流与碰撞耗散. $n\sigma_{ij}^2$ 是随机运动的应力, 可各向异性, 不要求有普通气体热平衡温度. 这些关系是精确的低阶矩方程, 但二次矩的演化会依赖三次矩; 仅有它们并没有闭合求出完整 $f$.

\originalsection{球对称稳态的径向平衡}
对定常、球对称且无平均流的示踪种群, 有 $u_r=u_\theta=u_\varphi=0$, 两个切向弥散相同, 记单一切向分量为 $\sigma_t^2$. 球坐标应力散度的径向部分给
\[
\frac{\dd(n\sigma_r^2)}{\dd r}
+\frac{2n(\sigma_r^2-\sigma_t^2)}r
=-n\frac{GM(<r)}{r^2}.
\]
定义 $\beta=1-\sigma_t^2/\sigma_r^2$（等价于 $1-(\sigma_\theta^2+\sigma_\varphi^2)/(2\sigma_r^2)$）, 得
\[
\frac{\dd(n\sigma_r^2)}{\dd r}+
\frac{2\beta n\sigma_r^2}r=-n\frac{GM(<r)}{r^2}.
\]
$\beta=0$ 等向, $\beta\to1$ 近径向, 大负 $\beta$ 为切向占优. $\sigma_r=0$ 的纯圆轨道极限不宜直接除以它来定义有限 $\beta$, 应回到原应力式取极限.

把导数写成对数斜率, 得质量推断式
\[
M(<r)=-\frac{r\sigma_r^2}{G}
\left[\frac{\dd\ln n}{\dd\ln r}
+\frac{\dd\ln\sigma_r^2}{\dd\ln r}+2\beta\right].
\]
它不是观测数值代入即唯一得到质量的定理, 因为真实三维 $n,\sigma_r,\beta$ 不都可直接测量. 且任意给定函数使此式得到 $M>0$ 只是必要一致性检查, 还未保证存在非负完整分布函数.

当 $\beta$ 固定, 可用积分因子 $r^{2\beta}$ 与外侧应力衰减边界得到
\[
n(r)\sigma_r^2(r)=r^{-2\beta}\int_r^\infty
n(s)GM(<s)s^{2\beta-2}\dd s.
\]
不同边界应力会增加齐次项, 因而截断系统须说明边界压力或轨道支持. 后面的薄盘题也要仔细处理锐利边缘, 不能默认为无任何边界效应.

\originalsection{投影为何产生质量—各向异性退化}
观测给投影半径 $R$ 的表面示踪密度 $\Sigma_n(R)$ 与视线弥散. 对球模型沿视线积分,
\[
\Sigma_n(R)=2\int_R^\infty\frac{n(r)r\dd r}{\sqrt{r^2-R^2}}.
\]
在空间点 $\sin\theta=R/r$, 视线速度二次矩为 $\sigma_r^2\cos^2\theta+\sigma_t^2\sin^2\theta=\sigma_r^2(1-\beta R^2/r^2)$. 因而
\[
\Sigma_n(R)\sigma_{\rm los}^2(R)=2\int_R^\infty
\left(1-\beta\frac{R^2}{r^2}\right)
\frac{n(r)\sigma_r^2(r)r\dd r}{\sqrt{r^2-R^2}}.
\]
同一投影速度数据可以由不同质量和轨道各向异性组合近似重现. 这就是退化的物理来源, 不是代数运算不够熟练. 自行、不同示踪种群、高阶速度分布或独立质量信息能提供额外约束, 但应分别验证模型假设.

\begin{example}
取一个半径区间内 $n\propto r^{-\gamma}$、常数 $\sigma_r,\beta$ 及平坦圆速度, 则 $v_c^2=(\gamma-2\beta)\sigma_r^2$. 若 $\gamma=3$, 同一三维径向弥散在等向假设下给 $v_c^2=3\sigma_r^2$, 在 $\beta=1/2$ 下给 $2\sigma_r^2$. 这个例子显示动力学推断依赖轨道分布; 它尚未要求两模型有完全相同的投影视线弥散, 后者要另用上式计算.
\end{example}

\originalsection{一个可检查的各向同性平衡模型}
令 $f=f_0\exp(-E/\sigma^2)$, $E=v^2/2+\Phi$. 它是轨道能量的函数, 故满足定常无碰撞方程. Gaussian 速度积分给 $n=f_0(2\pi\sigma^2)^{3/2}\exp(-\Phi/\sigma^2)$, 三个速度分量均有弥散 $\sigma^2$. 若其自身产生球势, 奇异等温球的一个形式解为
\[
\Phi=2\sigma^2\ln(r/r_0),\qquad
\rho=\frac{\sigma^2}{2\pi Gr^2},\qquad
v_c^2=2\sigma^2.
\]
由 $\nabla^2\ln r=1/r^2$ 和 $\rho\propto\exp(-\Phi/\sigma^2)$ 双向检查可知它自洽. 但中心奇异、无限总质量和无限延伸使它不是一个完整有限星系模型. 名称“等温”描述等向速度弥散常数, 不意味着无碰撞恒星靠两体热接触已经达到理想气体热平衡. 它把势、密度、分布函数和矩方程放在同一个算例里, 同时展示边界条件不能省略.

\originalchapter{Jeans 引力不稳定与旋转盘的稳定窗口}

本章对应 8.902 第7讲. 定常分布函数只说明一组轨道可以不变, 并未保证小扰动不会增长. 核心问题是随机运动或压力来得及抹平密度起伏吗, 旋转的回复作用又保护哪些波长? 先完整推导流体近似的色散关系, 再区分恒星盘的无碰撞结果. 原课以尺度估计引入 Toomre 条件, 其中的近似系数与严格动力学系数不能混为一谈.

\originalsection{均匀背景的近似及闭合假设}
对理想可压缩流体, 取 $\rho=\rho_0+\delta\rho$, 平均速度 $\boldsymbol u=\delta\boldsymbol u$, 线性压力扰动 $\delta P=c_s^2\delta\rho$. $c_s^2$ 是沿所采用热力学过程的导数 $\partial P/\partial\rho$, 绝热与等温值未必相同. 用流体的连续性与 Euler 方程,
\[
\partial_t\delta\rho+\rho_0\nabla\cdot\delta\boldsymbol u=0,
\qquad
\partial_t\delta\boldsymbol u=-\nabla\delta\Phi-
\frac{c_s^2}{\rho_0}\nabla\delta\rho,
\quad \nabla^2\delta\Phi=4\pi G\delta\rho.
\]
严格静态、无限均匀且非零密度的背景不能同时有 $\nabla\Phi_0=0$ 与完整 Poisson 方程. 上式只给扰动密度求引力, 将均匀背景力另行抵消, 是所谓 Jeans swindle. 它可以作为局部近似或适当周期模型, 不是已建立的无限孤立平衡宇宙. 后面在膨胀背景中写扰动方程时, 背景引力将通过 Friedmann 方程处理.

对无碰撞恒星, Jeans 矩方程中的应力尚无普适状态方程. 将其简化为同一等向 $\delta P=c_s^2\delta\rho$ 是额外流体闭合, 只能给速度弥散支持的尺度直觉; 完整动力学需解分布函数的线性响应, 还可能有相混合效应.

\originalsection{波长与增长率从色散关系读出}
对连续性再微分并用另外两式消去速度和势,
\[
\partial_t^2\delta\rho-c_s^2\nabla^2\delta\rho
-4\pi G\rho_0\delta\rho=0.
\]
代入 $\delta\rho\propto\exp[\ii(\boldsymbol k\cdot\boldsymbol x-\omega t)]$,
\[
\omega^2=c_s^2k^2-4\pi G\rho_0.
\]
高 $k$ 短波由压力支持, 低 $k$ 长波有增长模与衰减模. 令
\[
k_J=\frac{\sqrt{4\pi G\rho_0}}{c_s},\qquad
\lambda_J=\frac{2\pi}{k_J}=c_s\sqrt{\frac{\pi}{G\rho_0}},
\]
则 $\lambda>\lambda_J$ 时增长率 $s=\sqrt{4\pi G\rho_0-c_s^2k^2}$. 正好临界时 $\omega=0$, 通常解有常数项及线性时间项, 不能把它当作已有正频率稳定振动.

Jeans 质量的几何约定需说明. 本卷按直径 $\lambda_J$ 的球定义
\[
M_J=\frac{4\pi}{3}\rho_0\left(\frac{\lambda_J}{2}\right)^3
=\frac{\pi^{5/2}}6\frac{c_s^3}{G^{3/2}\rho_0^{1/2}}.
\]
原课把 $\lambda_J$ 当球半径写 $4\pi\rho_0\lambda_J^3/3$, 比本卷大8倍; 波长阈值相同, 质量只因定义不同. 比较两个文献的 $M_J$ 前须核对这个选择. 尺度上 $t_{\rm cross}\sim\lambda/c_s$, $t_{\rm ff}\sim1/\sqrt{G\rho_0}$, 长尺度在压力通信前发生塌缩, 这解释了上述符号, 不用时间尺度估计冒充所有 $\pi$ 因子的精确推导.

\originalsection{薄盘自引力的二维 Fourier 核}
在局部盘面取径向坐标 $x$, 对薄盘扰动 $\delta\Sigma\exp(\ii kx)$, 三维密度为 $\delta\Sigma\delta(z)\exp(\ii kx)$. 离盘面 $\delta\Phi\propto\exp(-|k||z|)$, Poisson 方程积分跨盘面给导数跳跃, 因而
\[
\delta\Phi_k(z)=-\frac{2\pi G}{|k|}\delta\Sigma_k\ee^{-|k||z|}.
\]
负号与引力吸引一致. 由此盘的引力项正比于 $|k|$, 与均匀三维介质的常数 $4\pi G\rho_0$ 不同. 这就是不能直接把三维 Jeans 长度换一个密度就当作严格盘判据的原因.

\originalsection{旋转气体盘的完整局部色散关系}
在导引中心旋转的局部参考系中, 设径向扰动速度为 $u$, 切向为 $v$, 只取轴对称、短波和线性扰动, 背景面密度 $\Sigma_0$ 近似常数. 含 Coriolis 与背景剪切的方程为
\begin{align*}
-\ii\omega\delta\Sigma+\ii k\Sigma_0u&=0,\\
-\ii\omega u-2\Omega v&=-\ii k\left(c_s^2\frac{\delta\Sigma}{\Sigma_0}+\delta\Phi\right),\\
-\ii\omega v+\frac{\kappa^2}{2\Omega}u&=0.
\end{align*}
先由第一式 $\delta\Sigma=k\Sigma_0u/\omega$, 第三式 $v=\kappa^2u/(2\ii\omega\Omega)$, 再代入径向式及薄盘势核, 得
\[
\omega^2=\kappa^2+c_s^2k^2-2\pi G\Sigma_0|k|.
\]
转动回复项 $\kappa^2$ 保护足够长的波, 压力项 $c_s^2k^2$ 保护短波, 中间尺度可能不稳定. 将 $k\ge0$ 时右端配方,
\[
\omega^2=c_s^2\left(k-\frac{\pi G\Sigma_0}{c_s^2}\right)^2
+\kappa^2-\frac{\pi^2G^2\Sigma_0^2}{c_s^2}.
\]
故全部局部轴对称波无指数增长的条件为
\[
Q_g\equiv\frac{c_s\kappa}{\pi G\Sigma_0}\ge1.
\]
若 $Q_g<1$, 不稳定波数区间由二次方程两根 $k_-<k<k_+$ 给出. 冷盘 $c_s=0$ 时只有低 $k$ 转动保护, 更短波仍不稳定. 随机运动与转动必须共同作用才可能封闭不稳定窗口.

\begin{center}
\begin{tikzpicture}
\begin{axis}[width=8.2cm,height=5cm,xlabel={$q=kc_s/\kappa$},ylabel={$\omega^2/\kappa^2$},xmin=0,xmax=3,ymin=-1,ymax=5,grid=major,legend style={font=\small}]
\addplot[structurecolor,thick,domain=0:3]{1+x^2-2*x/0.8};\addlegendentry{$Q_g=0.8$}
\addplot[black,dashed,domain=0:3]{1+x^2-2*x/1.2};\addlegendentry{$Q_g=1.2$}
\addplot[gray,domain=0:3]{0};
\end{axis}
\end{tikzpicture}
\end{center}
图按色散式直接绘制, 负的频率平方才是增长区, 不是频率为负. 有限厚度减弱小尺度引力响应; 多组分气体与恒星共同耦合, 又会改变单组分判据. $Q_g>1$ 只控制这里的局部轴对称线性模, 不保证没有非轴对称棒、螺旋或非线性演化.

\originalsection{恒星盘的系数来自分布函数响应}
无碰撞 Gaussian 径向速度分布的薄盘结果通常写为
\[
Q_*\equiv\frac{\sigma_R\kappa}{3.36G\Sigma_0}\ge1.
\]
这里以 Toomre 1964 原论文的临界结果作引用输入, 完整速度空间线性响应不由前面流体闭合证明. 在零频率边界, 其径向随机运动减弱因子含 $[1-\ee^{-\chi}I_0(\chi)]/\chi$, $\chi=(k\sigma_R/\kappa)^2$, $I_0$ 为修正 Bessel 函数. 临界常数可独立重算为
\[
C=2\pi\max_{q>0}\frac{1-\ee^{-q^2}I_0(q^2)}q
\simeq3.35828,
\quad q_{\max}\simeq0.97370.
\]
课程尺度估计 $\sqrt{32/3}\simeq3.266$ 接近这个数, 但其后将3.26写成 Toomre 精确计算的系数并不准确. 本卷区分近似估算、流体的 $\pi$ 与恒星 Gaussian 响应的3.36. 引用原论文只用于此事实核对及常数复算, 不收录论文原件或受限图; 详见来源补充清单.

从这一章回看前面两章, $\kappa$ 来自单星近圆轨道的回复作用, $\sigma_R$ 来自系综的随机运动, $\Sigma_0$ 决定集体自引力. 三者在稳定性中才相遇, 不能只从维里定理知道速度尺度就自动断言稳定. 后面的恒星形成需进一步加入冷却与反馈, 稳定性阈值只是其中的一环.

\originalchapter{恒星种群、形成史与化学守恒}

本章对应 8.902 第8讲和第9讲开头的 G 矮星问题. 前面把恒星视作有质量的轨道示踪体, 现在恢复它们的出生、寿命、光谱和元素回馈. 光谱合成把形成史映入观测, 化学演化则用质量守恒约束形成史; 两种信息可以互补, 却都需要初始质量函数、混合和流入流出的假设.

\originalsection{形成率的指示量不是无条件换算表}
恒星形成率 $\psi(t)$ 是单位时间新形成恒星的总初始质量, 常用 $\Ms\,\mathrm{yr^{-1}}$. 年轻大质量恒星产生较强紫外与电离光子, 后者使 H II 区产生 H$\alpha$; 尘埃吸收部分光再辐射于远红外. 因而这些光度能作为形成率指示量, 但响应的时间窗、尘埃和 IMF 会影响标定.

课程采用的数量级关系为
\begin{align*}
\frac{\psi_{\rm FIR}}{\Ms\,\mathrm{yr^{-1}}}&\sim\frac{L_{\rm FIR}}{5.8\times10^9\Ls},\\
\frac{\psi_{\rm H\alpha}}{\Ms\,\mathrm{yr^{-1}}}&\sim\frac{L_{\rm H\alpha}}{1.3\times10^{41}\,\mathrm{erg\,s^{-1}}},\\
\frac{\psi_{\rm UV}}{\Ms\,\mathrm{yr^{-1}}}&\sim\frac{L_{\nu,\rm UV}}{7.2\times10^{27}\,\mathrm{erg\,s^{-1}Hz^{-1}}}.
\end{align*}
UV 项须使用相应的频率光度密度; 原笔记只写 $\mathrm{erg/s}$ 会使量纲与该标定不符. 这些数是课程在给定模型背景下的近似例子, 不是所有金属丰度、突发形成史和 IMF 通用的常数. 老恒星或 AGN 也能加热尘埃, 吸收又可隐藏紫外和 H$\alpha$, 应比较不同指示量而非任取一个光度直接宣称唯一形成率.

冷而密的气体可在自由落体时间内聚集, 常用局部尺度模型 $\psi\sim\epsilon_{\rm ff}M_g/t_{\rm ff}$, $t_{\rm ff}\propto(G\rho)^{-1/2}$. 系数 $\epsilon_{\rm ff}$ 包含湍流、磁场和反馈等未由这一式求出的物理. 分子相比例还依赖屏蔽、金属丰度和团块结构. 原课受限论文的气体相图和观测散点图不重新分发; 这里解释模型依赖, 不把相图百分数当作普适宇宙常数.

\originalsection{IMF 归一化决定每个积分的单位}
用物理质量 $m$ 定义 $\phi(m)$, 使总初始质量 $M_{\rm new}$ 的一次形成产生
\[
\dd N=M_{\rm new}\phi(m)\dd m,\qquad
\int_{m_l}^{m_h}m\phi(m)\dd m=1.
\]
于是 $[\phi]=M^{-2}$, 右侧归一化是无量纲1. 若改以 $x=m/\Ms$ 表示, 可定义无量纲 $\varphi(x)$ 使 $\int x\varphi(x)\dd x=1$, 则 $\dd N=(M_{\rm new}/\Ms)\varphi(x)\dd x$. 原课把这两种单位约定交叉书写, 本卷统一使用前一种.

Salpeter 形式 $\phi=Cm^{-2.35}$ 在选定 $m_l,m_h$ 上由上述质量积分定 $C$. 同一幂律下数目与质量贡献分别按 $m^{-2.35}\dd m$ 和 $m^{-1.35}\dd m$ 加权. 低质量恒星多且长寿, 高质量恒星少却可能主导紫外光, 因而按光度选样会偏向后者. 课程的 $m_l\sim0.08\Ms$ 对应氢燃烧门槛尺度, $m_h\sim100\Ms$ 是模型截断; 不能只由 $L_{\rm Edd}\propto M$ 宣称100太阳质量是所有恒星绝不能超过的硬上限.

\originalsection{寿命与合成光谱的双重积分}
恒星的有效温度按 $L=4\pi R^2\sigma_{\rm SB}T_{\rm eff}^4$ 定义, 不要求真实频谱在每个频率上都是理想黑体. HR 图通常横轴温度向右降低, 纵轴光度增加. 固定半径线满足 $\log L=4\log T+\mathrm{const}$, 但主序的物理定义是核心氢燃烧, 各主序恒星半径不同. 原课将 $R=1$ 线直接叫主序, 此处纠正这一概念.

假定可燃料量与质量近似成正比, 则主序寿命 $\tau_{\rm MS}\propto M/L$. 在采用 $L\propto M^p$ 的质量区间, 有 $\tau_{\rm MS}\propto M^{1-p}$. 课程以 $p\sim4$ 得 $M^{-3}$, 低质量例子取 $p\sim3$ 得 $M^{-2}$. 这些是区间近似, 不可外推整个质量谱. 主序转折质量由 $\tau_{
\rm MS}(m_t)\simeq$ 种群年龄给出, 为老种群年龄提供线索.

设单星频率光度为 $\ell_\nu(m,Z,\tau)$, 对同年龄同初始金属丰度的种群, 每单位初始形成质量的光谱为
\[
\ell_\nu^{\rm SSP}(\tau,Z)=\int\phi(m)\ell_\nu(m,Z,\tau)\dd m.
\]
单星演化和残骸阶段含在 $\ell_\nu$ 内, 再叠加出生时刻,
\[
L_\nu(t)=\int_0^t\psi(t')\ell_\nu^{\rm SSP}(t-t',Z(t'))\dd t'.
\]
先选 IMF 和形成史, 再生成光谱、通过尘埃与滤镜映入数据, 比较后调整, 才是合成的工作流. 年龄、金属丰度、尘埃和 IMF 的效应可退化, 线特征与动力学质光比能提供独立限制. 本卷没有假称已运行某个完整种群合成软件或拟合真实星系全谱.

\originalsection{元素丰度与四条质量守恒式}
金属质量分数 $Z=M_{\rm heavy}/M$ 与元素数目比的对数
\[
[X/Y]=\log_{10}\frac{(n_X/n_Y)}{(n_X/n_Y)_\odot}
\]
不同. $[\mathrm{Fe/H}]=-1$ 表示铁氢数目比为相应太阳标定的十分之一, 不直接意味着所有重元素总质量分数都等于太阳的十分之一. 课程的 $Z_\odot\simeq0.02$ 是所用近似标度; 不同丰度标定需显式说明.

令气体质量 $M_g$, 恒星及残骸中锁定质量 $M_*$, 总重子质量 $M=M_g+M_*$. 形成率为 $\psi$, 恒星返还气体率为 $E$, 返还金属率为 $E_Z$, 流入率为 $F$、金属分数 $Z_F$, 流出率为 $O$、金属分数 $Z_O$. 均匀混合模型有
\begin{align*}
\dot M&=F-O, &\dot M_*&=\psi-E,\\
\dot M_g&=-\psi+E+F-O,
&\frac{\dd(ZM_g)}{\dd t}&=-Z\psi+E_Z+Z_FF-Z_OO.
\end{align*}
其中恒星形成先锁定原有金属, 返还包含原有金属和新核合成金属. 不能只保留新金属而漏掉原有返还部分. $Z_O$ 可不同于平均 $Z$, 选择性金属风与完全混合风的结果不同.

若质量 $m$ 的星在 $\tau_m$ 后返还 $m-w_m$, 则
\[
E(t)=\int(m-w_m)\psi(t-\tau_m)\phi(m)\dd m.
\]
积分只计当时已演化的恒星, 出生前取 $\psi=0$. 相应金属率可写
\[
E_Z(t)=\int\![(m-w_m)Z(t-\tau_m)+p_Z(m)]
\psi(t-\tau_m)\phi(m)\dd m,
\]
其中 $p_Z(m)$ 是新产生并释放的金属质量, 单个元素的净产生量也可能为负. 延迟时间使金属史并非仅由当前 $\psi(t)$ 决定.

\originalsection{瞬时回收近似与混合流出的抵消}
对寿命远短于气体演化时间的返还星, 近似 $\psi(t-\tau_m)\simeq\psi(t)$ 并立即混合. 定义返还比例和按锁定质量计的净产额
\[
\mathcal R=\int_{m_{\rm lim}}^{m_h}(m-w_m)\phi(m)\dd m,
\qquad y=\frac1{1-\mathcal R}\int_{m_{\rm lim}}^{m_h}p_Z(m)\phi(m)\dd m.
\]
于是 $E=\mathcal R\psi$, $E_Z=\mathcal R Z\psi+(1-\mathcal R)y\psi$. 注意锁定质量包括长寿恒星和残骸, 不只是残骸. 用 $\dd(ZM_g)/\dd t=M_g\dot Z+Z\dot M_g$, 两条守恒式相减得到
\[
M_g\dot Z=(1-\mathcal R)y\psi+(Z_F-Z)F-(Z_O-Z)O.
\]
若流出完全混合, $Z_O=Z$, 流出项在金属分数的瞬时导数中抵消; 它仍通过减少 $M_g$ 改变演化. 原讲义相应公式保留 $+Ze$ 项, 与前面的混合流出守恒式不一致. 不能因为符号展开失误把流出误认为直接提高金属分数.

瞬时回收对低质量长寿星和 Ia 型超新星延迟等不适用, 不应使用该近似解释所有元素比的时钟. 后面解析闭盒结果先限制在它的适用模型中.

\originalsection{闭盒解与恒星的平均金属丰度}
取 $F=O=0$, 初始全为无金属气体 $M_g(0)=M$, 则 $M$ 恒定,
\[
\dot M_g=-(1-\mathcal R)\psi,
\qquad M_g\dot Z=(1-\mathcal R)y\psi.
\]
消去时间得到 $\dd Z=-y\dd M_g/M_g$, 故以气体分数 $u=M_g/M$ 表示,
\[
Z=y\ln(1/u),\qquad M_*=(1-u)M.
\]
形成率决定到达某个 $u$ 需要多久, 却不改变这一 $Z(u)$ 的闭盒关系. 所有产生的金属在气体和锁定恒星中, 总金属质量为 $yM_*$. 因而恒星质量加权平均
\[
\langle Z_*\rangle M_*+ZM_g=yM_*,\qquad
\langle Z_*\rangle=y-Z\frac{u}{1-u}.
\]
当 $u\to0$ 平均趋 $y$, 而气体 $Z$ 形式上发散. 这只是极少量剩余气体中积累金属的模型极限, 到金属分数不再小或产额假设失效时不能继续物理外推. 当 $u\to1$ 时两项相消给初始平均接近零, 不是负金属丰度.

\originalsection{G 矮星问题约束的不只是一个平均值}
长寿恒星保留出生时的气体丰度. 闭盒单调增金属的模型中, 丰度不超过 $Z$ 时已锁定的质量为 $M[1-\exp(-Z/y)]$. 若现在气体分数为 $u_0$, 现有气体丰度 $Z_0=y\ln(1/u_0)$, 则
\[
\frac{M_*(\le Z)}{M_{*,0}}=
\frac{1-\exp(-Z/y)}{1-u_0},\qquad 0\le Z\le Z_0.
\]
微分分布正比 $\exp(-Z/y)$, 一直预测相当份额的低金属长寿星. 在恒定 IMF、样本寿命选择等条件下, 才能把这个质量比例映成某类 G 矮星的数目比例. 课程以实际低金属星较少的情况指出闭盒模型缺口, 候选修正包括流入、预增丰度与形成史假设; 不能仅靠平均 $\langle Z_*\rangle$ 替代整个分布检验.

作为可检验的模型扩展, 若完全混合流出 $O=\lambda\psi$, 无流入, 常 $\lambda$, 得
\[
Z=y_{\rm eff}\ln\frac{M_{g,0}}{M_g},\qquad
y_{\rm eff}=y\frac{1-\mathcal R}{1-\mathcal R+\lambda}.
\]
这说明混合流出虽在 $\dot Z$ 显式项中抵消, 却降低按剩余气体比推断的有效产额. 若有无金属流入, 则 $-ZF$ 直接稀释, 又是另一机制. 这些有限模型从守恒式推导, 不是假称已唯一拟合银河化学演化. 下一章的 AGN 将增加能量与反馈信息, 与本章质量、金属预算相联系.

\originalchapter{活动星系核的尺度、能源与吸积}

本章对应 8.902 第9讲的 AGN 部分. 活动星系核是星系中心显著的非恒星辐射源, 并非要求它在每个波段都压倒整个宿主. 要判断中央能源, 可以同时利用亮度、变化时间、谱线速度和辐射机制. 这条观测证据链比单凭一张有喷流的图片更具体; 本卷不复制原讲义中带许可例外的照片.

\originalsection{分类反映辐射与观测方向}
射电星系的喷流可把能量输送到远处射电瓣, 相对论电子的同步辐射形成明显非热谱. 类星体在光学呈明亮且紧致的核; 并非所有光学类星体都有很强射电喷流. 历史术语 quasar 与 QSO 的射电选择背景不同, 当代广义分类不能只按名称断言有无喷流. Seyfert 核通常在可见宿主中较易辨认, 核光度比最亮类星体低. BL Lac 等 blazar 的喷流接近视线, 相对论束射使连续谱增强、变化快且可有较高偏振.

课程所列类星体总光度约 $10^{45}$--$10^{48}\,\mathrm{erg\,s^{-1}}$, Seyfert 核约 $10^{42}$--$10^{45}\,\mathrm{erg\,s^{-1}}$, 是样本数量级而非严格分类边界. 原课射电项将 $10^{18}$ 的阈值标作 $\Ls$, 与随后的总能源尺度难以自洽; 该单位未能核定, 本卷不以它定义射电强弱或计算黑洞质量. 观测射电强度需说明频率、光谱光度密度和选择阈值, 不能擅自把原数改配一个猜测单位.

\originalsection{变化时间限制发光区大小}
若区域的相干变化受内部信息传播控制且无强束射, 本征时间窗 $\Delta t$ 给发光区尺度 $R\lesssim c\Delta t$. 对红移源, $\Delta t_{\rm obs}=(1+z)\Delta t$. 若喷流有 Doppler 因子 $\mathcal D$, 则本征区尺度上界近似
\[
R'\lesssim\frac{c\Delta t_{\rm obs}\mathcal D}{1+z},\qquad
\mathcal D=[\Gamma(1-\beta\cos i)]^{-1}.
\]
因而不能把所有快变对象都按静止源的 $c\Delta t_{\rm obs}$ 估大小. 无红移与束射时, 一天光程约 $8.4\times10^{-4}\,\mathrm{pc}$, 与课程给出的日尺度核尺寸相符.

黑洞 Schwarzschild 半径为 $r_s=2GM_\bullet/c^2\simeq9.57\times10^{-14}(M_\bullet/\Ms)\,\mathrm{pc}$. 因此百万太阳质量给 $r_s\sim10^{-7}\,\mathrm{pc}$, 十亿太阳质量给 $10^{-4}\,\mathrm{pc}$. 快变区可以在黑洞附近, 但大小上界本身既不等于视界半径, 也未唯一决定质量. 宽线气体的速度和响应时间可再提供质量尺度 $M_\bullet\sim f_{\rm geom}R_{\rm BLR}\Delta v^2/G$, 其中 $R_{\rm BLR}\sim c\tau_{\rm rest}$, 几何和投影因子需另标定.

高亮度集中在小区域, 使单靠普通恒星密集堆积或频繁超新星供能的模型面临质量、寿命和密度约束. 例如 $L=10^{46}\,\mathrm{erg\,s^{-1}}$ 一年释放 $3.16\times10^{53}\,\mathrm{erg}$, 若每次事件可辐射 $10^{51}\,\mathrm{erg}$, 要数百次每年且集中在极小核区. 这是数量级排查, 未把所有瞬变核都当作被同一单数值严格排除的模型. 对持续高功率活动核, 引力吸积提供更有效的能源通道.

\originalsection{吸积效率与薄盘加热}
气体落入深引力势可把部分静质量能转成辐射, 写 $L=\eta\dot M_{
\rm in}c^2$. $\eta$ 是辐射效率, 随吸积模式、内边界、自旋和能量损失改变, 不是所有黑洞同一常数. 课程取 $\eta\simeq1/16$ 示范, 与氢燃烧的约0.007静质量转化率比较; 原始核燃料中实际可燃烧比例又是不同因素.

气体有角动量时不会直接径向落下. 在薄、近 Kepler、稳态且内缘零力矩模型中, 角动量向外输运使气体向内流. 以运动黏性系数 $\nu_{\rm visc}$ 表示应力, 角动量守恒给
\[
\nu_{\rm visc}\Sigma=\frac{\dot M_{
\rm in}}{3\pi}\left[1-\sqrt{\frac{r_{
\rm in}}r}\right].
\]
每面局部耗散通量为 $F=\tfrac12\nu_{\rm visc}\Sigma(r\Omega')^2$. 用 $\Omega^2=GM_\bullet/r^3$, 得
\[
F(r)=\frac{3GM_\bullet\dot M_{
\rm in}}{8\pi r^3}
\left[1-\sqrt{\frac{r_{
\rm in}}r}\right].
\]
这说明热量来自剪切应力与引力能源的共同预算. 若光厚且局部辐射有效, $\sigma_{\rm SB}T_{
\rm eff}^4=F$, 远离内缘 $T_{
\rm eff}\propto r^{-3/4}$. 不同半径的热谱叠加形成多温度盘; X 射线可来自更热的日冕等额外成分, 不应把课程给出的 $10^8$ K 无条件当作所有超大质量黑洞薄盘的有效温度.

双面面积积分给牛顿薄盘 $L=\int_{r_{
\rm in}}^\infty4\pi rF(r)\dd r=GM_\bullet\dot M_{
\rm in}/(2r_{
\rm in})$. 该近似靠近视界时失效; Schwarzschild 的相对论薄盘效率约0.057, 来自 ISCO 结合能, 见已有 GR 卷前置结果. 本章不将 Newton 近似在 $r_{
\rm in}=6GM/c^2$ 的代入值冒充完整 GR 计算, 也未从一个有效黏性系数求出真实磁湍流机制.

\originalsection{Eddington 光度的假设与增长时间}
对近球对称辐射通量 $L/(4\pi r^2)$, 单位质量辐射力 $a_{
\rm rad}=\kappa_F L/(4\pi r^2c)$, 其中 $\kappa_F$ 是按实际光谱加权的不透明度. 与引力 $GM_\bullet/r^2$ 比较, 得
\[
L_{
\rm Edd}=\frac{4\pi GM_\bullet c}{\kappa_F}.
\]
完全电离氢且以 Thomson 散射为主时 $\kappa_F=\sigma_T/m_p$, 给 $L_{
\rm Edd}\simeq1.26\times10^{38}(M_\bullet/\Ms)\,\mathrm{erg\,s^{-1}}$. 电子受辐射力而通过等离子体耦合拖动质子, 所以配的是每质子质量, 不是直接把 $m_e$ 作为被引力束缚的全部质量.

称为“上限”时需保留几何、稳态、光子俘获和不透明度条件. 非球对称盘、束射或超临界流可使表观光度或质量流量与这份简单平衡不同, 黑洞也不会因辐射压力就自身被吹散. 定义给定效率的参考流入率 $\dot M_{
\rm Edd}=L_{
\rm Edd}/(\eta c^2)$,
\[
\dot M_{
\rm Edd}\simeq
2.22\times10^{-8}\left(\frac{0.1}{\eta}\right)
\frac{M_\bullet}{\Ms}\,\Ms\,\mathrm{yr^{-1}}.
\]
若 $\eta=1/16$, 系数约 $3.55\times10^{-8}$, 与原课的 $5\times10^{-10}$ 不相符; 此处从光度、$c^2$ 与秒到年的单位换算重新确定, 不保留错误系数.

若辐射拿走质量能的比例 $\eta$, 黑洞增长率为 $\dot M_\bullet=(1-\eta)\dot M_{
\rm in}$. 固定 Eddington 比 $\lambda=L/L_{
\rm Edd}$ 下, 质量指数增长, e 折叠时间
\[
t_{
\rm grow}=\frac{\eta}{(1-\eta)\lambda}
\frac{\kappa_Fc}{4\pi G}.
\]
对上述 Thomson 氢标定, 后因子约 $4.5\times10^8$ 年. $\eta=0.1$, $\lambda=1$ 给约五千万年. 这个模型比较高红移大黑洞形成所需时间时, 还应说明种子质量、间歇活动和效率史, 不将一次指数积分当作已唯一解释所有早期黑洞.

\originalsection{同步辐射与统一结构的用途}
Gaussian 单位下, 在磁场 $B$ 中运动的相对论电子, 典型同步频率为
\[
\nu_{
\rm syn}\sim\frac{3eB\gamma^2\sin\alpha}{4\pi m_ec},
\]
其中 $\alpha$ 为俯仰角, $\gamma$ 为电子 Lorentz 因子. 能量越高, 辐射频率越高, 冷却也更快. 对等向俯仰角平均, 同步损失功率 $P\simeq(4/3)\sigma_TcU_B\gamma^2\beta^2$, $U_B=B^2/(8\pi)$, 给 $t_{
\rm cool}\sim\gamma m_ec^2/P$. 非热谱、偏振和喷流方向可以约束这些量, 但仅由某个频率不能单独求出 $B$ 与 $\gamma$, 因为它们以乘积进入.

AGN 的简化统一图包含中央黑洞与盘、近核宽线区、较远窄线区、尘埃遮挡结构及可能的喷流. 宽线区轨道速度较高, 视线速度分布使谱线更宽; 窄线区较慢, 名称并非单纯由气体“厚”或“薄”决定. 尘埃吸收光并辐射红外, 遮挡方向使某些核的宽线难见. 喷流朝视线时束射增强, 对应 blazar 的一类表现.

\begin{center}
\begin{tikzpicture}[>=Latex]
\fill(0,0)circle(2pt)node[below]{黑洞};
\draw[structurecolor,thick](-0.6,0.08)--(0.6,0.08)node[right]{盘};
\draw[dashed](-1,-0.4)ellipse[x radius=.3,y radius=.5];
\draw[dashed](1,-0.4)ellipse[x radius=.3,y radius=.5];
\node at(1.75,-0.65){尘埃};
\draw[->,thick](0,.2)--(0,2)node[above]{喷流};
\draw[->](2.6,1.6)--(.25,.35)node[midway,above]{观测方向};
\node at(-1.4,1.1){宽线区};\node at(2.4,.65){窄线区};
\end{tikzpicture}
\end{center}
图只表示组成与方向, 不按长度比例绘制, 不替代某个真实活动核的几何测量. 统一模型解释许多方向相关的分类, 仍须允许不同吸积率、自旋和环境. AGN 的风与喷流还能向宿主输入能量, 改变前章气体流入流出与恒星形成; 反馈效能需独立能量预算, 不是仅看到黑洞就自动确定每个星系的化学史.

\originalchapter{从谱线位移到均匀宇宙的初始值问题}
本章对应8.286第3--8讲的观测和牛顿动力学入口，补充后面的相对论背景章。课程从可测量的频率开始，再问同一种膨胀律如何被不同位置的观测者共同描述。狭义相对论的一般原理见广义相对论卷；这里保留理解谱线和宇宙学坐标必需的推导。距离、速度和时间的定义必须先统一，才有可能比较数据和模型。

\originalsection{声波的Doppler效应与光的Doppler效应}
先在静止介质中考虑一维声波。波速为$u$，源周期为$T_s$；源向观察者运动，速率$v_s$。两次发射之间源前进$v_sT_s$，前一波峰却前进$uT_s$，所以前方波长为$\lambda=(u-v_s)T_s$。若观察者也向源运动，速率$v_o$，其遇到波峰的速率为
\[
 f_o=f_s\frac{u+v_o}{u-v_s}.
\]
这一区分源和观察者的公式依赖\emph{介质静止系}。使源远离观察者取$v_s<0$，观察者远离源取$v_o<0$。它不能直接令$u=c$后用于真空光波，因为光没有这一特殊介质系。

对相距渐增的光源与观察者，取观察者惯性系，光源退行速率$v=\beta c$，源的固有发射周期为$T_s$。该系里发射间隔是$\gamma T_s$，其中$\gamma=(1-\beta^2)^{-1/2}$；第二个波峰还需多走$v\gamma T_s$。故接收间隔为
\[
 T_o=\gamma(1+\beta)T_s,\qquad
 1+z=\frac{T_o}{T_s}=\sqrt{\frac{1+\beta}{1-\beta}}.
\]
这既含传播路程差，也含钟速差。反解$\beta=[(1+z)^2-1]/[(1+z)^2+1]$，小$z$时才有$v\simeq cz$。一般方向最清楚的定义是局部四动量测频：$E_{\rm obs}=-p_\mu u^\mu$，$\nu=E_{\rm obs}/h$，这样方向和参考系不会混淆。对于观测系中速度与光传播方向夹角$\theta$，有$\nu_o=\nu_s/[\gamma(1-\beta\cos\theta)]$。横向$\cos\theta=0$仍有红移$1+z=\gamma$。

两事件若在物体自己的静止系同地发生，Lorentz变换给$\Delta t=\gamma\Delta\tau$。测运动尺长却需在\emph{测量者系同时}取两端，才得$L=L_0/\gamma$；两个尺端在自身静止系同时的读数，变换后一般不同时。课程第3讲的分步波峰图正是在提醒这一区别。宇宙学红移来自沿途不同局部共动系的累计传播，不能把远处退行率$HD$无条件当作上述一个惯性系的$v$。

\originalsection{角度、视差与距离梯级}
周年视差$p$是地球轨道一AU基线在天体处张开的半角。小角关系$D=\mathrm{AU}/p$中$p$须用弧度；$p=1$角秒定义一秒差距，因此
\[
 D(\mathrm{pc})=\frac1{p(\mathrm{arcsec})},\qquad
 1\ \mathrm{pc}=206265\ \mathrm{AU}\simeq3.086\times10^{16}\ \mathrm m.
\]
这是几何测距，不需先假设恒星光度。再用已知距离的恒星校准光度或周期关系，才能推广到造父变星和超新星。若一个校准层的零点有偏移，它会传递到更远的层；单一光通量本身不能区分远而亮、近而暗。

谱线红移通过静止实验室波长比较得$z=(\lambda_o-\lambda_s)/\lambda_s$。近邻星系的红移与独立距离呈近线性关系，这提出均匀膨胀的模型。8.286第4讲分别讨论Lemaître于1927年的理论联系与Hubble于1929年的观测图；本卷保留这些年代，原历史照片和曲线不复制。现代拟合还需扣除本动速度；例如300 km/s的本动速度在1000 km/s退行速度上占30\%，在30000 km/s上占1\%。这说明近邻更利于几何校准，却不一定最适合单独估宇宙膨胀率。

\originalsection{均匀膨胀为什么没有特殊中心}
在平直近似下，共动网格位置记为$\mathbf x$，物理位置$\mathbf r=a(t)\mathbf x$。从原点O换到共动位置$\mathbf x_P$，相对坐标是$\mathbf r'=a(\mathbf x-\mathbf x_P)$，于是$\dot{\mathbf r}'=H\mathbf r'$。每个共动观察者都看到同样的Hubble律。气球表面的点是二维类比：表面上没有中心点，而球心属于额外的嵌入空间，并不对应表面上的观测者。三维宇宙无需假定真实存在一个外部空间。

尺度因子是网格间距的共同倍数；固定网格只是描述平均流动，不保证真实星系严格静止于网格。原课程Guth用$a$的单位“米／格”，$r$的单位“格”，曲率参数$k$的单位“格$^{-2}$”。本卷使用无量纲$a$和有长度的$\chi$。令本卷的$a=a_G/a_{G0}$、$\chi=a_{G0}r_G$，就有$K=k_G/a_{G0}^2$；两种写法描述同一几何。不能把其中一种的单位要求误报成另一种的错误。

\originalsection{无限均匀质量的引力为何需要边界条件}
牛顿理论对有限分布可由势积分定义引力。无限均匀静态分布则有边界问题：若将越来越大的球总取以P为中心，壳层定理会给P处$\mathbf g=0$；若总以Q为中心，P位于均匀球内部，得到$\mathbf g=-4\pi G\rho(\mathbf r_P-\mathbf r_Q)/3$。不同极限选择了不同的加速度零点。没有声明极限和边界条件，不能由“各方向相同”断言所有位置的引力同时为零。

Poisson方程给$\nabla\cdot\mathbf g=-4\pi G\rho$，处处为零的场在$\rho>0$时不满足该式。均匀模型在有限邻域内的一个解是$\mathbf g=-4\pi G\rho\mathbf r/3$，加上任意随时间变化的均匀加速度仍不改变散度。两个质点的\emph{相对}加速度
\[
 \ddot{\mathbf r}_{12}=-\frac{4\pi G\rho}{3}\mathbf r_{12}
\]
却不依赖这一共同加速度。第5--6讲用Newton与Bentley通信引出的难点，应理解为无限静态极限与动力学边界条件的问题；物理上可测的相对运动是均匀宇宙的可靠入口。

\originalsection{有限尘埃球怎样给出可求解的模型}
取初始半径内密度$\rho_i$均匀、初始流速$H_i\mathbf r_i$，忽略压强与外界潮汐。壳层不交叉时，$r(t)=a(t)r_i$，取$a_i=1$；包围质量$M=4\pi\rho_i r_i^3/3$恒定，外壳不影响其牛顿引力。故
\[
 \ddot a=-\frac{4\pi G\rho_i}{3a^2},\qquad
 \frac{\dot a^2}{2}-\frac{4\pi G\rho_i}{3a}=E_a.
\]
积分常数$E_a=H_i^2/2-4\pi G\rho_i/3$有$\mathrm{s^{-2}}$量纲。$E_a<0$存在最大尺度，$E_a=0$给$a\propto t^{2/3}$，$E_a>0$在晚期近线性膨胀。这些结论限定为无压强、无宇宙学常数的模型；加入真空能后空间正曲率不再保证再坍缩。

牛顿球是初始值模型和相对运动的解释工具。它没有推出整个宇宙的空间曲率、光线传播或压力引力作用。下面先建立常曲率的实际距离，再在相对论背景章明确这些缺口如何由度规和Einstein方程补上。

\originalchapter{常曲率空间的距离、体积与整体范围}
本章对应8.286第10--13讲的非欧几何部分。原幻灯片中部分介绍页署名Mustafa Amin，部分历史图片和科研曲线有逐件例外；本卷只重写数学论证和自绘几何。微积分、点积和球坐标足以理解这里的构造。更一般的曲率张量与坐标变换见广义相对论卷，本章专注宇宙测距实际用到的空间结构。

\originalsection{局部长度怎样定义曲面上的几何}
地球两点之间的最短航程沿球面大圆，而不是穿过地心的直弦。描述内在几何，需要规定相邻坐标点之间的实际长度
\[
 \dd\ell^2=g_{ij}(x)\dd x^i\dd x^j.
\]
二维正定条件是$g_{xx}>0$、$g_{xx}g_{yy}-g_{xy}^2>0$。它确保非零位移长度为正。任一点附近可把这个二次型化成欧氏形式；这只说明\emph{局部}可以选正交标架，不说明一个有限区域内处处都能用同一坐标消掉曲率。

半径$R$的二维球面参数为$(x,y,z)=R(\sin\theta\cos\varphi,\sin\theta\sin\varphi,\cos\theta)$。对$\theta,\varphi$求偏导，两向量互相正交，模长分别$R,R\sin\theta$，所以
\[
 \dd\ell^2=R^2(\dd\theta^2+\sin^2\theta\dd\varphi^2).
\]
经线与赤道构成的球面三角形可有三个直角，角和$3\pi/2$；这不违反平面几何，因为边是球面测地线。恒定Gaussian曲率$1/R^2$下，测地三角形面积$A$满足$A/R^2=\alpha+\beta+\gamma-\pi$。该式在此作为Gauss--Bonnet定理的球面特例，完整一般定理的证明不在本章范围。

\begin{center}\begin{tikzpicture}[scale=1.05]
\draw[structurecolor,thick] (0,0) circle (1.45);
\draw[structurecolor] (-1.45,0) arc[start angle=180,end angle=360,x radius=1.45,y radius=.45];
\draw[dashed] (-1.45,0) arc[start angle=180,end angle=0,x radius=1.45,y radius=.45];
\draw[thick] (0,1.45) to[out=255,in=65] (-1.45,0);
\draw[thick] (0,1.45) to[out=280,in=110] (0,-.45);
\draw[thick] (-1.45,0) arc[start angle=180,end angle=270,x radius=1.45,y radius=.45];
\node[above] at (0,1.45){北极};\node[left] at (-1.45,0){A};\node[below] at (0,-.45){B};
\node[align=left] at (4,0){球面上的大圆边\\角和可大于$\pi$};
\end{tikzpicture}\end{center}
图为自绘透视示意，A、B沿赤道相差$\pi/2$；投影纸面角度不是球面内在角度。

\originalsection{三维球面并非普通球的内部}
正曲率三维空间可辅助嵌入四维欧氏空间，定义
\[
 X_1^2+X_2^2+X_3^2+X_4^2=R^2,\quad
 (X_1,X_2,X_3)=R\sin\psi\,\mathbf n(\theta,\varphi),\quad X_4=R\cos\psi.
\]
这是三维\emph{表面}$S^3$，不是三维实心球。直接微分，$\mathbf n\cdot\dd\mathbf n=0$消掉混合项，得
\[
 \dd\ell^2=R^2[\dd\psi^2+\sin^2\psi(\dd\theta^2+\sin^2\theta\dd\varphi^2)],\quad0\leq\psi\leq\pi.
\]
四维旋转保持长度并可把任何一点送到极点，因此所有点等价：极点只是坐标选择。嵌入只是验证几何的手段，物理宇宙不必有一个可观测的第四空间方向。

设径向弧长$\chi=R\psi$，距离$\chi$的二球面积为$4\pi R^2\sin^2(\chi/R)$。从极点到反极点距离$\pi R$，过赤道之后二球面积反而减小。整个空间体积为
\[
 V_{S^3}=4\pi R^3\int_0^\pi\sin^2\psi\dd\psi=2\pi^2R^3.
\]
若用$r=R\sin\psi$，可写$\dd\ell^2=\dd r^2/(1-r^2/R^2)+r^2\dd\Omega^2$，但同一个$r$对应赤道两侧的两个$\psi$。$r=R$处分母为零是坐标图失效，非几何奇点；需要第二张图或保留$\psi$，不能由此认定宇宙在赤道终止。

\originalsection{负曲率的构造与均匀性}
负曲率空间可辅助使用四维符号为$(-,+,+,+)$的平坦空间，取上支双曲面
\[
 -X_0^2+\mathbf X^2=-R^2,\quad X_0=R\cosh\psi,\quad
 \mathbf X=R\sinh\psi\,\mathbf n,
\]
将环境空间的二次型限制到切向位移，得到正定空间度规
\[
 \dd\ell^2=R^2(\dd\psi^2+\sinh^2\psi\dd\Omega^2).
\]
这里辅助的$X_0$是构造坐标，不是宇宙时间。环境Lorentz型变换保持双曲面及其诱导长度，并可把任意上支点送到$(R,0,0,0)$，故负曲率空间也均匀；普通空间旋转则证明各向同性。

令$\chi=R\psi$，面积$4\pi R^2\sinh^2(\chi/R)$比欧氏面积增长快，$\psi$从0到无穷。球内体积为
\[
 V(\chi)=\pi R^3[\sinh(2\chi/R)-2\chi/R].
\]
负曲率的二维Klein圆盘用有限坐标圆盘表示整个无限空间。若圆盘坐标$\mathbf q$满足$|\mathbf q|<1$，映回双曲面$X_0=R/\sqrt{1-q^2}$、$\mathbf X=R\mathbf q/\sqrt{1-q^2}$。两点的距离由环境内积得到
\[
 \cosh\frac{d_{12}}R=\frac{1-\mathbf q_1\cdot\mathbf q_2}{\sqrt{(1-q_1^2)(1-q_2^2)}}.
\]
径向从原点到$q$为$d=R\operatorname{artanh}q$，$q\to1$时距离发散。圆盘里画出的直弦是测地线，圆盘上的欧氏角却通常不是空间实际角。不能看见有限圆盘就判其体积有限。

\originalsection{统一三种曲率并接到宇宙度规}
本卷尺度因子$a$无量纲，共动曲率$K$单位为长度$^{-2}$，统一定义$S_K$后空间度规为
\[
 \dd\ell^2=a^2[\dd\chi^2+S_K^2(\chi)\dd\Omega^2].
\]
今天的曲率半径$R_{K0}=|K|^{-1/2}$，其他时刻实际半径为$aR_{K0}$。$K>0$时上文$R=aR_{K0}$，$\psi=\chi/R_{K0}$；$K<0$以双曲函数代替。极小距离时
\[
 S_K(\chi)=\chi-\frac{K\chi^3}{6}+O(K^2\chi^5),
\]
所以几何偏差是曲率乘以距离平方的量级。局部实验接近欧氏，不排除宇宙尺度可测的曲率。

固定时间切片的体积元是$\dd V=a^3S_K^2(\chi)\dd\chi\dd\Omega$。距离章中的横向共动距离$D_M=S_K(\chi)$正是把实际面积$4\pi a^2D_M^2$写成欧氏形式时出现的量；它和径向弧长$\chi$只有平直时相同。加上共动钟测得的时间项$-c^2\dd t^2$便得到FLRW度规，而不是把三维球形图当作时空本身。

\originalsection{曲率与拓扑、空间闭合与再坍缩}
原第13讲在默认单连通模型下称正曲率有限、非正曲率无限。本卷明确这个默认：$S^3$是有限无边界，$\mathbb R^3$和$H^3$无限。但局部曲率零的三维环面也可有有限体积，负曲率空间也可取适当的紧致商。局部度规不唯一决定整体拓扑。

“闭合空间”描述空间范围与拓扑，“再坍缩”描述时间演化。无压尘埃且$\Lambda=0$时二者有特殊对应；一般组分下必须从Friedmann方程求未来轨迹，不能由一张球面图替代动力学判断。本卷只使用单连通常曲率模型计算常规距离，非平凡拓扑的观测约束和分类保留为进阶范围。

\originalchapter{均匀膨胀的几何与动力学}
本章对应8.902第10讲。星系动力学把一个星系当成系统，宇宙学则先把足够大尺度上的物质分布当成一个整体。我们先建立均匀背景，再在后面的章节讨论偏离均匀的扰动。这里需要牛顿引力、微积分和流体压强；Einstein方程的几何推导可参阅本仓库广义相对论卷。本章补出几何如何变成可求解的宇宙演化方程，并统一原笔记中有歧义的单位和符号。

\originalsection{为什么膨胀律必须呈线性}
宇宙学原理是在足够大尺度上的统计均匀性和各向同性，不要求每一点都没有星系、空洞或局部运动。均匀性指平移不改变大尺度统计，各向同性指旋转不改变它们。仅从我们一个位置看到各向同性，不能单独推断均匀性；还需没有特殊观测位置的假设。

在一个空间切片中，设相对位置为$\mathbf r$的共动粒子相对速度为$\mathbf f(\mathbf r)$。换到另一个共动原点后，同一相对位置应对应同一相对速度，故
\[
 \mathbf f(\mathbf r_1)-\mathbf f(\mathbf r_2)=\mathbf f(\mathbf r_1-\mathbf r_2).
\]
连续性使这一加法关系给出线性映射$\mathbf f=\mathsf H\mathbf r$。再要求它与所有空间旋转对易，就有$\mathsf H=H\mathsf1$。因此局部膨胀的速度场是$\mathbf v=H(t)\mathbf r$。均匀性负责不同原点看到相同规律，各向同性排除矩阵的特殊方向；原讲义把这两步的解释交叉排列，本卷按所用假设重写。

引入无量纲\textbf{尺度因子}$a(t)$，约定今天$a_0=1$。固定共动位置$\mathbf x$在平直切片上的物理位置为$\mathbf r=a\mathbf x$，于是
\[
 H=\frac{\dot a}{a},\qquad \dot{\mathbf r}=H\mathbf r+\mathbf u,
 \qquad\mathbf u=a\dot{\mathbf x}.
\]
$\mathbf u$称本动速度。距离很近且本动速度可忽略时，观测到的谱线速度满足$cz\simeq H_0D$。远处$HD$是特定宇宙时间切片上的退行率，可以大于$c$；它不是在同一惯性系中测得的局部粒子速度，不能代入狭义相对论的速度合成公式。

\originalsection{空间曲率进入距离，而尺度因子进入时间演化}
空间的常曲率和球对称允许把度规写成
\begin{equation}\label{eq:astro-flrw}
 \dd s^2=-c^2\dd t^2+a^2(t)\bigl[\dd\chi^2+S_K^2(\chi)\dd\Omega^2\bigr],
 \quad \dd\Omega^2=\dd\theta^2+\sin^2\theta\dd\varphi^2,
\end{equation}
其中$\chi$有长度量纲，$K$有$\mathrm{length}^{-2}$量纲，且
\[
 S_K(\chi)=\begin{cases}
 \sin(\sqrt K\chi)/\sqrt K,&K>0,\\
 \chi,&K=0,\\
 \sinh(\sqrt{-K}\chi)/\sqrt{-K},&K<0.
 \end{cases}
\]
正曲率空间的球面积比欧氏空间小，负曲率则较大。$K\to0$时两式都趋于$\chi$。切片的标量曲率为${}^{(3)}R=6K/a^2$，今天的曲率半径为$|K|^{-1/2}$。也可把曲率符号归一成$k=0,\pm1$而令尺度因子带长度；两套约定不能交叉使用。原笔记度规的一次$a$、曲率单位$L^2$均应分别为$a^2$和$L^{-2}$。

各向同性还限制物质的平均应力：共动观测者不看到净能流和剪切应力。以质量等价密度$\rho$表示能量密度$\epsilon=\rho c^2$，四速度满足$u^\mu u_\mu=-c^2$，完美流体张量为
\[
 T^{\mu\nu}=(\epsilon+p)\frac{u^\mu u^\nu}{c^2}+p g^{\mu\nu}.
\]
局部正交标架中$T^{\hat\mu\hat\nu}=\operatorname{diag}(\epsilon,p,p,p)$，混合张量才是$T^{\hat\mu}{}_{\hat\nu}=\operatorname{diag}(-\epsilon,p,p,p)$。不能把原讲义的混合分量符号当成两个上标的能量密度为负。

\originalsection{从局部守恒到密度的尺度律}
取一团边界跟随共动流动的体积$V\propto a^3$，忽略净热流和耗散。膨胀时流体做功，局部能量守恒给出
\begin{equation}\label{eq:astro-continuity}
 \dd(\rho c^2a^3)=-p\dd(a^3),\qquad
 \dot\rho+3H\left(\rho+\frac p{c^2}\right)=0.
\end{equation}
这是$\nabla_\mu T^{\mu\nu}=0$的共动时间分量，并不假定宇宙存在一个一般意义上守恒的全局总能量。若某组分独立守恒且$p=w\rho c^2$、$w$恒定，则积分得
\[
 \rho(a)=\rho_0 a^{-3(1+w)}.
\]
冷物质的静质量能占主导，$w\simeq0$，只因数密度稀释而$a^{-3}$。辐射的各向同性动量流给出$p=\epsilon/3$，其光子能量还因红移减小，故$a^{-4}$。宇宙学常数对应$w=-1$，密度恒定。发生衰变、能量交换或粒子自由度变化时，要在各组分方程中加入交换项；$T\propto a^{-1}$也需注明熵自由度不变。

\originalsection{Friedmann方程怎样由Einstein方程给出}
用式\eqref{eq:astro-flrw}计算联络，时间与空间混合项正比$H$，从而在$x^0=ct$约定下
\[
 R_{00}=-\frac{3}{c^2}\frac{\ddot a}{a},\qquad
 R_{ij}=\frac1{c^2}\left(\frac{\ddot a}{a}+2H^2+\frac{2Kc^2}{a^2}\right)g_{ij}.
\]
收缩后代入$G_{\mu\nu}=R_{\mu\nu}-Rg_{\mu\nu}/2$。时间分量和任一空间对角分量分别含$3(H^2+Kc^2/a^2)/c^2$与$-(2\dot H+3H^2+Kc^2/a^2)/c^2$。与
$G_{\mu\nu}+\Lambda g_{\mu\nu}=8\pi GT_{\mu\nu}/c^4$比较，得到
\begin{align}
 H^2&=\frac{8\pi G\rho}{3}-\frac{Kc^2}{a^2}+\frac{\Lambda c^2}{3},\label{eq:astro-friedmann}\\
 \frac{\ddot a}{a}&=-\frac{4\pi G}{3}\left(\rho+\frac{3p}{c^2}\right)+\frac{\Lambda c^2}{3}.\label{eq:astro-acceleration}
\end{align}
第一式约束初始膨胀率与密度，第二式给膨胀的加速度。对第一式求导并用式\eqref{eq:astro-continuity}可恢复第二式，所以三式只有两式独立。\textbf{正压强会加强减速}，这正是简单牛顿尘埃图像遗漏的相对论作用。

牛顿图像仍能帮助理解尘埃情形。半径$r=a x$的均匀球内质量$M=4\pi\rho r^3/3$恒定，壳层满足$\ddot r=-GM/r^2$。乘以$\dot r$并积分有$\dot r^2/2-GM/r=E$，于是$H^2=8\pi G\rho/3+2E/(a^2x^2)$；最后一项与$-Kc^2/a^2$对应。这里用的是壳层定理的牛顿类比。Birkhoff定理描述球对称\emph{真空外部}几何，不能把它表述成任意有压强宇宙球都等同一个独立系统的定理。

\originalsection{密度参数与可积的演化模型}
定义$\rho_{\rm crit}(t)=3H^2/(8\pi G)$，并取
\[
 \Omega_i(t)=\frac{\rho_i(t)}{\rho_{\rm crit}(t)},\qquad
 \Omega_K(t)=-\frac{Kc^2}{a^2H^2},\qquad
 \rho_\Lambda=\frac{\Lambda c^2}{8\pi G}.
\]
于是$\Omega_r+\Omega_m+\Omega_\Lambda+\Omega_K=1$。曲率可以用这一无量纲项记账，却不是一种真实物质密度。$\rho_{\rm crit,0}=1.878\times10^{-29}h^2\ \mathrm{g\,cm^{-3}}$，其中$H_0=100h\ \mathrm{km\,s^{-1}\,Mpc^{-1}}$。带正$\Lambda$的宇宙即使$K>0$也可能永远膨胀；把临界密度称为所有模型的再坍缩分界并不成立。

今天归一后的演化方程为
\begin{equation}\label{eq:astro-Ea}
 \frac{H^2(a)}{H_0^2}=E^2(a)=\Omega_{r,0}a^{-4}+\Omega_{m,0}a^{-3}
 +\Omega_{K,0}a^{-2}+\Omega_{\Lambda,0}.
\end{equation}
非零正系数的不同幂次解释早期辐射、中期物质、晚期暗能量优势；曲率项并不一定实际成为主导。平直单组分$w>-1$时，$\dot a\propto a^{-(1+3w)/2}$，积分得
\[
 a\propto t^{2/[3(1+w)]},\qquad Ht=\frac{2}{3(1+w)}.
\]
辐射、物质分别是$t^{1/2},t^{2/3}$；空的负曲率Milne模型为$a\propto t$；纯正$\Lambda$给$a\propto\exp(\sqrt{\Lambda c^2/3}\,t)$，没有有限过去的$a=0$。不能把单组分$H_0t_0$直接用于混合宇宙。

\begin{example}
\textbf{平直物质与正宇宙学常数的年龄。}
忽略辐射时，$\dot a=H_0\sqrt{\Omega_{m,0}/a+\Omega_{\Lambda,0}a^2}$。令$y=a^{3/2}$，年龄积分化为
\[
 H_0t(a)=\frac{2}{3\sqrt{\Omega_{\Lambda,0}}}
 \operatorname{arsinh}\left(\sqrt{\frac{\Omega_{\Lambda,0}}{\Omega_{m,0}}}a^{3/2}\right).
\]
当$\Omega_{\Lambda,0}\to0$时，用$\operatorname{arsinh}x\simeq x$恢复尘埃结果。这也给$ a(t)\propto\sinh^{2/3}(3H_0\sqrt{\Omega_{\Lambda,0}}t/2)$，连续接上早期幂律与晚期指数，而非拼接两条曲线。
\end{example}

课程引用Planck 2018模型参数作为\textbf{课程年代的示例}，取$\Omega_{m,0}=0.315$、$\Omega_{\Lambda,0}=0.685$、$H_0=67.4$，不是本卷声称的最新测量。物质与暗能量密度相等满足$a_{m\Lambda}=(\Omega_{m,0}/\Omega_{\Lambda,0})^{1/3}$，即$z\simeq0.296$；开始加速则满足$\rho_m=2\rho_\Lambda$，$z\simeq0.632$。原第10讲将相等红移写成约1.3，应改正。物质与辐射相等有$1+z_{\rm eq}=\Omega_{m,0}/\Omega_{r,0}$；数值取决于是否计入中微子及其是否仍相对论性，不能把不同组分定义所得的3200或3700当成同一组精确参数。

\originalsection{解的空间形状不等于其未来结局}
$K>0,0,<0$分别决定局部空间正、零、负曲率；是否再坍缩还取决于$p,\Lambda$和初始条件。反过来，平直空间也可有非平凡整体拓扑。由有限区域的观测推断曲率，不等于直接测出了整个宇宙的拓扑或总体积。下一章将保持同一套$K,\chi,a$约定，从光线传播建立实际距离观测量。

\originalchapter{尘埃宇宙的参数解与自由粒子的传播}
本章对应8.286第9--15讲的解析演化与测地线应用。Friedmann式已经建立，本章先把可积模型真正解出来，再从局部自由运动推出宇宙中的动量红移。测地线的一般理论及黑洞时空详见广义相对论卷；这里计算的是课程所需的特定背景和相应时间定义。

\originalsection{闭合尘埃模型的展开与收缩分支}
取$\Lambda=p=0$且$K>0$，记实际曲率半径$R=a/\sqrt K$。因为$\rho R^3$不变，定义常长度
\[
 A=\frac{4\pi G\rho R^3}{3c^2}.
\]
Friedmann式化为$\dot R^2=c^2(2A/R-1)$。令$R=A(1-\cos\eta)$，则$\dd R=A\sin\eta\dd\eta$；在展开分支选择正根，$\dd t=(A/c)(1-\cos\eta)\dd\eta$。以大爆炸为$t=0$积分得
\[
 R=A(1-\cos\eta),\qquad ct=A(\eta-\sin\eta),\quad0<\eta<2\pi.
\]
此参数式平滑经过$\eta=\pi$的最大尺度$2A$，并继续到收缩；直接只取$\dot R>0$会漏掉后一分支。$\eta\ll1$时$R\simeq A\eta^2/2$、$ct\simeq A\eta^3/6$，恢复$R\propto t^{2/3}$。

密度参数$\Omega=8\pi G\rho/(3H^2)$只含尘埃，代入得
\[
 \Omega=\frac2{1+\cos\eta},\qquad
 Ht=\frac{\sin\eta(\eta-\sin\eta)}{(1-\cos\eta)^2}.
\]
同一个$\Omega>1$可对应展开或收缩，$H$的符号区分两者。若已知$H>0$，应取$\eta=\arccos(2/\Omega-1)\in(0,\pi)$；若$H<0$，取$2\pi-\arccos(2/\Omega-1)$。这样比含$\arcsin$的分象限表达更不易漏分支。最大尺度处$H=0$使$\Omega$发散，但密度、$R,t$有限；$\Omega$不是那个时刻的良好演化变量。

\originalsection{开放与平直尘埃的年龄}
对$K<0$用$R=a/\sqrt{-K}$与同一定义的$A$，有$\dot R^2=c^2(2A/R+1)$。令$R=A(\cosh\eta-1)$，类似积分得
\[
 R=A(\cosh\eta-1),\quad ct=A(\sinh\eta-\eta),\quad
 \Omega=\frac2{1+\cosh\eta}.
\]
开放模型早期同样近尘埃平直解，晚期$R\simeq ct$，但它一直有物质，不能把晚期渐近性质说成所有时刻都为空的Milne模型。以$0<\Omega<1$表达年龄，得
\[
 Ht=\frac1{1-\Omega}-\frac{\Omega}{2(1-\Omega)^{3/2}}
 \operatorname{arcosh}\left(\frac2\Omega-1\right).
\]
闭合展开分支则为
\[
 Ht=\frac{\Omega}{2(\Omega-1)^{3/2}}
 \arccos\left(\frac2\Omega-1\right)-\frac1{\Omega-1}.
\]
两式$\Omega\to1$的极限都是$2/3$；数值计算近该点宜用级数或直接年龄积分，避免两个大数相减。$\Omega\to0$的开放极限为1。

8.286第8--10讲用2013年例值$H_0\simeq67.3$ km/s/Mpc，平直纯尘埃给$t_0\simeq9.69$ Gyr。若另外的恒星年龄证据大于此值，冲突指向\emph{这一特定组分模型}，并非所有膨胀宇宙。正真空能延长相同今日$H_0$下的年龄，具体混合解见背景章。三条尘埃曲线也可独立比较：$K>0$回转，$K=0$永远减速，$K<0$逐渐近匀速。

\begin{center}\begin{tikzpicture}\begin{axis}[width=.7\textwidth,height=4.8cm,xlabel={$ct/A$},ylabel={$R/A$},xmin=0,xmax=6.3,ymin=0,ymax=6,legend pos=north west]
\addplot[structurecolor,domain=.001:359.9,samples=200,variable=\t] ({rad(\t)-sin(\t)},{1-cos(\t)});\addlegendentry{闭合尘埃}
\addplot[black,dashed,domain=.001:2.9,samples=150,variable=\t] ({(exp(\t)-exp(-\t))/2-\t},{(exp(\t)+exp(-\t))/2-1});\addlegendentry{开放尘埃}
\addplot[gray,domain=.001:6.3,samples=100] {(9/2)^(1/3)*x^(2/3)};\addlegendentry{平直尘埃}
\end{axis}\end{tikzpicture}\end{center}
平直曲线用$\dot R^2=2Ac^2/R$定义同一质量尺度$A$，不把不存在的有限曲率半径当作其半径。

\originalsection{从驻值路径到有明确定义的参数}
空间曲线长度是$\int\sqrt{g_{ij}\dot x^i\dot x^j}\dd\lambda$。改变路径而固定两端，积分分部后的Euler方程包含$1/\sqrt{g_{ij}\dot x^i\dot x^j}$。若用弧长作参数，这个因子成为常数，简化为
\[
 \frac{\dd}{\dd\lambda}(g_{\mu\nu}x'^\nu)
 =\frac12\partial_\mu g_{\alpha\beta}x'^\alpha x'^\beta.
\]
对类时自由粒子用固有时$\tau$，满足$g_{\mu\nu}x'^\mu x'^\nu=-c^2$；对光则用仿射参数，不能用处处为零的固有时参数化。上式也由$L=g_{\mu\nu}x'^\mu x'^\nu/2$直接变分得出。展开导数并乘逆度规可得$x''^\mu+\Gamma^\mu_{\alpha\beta}x'^\alpha x'^\beta=0$。一般重参数化会增加沿切线的项；一个坐标加速度不能脱离参数定义比较。

\originalsection{平直FLRW里物理动量为何随尺度因子衰减}
取$\dd s^2=-c^2\dd t^2+a^2\delta_{ij}\dd x^i\dd x^j$。空间坐标不显含于度规，变分式的右侧为零，因此$a^2\dd x^i/\dd\tau=C_i$恒定。共动观察者测得粒子的本动速度为$v^i=a\dd x^i/\dd t$，$\dd t/\dd\tau=\gamma$，于是
\[
 p_{\hat i}=m\gamma v^i=\frac{mC_i}{a},\qquad
 E=\sqrt{m^2c^4+p^2c^2}.
\]
因而无碰撞粒子的\emph{物理动量}随$a^{-1}$，是协变空间动量$p_i=ma^2x'^i$守恒的局部标架表达。光的仿射测地线也有同样的尺度律，$E=pc$给$1+z=a_o/a_s$。曲空间可用空间等距对称性作相同的局部推论；不能假定三个普通平移坐标在所有$K$都全局存在。

重粒子非相对论时$v\simeq p/m\propto a^{-1}$；相对论时$v\simeq c$而不是同样按$a^{-1}$缩小。若存在碰撞、束缚引力或空间扰动，上述常量不再保留。星系内部的轨道和恒星半径不会因把背景$a$写进坐标就自动同步膨胀。

\originalsection{黑洞径向落体为何看似牛顿式却有不同的钟}
第15讲以Schwarzschild径向运动说明固有时和远处坐标时的区别；详细几何留给广义相对论卷。对外部真空取$f=1-r_s/r$，$r_s=2GM/c^2$，类时粒子有守恒$e=f\dd t/\dd\tau$。规范条件给
\[
 \left(\frac{\dd r}{\dd\tau}\right)^2=c^2(e^2-f),\qquad
 \frac{\dd^2r}{\dd\tau^2}=-\frac{GM}{r^2}.
\]
第二式虽与牛顿形式相同，$r$是面积半径，$\tau$是落体自己的钟，不能据此否定相对论修正。由$r_0$处相对静止释放，$e^2=f(r_0)$，以$r=r_0\cos^2\theta$积分得
\[
 \tau=\sqrt{\frac{r_0^3}{2GM}}(\theta+\sin\theta\cos\theta),\quad
 0\leq\theta\leq\frac\pi2.
\]
到$r=0$的理想点粒子固有时有限，实际潮汐破坏和经典奇点处理论失效须另外讨论。接近$r_s$时$\dd t/\dd\tau=e/f$使Schwarzschild坐标时发散；这不等于落体自己的钟在视界停下。课程所需的具体轨道数值应用另见恒星与致密天体习题，不复制GR卷整套证明。

\originalchapter{红移、光度距离与宇宙学观测}
本章对应8.902第11讲及第12讲开头的共动体积。需要区分两个问题：宇宙中相隔多远，以及我们从收到的光推断出哪一种距离。膨胀空间里这两个问题不能用一个欧氏距离同时回答。沿用式\eqref{eq:astro-flrw}的无量纲$a_0=1$、长度坐标$\chi$和曲率$K$。

\originalsection{红移如何从两束相邻光波得到}
定义光谱红移$1+z=\lambda_{\rm obs}/\lambda_{\rm em}=\nu_{\rm em}/\nu_{\rm obs}$。先考虑固定共动坐标的光源与观测者，忽略其本动速度和局部引力势。径向光线沿零测地线传播，$\dd s^2=0$给出$c\dd t=a\dd\chi$。相邻两个波峰跨过同一段共动距离，故
\[
 \int_{t_e}^{t_o}\frac{c\dd t}{a(t)}
 =\int_{t_e+\delta t_e}^{t_o+\delta t_o}\frac{c\dd t}{a(t)}.
\]
当一周期远短于膨胀时间时，端点变化给$\delta t_o/a_o=\delta t_e/a_e$，于是
\begin{equation}\label{eq:astro-redshift}
 1+z=\frac{a_o}{a_e},\qquad a_e=\frac1{1+z}\quad(a_o=1).
\end{equation}
同一关系既降低光子能量，又拉长源的变化时间：$E_o=E_e/(1+z)$、$\delta t_o=(1+z)\delta t_e$。源的超新星光变因此发生时间伸长。光子沿零测地线传播并不表示空间切片中的路径总是普通意义的最短路。

狭义相对论的径向退行速度$\beta c$给$1+z_{\rm Dop}=\sqrt{(1+\beta)/(1-\beta)}$：在源的静止系两次发射相隔$\Delta\tau$，观测者系相隔$\gamma\Delta\tau$，第二个波峰又需多走$\beta c\gamma\Delta\tau$，总接收间隔为$\gamma(1+\beta)\Delta\tau$。这解释近距离$z\simeq\beta$。宇宙学红移则积分了光传播期间不断改变的几何；远距离不能由一个全局惯性系的$\beta$解释。存在本动与引力红移时，局部频率比与背景膨胀的频率比相乘，而不是把大红移简单相加。

\originalsection{径向共动距离与回溯时间}
从$\dot a=aH$和式\eqref{eq:astro-redshift}可得
\[
 \frac{\dd z}{\dd t}=-(1+z)H(z),\quad
 E^2(z)=\Omega_{r,0}(1+z)^4+\Omega_{m,0}(1+z)^3
 +\Omega_{K,0}(1+z)^2+\Omega_{\Lambda,0}.
\]
发射与接收相隔的\textbf{回溯时间}为
\[
 t_0-t_e=\frac1{H_0}\int_0^z\frac{\dd z'}{(1+z')E(z')},
\]
径向共动距离则是
\begin{equation}\label{eq:astro-comoving-distance}
 D_C=\chi=\frac c{H_0}\int_0^z\frac{\dd z'}{E(z')}.
\end{equation}
额外的$1+z$有明确来源：回溯时间积分的是物理时间，共动距离把每段光程除以当时的$a$。所以$c(t_0-t_e)$一般不是$D_C$。今天同一宇宙时间切片上的径向固有距离为$a_0\chi=\chi$，发射时则为$a_e\chi$；我们不能在同一时刻同时看见这条切片上的所有点。

共动距离由观测和模型推断，而不是禁止观测的量。其优势是理想共动物体的$\chi$不随宇宙膨胀改变。非零曲率时，径向距离$D_C$和横向共动距离
\[
 D_M=S_K(D_C)
\]
不同；下列角度和通量关系使用$D_M$。将$\Omega_K$与$K$的符号混淆会把正弦和双曲正弦用反。

\originalsection{用尺寸与通量分别定义距离}
设源在发射切片上的横向物理长度为$l$，小角尺寸为$\vartheta$。度规给$l=a_eS_K(\chi)\vartheta$，所以\textbf{角直径距离}定义及结果是
\[
 D_A=\frac l\vartheta=\frac{D_M}{1+z}.
\]
$l$必须指源当时的实际物理尺度，不能误用今日共动尺度。远处角直径距离可能随红移增加而减小；物体更远并不总是更小。

\textbf{光度距离}由源静止系的总光度$L$与观测总辐射通量$F$定义$F=L/(4\pi D_L^2)$。观测光束铺在今日面积$4\pi D_M^2$上；每个光子的能量和单位接收时间的光子数各减小一个$1+z$，故
\begin{equation}\label{eq:astro-distance-duality}
 F=\frac{L}{4\pi D_M^2(1+z)^2},\qquad
 D_L=(1+z)D_M=(1+z)^2D_A.
\end{equation}
后一个式子称距离对偶关系。这里假定光子数沿传播守恒、光走度规的零测地线，且用一致的源与光束面积定义；尘埃吸收需单独校正。引力透镜改变单源视通量和视尺寸，不能无视它给单个源距离估计造成的误差。

式\eqref{eq:astro-distance-duality}使用\emph{全频率积分}的通量。固定观测波段对应源不同的静止系波段，需星系章节讨论的谱形与$K$修正。例如以每单位频率表示，在相应频率$\nu_e=(1+z)\nu_o$处
\[
 F_{\nu_o}=\frac{(1+z)L_{\nu_e}}{4\pi D_L^2}.
\]
其积分$\dd\nu_e=(1+z)\dd\nu_o$恢复总通量式；写成每单位波长时因变量雅可比不同而有相反的额外因子。这是检验谱通量红移系数的简便办法。

\begin{example}
\textbf{Einstein--de Sitter宇宙的三个距离。}
取$\Omega_{m,0}=1$且其他参数为零，$E(z)=(1+z)^{3/2}$。直接积分得
\[
 \frac{H_0D_C}{c}=2\left[1-(1+z)^{-1/2}\right],\qquad
 D_A=\frac{D_C}{1+z},\qquad D_L=(1+z)D_C.
\]
当$z\ll1$三者都趋于$cz/H_0$。令$y=1+z$，$D_A\propto y^{-1}-y^{-3/2}$，导数为零要求$\sqrt y=3/2$，即$z=5/4$。这给角直径距离的最大点，而$D_L$仍随$z$单调增加。
\end{example}
\begin{center}
\begin{tikzpicture}
\begin{axis}[width=.76\textwidth,height=5.3cm,xmin=0,xmax=4,ymin=0,ymax=6,
 xlabel={$z$},ylabel={$H_0D/c$},legend pos=north west,grid=major]
\addplot[structurecolor,thick,domain=0:4,samples=100]{2*(1-1/sqrt(1+x))/(1+x)};
\addlegendentry{$D_A$}
\addplot[black,dashed,domain=0:4,samples=100]{2*(1-1/sqrt(1+x))};\addlegendentry{$D_C$}
\addplot[structurecolor,densely dotted,thick,domain=0:4,samples=100]{2*(1+x)*(1-1/sqrt(1+x))};\addlegendentry{$D_L$}
\end{axis}
\end{tikzpicture}
\end{center}
图由上述解析式自绘，展示的是物质单组分玩具模型，不是将课程中的第三方观测图重新嵌入。

\originalsection{天体计数需要的共动体积}
在红移宽度$\dd z$与立体角$\dd\Omega$内，共动横截面积是$D_M^2\dd\Omega$，径向厚度由式\eqref{eq:astro-comoving-distance}给出$c\dd z/H(z)$，所以
\[
 \frac{\dd V_c}{\dd z\dd\Omega}=\frac{cD_M^2}{H(z)}
 =\frac{c(1+z)^2D_A^2}{H(z)}.
\]
发射时固有体积再乘$a_e^3$。原笔记体积式中的个别$d\chi/dz$和额外径向因子排印混乱，按面积乘径向厚度即可辨认。若共动数密度$n_c(z)$已知，理想计数为$\dd N=n_c\dd V_c$；实际还需通量选择函数、遮挡、红移误差和天体演化。源静止系发生率$\mathcal R_c$给观测事件率时，还要除以$1+z$，不能把对象计数与事件率混为一式。

\originalsection{参数如何变成可比较的数据预测}
标准烛光先给距离模数$\mu=5\log_{10}(D_L/10\ \mathrm{pc})$，标准尺先给$\vartheta=l/D_A$。将测得的$z$与模型预测的$D_L,D_A$对比，可以约束式\eqref{eq:astro-Ea}的参数。绝对光度标定与$H_0$有简并，演化和消光能模仿距离变化，故单类数据不能未经系统误差处理就给宇宙结论。

低红移时定义减速参数$q_0=-\ddot a_0/(a_0H_0^2)$，Taylor展开给
\[
 D_L=\frac c{H_0}\left[z+\frac{1-q_0}{2}z^2+O(z^3)\right].
\]
一阶只能测膨胀率，二阶才携带加速信息；曲率在这一展开中从更高阶进入。高红移应积分完整方程，不能把二阶展开当成适用于$z\sim1$的精确关系。这一逻辑把均匀背景与实际超新星观测连接起来。

\originalchapter{真空能、加速膨胀与参数推断}
本章对应8.286第18--21讲，并接回第1--2讲中的真空能问题。背景章给出了含$\Lambda$的演化式和解析年龄，距离章给出了光线测距。本章说明“负压强”的物理含义，以及观测为何需要多种证据共同约束。量子场论的正规化推导见QFT卷；这里不把一个发散的零点能积分当成可直接读出的宇宙常数。

\originalsection{恒定能量密度怎样产生负压强}
取真空能密度$\epsilon_v$不随体积变化，区域能量$U=\epsilon_vV$。若准静态改变体积而无热流，$\dd U=-p\dd V$给$p_v=-\epsilon_v$。假想活塞的外面真空压强取零，拉开活塞需对内部系统做正功；这解释了为何负压强并非“负能量密度”。真实宇宙无这样的活塞边界，但局部共动能量守恒给同一关系。

Lorentz不变真空不能挑选一个静止速度方向，故其应力必须与度规成比例：$T_{\mu\nu}^{v}=-\epsilon_vg_{\mu\nu}$。把它移到Einstein方程的左边可写成
\[
 \Lambda=\frac{8\pi G\epsilon_v}{c^4},\qquad
 \rho_v=\frac{\epsilon_v}{c^2},\quad p_v=-\rho_vc^2.
\]
加速方程中的$\rho+3p/c^2$因此贡献$-2\rho_v$。这仍是正$G$和正能量密度的引力理论，不是把Newton常数改成负数。在真空能主导且展开的阶段，$a\sim\ee^{H_vt}$，$H_v=\sqrt{8\pi G\rho_v/3}$。加入其他组分或曲率时不是从一开始就精确指数。

膨胀时$U$可增大而$\epsilon_v$不变，满足$\dd U=-p_v\dd V$。这不授权宣称整个宇宙有一个定义无歧义且值为零的全局总能量。一般FLRW时空缺少全局时间平移对称性；第1--2讲用负引力能与正物质能抵消作直观图像，本卷将其保留为特定模型和能量定义下的启发，不作为已证明的普遍定理。

\originalsection{由现象到组分：三种证据各约束什么}
远处Ia超新星在相同红移下偏暗，意味着在标准化光度、传播和选择效应受控后光度距离较大。距离是$H(z)$的累计积分，不能由某一个超新星直接测出瞬时$\ddot a$。原第21讲所述20--30\%偏暗是1998年研究的量级描述，具体依赖比较模型和红移，不是任何$z$都保持同一差值。

CMB声学角尺度主要比较最后散射时的声学标准尺与角直径距离，帮助限制空间曲率和组分，但存在几何退化。其他观测约束物质密度后，近空间平坦而物质不足的组合指向额外组分。恒星年龄则限制宇宙年龄不能比其中已形成的恒星更小。三者从距离、早期标准尺和时钟给互补信息；任何一个结果都需说明校准与模型依赖。

用$w(a)=p_{\rm de}/\epsilon_{\rm de}$描述可变化暗能量，独立守恒给
\[
 \rho_{\rm de}(a)=\rho_{{\rm de},0}
 \exp\left[-3\int_1^a(1+w(a'))\dd\ln a'\right].
\]
$w=-1$恢复宇宙学常数；$w$常数给$a^{-3(1+w)}$。正能量组分的$w<-1/3$可以促加速，却不一定就是严格真空。改变引力理论也可能改变从观测到组分的推断。本卷以GR和FLRW为计算前提，不将“暗能量”这个现象名称等同于唯一已确定的微观粒子。

\originalsection{超新星怎样从标准烛光变为实际拟合}
距离模数是$\mu=m-M=5\log_{10}[D_L/(10\ \mathrm{pc})]$。Ia超新星并非天然相同光度，要根据光变宽度与颜色作标准化。例如8.902 PS5采用$\mu=m-M+\alpha s-\beta c$，其$M$定义与写$\alpha(s-1)$的常见记法相差零点。红移$z$、色参数$c$及速度光速$c$用同一个字母时尤其需在语境中辨清。

因为$D_L\propto H_0^{-1}$，若没有外部绝对光度校准，改变$H_0$可由改变$M$抵消。此时超新星主要限制距离曲线\emph{形状}和组分，不能单独测出$H_0$。更完整分析还需校准误差、相关光变参数协方差、选择效应、宿主环境、吸收与模型内禀散布。后面真实SNLS习题已实际读取数据，按题给简化模型拟合，并明确两种散布解释；那是课程计算练习，不冒充发表论文的完整联合似然。

\originalsection{年龄与距离的积分边界不可省略}
对今天$a_0=1$的展开模型，含辐射、物质、曲率、真空能时
\[
 H_0t_0=\int_0^1\frac{a\dd a}
 {\sqrt{\Omega_{r,0}+\Omega_{m,0}a+\Omega_{K,0}a^2+\Omega_{v,0}a^4}}.
\]
这对应第20讲的年龄积分；只有在积分区间内根号为正且存在所选大爆炸展开分支才适用。允许负$\Lambda$、回转或反弹时须先查分支，不能把所有密度参数直接丢进一个数值积分再解释为年龄。

光度距离还含空间面积因子。正曲率时$D_M=R_{K0}\sin(\chi/R_{K0})$，负曲率用$\sinh$，$\chi=c\int_0^z\dd z'/H(z')$。红移光子的能量与到达率各降一个$1+z$，故$D_L=(1+z)D_M$。正曲率反极点附近需处理焦点和可能的多路径，普通近邻资料的模型适用区间与整个封闭空间的光线拓扑不能混成同一个简单距离分支。

\originalsection{为何真空能量级仍是独立的理论难题}
在无引力的实验里，势能加常数不改变相对动力学；有引力时常数进入$T_{\mu\nu}$，影响宇宙膨胀。量子自由场的每个谐振模式有零点能，形式上玻色场贡献$+\hbar\omega/2$，费米场的真空贡献符号不同；Higgs等场在最低势点的常数也参与记账。

以动量截止$\Lambda_{\rm UV}$粗估，一种无质量场的零点能密度随$\Lambda_{\rm UV}^4$增长。截止取Planck能标会与宇宙所需小尺度相差约$10^{120}$量级。这说明参数自然性和不同来源之间的抵消问题；它不等于用不受控截止已精确预测了$\Lambda$。有效场论中还存在允许的宇宙学常数反项，最终可观测的是重整化后的总量，而不是某一单项的发散。

第1--2、21讲讨论多真空与环境选择：若真空能过大而正，结构增长过早停止；若过大而负，宇宙可能过早坍缩。因此观察者所在区域可能有选择偏向。但给出可检验概率仍需真空分布、形成权重、观察者代理量与测度。课程的$10^{500}$景观数量是特定理论讨论中的估计，不是已观测的宇宙数；本卷将该框架和争议保留为进阶方向，不据它宣称已解决真空能问题。

\originalchapter{因果视界、平坦性与暴胀}
本章对应8.902第11讲末至第12讲的暴胀部分。Friedmann方程描述给定组分后怎样演化，却尚未解释为什么早期密度如此均匀、曲率如此小，以及结构的初始扰动从何而来。暴胀是一类早期加速膨胀模型，必须同时说明加速的条件、持续时间及结束机制；它不是在时间轴上增加一个没有动力学解释的箭头。

\originalsection{复合为什么发生在远低于电离能的温度}
宇宙冷却后，反应$p+e\rightleftarrows\mathrm H+\gamma$逐渐形成中性氢。13.6 eV对应约$1.6\times10^5$ K，但每个重子对应大量光子，高能尾部仍可电离氢。热平衡、非简并粒子近似下，令$x=n_e/n_b=n_p/n_b$且忽略氦，Saha关系为
\[
 \frac{x^2}{1-x}=\frac1{n_b}\left(\frac{m_e\kb T}{2\pi\hbar^2}\right)^{3/2}
 \exp\left(-\frac{\chi_{\rm H}}{\kb T}\right),\qquad\chi_{\rm H}=13.6\ \mathrm{eV}.
\]
这里的氢基态、电子自旋统计因子采用常见简化约定；它们不改变指数导致的温度尺度。光子数密度$n_\gamma=2\zeta(3)(\kb T)^3/(\pi^2\hbar^3c^3)$、$\eta_b=n_b/n_\gamma\ll1$。将$n_b=\eta_b n_\gamma$代入，可看出小$\eta_b$要求温度进一步下降才使$x$显著小于1，因而复合温度约为几千K。

\textbf{复合与光子脱耦不是同一个定义}。中性比例涉及化学反应，脱耦涉及散射率$n_e\sigma_Tc$与膨胀率及沿途光学深度的竞争。Saha平衡在实际复合后期失效，精确最后散射面还需非平衡原子过程；$z_*\sim1100$在此作为课程所用尺度，并不是由上述一个平衡式精确算出的观测红移。宇宙微波背景的均匀性因此提出一个因果问题，而非仅是温度计算问题。

\originalsection{三种视界必须分别定义}
从某个初始时刻$t_i$起，光能传播的最大共动径向距离定义\textbf{粒子视界}
\[
 \chi_p(t)=\int_{t_i}^t\frac{c\dd t'}{a(t')},\qquad D_p=a(t)\chi_p(t).
\]
它取决于模型的过去历史和积分是否收敛。\textbf{事件视界}$\chi_e(t)=\int_t^{t_f}c\dd t'/a(t')$则问现在发出的信号将来最多到哪里，只有该未来积分有限时才存在。\textbf{共动Hubble半径}$R_H=c/(aH)$是一个局部膨胀尺度，不一般等于前两者。

若$a\propto t^p$且$0<p<1$，从$t_i=0$积分有
\[
 D_p=\frac{ct}{1-p},\qquad\frac cH=\frac{ct}{p}.
\]
辐射$p=1/2$时$D_p=c/H$；物质$p=2/3$时$D_p=2c/H$。纯de Sitter模型的$H$为常数，有未来事件视界$c/H$，但其粒子视界还取决于选取的过去边界。用$1/(aH)$讨论暴胀的模出入是方便近似，不应将它直接叫作所有模型的因果视界。

复合前已有辐射与物质两组分。忽略曲率和暗能量，可把粒子视界积分精确做到
\[
 \chi_p(a)=\frac{2c}{H_0\Omega_{m,0}}
 \left[\sqrt{\Omega_{r,0}+\Omega_{m,0}a}-\sqrt{\Omega_{r,0}}\right].
\]
这由$H=H_0a^{-2}\sqrt{\Omega_{r,0}+\Omega_{m,0}a}$直接积分得到。复合时的角尺度约为$\vartheta_p=a_*\chi_p(a_*)/D_A(z_*)=\chi_p(a_*)/D_M(z_*)$。选取课程示例参数后它是度数量级，远小于整个天空。原笔记在复合时把辐射与物质之比估为0.05；对$\Omega_{r,0}\simeq9.2\times10^{-5}$、$\Omega_{m,0}=0.315$、$z_*=1090$，该比为$\Omega_{r,0}(1+z_*)/\Omega_{m,0}\simeq0.32$，纯物质近似已有可见误差。

在从减速热大爆炸外推的模型里，大尺度温度相近需要一个初始条件解释：相距远的最后散射区域没有充分的早期因果联系，为什么还具有相近热状态？天空温度近各向同性并非数学上不可能，它指出了模型本身没有产生这一初始条件的机制。暴胀也仍需相应初始区域、势和再加热条件，不能宣称它无条件解决所有初始条件问题。

\originalsection{为什么加速膨胀能改善平坦性与因果联系}
包括真空能在内的总密度参数满足
\[
 \Omega_{\rm tot}-1=\frac{Kc^2}{a^2H^2},\qquad
 \frac{\dd}{\dd t}\left(\frac c{aH}\right)=-\frac{c\ddot a}{\dot a^2}.
\]
因此$\ddot a>0$使共动Hubble半径减小，也压低曲率相对贡献。对恒定$w>-1/3$的平直主导组分，$|\Omega_{\rm tot}-1|\propto a^{1+3w}$：辐射时期随$a^2$增长，物质时期随$a$增长。越向过去外推，平坦性要求越紧；加速期将这个趋势反转。

若暴胀前一片足够大、已具有因果联系的区域被拉伸，后来的可观测区域可来自其内部。扰动波长$\lambda_{\rm phys}=a\lambda_c$增长，而$H$近恒定，于是固定共动模可以从Hubble半径内被拉到外。暴胀结束后的减速阶段$1/(aH)$重新增长，模再进入。在这一叙述中\emph{粒子视界的累计积分仍增长}；减小的是共动Hubble半径，二者不能互换。

无相互作用的重遗迹数密度随$a^{-3}$，$N$个e折使它降低$\ee^{-3N}$。这可稀释在较早相变形成的磁单极，但再加热不能又大量重造同类遗迹。磁单极的产生来自特定大统一对称性破缺，而不是所有量子场论的必然预言。课程中的单极问题是这类模型与观测稀少性之间的张力。

\originalsection{标量场怎样实现暂时的负压强}
本节取$\hbar=c=1$并定义约化Planck质量$\bar M_P^2=1/(8\pi G)$。对负时间号度规，规范标量场作用量为
\[
 S_\phi=\int\sqrt{-g}\dd^4x\left[-\frac12g^{\mu\nu}\partial_\mu\phi\partial_\nu\phi-V(\phi)\right].
\]
均匀场的能量密度和压强为$\epsilon_\phi=\dot\phi^2/2+V$、$p_\phi=\dot\phi^2/2-V$。代入加速方程，忽略其他组分，有$\ddot a>0$当且仅当$\dot\phi^2<V$。缓慢运动的场有接近$-1$的$w$，但仅仅有势能不保证它持续缓慢运动。

变分场作用量，或把上述$\epsilon,p$代入守恒式，得
\begin{equation}\label{eq:astro-inflaton}
 \ddot\phi+3H\dot\phi+V'(\phi)=0,\qquad
 H^2=\frac{\dot\phi^2/2+V}{3\bar M_P^2},\qquad
 \dot H=-\frac{\dot\phi^2}{2\bar M_P^2}.
\end{equation}
$3H\dot\phi$是均匀膨胀导致的阻尼，不是把能量自动转移到热辐射；实际再加热还需场与其他粒子的耦合。直接令$\Lambda$在早期占主导不能说明结束，因为真正恒定的$\Lambda$不会自行衰变。

\originalsection{慢滚条件与持续时间}
当$\dot\phi^2\ll V$且$|\ddot\phi|\ll|3H\dot\phi|$，有
\[
 H^2\simeq\frac V{3\bar M_P^2},\qquad 3H\dot\phi\simeq-V',\qquad
 \epsilon_V=\frac{\bar M_P^2}{2}\left(\frac{V'}V\right)^2,\quad
 \eta_V=\bar M_P^2\frac{V''}V.
\]
$\epsilon_V\ll1,|\eta_V|\ll1$通常保证慢滚近似自洽，精确加速条件则是$\epsilon_H=-\dot H/H^2<1$。将慢滚方程代回去可得$\dot\phi^2/V\simeq2\epsilon_V/3$；原笔记$1/(24\pi G)(V'/V)^2$实际上是这个动能比的近似，和本卷标准$\epsilon_V=1/(16\pi G)(V'/V)^2$不是同一定义。

从起始场值滚到结束场值的e折数为
\begin{equation}\label{eq:astro-efolds}
 N=\ln\frac{a_f}{a_i}=\int_{t_i}^{t_f}H\dd t
 \simeq\frac1{\bar M_P^2}\int_{\phi_f}^{\phi_i}\frac V{V'}\dd\phi.
\end{equation}
减小共动Hubble半径的比值为$(a_iH_i)/(a_fH_f)$，近恒定$H$时约为$\ee^{-N}$。常说的50--60个e折是把某个今日观测尺度接到暴胀和再加热历史后的数量级，并非所有模型统一的固定数。再加热温度、暴胀能标和后续状态方程都会改变它。

\begin{example}
\textbf{二次势给出的慢滚练习。}
取$V=m^2\phi^2/2$，则$\epsilon_V=\eta_V=2\bar M_P^2/\phi^2$。用$\epsilon_V\simeq1$估结束场值$\phi_f\simeq\sqrt2\bar M_P$，积分式\eqref{eq:astro-efolds}得$N\simeq(\phi_i^2-\phi_f^2)/(4\bar M_P^2)$。这是一项解析练习，说明持续时间取决于场值区间，而不能只用势能高度判断。本卷不据它宣称该势符合当前观测。
\end{example}

\originalsection{结构种子、结束与热史的连接}
近de Sitter阶段一个轻的规范标量场的真空涨落在Hubble尺度冻结，典型$\delta\phi\sim H/(2\pi)$。不同区域的场值差意味着滚到同一结束面所需时间差$\delta t\sim-\delta\phi/\dot\phi$，进而产生曲率扰动$\mathcal R\sim-H\delta\phi/\dot\phi$。在单场慢滚、绝热真空与领先阶近似下，其无量纲功率为
\[
 \mathcal P_{\mathcal R}\simeq\frac{H^2}{8\pi^2\bar M_P^2\epsilon_H}.
\]
这里说明物理联系并给出领先阶量子结果；其完整规范扰动量子化超出原第12讲的经典推导，不将这一估计冒充已证明的一般初始谱。多场、非规范动能和非平凡初态均可改变结果。

场接近势的最低点后可能振荡并将能量转移给普通粒子，形成热辐射；暴胀结束、能量转移和达到热平衡可以是不同阶段。以后随膨胀冷却，经历辐射—物质相等、复合、星系形成与晚期暗能量优势。原讲义时间轴的$10^{-34}$秒是示意的模型尺度，而不是测得的统一暴胀时刻。接下来先研究扰动如何在这一背景中增长，再回到可直接由热平衡与反应率计算的早期热史。

\originalchapter{规范对称性、Higgs场与宇宙磁单极问题}
本章对应8.286第21--22讲，特别采用第22讲Scott A. Hughes代课逐字稿，并补齐其示意讨论中的数学条件。原课程教师为Alan H. Guth，但不能把本讲口述全部署为Guth。QFT卷给出规范理论的一般结构，本章只解释热宇宙中对称性变化为何关系到物质质量与遗迹；电弱实验机制和假设的大统一相变分别说明。

\originalsection{从电磁势的自由度到局部相位}
取$\hbar=c=1$，电荷为$q$的场满足$\psi'=\ee^{\ii q\Lambda(x)}\psi$。若$D_\mu=\partial_\mu-\ii q A_\mu$且$A'_\mu=A_\mu+\partial_\mu\Lambda$，直接展开得
\[
 D'_\mu\psi'=\ee^{\ii q\Lambda}D_\mu\psi.
\]
所以由$|D\psi|^2$建立的动力学不因任意局部相位改变。$F_{\mu\nu}=\partial_\mu A_\nu-\partial_\nu A_\mu$也保持不变。这把经典势的规范自由度与量子场局部相位联系起来；改变描述不会改变电场、磁场或能量密度。

$U(1)$由单位模复数$\ee^{\ii\theta}$组成，它有一个连续参数。$SU(n)$由$n\times n$酉矩阵且行列式为1组成，连续参数数目$n^2-1$。\emph{$SU(1)$只有单位元素，不能代替$U(1)$}。非Abel群的生成元不对易，导致规范场强含$[A_\mu,A_\nu]$并允许规范玻色子自相互作用；这本身不让玻色子获得质量。没有破缺的QCD仍有无质量的拉格朗日胶子，而Abel规范场也可通过Higgs机制获得质量。原代课口述将质量与非Abel性联系过紧，本卷按动力学区别重写。

标准模型局部规范结构是$SU(3)_c\times SU(2)_L\times U(1)_Y$，局部规范场共$8+3+1$种。弱$SU(2)$在内部空间作用，不等同于真实三维空间旋转；$SU(2)$是$SO(3)$的双覆盖而不是同一个群。Higgs的“双重态”指弱同位旋的两个复分量，两分量都是Lorentz标量，不能称为自旋$1/2$的时空旋量。

\originalsection{电弱Higgs机制中的势与质量}
设Higgs弱双重态$H=(H_1,H_2)^T$，超荷$Y=1/2$。最简单势为
\[
 V(H)=-\mu_H^2 H^\dagger H+\lambda(H^\dagger H)^2+V_0,
 \quad \mu_H^2,\lambda>0.
\]
令$x=H^\dagger H$，驻点$-\mu_H^2+2\lambda x=0$给$x=v^2/2$，其中$v^2=\mu_H^2/\lambda$。可选规范使最低能背景$H=(0,v/\sqrt2)^T$，并将径向涨落写为$(0,(v+h)/\sqrt2)^T$。二阶势给$m_h^2=2\lambda v^2$。

协变导数$D_\mu=\partial_\mu-\ii g\tau^aW^a_\mu/2-\ii g'YB_\mu$，将背景代入$|DH|^2$，得到
\[
 -\mathcal L_{\rm mass}=
 \frac{g^2v^2}{4}W^+_\mu W^{-\mu}
 +\frac{v^2}{8}(gW^3_\mu-g'B_\mu)^2,
\]
这里使用负时间号度规，质量拉格朗日项为所列二次型的负值；质量平方由系数确定。故$m_W=gv/2$、$m_Z=v\sqrt{g^2+g'^2}/2$；正交组合$A_\mu=(g'W^3_\mu+gB_\mu)/\sqrt{g^2+g'^2}$无质量。残余电磁生成元是$Q=T_3+Y$，对所选背景作用为零。准确结构是$SU(2)_L\times U(1)_Y$到$U(1)_{\rm em}$，不能把超荷$U(1)$直接说成未混合的光子。

带Yukawa耦合的费米子有$m_f=y_fv/\sqrt2$，其相对论能谱的静止能不再为零。电荷量子数由群表示规定，不是在Higgs场冷却时从无到有随意创造；左手电子与中微子同处一个弱双重态，也不表示所有场和相互作用完全相同。最小标准模型没有右手中微子场，通常的可重整化Higgs耦合因而不给活跃中微子质量；振荡说明需要额外结构。本卷不把特定质量机制当成已经实验选定。

\originalsection{热相变与大统一假设的区别}
高温有效势会含热修正，例如$c_TT^2H^\dagger H$。当正热项大于零温负二次项时，最低点可接近零；降温后非零径向背景成为有利状态。这是一个解释结构的简化势，电弱转变的实际阶数与模型参数、涨落和规范处理有关，不能由图中一条曲线认定真实宇宙必然先形成尖锐泡壁。

第21--23讲讨论$SU(5)$大统一。其生成元数目24，标准模型可嵌入其中；超荷方向可写成与
\[
 \operatorname{diag}(-1/3,-1/3,-1/3,1/2,1/2)
\]
成比例的无迹矩阵。这展示颜色三维块、弱二维块和一个Abel方向怎样共存于更大的群，不是大统一已存在的证明。最小$SU(5)$的高能破缺通常涉及与电弱双重态不同的Higgs表示。规范耦合随能标的演化可能提示统一；超对称粒子会改变演化斜率，但是否存在相应粒子和高能理论须独立检验。课程$10^{16}$ GeV指假设的大统一尺度，不能误写成已确立电弱Higgs转变的温度。

\originalsection{什么样的真空空间允许拓扑缺陷}
若一套场的势最低点组成空间$\mathcal M$，不同区域可选择不同最低能配置。场往往能连续调整以消除差异，只有不能连续收缩的边界配置才强迫留下稳定缺陷。离散不连通最低点可给畴壁；绕圆一周不能缩掉的环路可给弦；包围点缺陷的二球映射若不能缩为一点，可给单极，分别由$\pi_0(\mathcal M),\pi_1(\mathcal M),\pi_2(\mathcal M)$描述。

这些拓扑条件决定是否\emph{容许}缺陷，生成丰度还需相变动力学。两个区域随意有不同Higgs相位，并不自动证明其中夹着单极。电弱Higgs双重态最低点的情形也不能无条件套用大统一单极结论。课程示意中的畴壁、宇宙弦与磁单极是不同破缺模型的可能结果，不是所有转变的统一三步演化过程。

一个可看清单极的例子是$SU(2)$规范场与实伴随三重态$\Phi^a$，势最低点$\Phi^a\Phi^a=v_G^2$。远离核心的“刺猬”边界$\Phi^a\simeq v_Gx^a/r$把空间二球一对一映到内部二球，绕数为1；若场处处连续且从不离开最低点，就无法将此边界延伸到实心球中心，故核心需偏离最低势。这个绕数条件才是非平凡结的原因。

单独标量的角方向梯度$\partial\Phi\sim v_G/r$，远处能量$\int r^2|\partial\Phi|^2\dd r\sim v_G^2\int\dd r$发散。$1/r$是\emph{角向}梯度，不意味着场的模长必随$\ln r$增长。规范场可使$D_i\Phi$在远处趋零，并留下一条$B_r\propto r^{-2}$的残余磁场；其远区能量$\int r^2B^2\dd r$收敛。完整有限核心轮廓与非Abel单极解需解耦合场方程，这里给出拓扑和有限能量的入口，不声称已证明一般GUT的全部解。

\originalsection{磁荷与质量的量级}
Gaussian单位下，若最小相关电荷为$q$，Dirac量子化给$qg_m=n\hbar c/2$。取$q=e,n=1$作为课程约定，则
\[
 \frac{g_m}{e}=\frac1{2\alpha_{\rm em}}\simeq68.5.
\]
两粒子在同一距离下磁Coulomb力的大小比单位电荷力约大$68.5^2$；异号吸引，同号排斥。含分数电荷、具体群的整体结构和磁荷谱时须重新检查最小允许$n$，不能不加条件地把所有GUT单极都等同于这一最小Dirac荷。

典型平滑GUT单极的质量能量量级为$M_Xc^2/\alpha_G$，其中$M_X$是相关重规范玻色子的质量，$\alpha_G$是高能规范耦合；它不一定等于今日的电磁精细结构常数。课程以$10^{16}$ GeV和约$10^2$倍作粗估，得$10^{18}$ GeV的质量能量，相当于约$1.8$微克。它强调“一个粒子也可宏观重”的可能性，不是已测到这种粒子的质量。

\originalsection{过多遗迹如何成为宇宙学问题}
在非暴胀的热相变里，场的关联长度$\xi$不能远超过当时因果尺度。若按每关联体积约一个单极估计，$n_M\sim\xi^{-3}$；取最大的$\xi\sim c/H$给粗略的低丰度估计。自然单位下，辐射时期$H=1.66\sqrt{g_*}T^2/M_P$、$s=2\pi^2g_{*s}T^3/45$，故
\[
 Y_M=\frac{n_M}{s}\sim10.4\frac{g_*^{3/2}}{g_{*s}}
 \left(\frac T{M_P}\right)^3.
\]
忽略后续湮灭且保持熵史时，$\Omega_Mh^2\simeq2.74\times10^8(M_Mc^2/\mathrm{GeV})Y_M$。假设$T=10^{16}$ GeV、$g_*=g_{*s}=100$和$M_Mc^2=10^{18}$ GeV，量级远大于1，约$10^{19}$。具体数值随模型和关联长度改变；原课$10^{20}$是数量级警示，并非精密丰度预测。

单极形成后若经历$N$个e折，数密度被稀释$\ee^{-3N}$；场的势能却可近恒定，并在结束时转为新的热浴。它解释暴胀为何可移走旧遗迹又保留普通物质。不过再加热若再次超过同一相变的产生阈值，问题可能重现，需规定模型热历史和残余湮灭。下面将这一约束与暴胀的结束机制合起来讨论。

\originalchapter{暴胀的退出、再加热与初始条件边界}
本章对应8.286第1--2、23讲，补充暴胀章的慢滚方程。第23讲的实际逐字稿已读完；幻灯片中Max Tegmark为Guth与Kaiser制作的WMAP曲线属于Science的例外图片，本卷不复制，也不把逐字稿误称的“Planck七年数据”沿用为图的来源。本章把模型机制、数量级例值和观测判断分开。

\originalsection{旧暴胀为何有加速却难以顺利结束}
设标量势在$\phi=0$有高于全局最低点的局部极小，场暂时陷于其间。若梯度和动能小，$p\simeq-V$、$H^2\simeq V/(3\bar M_P^2)$，可产生指数展开。经典场不能从最低的局部点自动爬过势垒；量子隧穿却可在局部形成真空泡，随后泡壁膨胀。隧穿速率来自势形和量子场论的反弹计算，不能只用势垒的高度唯一确定。

旧模型的困难是需要把能量转成足够均匀的热宇宙。若泡形成太慢，背景假真空继续迅速膨胀，泡难以并合填满空间；若太快，又难持续足够多的e折。即使泡壁碰撞可转化能量，空间强烈非均匀也未自动消失。这称平滑退出问题，指这一具体机制失败，而不是指数膨胀的微分方程本身无解。

\begin{center}\begin{tikzpicture}\begin{axis}[width=.47\textwidth,height=4.3cm,xlabel={$\phi$},ylabel={$V$},xmin=-1.25,xmax=1.25,ymin=0,ymax=.5,title={假真空与势垒}]
\addplot[structurecolor,domain=-1.25:1.25,samples=160]{(x^2-1)^2*(x^2+.2)};
\addplot[only marks,mark=*] coordinates {(0,.2)};
\end{axis}\end{tikzpicture}\hfill
\begin{tikzpicture}\begin{axis}[width=.47\textwidth,height=4.3cm,xlabel={$\phi$},xmin=-1.25,xmax=1.25,ymin=0,ymax=1.2,title={平缓山顶与连续滚动}]
\addplot[structurecolor,domain=-1.25:1.25,samples=160]{(1-x^2)^2};
\addplot[only marks,mark=*] coordinates {(0,1)};
\end{axis}\end{tikzpicture}\end{center}
图中两势是自绘多项式示意，纵横尺度任意，右图在原点的曲率是否足够小仍取决于有量纲参数。图形“较平”本身不能代替$|\eta_V|\ll1$的检查。

\originalsection{新暴胀把连续滚动与膨胀同时利用}
若势在一段区间足够平缓，场在小涨落推动后缓慢滚动，许多e折可发生在已近均匀的一个关联区域形成之后。区域被拉大，而退出通过场滚到陡区和最低点完成，避免完全依靠独立真空泡的并合。山顶是经典不稳定点，称它“暂时近真空”是在描述停留时间长，不意味着它是严格局部稳定最低点。

近山顶写$V\simeq V_f-m_t^2\phi^2/2$，$H$近常数。线性场方程为$\ddot\phi+3H\dot\phi-m_t^2\phi=0$，增长指数为
\[
 s_+=\frac{-3H+\sqrt{9H^2+4m_t^2}}2\simeq\frac{m_t^2}{3H}
 \quad(m_t^2\ll H^2).
\]
所以离开山顶时间$\sim3H/m_t^2$可远长于$H^{-1}$。这一推导具体说明“小二阶导数”何以带来长停留；仅知道$V_f$大不能保证持续时间。实际场不精确停在$\phi=0$，初始涨落和之后势形决定总e折。

\originalsection{慢滚结束后怎样产生热宇宙}
规范场在二次最低势$V\simeq m^2\phi^2/2$附近振荡。若$H\ll m$，一周期平均的动能和势能相等，故平均$w\simeq0$；在此阶段能量随$a^{-3}$稀释，而非立即变成$a^{-4}$辐射。更一般$V\propto|\phi|^p$且周期平均有效时，场的virial关系给$\langle w\rangle=(p-2)/(p+2)$。

一个简化扰动衰变模型引入速率$\Gamma_\phi$：
\[
 \dot\epsilon_\phi+3H\epsilon_\phi=-\Gamma_\phi\epsilon_\phi,
 \qquad \dot\epsilon_r+4H\epsilon_r=\Gamma_\phi\epsilon_\phi.
\]
交换项相消保证两部门总局部守恒。$H\sim\Gamma_\phi$附近若辐射已主导并快速热化，令$\epsilon_r=\pi^2g_*T_{\rm reh}^4/30$，得到
\[
 T_{\rm reh}\sim\left(\frac{90}{\pi^2g_*}\right)^{1/4}
 \sqrt{\Gamma_\phi\bar M_P}.
\]
共振预热、非扰动产生和热化延迟可能使此模型不足；最大瞬时温度也不一定等于这里定义的再加热温度。再加热不能超过可用势能预算，不能仅因课程画$10^{16}$ GeV能标便断言真实再加热温度已知。

\originalsection{课程的大统一数量级例值怎样计算}
恢复常数，用假设能标$E_G=10^{16}$ GeV作维度分析，能量密度$\epsilon_f\sim E_G^4/(\hbar c)^3$，质量等价密度$\rho_f=E_G^4/(\hbar^3c^5)$。代入$H_f=\sqrt{8\pi G\rho_f/3}$得
\[
 \rho_f\sim2.3\times10^{81}\ \mathrm{g\,cm^{-3}},\quad
 H_f^{-1}\sim2.8\times10^{-38}\ \mathrm s,\quad
 \frac c{H_f}\sim8.4\times10^{-28}\ \mathrm{cm}.
\]
第23讲一度把最后的时间写作$H_f$，随课堂提问已改为$H_f^{-1}$；本卷沿用修正后的单位。它们是\emph{假设该能标}所得样例，暴胀能标和初始时刻不能由这些数字独立测定。

从$10^{-27}$ cm扩大到1 cm需要$N\simeq\ln10^{27}\simeq62$个e折，若到10 cm约64。与原导论“约100次加倍”是同一量级，因为$100\ln2\simeq69$；具体尺寸选择不同不应冒充矛盾或精密预测。暴胀后还要接热膨胀：无净熵产生时
\[
 \frac{a_0}{a_{\rm reh}}=
 \frac{T_{\rm reh}}{T_{\gamma0}}
 \left(\frac{g_{*s,\rm reh}}{g_{*s,0}}\right)^{1/3}.
\]
不能漏掉自由度变化，更不能把$T_{\rm reh}$固定为$E_G/\kb$。这样才把一个初始因果区、e折数与今日可见区域连接起来。前面视界章的50--60 e折是指定观测模的常用估计，和任意样例区域的总e折定义并不完全相同。

\originalsection{无毛图像有明确的适用条件}
在近de Sitter的展开期，普通物质密度$a^{-3}$、辐射$a^{-4}$、空间曲率相对项$a^{-2}$均快速减小。若没有显著各向异性应力，小的均匀剪切满足$\dot\sigma+3H\sigma=0$，其能量尺度随$a^{-6}$。这说明为何许多已开始展开、真空能占优的区域趋于局部均匀和各向同性。

这种论证不证明任意初始几何必然暴胀。过大的梯度、黑洞式坍缩、特定闭合初态或不满足所用能量条件的物质都可能改变结局。正曲率切片被拉大后局部越来越平，却不因此改变其整体拓扑。de Sitter还可有不同空间切片的坐标描述，“空间平直的指数模型”与“时空完全无曲率”尤其不能混淆。

若未来始终指数展开，光从$t_i$出发的最大共动距离是$c/(a_iH)$，实际事件视界为$c/H$。它可限制外界信号对已给定de Sitter区域的影响。但有限暴胀后若接上永远减速模型，完整未来积分可能发散，从而没有永久事件视界；不能拿有限阶段的指数近似证明该区域永远隔绝任意外部变化。

\originalsection{永恒暴胀与观测成功之间还隔着什么}
局部退出不一定使\emph{全局}膨胀体积减少。作为示意，若假真空在单位时间的有效退出概率率为$\kappa$，粗略未退出体积$V_f\propto\ee^{(3H-\kappa)t}$；$3H>\kappa$可使其平均增大。这未包括泡间相关、边界和随机几何，是说明竞争的模型而非一般证明。慢滚场若一个Hubble时间的量子涨落$H/(2\pi)$足以和经典滚动$|\dot\phi|/H$竞争，也可能产生自再生区，但需检查势和随机过程。

因此永恒暴胀是某些模型及条件下的可能结果，不由“足够e折”这个条件自动推出。多真空、初始波函数、从无隧穿和景观测度均是第1--2、23讲提出的高阶方向；本卷给其物理问题和前置关系，保留完整量子宇宙学与随机测度为待补范围。

CMB的近标度不变、近绝热扰动以及峰结构与一类暴胀初始条件相容。位置与相对高度还依赖重子、暗物质、几何和散射史；归一振幅往往需拟合势的参数。课程展示的六参数拟合不能说成所有点均为无参数预测，也不能由平均谱与角功率谱的成功选出唯一暴胀势或证明多重宇宙。真实数据与模型对应关系，须和CMB章、热史章及距离拟合一同理解。

\originalchapter{密度扰动为何增长，以及何时停止增长}
本章对应8.902第12讲末、第13讲和第14讲的增长部分。星系不是均匀Friedmann解本身的产物，而来自小的密度起伏。我们先在适用的尺度上从流体方程推导增长，再说明辐射、本动速度和规范的限制。原笔记把一个启发式辐射方程延伸到视界外，本卷保留其物理主题，改用有明确适用条件的论述。

\originalsection{在膨胀背景中定义扰动}
令$\rho(\mathbf x,t)=\bar\rho(t)[1+\delta(\mathbf x,t)]$，$\mathbf r=a\mathbf x$，物理速度$\mathbf v=H\mathbf r+\mathbf u$。$\delta$是相对于同一时刻背景的密度对比，不等于绝对密度差；$\mathbf u$是本动速度。对非相对论、亚Hubble尺度、弱引力势的流体，原物理坐标方程变换到共动坐标并减去背景后给出
\begin{align}
 \dot\delta+\frac1a\boldsymbol\nabla\!\cdot[(1+\delta)\mathbf u]&=0,\label{eq:astro-pert-cont}\\
 \dot{\mathbf u}+H\mathbf u+\frac1a(\mathbf u\cdot\boldsymbol\nabla)\mathbf u
 &=-\frac1a\boldsymbol\nabla\phi-\frac{\boldsymbol\nabla p}{a\rho},\label{eq:astro-pert-euler}\\
 \nabla^2\phi&=4\pi Ga^2\bar\rho\delta.\label{eq:astro-pert-poisson}
\end{align}
梯度作用于$\mathbf x$。$H\mathbf u$可理解为自由粒子的本动物理动量随$a^{-1}$衰减。引力加速度是$-\boldsymbol\nabla\phi/a$；原笔记的正号与其$\nabla^2\phi=4\pi G\rho$不一致。

暗物质原则上由无碰撞Boltzmann方程描述；上述流体式是其矩方程，在速度分布足够冷、尚未壳层交叉时可近似为单流尘埃。发生多流时，平均速度仍存在，但不能忽略速度弥散应力，也不能把不同流简单当成一个无压强速度场。

\originalsection{线性方程的两个Hubble阻尼因子}
保留$|\delta|\ll1$和$\mathbf u$的一阶项，取绝热压强扰动$\delta p=c_s^2\bar\rho\delta$，有$\dot\delta=-\nabla\cdot\mathbf u/a$。求时间导数时，一个$H$来自$1/a$，另一个来自Euler方程中的$H\mathbf u$，代入Poisson式得
\begin{equation}\label{eq:astro-growth}
 \ddot\delta+2H\dot\delta-4\pi G\bar\rho\delta
 =\frac{c_s^2}{a^2}\nabla^2\delta.
\end{equation}
这一步不能漏掉共动梯度前的$a^{-1}$，否则会错误地得到只有一次$H$阻尼。对
$\delta(\mathbf x,t)=\int\delta_{\mathbf k}(t)\ee^{\ii\mathbf k\cdot\mathbf x}\dd^3k/(2\pi)^3$，方程成为
\[
 \ddot\delta_{\mathbf k}+2H\dot\delta_{\mathbf k}
 +\left(\frac{c_s^2k^2}{a^2}-4\pi G\bar\rho\right)\delta_{\mathbf k}=0.
\]
物理波数为$k/a$。瞬时Jeans分界$k_{J,\rm phys}=\sqrt{4\pi G\bar\rho}/c_s$，物理波长$\lambda_J=c_s\sqrt{\pi/(G\bar\rho)}$。膨胀使系数随时间变化，不能把增长时间始终当成静态指数$\ee^{t\sqrt{4\pi G\rho}}$。

\originalsection{物质主导时期的两个独立解}
Einstein--de Sitter背景有$H=2/(3t)$、$\bar\rho=1/(6\pi Gt^2)$。忽略压强，试$\delta\propto t^n$，式\eqref{eq:astro-growth}给
\[
 n(n-1)+\frac43n-\frac23=0,
 \qquad n=\frac23,-1.
\]
所以增长模$D_+\propto a$，衰减模$D_-\propto t^{-1}=a^{-3/2}$。线性叠加的系数由初始密度与速度共同决定；只指定$\delta(t_i)$不足以唯一决定增长。常用增长因子$D(a)$约定$D(1)=1$，不是把每个模的今天振幅都设为1。

暗能量或曲率开始影响背景时，冷物质自引力项仍含$\bar\rho_m$，而$H$包含所有背景组分。用$a$作自变量，记$D_a=\dd D/\dd a$，得
\[
 D_{aa}+\left(\frac3a+\frac{E_a}{E}\right)D_a
 -\frac{3\Omega_{m,0}}{2a^5E^2}D=0.
\]
在忽略辐射、包含物质、曲率和恒定$\Lambda$的模型中，$E(a)$本身是一个独立解。降阶法使另一个解为$E(a)\int\dd a/(a^3E^3)$，取早期$D_+\simeq a$的归一化有
\[
 \widetilde D_+(a)=\frac{5\Omega_{m,0}}2E(a)\int_0^a\frac{\dd a'}{a'^3E^3(a')},
 \qquad D(a)=\frac{\widetilde D_+(a)}{\widetilde D_+(1)}.
\]
正$\Lambda$占优势时，密度下降而$H$趋向常数，增长逐渐冻结。这一积分在有不可忽略辐射或一般时变暗能量时不直接适用，应解原方程。

\originalsection{光子与重子为何会振荡}
复合前，Thomson散射将电子与光子紧耦合，Coulomb作用又将电子与离子耦合。对绝热扰动$\delta\rho_b/\rho_b=(3/4)\delta\rho_\gamma/\rho_\gamma$，定义$R_b=3\rho_b/(4\rho_\gamma)$，总压强由辐射给出$\delta p=c^2\delta\rho_\gamma/3$，于是有效声速
\[
 c_{s,\gamma b}^2=\frac{c^2}{3(1+R_b)}.
\]
只有重子惯性可以忽略时才是$c/\sqrt3$。在因果联系内，光子压强与引力竞争，使扰动进行声学振荡；重子不再能像冷尘埃那样自由增长。复合解除光子压强支持后，气体声速下降到公里每秒量级，能进入已有暗物质势阱。

若把Jeans质量定义为\emph{直径为}$\lambda_J$的均匀球质量，则$M_J=(\pi/6)\rho\lambda_J^3$。原课以半径$\lambda_J$定义，故相差8倍；课程的$10^{16}\Ms$和$10^5\Ms$应理解为相应物理阶段的尺度估计，不能在未统一球半径和温度的情况下比较系数。复合后气体在一段时间内仍通过残余电子的Compton散射与CMB保持热耦合，后来绝热冷却趋于$T_b\propto a^{-2}$，并非一直$T_b=T_\gamma\propto a^{-1}$。

在紧耦合不是无限强时，光子走随机游走，扩散物理距离约为$\sqrt{\lambda_{\rm mfp}ct}$；越近脱耦，平均自由程越大。光子携带热量和动量离开密度峰，阻尼小尺度声学扰动，这称Silk阻尼。更精确的共动阻尼积分要同时计入剪切黏性、热传导和随时间变化的电离率；课程的Silk质量约$10^{12}\Ms$只是数量级，不表示所有更小质量对象永远不能后来形成。

\originalsection{冷暗物质在辐射时期的增长为何受抑制}
冷暗物质不受光子压强直接支持，但其密度在辐射主导背景中较小。对亚Hubble冷物质扰动，若在声学振荡平均后忽略辐射密度扰动的引力源，仍用式\eqref{eq:astro-growth}，只是$H^2\propto\bar\rho_r+\bar\rho_m$而自引力只含$\bar\rho_m$。令$y=a/a_{\rm eq}$可整理为
\begin{equation}\label{eq:astro-meszaros}
 D_{yy}+\frac{2+3y}{2y(1+y)}D_y-\frac{3}{2y(1+y)}D=0.
\end{equation}
一个解是$D_1=1+3y/2$，另一个可由降阶得到
\[
 D_2=\left(1+\frac32y\right)
 \ln\frac{\sqrt{1+y}+1}{\sqrt{1+y}-1}-3\sqrt{1+y}.
\]
在$y\ll1$时两解约为常数与$\ln(4/y)-3$；$y\gg1$时增长解趋于$a$。进入视界后的辐射势源与初始匹配通常造成冷暗物质的缓慢对数增长；原课的“停滞”是相对于物质时期$a$增长的概括，不是精确地处处常数。热或暖暗物质还会自由流动擦除一定尺度以下的起伏，不能把所有暗物质都当成无任何小尺度截止的CDM。

\originalsection{视界外的密度增长为何依赖规范}
在视界外，牛顿Poisson方程与瞬时引力解释不能直接外推。改变时间切片$t\mapsto t+T(\mathbf x)$会把$\delta\rho$变为$\delta\rho-\dot{\bar\rho}T$，所以$\delta$自身依赖规范。课程所写的辐射时期$\delta\propto a^2$、物质时期$\delta\propto a$可用于合适的共动密度变量；它们不表示远处区域在超光速交换信息，也不表示任意规范的密度都增长。

以$c=1$的共形Newton规范为例，$\dd s^2=a^2[-(1+2\Psi)\dd\eta^2+(1-2\Phi)\dd\mathbf x^2]$，$\mathcal H=a'/a$，$\theta=\ii\mathbf k\cdot\mathbf v$。无压强冷物质的守恒给
\[
 \delta_c'=-\theta_c+3\Phi',\qquad
 \theta_c'+\mathcal H\theta_c=k^2\Psi.
\]
因此辐射与势的时间变化不能通过把牛顿式的$4\pi G\rho$换成$32\pi G\rho/3$而全部处理。无剪切的自由辐射还含$\delta_\gamma'=-4\theta_\gamma/3+4\Phi'$、$\theta_\gamma'=k^2(\delta_\gamma/4+\Psi)$；有Thomson散射则加入拖曳，自由流中微子还需高阶矩。这里仅用于说明遗漏的项，未声称完成一般多组分Boltzmann数值演化。符号与规范限制另核对\href{https://arxiv.org/abs/astro-ph/9506072}{Ma与Bertschinger（1995）}的守恒式及CDM式，原论文不收录到源码包。

\originalsection{增长对星系形成的启示与限度}
只让重子从$z\sim1100$开始、以$D\propto a$增长，放大倍数只有约$10^3$。CMB的温度起伏约$10^{-5}$说明一个纯重子、标准初始条件的简单模型难以产生足够的晚期结构，已有冷暗物质势阱改善这一问题。但温度起伏不等于所有尺度上的物质$\delta$，还涉及传递函数、声学效应和投影。原课“没有CDM就不可能有结构”应限定在所讨论的引力、初始谱和背景模型中，不是关于任何替代理论的无条件定理。下一章把不同尺度的增长转成可与巡天比较的统计量。

\originalchapter{从随机密度场到功率谱与星系相关}
本章对应8.902第14讲后半和第15讲的统计部分。单个密度峰的位置由具体初始条件决定，理论与巡天更常比较峰的统计分布。学习本章需要Fourier变换、期望值和方差。为避免原笔记把无限体积的Fourier振幅平方当成已归一的谱，我们先固定约定，再连接尺度不变、传递函数和偏倚。

\originalsection{功率谱为何要定义为二点期望值}
取$\delta(\mathbf x)=\int\delta_{\mathbf k}\ee^{\ii\mathbf k\cdot\mathbf x}\dd^3k/(2\pi)^3$，其逆变换为$\delta_{\mathbf k}=\int\delta(\mathbf x)\ee^{-\ii\mathbf k\cdot\mathbf x}\dd^3x$。实场满足$\delta_{-\mathbf k}=\delta_{\mathbf k}^*$。均匀性使不同Fourier波数的相关具有Dirac函数，各向同性使相关系数只依赖$k=|\mathbf k|$，故定义
\begin{equation}\label{eq:astro-power-definition}
 \langle\delta_{\mathbf k}\delta_{\mathbf k'}^*\rangle
 =(2\pi)^3\delta_D^{(3)}(\mathbf k-\mathbf k')P(k).
\end{equation}
括号是随机场的集合平均；在遍历性条件下可用足够大区域的空间平均估计。$\delta$无量纲，$\delta_{\mathbf k}$有体积量纲，故$P$也有体积量纲。有限周期盒内，对上述未除体积的离散变换，有$P(k)=\langle|\delta_{\mathbf k}|^2\rangle/V$。原笔记省略$V$或Dirac函数所写的$P=|\hat\delta|^2$不能跨不同体积比较。

若平均密度定义正确，$\langle\delta\rangle=0$；不是不除体积的$\int\delta\dd^3x$在任何有限巡天里必然等于0。有限区域的平均本身受大尺度起伏影响，这是估计相关函数时的积分约束问题。

\originalsection{方差与相关函数是一套谱的两种读法}
同一位置的方差由式\eqref{eq:astro-power-definition}得
\[
 \langle\delta^2\rangle=\int\frac{\dd^3k}{(2\pi)^3}P(k)
 =\frac1{2\pi^2}\int_0^\infty k^2P(k)\dd k.
\]
定义每单位对数波数的无量纲功率$\Delta^2(k)=k^3P(k)/(2\pi^2)$，则方差为$\int\Delta^2\dd\ln k$，若积分发散还须指定平滑或物理截止。这个$\Delta^2$与密度对比$\delta$不同，不能把$P$本身叫作无量纲振幅。

二点相关函数$\xi(r)=\langle\delta(\mathbf x)\delta(\mathbf x+\mathbf r)\rangle$为
\begin{equation}\label{eq:astro-correlation}
 \xi(r)=\frac1{2\pi^2}\int_0^\infty k^2P(k)\frac{\sin kr}{kr}\dd k.
\end{equation}
这里角积分$\int\ee^{\ii kr\cos\theta}\dd\Omega=4\pi\sin(kr)/(kr)$把三维Fourier变换化成一维积分。$\xi$可以为负，表示该距离的配对比随机更少；非负$P$则来自随机场协方差的正定性，两者并不矛盾。

若星系数密度场$n_g=\bar n_g(1+\delta_g)$，对相互不重合的小体积，联合计数概率为
\[
 \dd\mathcal P=\bar n_g^2[1+\xi_{gg}(r)]\dd V_1\dd V_2.
\]
原式里的一阶项是在平均后消失，而不是每个位置的$\delta$都为零。零间距离散对象还带Poisson接触项；现实估计需要随机目录与相同选择函数。

\originalsection{指定质量尺度就必须指定平滑窗}
把密度场与归一窗$W_R$卷积，$\delta_R(\mathbf x)=\int W_R(\mathbf y)\delta(\mathbf x-\mathbf y)\dd^3y$，$\int W_R\dd^3y=1$。Fourier空间乘上$\widetilde W(kR)$，于是
\begin{equation}\label{eq:astro-smoothed-variance}
 \sigma_R^2=\frac1{2\pi^2}\int_0^\infty k^2P(k)|\widetilde W(kR)|^2\dd k.
\end{equation}
半径$R$的实空间均匀球窗为$3/(4\pi R^3)$，球外零，积分角度与半径可得
\[
 \widetilde W_{ \rm TH}(q)=3\frac{\sin q-q\cos q}{q^3},\qquad
 \widetilde W_{ \rm TH}(0)=1.
\]
Gaussian窗$W_R\propto\exp[-r^2/(2R^2)]$则给$\widetilde W=\exp(-k^2R^2/2)$。不同窗的$R$不是同一种有效质量，需声明映射。采用球窗时通常定义$M=4\pi\bar\rho_{m,0}R^3/3$，$R$是共动半径。

若$P=Ak^n$，令$q=kR$，有$\sigma_R^2\propto R^{-(n+3)}$，故
\[
 \sigma_R\propto R^{-(n+3)/2}\propto M^{-(n+3)/6}.
\]
球窗需$-3<n<1$保证红外与紫外均收敛；$n=1$的纯幂律仍有紫外对数问题。课中的硬截止$k_{\max}\sim1/R$是尺度估计，不能代替指定窗的精确系数。实际传递函数和非线性尺度改变全范围纯幂律。

\originalsection{Harrison--Zel'dovich谱的尺度不变究竟指什么}
“所有尺度相同振幅”须说明变量和比较时刻。原课的论证比较的是\emph{每个尺度进入Hubble半径时}的密度振幅，不是固定时刻$P_\delta(k)$为常数。对在长波共动规范下$P_\delta\propto k^n$的共同初始切片，$\sigma_i(M)\propto M^{-(n+3)/6}$。

辐射时期，物理Hubble尺度$\propto a^2$，其中物质质量$M_H\propto\rho_m(a)(a^2)^3\propto a^3$。相应长波共动密度扰动增长$\propto a^2\propto M_H^{2/3}$。物质时期，物理Hubble尺度$\propto a^{3/2}$，$M_H\propto a^{3/2}$，长波增长$\propto a\propto M_H^{2/3}$。两种情况下，进入时振幅的质量依赖都是
\[
 \sigma_H(M)\propto M^{-(n+3)/6}M^{2/3}=M^{(1-n)/6}.
\]
要求它近乎不依赖$M$给$n\simeq1$。原笔记中物质时期$\sigma\propto a^{3/2}$与随后的$(1+z)^{-1}$不一致，正确增长为$a$。该论证要与上一章的规范说明一起读，不能说视界外因果作用创造了增长。

更直接的现代约定是对原始曲率扰动定义
\[
 \mathcal P_{\mathcal R}(k)=A_s(k/k_*)^{n_s-1}.
\]
$n_s=1$指每单位$\ln k$的曲率功率相同。领先阶、物质主导大尺度关系$\delta\propto k^2\mathcal R$使$P_\delta\propto k^4P_\mathcal R\propto k^{n_s}$。所以“曲率尺度不变”与“密度谱近$k^1$”可以同时成立，并不矛盾。$A_s$是特定枢轴$k_*$处的曲率振幅，不能直接当成$\sigma_8$。

\originalsection{传递函数留下相等时刻的尺度}
不同波数进入Hubble半径的时刻不同。较大$k$在辐射时期进入，冷暗物质增长受抑制；较小$k$较晚进入，少受这一抑制。定义长波极限$T(k)\to1$的传递函数，将已规定的物质谱基准写成
\[
 P_m(k,a)=A k^{n_s}T^2(k)D^2(a).
\]
$A$含曲率到物质的转换常数与选取的增长归一。转折波数约为$k_{\rm eq}=a_{\rm eq}H_{\rm eq}/c$。写$q=k/k_{\rm eq}$，无重子冷暗物质的短波渐近大致为$T\propto\ln q/q^2$，课程忽略对数后的$k^{-2}$仅是简化尺度律，必须有适当波数单位或比值。因而物质谱在大尺度近$k^{n_s}$，小尺度近$k^{n_s-4}\ln^2 k$；“暴胀幂律”并不直接等于今日密度谱。

重子的声学振荡在传递函数中产生重子声学振荡特征，中微子自由流和暗物质的微观速度也会改变它。线性阶段的尺度独立$D$并不适用于所有组分和所有模型；进入非线性后，模之间耦合，不能只乘上一个$D^2$。

\originalsection{振幅、偏倚与数据的年代边界}
定义$\sigma_8=\sigma_R$在今日球窗$R=8h^{-1}\ \mathrm{Mpc}$处的物质方差。若在某尺度上采用确定性线性偏倚$\delta_g=b\delta_m$，则$P_{gg}=b^2P_{mm}$、$\sigma_{8,g}=|b|\sigma_{8,m}$。实际星系谱还可有离散噪声与随机偏倚，Poisson近似的附加项为$1/\bar n_g$，多类对象的互相关能帮助区分偏倚与物质振幅。

课程举出的$\xi_{gg}\simeq(r/r_0)^{-1.8}$、$r_0\sim5h^{-1}\ \mathrm{Mpc}$是一定星系样本和距离范围内的经验近似；星系团更强聚集说明不同示踪物有不同偏倚，不是任意星系团样本都恒满足$\xi_{cc}=20\xi_{gg}$。原笔记给WMAP与SDSS时代的$n=0.953\pm0.016$、$\sigma_8=0.756\pm0.035$，本卷保留其\textbf{历史例示身份}，不称最新测量。测量功率谱还需处理红移空间畸变、巡天窗口和选择效应；下一章将线性方差连接到坍缩阈值及暗晕质量函数。

\originalchapter{球形坍缩、暗晕与质量函数}
本章对应8.902第15讲末与第16讲。线性增长在$|\delta|\ll1$时可靠，暗晕却处在高度非线性状态。均匀球形过密区提供一个能完整求解的桥梁：由其真实非线性演化确定坍缩时刻，再把早期扰动按线性理论外推到这个时刻，建立用于统计的阈值。阈值与真实密度必须始终分开。

\originalsection{均匀过密球的参数解}
先取Einstein--de Sitter背景、无压强尘埃、球对称且壳层尚未交叉。半径$R$内总质量$M$恒定，边缘满足
\[
 \ddot R=-\frac{GM}{R^2},\qquad \frac12\dot R^2-\frac{GM}{R}=E.
\]
过密增长模最终转向，取$E<0$、$A=GM/(2|E|)$。令$R=A(1-\cos\vartheta)$，代入能量式选膨胀分支，可得$\dd t=B(1-\cos\vartheta)\dd\vartheta$，其中$A^3=GMB^2$。因此
\begin{equation}\label{eq:astro-top-hat-param}
 R=A(1-\cos\vartheta),\qquad t=B(\vartheta-\sin\vartheta).
\end{equation}
把大爆炸时刻设为背景与扰动共同的$t=0$排除这一模型的衰减模；若另设不同起始时间，不能继续使用相同阈值。转向发生在$\vartheta=\pi$，$R_{\rm ta}=2A$、$t_{\rm ta}=\pi B$；形式坍缩在$\vartheta=2\pi$、$t_c=2t_{\rm ta}$。

背景$\bar\rho=1/(6\pi Gt^2)$，球内$\rho=3M/(4\pi R^3)$。用参数解与$A^3=GMB^2$消去维度常数得
\begin{equation}\label{eq:astro-nonlinear-density}
 1+\delta_{\rm nl}=\frac92\frac{(\vartheta-\sin\vartheta)^2}{(1-\cos\vartheta)^3}.
\end{equation}
原笔记归一变量$x=a/a_{\rm ta}$、$y=R/R_{\rm ta}$给相同结果$1+\delta=\xi(x/y)^3$，其中$\xi=9\pi^2/16$；个别原式漏掉立方。牛顿运动式中的$\xi$也不能在尚未用边界条件求出前直接替换为3。

\originalsection{1.686阈值是怎样算出的}
早期$\vartheta\ll1$时，展开
\[
 1-\cos\vartheta=\frac{\vartheta^2}{2}-\frac{\vartheta^4}{24}+\cdots,\quad
 \vartheta-\sin\vartheta=\frac{\vartheta^3}{6}-\frac{\vartheta^5}{120}+\cdots.
\]
代入式\eqref{eq:astro-nonlinear-density}得$\delta_{\rm nl}=3\vartheta^2/20+O(\vartheta^4)$。线性增长模$\delta_L\propto t^{2/3}$，用这个早期系数定其归一化后，任意参数时刻的线性外推值为
\[
 \delta_L(\vartheta)=\frac3{20}[6(\vartheta-\sin\vartheta)]^{2/3}.
\]
所以转向时$\delta_{L,\rm ta}=\frac35(3\pi/4)^{2/3}\simeq1.062$，形式坍缩时
\begin{equation}\label{eq:astro-delta-collapse}
 \delta_c=\frac35\left(\frac{3\pi}{2}\right)^{2/3}\simeq1.68647.
\end{equation}
与此相对，转向时真实密度比$1+\delta_{\rm nl}=9\pi^2/16\simeq5.55$，形式坍缩时则发散。\textbf{$\delta_c$不是暗晕的真实密度对比}，它是把早期同一扰动沿线性增长延伸所得的统计标签。

\originalsection{为什么估计的维里半径是转向半径的一半}
均匀球的势能$U=-3GM^2/(5R)$。转向时宏观径向动能为零，故$E=U_{\rm ta}$。理想化最终稳定暗晕若没有净能量流出且满足$2T+U=0$，则$E=U_{\rm vir}/2$，给
\[
 R_{\rm vir}=\frac12R_{\rm ta}.
\]
将维里化时刻近似设为形式坍缩时刻$2t_{\rm ta}$，背景密度已降低4倍，球的密度则比转向时增大8倍。因此
\[
 \Delta_{\rm vir}\equiv\frac{\rho_{\rm vir}}{\bar\rho(t_c)}
 =32\frac{9\pi^2}{16}=18\pi^2\simeq177.65.
\]
严格无压强均匀球本身在$R=0$奇点终止，\emph{并未自行产生维里化}。真实体系的壳层交叉、相混合与非球对称扰动提供松弛；上面的半径与时刻是能量加维里条件的模型估计。有限压强、表面压强、潮汐、质量交换与$\Lambda$会改变关系，不能把它当成任意暗晕定理。

常用$M_{200c}$定义平均密度为$200\rho_{\rm crit}$，$M_{200m}$则为$200\bar\rho_m$。Einstein--de Sitter里二者参考密度相同，其他背景下不同。给出$M_{200}$时须说明后缀与红移，再用$M=(4\pi/3)R^3\Delta\rho_{\rm ref}$换算，不能把200倍平均密度和200倍临界密度任意互换。

\originalsection{Gaussian尾部如何估计已坍缩的质量比例}
选定共动球窗半径$R(M)=[3M/(4\pi\bar\rho_{m,0})]^{1/3}$，将\emph{线性}密度场平滑到$M$尺度，方差$\sigma^2(M,z)=D^2(z)\sigma^2(M,0)$由式\eqref{eq:astro-smoothed-variance}计算。若初始场为Gaussian，平滑密度的概率密度为
\[
 p(\delta_M)=\frac1{\sqrt{2\pi}\sigma}\exp\left(-\frac{\delta_M^2}{2\sigma^2}\right).
\]
单点超过球形阈值的概率是
\[
 F_{\rm naive}(>M,z)=\int_{\delta_c}^{\infty}p(\delta_M)\dd\delta_M
 =\frac12\operatorname{erfc}\left(\frac{\delta_c}{\sqrt2\sigma(M,z)}\right).
\]
按初始近均匀的Lagrange区域质量来解释，这估计的是\emph{质量比例}；不是高度非均匀、已坍缩Euler空间中的几何体积分数。一个在小尺度平滑后低于阈值的点，仍可能属于更大尺度已经超过阈值的区域，这称云中之云问题。

当$M\to0$且$\sigma\to\infty$时，上式趋于$1/2$。Press--Schechter处方将它乘2，取$F=\operatorname{erfc}(\nu/\sqrt2)$，$\nu=\delta_c/\sigma$。对于锐$k$窗产生的Markov随机游走和恒定吸收边界，反射原理可导出同一因子；实空间球窗有相关步长时不再自动成立。扩展处方的具体边界与首次通过计算在公开PS6题解中另作推导。这一模型不表示所有暴胀都精确产生Gaussian场，也不意味着所有暗晕遵守单一球形阈值。

\originalsection{质量比例转为数密度时的符号与单位}
$F(>M)$随$M$增大而减小，所以$\dd F/\dd M<0$。质量在$[M,M+\dd M]$中的比例是$-F_M\dd M$，乘今日平均物质密度并除以对象质量，得\emph{共动}质量函数
\begin{equation}\label{eq:astro-press-schechter}
 \frac{\dd n_c}{\dd M}=\sqrt{\frac2\pi}\frac{\bar\rho_{m,0}}{M^2}
 \frac{\delta_c}{\sigma(M,z)}\left|\frac{\dd\ln\sigma}{\dd\ln M}\right|
 \exp\left[-\frac{\delta_c^2}{2\sigma^2(M,z)}\right].
\end{equation}
若用$\bar\rho_m(z)$代替$\bar\rho_{m,0}$，得到的是物理体积数密度，相差$(1+z)^3$。原笔记同时使用$\rho_c(z)\Omega_m(z)$和“共动”名称，需作这一转换；其$F_M$还漏负号。窗的方差中必须是$|\widetilde W|^2$并带$1/(2\pi^2)$，不能用单次$\widetilde W$或$1/(2\pi)$。

当小质量处$\sigma\to\infty$且大质量处$\sigma\to0$，积分$\int_0^\infty M(\dd n_c/\dd M)\dd M=\bar\rho_{m,0}$。有微观截止或不满足这些边界时，不能继续说该处方必然把所有质量放入暗晕。$\nu\gg1$的稀有高质量暗晕受到指数抑制，对增长因子和振幅很敏感；小质量部分的形状则还依赖$\sigma(M)$的斜率。

球形模型是解析基准。椭球坍缩、合并、环境与实际暗晕识别方法改变精确质量函数，课程展示的Sheth--Tormen、Jenkins、Tinker等曲线是各自条件下的改进或模拟标定，并非从本章简单积分直接推得。比较模型必须统一质量定义、背景宇宙学、红移和体积约定，才能判断真正的差异。

\originalchapter{微波背景的温度谱与宇宙学参数}
本章对应8.902第17--18讲。CMB近黑体的平均谱说明早期热史，天空温度起伏则同时记录初始扰动、光子重子流体和光传播几何。要从图中的一个峰解释参数，必须先明确温度场、角功率谱和声学尺度的定义；下述领先阶模型解释机制，不替代完整Boltzmann计算与观测似然分析。

\originalsection{收到的温度起伏包含哪些物理项}
令$T(\mathbf n)=\bar T[1+\Theta(\mathbf n)]$，$\mathbf n$指观测者看向天空的方向。最后散射附近，光子的本征温度、局部势引起的引力红移及流体本动速度都影响接收频率。用$c=1$、上一章的共形Newton规范作示意，并扣除观测者处的单极和偶极，有
\[
 \Theta_{\rm obs}\simeq[\Theta_0+\Psi-\mathbf n\cdot\mathbf v_b]_* 
 +\int_{\eta_*}^{\eta_0}(\Phi'+\Psi')\dd\eta.
\]
这是瞬时最后散射和无后续散射近似；完整式要对可见度函数和沿光线位置积分。$\Theta_0$是局部光子温度单极，$\Psi$是引力势项，速度项的符号随光线方向约定改变，积分项称积分Sachs--Wolfe效应。静态势中进入与离开时的频移相消，演化势则不能完全相消。

在物质主导、绝热长波、$\Phi=\Psi$且忽略衰减模时，Newton规范密度满足$\delta_m\simeq-2\Psi$、$\delta_\gamma=4\delta_m/3$，故$\Theta_0=\delta_\gamma/4=-2\Psi/3$，普通Sachs--Wolfe项为$\Theta_0+\Psi=\Psi/3$。势阱的本征温度与引力红移部分抵消；不能将其说成任意尺度上Sachs--Wolfe与Doppler两者自动抵消。光传播的时间延迟也是几何效应，但不能把所有时间延迟都当成上述温度抵消的充分解释。

复合后仍有多个二次效应。再电离产生新的Thomson散射，衰减原有较小尺度温度起伏的振幅，并生成大尺度偏振；散射具有角分布，不是精确各向同性。变化的势产生积分Sachs--Wolfe项，透镜改变天空方向并重新分配角功率。热Sunyaev--Zel'dovich效应是\emph{热电子对光子的逆Compton散射}，可使低频强度降低、高频升高，改变谱形；运动SZ来自电子整体速度。它们不能都当成最后散射面上的初始密度变化。

\originalsection{用球谐而不是平面波描述整个天空}
取归一球谐$\int Y_{\ell m}^*Y_{\ell'm'}\dd\Omega=\delta_{\ell\ell'}\delta_{mm'}$，有
\[
 \Theta(\mathbf n)=\sum_{\ell=0}^{\infty}\sum_{m=-\ell}^{\ell}a_{\ell m}Y_{\ell m}(\mathbf n),
 \qquad a_{\ell m}=\int\Theta Y_{\ell m}^*\dd\Omega.
\]
统计各向同性使$\langle a_{\ell m}a_{\ell'm'}^*\rangle=C_\ell\delta_{\ell\ell'}\delta_{mm'}$。对一张天空估计$\widehat C_\ell=(2\ell+1)^{-1}\sum_m|a_{\ell m}|^2$。$C_\ell$是随机场的期望谱，估计器的$m$平均不等于多个独立宇宙的平均。

$\ell=0$是平均温度，$\ell=1$的主导项来自太阳系相对CMB的运动，并有地球轨道的时变贡献；$\ell=2$为四极。特征角尺度大约$\vartheta\sim\pi/\ell$，只是数量级映射，不是精确的一一波长对应。

加法定理给角相关函数
\[
 C(\vartheta)=\langle\Theta(\mathbf n)\Theta(\mathbf n')\rangle
 =\sum_\ell\frac{2\ell+1}{4\pi}C_\ell P_\ell(\cos\vartheta).
\]
常绘$\mathcal D_\ell=\ell(\ell+1)C_\ell/(2\pi)$；若温度系数以$\mu\mathrm K$而不是$\Theta$表示，就还乘$\bar T^2$并换单位。课程图中的$\mu\mathrm K^2$不能与无量纲$C_\ell$直接相等。全空、Gaussian、无噪声时
\[
 \frac{\operatorname{Var}(\widehat C_\ell)}{C_\ell^2}=\frac{2}{2\ell+1}.
\]
这来自有限的独立$m$自由度，称宇宙方差；低$\ell$即使仪器完美也不能任意精确，掩膜和噪声则进一步改变误差。

\originalsection{为什么尺度不变谱产生大角度平台}
在普通Sachs--Wolfe近似下，曲率模投影到半径$D_M(z_*)$的天空，温度传递函数约为$\Delta_\ell(k)=j_\ell(k\chi_*)/5$，这里取平直空间和对应曲率符号约定。角谱
\[
 C_\ell=4\pi\int\mathcal P_{\mathcal R}(k)|\Delta_\ell(k)|^2\dd\ln k.
\]
若$\mathcal P_{\mathcal R}=A_s$恒定，利用$\int_0^\infty j_\ell^2(x)\dd x/x=1/[2\ell(\ell+1)]$（$\ell\ge1$），则
\[
 C_\ell\simeq\frac{2\pi A_s}{25\ell(\ell+1)},\qquad \mathcal D_\ell\simeq\frac{A_s}{25}.
\]
这解释大角度近平台。晚期积分Sachs--Wolfe、有限最后散射宽度、张量模和略倾斜初始谱会改变它，不应把低$\ell$平台当成在所有背景下精确成立。

\originalsection{声学峰与声视界的积分}
紧耦合流体的声速为$c/\sqrt{3(1+R_b)}$。忽略驱动和系数慢变，Fourier扰动满足声学振子方程，解的相位为$kr_s$，其中\emph{共动}声视界
\begin{equation}\label{eq:astro-sound-horizon}
 r_s(a_*)=\int_0^{a_*}\frac{c_s(a)\dd a}{a^2H(a)}
 =\int_{z_*}^{\infty}\frac{c_s(z)\dd z}{H(z)}.
\end{equation}
原课用$t_*c/\sqrt3$估物理声程，忽略了膨胀期间速度与尺度因子的积分；不能用于精确峰位。观察角为$\theta_s=r_s/D_M(z_*)$，基准声学多极$\ell_A=\pi/\theta_s$。各峰近乎位于整数$kr_s=n\pi$的压缩或稀疏状态，但势演化和相位移使第一温度峰\emph{不精确位于}$\ell_A$。课堂常说第一峰约$\ell\sim200$，而常用模型的$\ell_A$约300，两者指的是不同量。

绝热初态让不同波数的振子有相干初始相位，故投影后保留一系列峰。压缩与稀疏对应奇偶峰；速度项与密度项相差约四分之一周期，会改变实际温度谱。无限小尺度不能无限振荡，光子扩散产生近似$\exp[-(k/k_D)^2]$的振幅阻尼，有限最后散射宽度也平滑起伏。课程的约5角分是特定尺度估计，不是普遍锐截止。

\originalsection{几何、物质和重子各怎样影响谱}
固定早期物理组分和声视界时，空间几何主要通过$D_M$改变峰的角尺度：负曲率使相同尺看起来较小，峰向更大$\ell$移动；正曲率相反。但$H_0$、暗能量和曲率都影响距离，CMB存在几何简并，不能仅凭第一峰一项精确分别测出全部参数。

早期方程更直接依赖物理密度$\omega_m=\Omega_mh^2$、$\omega_b=\Omega_bh^2$。$\omega_m$改变辐射—物质相等时刻和势阱衰减：辐射占比更高时，变化的势驱动声学振子，产生辐射驱动和早期积分Sachs--Wolfe贡献。原课用“同一红移下年龄较小、来不及形成结构”解释峰高变化，是直观概括；其结论还依赖固定哪些参数，不能直接作为任意改变$\Omega_m$的单调律。

提高$\omega_b$增加$R_b$，改变声速和振子惯性。在近恒定势阱中，它还移动平衡位置：压缩加深，稀疏不能对称地离开原势阱，因而相对增强奇峰、抑制偶峰。第一与第二峰的比值提供重子加载信息，结合其余峰、倾斜谱和阻尼尾才可区分其他效应。BBN的重子约束提供独立交叉检验。

再电离光学深度$\tau$使原温度谱大部分乘$\ee^{-2\tau}$，与初始幅度$A_s$简并；偏振可帮助打破这一简并。透镜、前景和仪器波束也参与实际似然。本卷不把示意机制宣称为已实施完整参数估计。课程引用的Planck 2018温度谱是2020年发表的CC BY图，本卷仅保留来源说明，不复制图；这里的公式与机制为独立中文推导。

\originalchapter{宇宙热史、脱耦与热遗迹}
本章对应8.902第18讲末和第19讲。膨胀给一个时间尺度$H^{-1}$，反应给另一个时间尺度$\Gamma^{-1}$：比较它们，可以判断热平衡是否能跟上宇宙演化。本章先说明温度、熵和相对论自由度，再推导冷暗物质的冻结。每条方程都声明物理数密度还是共动量，以修正原课相关定义中的尺度因子问题。

\originalsection{黑体辐射的能量、数目与熵}
在频率$\nu$与$\nu+\dd\nu$内，单位体积含两个光子极化的态密度为$8\pi\nu^2\dd\nu/c^3$。每态平均占据数为$[\exp(h\nu/\kb T)-1]^{-1}$，故
\[
 u_\nu(T)=\frac{8\pi h\nu^3}{c^3}\frac1{\exp(h\nu/\kb T)-1},\quad
 u=\int u_\nu\dd\nu=\frac{\pi^2\kb^4}{15\hbar^3c^3}T^4.
\]
积分用$x=h\nu/(\kb T)$和$\int x^3/(\ee^x-1)\dd x=\pi^4/15$。各向同性光子动量通量给$p=u/3$。类似数目积分得
\[
 n_\gamma=\frac{2\zeta(3)}{\pi^2}\left(\frac{\kb T}{\hbar c}\right)^3,
 \qquad \frac{u}{n_\gamma}=\frac{\pi^4}{30\zeta(3)}\kb T\simeq2.701\kb T.
\]
光子的化学势在充分产生与吸收的热平衡中为零；对这种辐射，熵密度$s=(u+p)/T=4u/(3T)$。这里$s$有熵单位，后文取自然单位时可把$\kb$吸收入定义。

按频率与波长绘谱时最大点不同：$u_\nu$的峰满足$3(1-\ee^{-x})=x$，$x\simeq2.821$；每单位波长谱因雅可比$c/\lambda^2$变成$\lambda^{-5}/[\exp(hc/(\lambda\kb T))-1]$，峰满足$5(1-\ee^{-y})=y$，$y\simeq4.965$。不能用$\lambda=c/\nu_{\max}$把前者直接当成后者。

脱耦后无碰撞光子的相空间占据数沿轨道保持不变，物理动量$p\propto a^{-1}$。初始Planck分布$[\exp(pc/(\kb T_i))-1]^{-1}$仍可写成温度$T_i a_i/a$的Planck分布。因此\textbf{红移保持黑体形状并不要求每个时期仍在局部热平衡}。新的能量注入或不同温度混合可产生谱畸变；发生吸收与发射时也需另算碰撞项。

\originalsection{相对论自由度改变温度与时间关系}
本节起取$\hbar=c=\kb=1$，温度直接以能量表示，$1\ \mathrm{eV}$对应约$1.1605\times10^4$ K。热平衡物种的占据数为$f=(\exp[(E-\mu)/T]\mp1)^{-1}$，玻色取负、费米取正。化学势可忽略且$T\gg m$时
\[
 \epsilon=\frac{\pi^2}{30}g_*(T)T^4,\qquad
 s=\frac{2\pi^2}{45}g_{*s}(T)T^3.
\]
相同温度的每个费米自由度在能量与熵积分中相对玻色贡献$7/8$；不同温度物种在$g_*$中带$(T_i/T)^4$，在$g_{*s}$中带$(T_i/T)^3$。某组分变成非相对论性并湮灭后，它的熵转给仍耦合的粒子，不能维持所有物种的同一温度。

若没有净熵产生，$sa^3$守恒，故$T\propto g_{*s}^{-1/3}a^{-1}$。自由度不变时才简化为$T\propto a^{-1}$。辐射主导的Friedmann式给
\[
 H=1.66\sqrt{g_*}\frac{T^2}{M_P},\qquad M_P=G^{-1/2},
\]
这里$M_P$是\emph{非约化}Planck质量，与暴胀章$\bar M_P=M_P/\sqrt{8\pi}$区别明确。若$g_*$近恒定，$H=1/(2t)$，则
\[
 t=0.301\frac{M_P}{\sqrt{g_*}T^2}
 \simeq0.738\left(\frac{g_*}{10.75}\right)^{-1/2}
 \left(\frac{T}{\mathrm{MeV}}\right)^{-2}\mathrm s.
\]
自由度变化快或有非辐射主导阶段时不能用这一简单系数精确计时。

\originalsection{脱耦意味着反应不再能维持平衡}
每粒子的有效反应率$\Gamma$通常为$n\langle\sigma v\rangle$。当$\Gamma\gg H$且反应能改变相关变量时，系统可追踪平衡；当$\Gamma\ll H$，原有分布被膨胀带走而难以恢复。\emph{化学脱耦}冻结粒子总数，\emph{动理学脱耦}冻结与热浴的动量交换，两者可发生在不同时刻。

非相对论且非简并的平衡粒子有
\[
 n_{ \rm eq}=g\left(\frac{mT}{2\pi}\right)^{3/2}\exp[-(m-\mu)/T].
\]
当可与热浴湮灭且无不对称荷时取$\mu\simeq0$。降温后该密度随$\ee^{-m/T}$迅速降低，若湮灭一直足够快，粒子几乎消失；反应提前跟不上则留下遗迹。$\Gamma\simeq H$给数量级交叉点，真实偏离平衡时刻还取决于平衡丰度变化速率，不能将冻结当成瞬间关闭碰撞。

\originalsection{冷遗迹的Boltzmann方程与共动变量}
对稳定、有效自共轭的粒子，令$n$为物理数密度，反应截面约定使方程为
\begin{equation}\label{eq:astro-relic-boltzmann}
 \dot n+3Hn=-\langle\sigma v\rangle(n^2-n_{ \rm eq}^2).
\end{equation}
右侧的反向产生率由详细平衡确定。有粒子与反粒子分别计数、或并存多种共湮灭物种时，需明确相应有效截面和数目因子；本式不是任意粒子模型的全部动力学。

真正的共动粒子数密度是$N_c=na^3$。直接求导得
\[
 \dot N_c=-\frac{\langle\sigma v\rangle}{a^3}(N_c^2-N_{c,\rm eq}^2).
\]
原笔记的$n_c=n a^{-3}$及缺$a^{-3}$的共动碰撞式应修正。更便于热史的量是熵归一丰度$Y=n/s$；若$sa^3$守恒，膨胀稀释在$\dot Y$里自动消掉。

令$x=m/T$。在$g_{*s}$近恒定、辐射主导时$\dot x=Hx$，式\eqref{eq:astro-relic-boltzmann}变成
\begin{equation}\label{eq:astro-relic-y}
 \frac{\dd Y}{\dd x}=-\frac{s\langle\sigma v\rangle}{Hx}(Y^2-Y_{ \rm eq}^2),\qquad
 Y_{ \rm eq}\simeq0.145\frac g{g_{*s}}x^{3/2}\ee^{-x}.
\end{equation}
若熵自由度随温度变化，右侧还乘$1+Tg_{*s}'/(3g_{*s})$，因为$\dot x$相应减慢。额外熵注入或非标准膨胀则需重新建立两条背景方程。

\originalsection{冻结后的积分为何给出反比于截面的遗迹}
取$s$波湮灭、$\langle\sigma v\rangle$近恒定，定义
\[
 \lambda=0.264\frac{g_{*s}}{\sqrt{g_*}}M_Pm\langle\sigma v\rangle.
\]
则$s\langle\sigma v\rangle/(Hx)=\lambda/x^2$。冻结后$Y_{ \rm eq}$很小，式\eqref{eq:astro-relic-y}可积分为
\[
 \frac1{Y_\infty}=\frac1{Y(x_f)}+\frac\lambda{x_f}.
\]
若后续湮灭项主导，$Y_\infty\simeq x_f/\lambda$。所以更大的截面使粒子跟踪平衡更久、冻结温度更低，留下更少遗迹。只用$Y_\infty=Y_{ \rm eq}(x_f)$会漏掉冻结后的残余湮灭，通常不足以给可靠丰度。

常规弱相互作用量级的模型会给$x_f$约20--30，但值依赖质量、内部自由度和截面。今日密度为$\rho_{\chi,0}=m s_0Y_\infty$，对标准熵史换算有
\[
 \Omega_\chi h^2\simeq2.74\times10^8\left(\frac m{\mathrm{GeV}}\right)Y_\infty.
\]
数量级上，电弱尺度有效湮灭截面可产生所需宇宙学遗迹，这称WIMP机制的巧合；单个丰度不能唯一反推出$m$和$\sigma$，更不是已发现WIMP的证据。该例也不能排除非热产生、不对称暗物质或其他候选。

\originalsection{热、冷遗迹与自由流的尺度}
相对论性脱耦的粒子，在固定熵自由度下$Y$近恒定而不受Boltzmann指数抑制；非相对论性脱耦则受上述冻结控制。但热暗物质、暖暗物质、冷暗物质的结构效应最终取决于相关时期的\emph{速度分布和自由流尺度}，不能只看粒子第一次化学脱耦时是否相对论性。

无碰撞粒子的动量$p\propto a^{-1}$。速度$v=pc^2/\sqrt{m^2c^4+p^2c^2}$在相对论时期约为$c$，非相对论时期才随$a^{-1}$下降。共动自由流距离应积历史速度
\[
 \lambda_{\rm fs,c}(t)=\int_{t_{ \rm kd}}^t\frac{v(t')}{a(t')}\dd t'.
\]
原课用$v\sqrt{\pi/(G\rho)}$作Jeans类比给局部尺度，不能代替这个累计积分；其keV质量到几十Mpc的估计需指定热分布、组分温度和质量，不能推广到所有keV粒子。自由流擦除小尺度初始结构，大尺度可较早形成；冷暗物质则容许较小暗晕先出现并合并，但也存在微观截止，后期重子过程会改变可见星系。

\originalsection{热史中可靠计算与模型假设的分界}
核合成前后MeV温度的弱反应、中微子脱耦、电子正电子湮灭属于标准模型可计算过程。夸克到强子的转变在百MeV尺度，具体热力学不由理想气体式决定。更高温的大统一相变、Peccei--Quinn相变和Planck尺度量子引力阶段依赖相应理论是否正确，课程时间轴中的$10^{16}$、$10^{12}$ GeV不是已观测的宇宙事件。

光子热化还需区分弹性Compton散射与能改变光子数的双Compton、制动辐射等过程。原课“keV时光子脱离平衡并数目守恒”不是一个可由$T=m$判定的普遍瞬时事件。辐射—物质相等、复合及第一批天体形成也分别依赖背景密度、原子反应和结构演化。下一章利用这些明确的反应尺度推导轻元素核合成，而不把所有热史节点压成未经解释的时间表。

\originalchapter{辐射自由度、中微子温度与热史的分段}
本章补齐8.286第15--18讲的量子气体和中微子内容。热遗迹章已推导黑体能量、数密度、熵和冻结；这里解释不同统计权重怎样从积分产生，以及脱耦后为什么不能把所有粒子放进同一个温度。课程2013年的中微子参数和历史CMB谱图只作年代明确的背景，原第三方图片不收录。

\originalsection{同样多的自旋态为什么贡献不同的能量}
先取$\hbar=c=\kb=1$。具有$g$个内部态的自由粒子，相空间计数为$g\dd^3p/(2\pi)^3$。若化学势可忽略，平衡分布是$f_B=(\ee^{E/T}-1)^{-1}$或$f_F=(\ee^{E/T}+1)^{-1}$。对无质量粒子，$E=p$，故
\[
 n=\frac{gT^3}{2\pi^2}\int_0^\infty\frac{x^2\dd x}{\ee^x\mp1},\qquad
 \epsilon=\frac{gT^4}{2\pi^2}\int_0^\infty\frac{x^3\dd x}{\ee^x\mp1}.
\]
用$1/(\ee^x-1)=\sum_{j\geq1}\ee^{-jx}$，逐项积分$\int x^{s-1}\ee^{-jx}\dd x=\Gamma(s)j^{-s}$；费米展开则有交替符号。于是
\[
 \int_0^\infty\frac{x^{s-1}\dd x}{\ee^x+1}
 =(1-2^{1-s})\Gamma(s)\zeta(s).
\]
能量用$s=4$，权重$7/8$；数目用$s=3$，权重$3/4$。它们来自不同积分，不是可以随意互换的“费米修正”。压强从各向同性动量流$\int p_iv_i f\dd^3p$给$\epsilon/3$，熵$ s=(\epsilon+p)/T$因而与能量有同样的$7/8$权重。

原8.286把能量／熵有效数记作$g$，把数目有效数记作$g^*$；本卷采用$g_*$表示能量、$g_{*s}$表示熵，另记$g_{*n}$表示数目。不同温度物种贡献分别带$(T_i/T)^4$、$(T_i/T)^3$、$(T_i/T)^3$。符号转换后关系为
\[
 \epsilon=\frac{\pi^2}{30}g_*T^4,\quad
 s=\frac{2\pi^2}{45}g_{*s}T^3,\quad
 n=\frac{\zeta(3)}{\pi^2}g_{*n}T^3.
\]
质量不小于温度时需要保留$E=\sqrt{p^2+m^2}$积分，不能以阶跃的自由度表作精确热史。

\originalsection{MeV时期的自由度为什么是10.75}
光子有两个横向极化。电子和正电子各有两个自旋态，在$T\gg m_e$时共四个费米自由度。标准热背景中每个活跃中微子味有一个中微子螺旋态和一个反中微子螺旋态，三个味合计六个，所以
\[
 g_*=2+\frac78\,4+\frac78\,6=\frac{43}{4}=10.75,
 \qquad g_{*n}=2+\frac34\,4+\frac34\,6=9.5.
\]
这里假定这些物种都在近共同温度，且未额外热化其他态。中微子有质量，不意味着Dirac模型的其他螺旋态必定在早期热浴中以同样丰度产生；相互作用率仍须检查。Majorana与Dirac的粒子解释不同，但通常的活跃热中微子计数在上述近似下相同。

四费米弱作用在$T\ll m_W$时给数量级$\sigma\sim G_F^2T^2$，背景数密度$\sim T^3$，故$\Gamma\sim G_F^2T^5$。膨胀率$H\sim T^2/M_P$；随降温$\Gamma/H\propto T^3$下降，交叉在MeV量级。精确脱耦不是一个瞬间，三个味也非完全同步；这里的突变近似用于解释主效应。

\originalsection{电子湮灭怎样使光子比中微子热}
中微子脱耦后，其动量分布随$a^{-1}$平移，温度参数$T_\nu\propto a^{-1}$。电子与正电子仍和光子耦合，湮灭时将熵主要传给这一部门。湮灭前该部门的熵自由度是$2+7\times4/8=11/2$，湮灭后为2。若两端充分相对论／完全湮灭且无额外熵交换，部门熵守恒给
\[
 \frac{11}{2}(a_iT_i)^3=2(a_fT_{\gamma,f})^3,
 \qquad a_fT_{\nu,f}=a_iT_i,
 \qquad \frac{T_\nu}{T_\gamma}=\left(\frac4{11}\right)^{1/3}.
\]
因而$T_\gamma=2.725$ K时$T_\nu\simeq1.945$ K。这个“温度”在中微子变成非相对论性后是描述遗留动量分布的参数，不能再把它说成与同温物质反复碰撞形成的平衡温度。

每个味含粒子和反粒子的数密度相对光子为
\[
 \frac{n_{\nu_i}+n_{\bar\nu_i}}{n_\gamma}
 =\frac{(3/4)\,2}{2}\left(\frac4{11}\right)=\frac3{11}.
\]
光子$n_\gamma\simeq411\ \mathrm{cm^{-3}}$，所以每味约112$\ \mathrm{cm^{-3}}$。若仍相对论性，三味能量使
\[
 \frac{\epsilon_r}{\epsilon_\gamma}
 =1+\frac78\left(\frac4{11}\right)^{4/3}3\simeq1.681.
\]
实际有限时长的脱耦及有限温度修正通常用$N_{\rm eff}$写为$1+(7/8)(4/11)^{4/3}N_{\rm eff}$，而不是把非整数$N_{\rm eff}$当作已经发现了非整数个粒子种类。本卷不在突变模型中冒充已算出高精度修正。

\originalsection{中微子质量与背景动力学的联系}
当$p\ll m_\nu$时，$\epsilon_\nu\simeq\sum_i m_i n_i$，以标准遗留数密度除以今日临界密度可得近似
\[
 \Omega_\nu h^2\simeq\frac{\sum_i m_ic^2}{94\ \mathrm{eV}}.
\]
系数取决于标准温度和脱耦约定。这不是质量测量；它把某个假定质量和热史转成宇宙密度。振荡测得质量平方差，不能独自确定绝对质量和最轻本征态。第18讲所列约$7.5\times10^{-5}$与$2.3\times10^{-3}\ \mathrm{eV^2}$为2013年例值，不作当前最佳拟合声明。

非相对论背景中微子不应同时算在$a^{-4}$的辐射密度和$a^{-3}$的物质密度里。精确转换期应计算
\[
 \epsilon_\nu(a)=\frac{g}{2\pi^2a^4}\int_0^\infty
 q^2\sqrt{q^2+a^2m^2}\,f(q)\dd q,
\]
其中$q=ap$为共动动量。高温与低温极限分别给两种尺度律，过渡连续。自由流还能抑制小尺度物质增长，密度记账与扰动效应是两个相关却不同的计算。

\originalsection{时间、温度与可见的历史证据}
在$g_*=10.75$近不变的辐射时期，$t\simeq0.738(T/\mathrm{MeV})^{-2}$ s，即$T\simeq0.859/\sqrt{t/\mathrm s}$ MeV，与第18讲的0.860一致。电子湮灭后$g_*$与光子温度关系改变，需要另算系数；复合附近又已不是纯辐射主导。把一个温度—时间式从MeV一直外推到今天会混淆这些阶段。

第18--19讲回顾1975年的早期谱数据、1987年火箭结果与1990年COBE/FIRAS黑体谱。其物理意义是宽频谱与同一黑体温度相容，不是所有早期测量都同等可靠。本卷的Planck曲线由公式生成，不复制原测量图，也不编造原误差棒。黑体\emph{平均谱}支持热起源，角方向的微小\emph{涨落谱}检验几何和结构种子；后者须联系CMB章的$C_\ell$，不能用前者的温度拟合取代。

\originalchapter{轻元素核合成与重子密度}
本章对应8.902第19讲末与第20讲。BBN把宇宙膨胀、弱反应和核反应放在同一时间尺度上，能预测氦的大致质量分数，并以氘丰度约束重子密度。解析估计足以看见机制，却不能称作已经执行了完整反应网络的精确丰度计算。下文保留课程使用的近似参数，并指出原笔记把中子平均寿命称为半衰期等会影响结果的区别。

\originalsection{弱反应控制中子与质子的比例}
核子形成后，弱反应$n+\nu_e\rightleftarrows p+e^-$、$n+e^+\rightleftarrows p+\bar\nu_e$以及中子衰变与其逆过程持续改变$n/p$。在充分热平衡、轻子化学势可忽略时，非相对论核子数密度比为
\[
 r_{ \rm eq}\equiv\frac{n_n}{n_p}
 =\left(\frac{m_n}{m_p}\right)^{3/2}
 \exp\left[-\frac{(m_n-m_p)c^2}{\kb T}\right]
 \simeq\exp[-1.293\ \mathrm{MeV}/(\kb T)].
\]
核子自旋自由度相同。温度下降使较重的中子相对减少。弱相互作用的数量级率$\Gamma_w\sim G_F^2T^5$，而辐射时期$H\propto T^2/M_P$，所以$\Gamma_w/H\propto T^3$，最终在MeV量级不再足以维持平衡。

课程取冻结温度约0.8 MeV，瞬时平衡估计$r_f\simeq0.20$。这与中微子逐渐脱耦发生在相近时期，但\textbf{弱反应停止与中微子脱耦不是一个完全相同的瞬间}。真实弱冻结持续一段时间，且需计入中微子温度、电子正电子分布与衰变；后面的氦数字不能由一个固定温度精确确定。

\originalsection{氘瓶颈为何延迟核合成}
第一步$p+n\rightleftarrows\mathrm D+\gamma$形成结合能$B_D\simeq2.22$ MeV的氘。即使$\kb T<B_D$，每个重子对应的光子太多，谱的高能尾部仍可破坏氘。没有足够稳定的氘，快速通向更紧束缚的氦的反应链无法开始，这称氘瓶颈。

在理想非相对论化学平衡中，$\mu_D=\mu_p+\mu_n$。将Maxwell--Boltzmann密度式相除可得
\[
 \frac{n_D}{n_pn_n}=\frac{g_D}{g_pg_n}
 \left(\frac{2\pi\hbar^2}{\kb T}\frac{m_D}{m_pm_n}\right)^{3/2}
 \exp\left(\frac{B_D}{\kb T}\right).
\]
$g_D=3$、$g_p=g_n=2$，忽略核子小质量差而取$m_D\simeq2m_N$，则右侧为$6[\pi\hbar^2/(m_N\kb T)]^{3/2}\exp[B_D/(\kb T)]$。正指数来自氘总质量低于两个自由核子的质量，不能误写成单粒子平衡密度的负Boltzmann指数。

为估算瓶颈解除，仿课程令自由氘数与自由中子数相当，$n_D/n_n\sim1$。这是把约一半尚未进入更重核的中子结合进氘的启发式条件，随后真实核反应网络会立即改变氘丰度。取$n_p\simeq(5/6)\eta_b n_\gamma$、$\eta_b=n_b/n_\gamma$，用黑体光子数公式得到
\begin{equation}\label{eq:astro-deuterium-bottleneck}
 1\sim\frac{10\zeta(3)}{\sqrt\pi}\eta_b
 \left(\frac{\kb T_D}{m_Nc^2}\right)^{3/2}
 \exp\left(\frac{B_D}{\kb T_D}\right).
\end{equation}
系数约6.78，与课程约6.9一致。小$\eta_b\sim10^{-9}$使所需温度远低于$B_D$；反应通常在$\kb T$约0.06--0.1 MeV的范围迅速展开。增加$\eta_b$使同样的结合比例更早、在更高温度达到。

将式\eqref{eq:astro-deuterium-bottleneck}具体取$\eta_b=5\times10^{-10}$、$m_Nc^2=939$ MeV时，可以数值解得约0.066 MeV。电子湮灭后$g_*\simeq3.36$的简化辐射年龄式给约300秒，而较高的0.08 MeV对应约200秒。原课的“$8\times10^8$ K、200秒”是数量级近似，不能与一个特定Saha半结合条件共同声称成精确解；精确时间还依赖自由度演化和非平衡反应。

\originalsection{中子衰变必须同时计入生成的质子}
在弱反应已可忽略、核结合尚少的近似阶段，令$X_n=n_n/(n_n+n_p)$，重子数守恒。中子平均寿命为$\tau_n$，则
\[
 X_n(t)=X_{n,f}\exp[-(t-t_f)/\tau_n],\qquad
 r(t)=\frac{r_f\exp[-\Delta t/\tau_n]}
 {1+r_f[1-\exp(-\Delta t/\tau_n)]}.
\]
分母增加来自衰变产生的质子。只有$r_f$很小时才近似$r\simeq r_f\exp(-\Delta t/\tau_n)$。课程所用890秒应理解为\emph{平均寿命}；其对应半衰期为$\tau_n\ln2\simeq617$秒。若真的把890秒当成半衰期，应使用$2^{-\Delta t/890\mathrm s}$，会给不同答案。

例如$r_f=0.20$、$\Delta t=200$秒、$\tau_n=890$秒给$r\simeq0.154$，并非由此精确推出0.25的氦分数。有效冻结温度、瓶颈结束时刻和残余弱反应的改进均重要，原讲义把后面的步骤称“精确计算”超过了这些近似所能支持的精度。本卷采用它说明数量级和参数方向，不伪造一个没有运行的精确网络结果。

\originalsection{氦质量分数的简洁估计}
瓶颈解除后，氘、氚与${}^3\mathrm{He}$反应迅速将大部分剩余中子锁入强束缚的${}^4\mathrm{He}$。每个氦核需两个中子，若忽略少量其他核和结合能对质量的修正，有$n_{\rm He}=n_n/2$、$n_{\rm H}=n_p-n_n$，因此
\begin{equation}\label{eq:astro-helium-fraction}
 Y_p=\frac{4n_{\rm He}}{n_{ \rm H}+4n_{\rm He}}
 =\frac{2n_n}{n_p+n_n}=2X_n=\frac{2r}{1+r}.
\end{equation}
若$r=1/7$，便得$Y_p=1/4$。若所有冻结时$r_f=1/5$的中子瞬间结合，则给$1/3$；延迟期间的中子衰变使结果下降。上面200秒例的$Y_p\simeq0.266$直观表明，只改变最后一位四舍五入不能把近似模型变成精确预测。

大部分核合成停止后，仍留下少量氘、${}^3\mathrm{He}$和${}^7\mathrm{Li}$或随后衰变成锂的${}^7\mathrm{Be}$。缺乏稳定的质量数5、8核，加上密度迅速下降与Coulomb势垒，使普通BBN难以生产大量更重元素；恒星内部的高密度、长时间和三氦过程才有不同机会。不能由本章一个氦式推算全部轻元素丰度。

\originalsection{从丰度到重子参数，需要哪些额外信息}
较高$\eta_b$提高核反应相对于膨胀的速率，让氘更有效地被烧成氦，故剩余D/H通常下降；氦分数对$\eta_b$依赖较弱。额外辐射自由度提高$H$，弱反应更早冻结且核合成前时间缩短，通常增加中子存留和氦分数。中子寿命、轻子化学势与反应截面也进入定量推断，不能将$Y_p$只看成$\Omega_b$的单参量函数。

对既定今日CMB温度，由$n_{b,0}=\Omega_b\rho_{ \rm crit,0}/m_p$与黑体$n_{\gamma,0}$得
\[
 \eta_b\simeq2.74\times10^{-8}\,\Omega_bh^2
 \left(\frac{2.7255\ \mathrm K}{T_0}\right)^3.
\]
因此氘等轻元素约束的是物理重子密度$\omega_b=\Omega_bh^2$；若要单独给$\Omega_b$，还需$h$。将BBN与CMB的重子加载比较，是不同物理过程的交叉核验，不是用同一份观测重新作两次拟合。

观测原初氦需从低金属丰度区域外推并控制恒星污染，氘常用远处吸收系统中的同位素谱线，需排除氢线混叠和化学演化。标准核网络与贫金属恒星锂丰度的差异是需要额外物理或天体系统误差解释的难题，本卷不凭一段解析估计声称已解决。反应网络、误差传播及真实数据拟合未实施，范围在覆盖清单单列。

\originalchapter{Lyman吸收森林与星系际气体}
本章对应8.902第21讲。远方类星体提供亮的背景光，沿途气体在自己的静止系吸收特定波长；我们因此能从一条观测光谱识别不同红移处的气体。吸收线既依赖密度，又依赖电离、温度和速度。本章从频率相关截面推导光学深度，并说明Gunn--Peterson暗槽能告诉我们什么以及不能单独告诉什么。

\originalsection{为什么森林位于类星体Ly\texorpdfstring{$\alpha$}{alpha}发射线的蓝侧}
氢$1s\leftrightarrow2p$的Lyman-$\alpha$静止波长约$\lambda_\alpha=1215.67$埃。红移$z_q$的类星体发射线位于$\lambda_\alpha(1+z_q)$；沿途红移$z_a<z_q$的吸收体在$\lambda_{ \rm obs}=\lambda_\alpha(1+z_a)$处吸收，因此该线在类星体发射线的蓝侧。光子发射时可来自类星体更短波长的\emph{连续谱}，并非必须先在类星体处是一个Ly$\alpha$线光子，再沿途被同一条线重复吸收。

一连串不同$z_a$的气体区域产生森林。谱中的金属双线还可帮助识别独立吸收系统，例如Mg II的2796、2803埃双线，但一条森林线不能仅凭深浅唯一指定一个孤立球云。实际星系际气体沿宇宙网连续分布，本动速度、温度与密度共同塑造观测线形。

\originalsection{柱密度与三个吸收尺度}
中性氢柱密度$N_{\rm HI}=\int n_{\rm HI}\dd l$的单位为$\mathrm{cm}^{-2}$。在Lyman极限912埃附近，氢光电离截面约$6.3\times10^{-18}\ \mathrm{cm}^2$，故$N_{\rm HI}\sim1.6\times10^{17}\ \mathrm{cm}^{-2}$已有连续吸收光学深度约1。较低柱密度通常表现为森林线，足够高时成为Lyman极限系统。

常用阻尼Ly$\alpha$系统定义$N_{\rm HI}\ge2\times10^{20}\ \mathrm{cm}^{-2}$，高柱密度使自然线宽的翼显著。原课印成$10^{32}\ \mathrm{cm}^{-2}$，不能原样用作物理分类。该定义另核对\href{https://arxiv.org/abs/astro-ph/0509481}{Wolfe、Gawiser与Prochaska（2005）}的开篇分类；只引用定义，不重印受限综述或其图。柱密度分类是观测约定，具体线宽仍受气体速度结构影响。

\originalsection{从归一线型到光学深度}
本节用Gaussian电磁单位。某共振线的吸收截面为
\[
 \sigma_\nu=S_\alpha\varphi(\nu-\nu_\alpha),\qquad
 S_\alpha=\frac{\pi e^2}{m_ec}f_\alpha,\qquad\int\varphi(\nu)\dd\nu=1.
\]
振子强度$f_\alpha$是无量纲矩阵元参数，不是逐次吸收概率。$\varphi$单位为$\mathrm{Hz}^{-1}$，$S_\alpha$单位为$\mathrm{cm}^2\mathrm{Hz}$，乘积$\sigma_\nu$才有面积单位。Ly$\alpha$的$f_\alpha\simeq0.416$给$S_\alpha\simeq1.10\times10^{-2}\ \mathrm{cm}^2\mathrm{Hz}$；原课将系数与线型的单位混在一起，应改正。

热速度使线心近Gaussian，自然有限寿命给Lorentz翼，两者卷积为Voigt型。热Doppler参数$b=\sqrt{2\kb T/m_H}$，$\Delta\nu_D=(b/c)\nu_\alpha$；若另有Gaussian非热速度，平方宽度相加。因此$b=20\ \mathrm{km\,s^{-1}}$在\emph{纯热}假设下给$T\simeq2.4\times10^4$ K，非热展宽存在时给的是温度上限。Lorentz翼按频差的平方反比衰减，不能把整个Voigt型都称为Gaussian。

沿观测频率$\nu_o$的光路，当地光子频率是$\nu_o(1+z)$，当地物理路径元为
\[
 \dd l=c|\dd t|=\frac{c\dd z}{(1+z)H(z)}.
\]
因此均匀背景近似下
\begin{equation}\label{eq:astro-lya-opacity}
 \tau(\nu_o)=\int_0^{z_q}\frac{c\,n_{\rm HI}(z)}{(1+z)H(z)}
 S_\alpha\varphi[\nu_o(1+z)-\nu_\alpha]\dd z.
\end{equation}
$n_{\rm HI}$是\emph{物理}数密度。射出原窄光束的共振散射也表现为消光，若忽略重新散入该光束的光，$F_\nu=F_{\nu,\rm cont}\ee^{-\tau}$。这不是将未经距离和红移变换的源静止系$F$直接代到观测谱中。

\originalsection{Gunn--Peterson光学深度与饱和的限制}
若线型很窄、沿共振区域的背景变化缓慢，用$\varphi\simeq\delta_D$，令$z_r=\nu_\alpha/\nu_o-1$。Dirac变量变换给$\delta_D[\nu_o(1+z)-\nu_\alpha]=\delta_D(z-z_r)/\nu_o$，代入式\eqref{eq:astro-lya-opacity}，有
\begin{equation}\label{eq:astro-gp-opacity}
 \tau_{ \rm GP}(z_r)=\frac{S_\alpha\lambda_\alpha n_{\rm HI}(z_r)}{H(z_r)}.
\end{equation}
在物质主导期$H\simeq H_0\sqrt{\Omega_{m,0}}(1+z)^{3/2}$，若中性比例$x_{\rm HI}$不变而总氢密度随$(1+z)^3$增长，便有$\tau_{ \rm GP}\propto x_{\rm HI}(1+z)^{3/2}$。注意$S_\alpha\lambda_\alpha$单位是$\mathrm{cm}^3\mathrm{s}^{-1}$，所以最后结果无量纲。

课程用$z=3$、$n_H\sim10^{-5}\ \mathrm{cm}^{-3}$、$H_0\sim70\ \mathrm{km\,s^{-1}\,Mpc^{-1}}$，完全中性时得到$\tau\sim10^5$。要透过显著光，气体的中性比例可需低到$10^{-5}$量级。反过来，$x_{\rm HI}$即使远小于1也能造成饱和黑暗槽。\textbf{黑暗槽不能单独证明该处气体“大部分中性”}，它往往只给中性比例的下限；课程以$z\sim6$暗槽直接推断之前大部分中性，需保留这项限制。定量再电离分析还用阻尼翼、空间起伏和其他数据。

\originalsection{光电离如何维持少量中性氢}
若背景紫外辐射场的角平均强度为$J_\nu$，每个中性原子的光电离率为
\[
 \Gamma_{ \rm HI}=4\pi\int_{\nu_{ \rm LL}}^\infty
 \frac{J_\nu\sigma_{ \rm ion}(\nu)}{h\nu}\dd\nu.
\]
它单位为$\mathrm{s}^{-1}$，不是整个宇宙每秒仅发出多少光子的总率。忽略碰撞电离、取适当复合系数$\alpha(T)$时，局部电离平衡要求
\[
 \Gamma_{ \rm HI}n_{\rm HI}=\alpha(T)n_en_{\rm HII}.
\]
高度电离氢气$n_e\simeq n_{\rm HII}\simeq n_H$，于是$x_{\rm HI}\simeq\alpha n_H/\Gamma_{ \rm HI}$。以课程例$\Gamma\sim10^{-12}\ \mathrm{s}^{-1}$、$\alpha\sim4\times10^{-13}\ \mathrm{cm^3\,s^{-1}}$、$n_H\sim10^{-5}\ \mathrm{cm}^{-3}$，给$x_{\rm HI}\sim4\times10^{-6}$，与透明度要求相容。

每个中性原子的电离时间约$\Gamma^{-1}\sim3\times10^4$年，而每个离子的复合时间$1/(\alpha n_e)\sim2.5\times10^{17}$秒，约$8\times10^9$年。原课将$2\times10^{17}$秒换成$3\times10^6$年是单位换算错误。二者差异由平衡时\emph{中性原子很少}调节，不能靠把两个不同比例的粒子时间强行设成相同而推导平衡。

线宽对应的$\kb T$小于13.6 eV，不足以单独排除所有碰撞电离，因为Maxwell分布仍有高能尾部；必须比较碰撞电离与复合率。所需极低中性比例和强紫外背景使光电离成为低密度森林气体的合适解释。复合系数用Case A还是Case B取决于再发出的电离光子是否逃逸；气体团聚又使$\langle n_e n_p\rangle$不同于平均密度平方。

\originalsection{光谱能推断什么，以及还缺什么}
由$x_{\rm HI}\propto\alpha(T)n_H/\Gamma$和式\eqref{eq:astro-gp-opacity}可看出，吸收深度对密度近乎二次敏感，也依赖紫外背景和温度。相同$\tau$可以来自不同密度与电离条件；线宽有热和速度简并，饱和线又失去线心柱密度信息。金属双线、氘线、阻尼翼和不同Lyman线可帮助打破部分简并。

从森林通量谱推到物质功率谱，还需非线性流体、热史、速度畸变和辐射转移模型。原课提到的THESAN和HERA说明研究方法，但本卷没有运行相应模拟或处理实际望远镜数据，不能把课堂示意图当作已完成的模型核验。星系际介质、再电离及暗物质结构因此通过同一条吸收光路联系起来，又保留明确的反演边界。

\addtocontents{toc}{\protect\Needspace{3\baselineskip}}
\originalchapter{辐射、测光与物质状态}\label{ch:radiation}
\originalsection{强度与可观测量}
天体通常不能送入实验室. 光的能量、频谱和方向分布便是推断温度、运动与组成的入口. 比强度 $I_\nu$ 定义为沿方向 $\boldsymbol n$ 的射线在时间 $\dd t$ 内, 穿过垂直射线的面积 $\dd A_\perp$, 在立体角 $\dd\Omega$ 和频段 $\dd\nu$ 中携带的能量除以这四个增量. 若使用固定面元 $\dd A$, 则 $\dd A_\perp=\cos\theta\dd A$. 因而通量为 $F_\nu=\int I_\nu\cos\theta\dd\Omega$. 各向同性强度在整球上的净通量为零, 在向外半球上的通量为 $\pi I_\nu$. 强度与通量的区别在恒星表面与观测者处尤其重要.

半径为 $R$ 的球对称光源具有光度 $L$, 真空中能量守恒给出 $F=L/(4\pi d^2)$. 没有吸收时同一条射线的 $I_\nu$ 不随距离改变: 物体所占立体角与接收到的通量都按 $d^{-2}$ 缩小. 宇宙膨胀另使 $I_\nu/\nu^3$ 沿自由传播保持不变, 积分后的面亮度按 $(1+z)^{-4}$ 衰减. 后一结论有两个来自光子能量和到达率的因子, 再有两个来自角距离和光度距离之比.

星等是对数亮度: $m_1-m_2=-2.5\log_{10}(F_1/F_2)$. 绝对星等 $M$ 定义在 $10\,\mathrm{pc}$, 故无消光且低红移时 $m-M=5\log_{10}(d/10\,\mathrm{pc})$. 波段星等不能直接等同于总光度; 用不同温度的恒星比较时须保留波段响应和总光度校正. 消光 $A_\lambda$ 使 $m_\lambda$ 增加, 色余 $E(B-V)=A_B-A_V$ 则刻画相对变红.

\originalsection{黑体与输运方程}
频率为 $\nu$ 的光子振子平均能量为 $h\nu/(\ee^{h\nu/\kb T}-1)$. 每单位体积、频率有 $8\pi\nu^2/c^3$ 个含两偏振的模, 故
\begin{equation}\label{eq:planck}
B_\nu(T)=\frac{2h\nu^3}{c^2}\frac1{\ee^{h\nu/\kb T}-1},\qquad
u_\nu=\frac{4\pi}{c}B_\nu.
\end{equation}
这个分布需要热平衡和光子数可改变, 于是光子化学势为零. $h\nu\ll\kb T$ 时得到 Rayleigh--Jeans 极限 $B_\nu\simeq2\nu^2\kb T/c^2$; 高频端由指数抑制控制. 由 $B_\lambda\dd\lambda=-B_\nu\dd\nu$ 可知两个峰值并不以 $\lambda\nu=c$ 对应, 因为谱密度的横轴不同.

沿射线路径 $s$ 的发射系数 $j_\nu$ 和吸收系数 $\alpha_\nu$ 给出 $\dd I_\nu/\dd s=-\alpha_\nu I_\nu+j_\nu$. 定义 $\dd\tau_\nu=\alpha_\nu\dd s$、源函数 $S_\nu=j_\nu/\alpha_\nu$, 解一阶线性方程得
\begin{equation}
I_\nu(\tau)=I_\nu(0)\ee^{-\tau}+\int_0^\tau S_\nu(t)\ee^{-(\tau-t)}\dd t.
\end{equation}
局域热平衡且真吸收主导时 $S_\nu=B_\nu$. 常源函数解为 $I_\nu=I_\nu(0)\ee^{-\tau}+B_\nu(1-\ee^{-\tau})$. 光学厚介质抹掉背景, 光学薄介质保留背景并加上线积分发射. 散射主导时源函数含辐射场自身, 不能直接以黑体代替.

\originalsection{电离平衡}
在稀薄、非简并、局域热平衡气体中, 中性原子、电离离子和电子的化学势满足 $\mu_0=\mu_1+\mu_e$. 非相对论粒子的配分因子含热德布罗意体积 $\lambda_T^3=(2\pi\hbar^2/m\kb T)^{3/2}$. 对重原子与离子近似等质量, 消去化学势可得
\begin{equation}\label{eq:saha}
\frac{n_en_1}{n_0}=\frac{2U_1}{U_0}
 \left(\frac{2\pi m_e\kb T}{h^2}\right)^{3/2}\ee^{-\chi/\kb T}.
\end{equation}
$U_i$ 是不含自由电子自旋的内部配分函数, $\chi$ 为电离能, 前面的 2 是电子自旋简并. 高温不是唯一的电离条件; 低密度也提高电离比例, 因为电离增加粒子可用的平动状态. 公式不能用于强辐照且碰撞不足的光电离气体, 这种情形须比较实际电离与复合速率.

\originalsection{电离关系的统计推导}
为了把 Saha 式与可测密度联系起来, 需说明能量零点. 对非简并粒子 $i$, 单粒子能量写为 $E_i=E_{i0}+p^2/(2m_i)$, Maxwell--Boltzmann 数密度为 $n_i=U_i\lambda_i^{-3}\exp[(\mu_i-E_{i0})/\kb T]$. 反应 $H\leftrightarrow p+e$ 的化学势守恒并不等于三个 $E_{i0}$ 之和为零. 采用中性基态零点, 电离产物的最低总能比它高 $\chi$, 因而消去化学势后有 $\exp(-\chi/\kb T)$. 重离子与原子的热波长比近 1, 留下电子平动因子及电子自旋权重, 得式\eqref{eq:saha}.

这一步也说明为什么没有把原子质量放到最终的巨大 $T^{3/2}$ 因子中. 电离额外创造一个可以在空间独立平动的轻电子, 可用状态数正是由它的热波长控制. 在给定密度下半电离路径不是 $\kb T=\chi$, 而是指数与电子平动状态数共同竞争. 对高密度气体, 电荷屏蔽、能级截断与简并会使这里的理想配分函数失效.

辐射扩散同样可以从频率输运明确构造平均值. 局域黑体且光学厚时 $F_\nu=-(4\pi/3\kappa_\nu\rho)(\partial B_\nu/\partial T)\nabla T$. 对频率积分后定义
\[
\frac1{\kappa_R}=\frac{\int_0^\infty\kappa_\nu^{-1}(\partial B_\nu/\partial T)\dd\nu}
{\int_0^\infty(\partial B_\nu/\partial T)\dd\nu}.
\]
低不透明度频段更容易运走能量, 因而是调和权重. 分母由 $\int B_\nu\dd\nu=\sigma T^4/\pi$ 求导, 恢复正文的扩散式. 频率平均不是单纯的数学便利, 它对应实际的能量输运方向.

\addtocontents{toc}{\protect\Needspace{3\baselineskip}}
\originalchapter{恒星结构、燃烧与演化}\label{ch:stellar}
\originalsection{静水平衡与能量预算}
恒星长期发光却没有立即坍缩, 说明压力梯度平衡重力, 同时有能量来源补充辐射损失. 球对称、慢于动力学时间变化的恒星以 $r$ 为半径, $m_r$ 为内部质量, $L_r$ 为穿过球面的光度. 质量守恒、力平衡及能量守恒分别为
\begin{equation}\label{eq:stellar}
\frac{\dd m_r}{\dd r}=4\pi r^2\rho,\qquad
\frac{\dd P}{\dd r}=-\frac{Gm_r\rho}{r^2},\qquad
\frac{\dd L_r}{\dd r}=4\pi r^2\rho\epsilon.
\end{equation}
$\epsilon$ 是单位质量净产热率, 演化模型中还须计中微子损失与热储能变化. 这三式尚不封闭: 需要状态方程、能量输运规律以及不透明度和核反应率. 中心条件为 $m_r,L_r\to0$, 表面则与大气边界匹配; 不能把给定总质量与任意中心温度当作互相独立的自由参数.

粗略地 $P_c\sim GM^2/R^4$、$\rho\sim M/R^3$. 理想气体 $P\sim\rho\kb T/(\mu m_p)$ 于是给 $T_c\sim GM\mu m_p/(\kb R)$. 这解释引力收缩为什么提高温度. 引力能 $W\sim-GM^2/R$, 热平衡下维里关系 $2K+W=0$ 使总能 $E=W/2<0$. 星体丢失能量后 $E$ 更负, $K$ 反而增大; 自引力系统具有这种负热容行为. 仅靠收缩供能的时间 $t_{\rm KH}\sim GM^2/(RL)$ 远短于太阳的核燃烧寿命.

\originalsection{辐射扩散与对流}
当平均自由程 $\ell=(\kappa_R\rho)^{-1}$ 远小于温度变化尺度, 各方向强度仅有小各向异性. 将 $B_\nu[T(s)]$ 按一个自由程展开并作角平均, $\int\cos^2\theta\dd\Omega=4\pi/3$, 得 $F=-c\nabla u/(3\kappa_R\rho)$. 使用 $u=a_RT^4$ 后,
\begin{equation}
F_r=-\frac{4a_RcT^3}{3\kappa_R\rho}\frac{\dd T}{\dd r},\qquad
\nabla_{\rm rad}\equiv\frac{\dd\ln T}{\dd\ln P}
=\frac{3\kappa_RL_rP}{16\pi a_RcGm_rT^4}.
\end{equation}
Rosseland 平均是按输运权重定义的不透明度调和平均, 不是任意频率不透明度的算术平均. 若气团上移时热交换慢, 它按绝热梯度 $\nabla_{\rm ad}$ 冷却. 环境在同一压强下温度下降更快, 即 $\nabla_{\rm rad}>\nabla_{\rm ad}$, 则气团较热、密度较低并继续上升. 这就是均匀组成下的 Schwarzschild 对流判据; 有平均分子量梯度时需加入组成稳定项.

\originalsection{核反应与演化尺度}
质子间库仑势垒使低能粒子难以接近, 量子隧穿使反应仍可发生. 非共振带电粒子截面可写为\[\sigma(E)=\frac{S(E)}E\exp\left[-\sqrt{\frac{E_G}E}\right].\] Maxwell 分布的 $\exp(-E/\kb T)$ 与隧穿因子竞争, 共同指数 $f(E)=E/\kb T+\sqrt{E_G/E}$ 在
$E_0=[E_G(\kb T)^2/4]^{1/3}$ 处最小. 对 $f$ 的二阶展开给出反应率主要来自 $E_0$ 附近的有限能段, 即 Gamow 窗口. 这里不将未知的 $S(E)$ 当作已推导常数, 详细核反应网络和实测截面不在本卷范围.

氢转为氦约释放反应质量的 $0.007c^2$. 若可燃质量占总质量的 $f$、氢质量分数为 $X$, 主序寿命估计 $t_{\rm nuc}\sim0.007fXM c^2/L$. 全对流低质量恒星可把更多氢带到燃烧区; 大质量恒星光度增长比质量快, 因而寿命短. 主序离开后的壳层燃烧、核心收缩和外层膨胀可在 H--R 图上区分, 但轨道的精确数值依赖完整输运、质量损失及组成演化, 不能用这一时间估计替代恒星演化程序.

\originalsection{平衡尺度与稳定性的区别}
估算恒星中心温度时, 必须区分维里总能与核燃料能. 对均匀球 $W=-3GM^2/(5R)$, 理想单原子气体总热动能 $K=3M\kb\bar T/(2\mu m_p)$. 维里关系给 $\bar T=GM\mu m_p/(5\kb R)$, 它是质量加权平均温度, 不等于中心温度. 中心温度通常更高, 要由具体密度与状态方程求出. 将此平均值的系数直接当作任意恒星中心公式会引入无依据的精度.

球的收缩和稳定也需要新的条件. 设均匀缩放时绝热 $P\propto\rho^\gamma$, 质量固定, 内能 $U\propto R^{-3(\gamma-1)}$, 引力能 $W\propto-R^{-1}$. 写 $E(R)=A R^{-q}-B/R$, $q=3(\gamma-1)$、$A,B>0$. 平衡 $E'(R)=0$ 给 $B=qA R^{1-q}$, 再代二阶导数得 $E''=q(q-1)A R^{-q-2}$. 因而这个同质缩放模型稳定要求 $\gamma>4/3$, 临界 $\gamma=4/3$ 对整体缩放边缘稳定. 这为相对论简并星质量极限提供物理解释, 但不是完整径向正常模分析, 对非均匀组成和强广义相对论星体须重新处理.

核燃烧还可给定温度反馈的直觉. 理想气体星体核心温度升高使反应率提高, 增加的热能推动膨胀, 密度与温度下降, 形成热调节. 简并核压力对温度不敏感, 热输入不能立即以相同方式膨胀调节, 可导致燃烧不稳定. 这解释氦闪的机制方向; 具体发生时刻和峰值必须使用相应反应率、热容和输运, 本卷没有编造未运行的氦闪曲线.

\addtocontents{toc}{\protect\Needspace{3\baselineskip}}
\originalchapter{简并星、吸积与恒星形成}\label{ch:compact}
\originalsection{电子简并与多方球}
白矮星不依靠高温维持压力. Pauli 原理使电子必须占据不同动量态. 零温时填满 $p\le p_F$, 两自旋态给 $n_e=p_F^3/(3\pi^2\hbar^3)$. 各向同性动量流的压强为
\begin{equation}
P=\frac1{3\pi^2\hbar^3}\int_0^{p_F}p^3v(p)\dd p.
\end{equation}
非相对论时 $v=p/m_e$, 因而 $P=K_{\rm NR}\rho^{5/3}$; 极相对论时 $v=c$, 因而 $P=K_{\rm R}\rho^{4/3}$. 其中 $\rho=\mu_em_pn_e$, $\mu_e$ 为每电子对应的质子质量倍数, 并非恒星气体平均分子量 $\mu$.

令 $P=K\rho^{1+1/n}$、$\rho=\rho_c\theta^n$, 取
$r=\alpha\xi$、$\alpha^2=(n+1)K\rho_c^{1/n-1}/(4\pi G)$.
把静水平衡式乘 $r^2/\rho$ 再求导, 使用质量守恒消去 $m_r$, 得
\begin{equation}\label{eq:lane}
\frac1{\xi^2}\frac{\dd}{\dd\xi}\left(\xi^2\theta'\right)=-\theta^n,
\qquad\theta(0)=1,\quad\theta'(0)=0.
\end{equation}
中心看似奇异, 正则展开却给 $\theta=1-\xi^2/6+O(\xi^4)$, 可由小正 $\xi$ 启动数值积分. 第一零点 $\xi_1$ 定义表面, $R=\alpha\xi_1$、$M=4\pi\alpha^3\rho_c[-\xi_1^2\theta'(\xi_1)]$. $n=3/2$ 时 $R\propto M^{-1/3}$; $n=3$ 时质量与 $\rho_c$ 无关, 给出 Chandrasekhar 质量尺度. 真实白矮星还受有限温度、库仑修正和逆 $\beta$ 过程影响.

\originalsection{吸积盘}
对质量 $M$ 的中心天体, 半径 $r$ 的 Kepler 速度为 $v_K=(GM/r)^{1/2}$. 圆轨道单位质量总能为 $-GM/(2r)$, 进入更小半径必须失去能量和角动量. 因此吸积不仅是物质向内掉落, 还需要应力把角动量向外输运. 薄盘中垂直重力近似 $-\Omega_K^2z$, 静水平衡给尺度高 $H\sim c_s/\Omega_K$; $H/r\ll1$ 要求 $c_s\ll v_K$.

对稳态、Kepler 薄盘, 以 $\dot M>0$ 为向内质量流率, 角动量守恒使应力矩满足 $\mathcal G=\dot M(\sqrt{GMr}-\sqrt{GMr_{\rm in}})$, 此处使用内缘零应力条件. 每环机械功率来自相邻环的相对角速度, 两面各自辐射通量为
\begin{equation}
F(r)=\frac{3GM\dot M}{8\pi r^3}\left(1-\sqrt{\frac{r_{\rm in}}r}\right).
\end{equation}
积分两面得到 $L=GM\dot M/(2r_{\rm in})$. 这个结果不适用于厚盘、强风、内缘有应力或强相对论区域. 辐射作用于自由电子, 通过电荷耦合带动离子; 比较 $\sigma_TL/(4\pi r^2c)$ 与质子引力 $GMm_p/r^2$ 得 $L_{\rm Edd}=4\pi GMm_pc/\sigma_T$. 这是球对称散射模型的力平衡尺度, 不表示所有盘必须严格低于它.

\originalsection{云的引力不稳定}
均匀气体中取小扰动 $\delta\rho,\boldsymbol v,\delta\Phi\propto\ee^{\ii(\boldsymbol k\cdot\boldsymbol x-\omega t)}$. 连续方程给 $-\omega\delta\rho+\rho_0\boldsymbol k\cdot\boldsymbol v=0$, Euler 方程给 $\omega\boldsymbol v=\boldsymbol k(c_s^2\delta\rho/\rho_0+\delta\Phi)$, Poisson 方程给 $-k^2\delta\Phi=4\pi G\delta\rho$. 消去两未知量得 $\omega^2=c_s^2k^2-4\pi G\rho_0$. 当 $k<k_J=\sqrt{4\pi G\rho_0}/c_s$ 时扰动增长. 本卷定义 $\lambda_J=2\pi/k_J$, Jeans 半径为 $R_J=\lambda_J/2$, 故 $M_J=\pi^{5/2}c_s^3/(6G^{3/2}\rho_0^{1/2})$. 不同文献的 Jeans 质量常数可能不同, 必须先声明约定.

这项无限均匀介质线性化以减去背景引力为理想化, 不是有限球的精确边界值解. 坍缩后若近似等温, $P$ 项相对于重力会越来越弱; 角动量和磁场却可阻止简单球对称坍缩, 所以形成盘、碎裂和角动量输运是恒星形成不可忽略的后续步骤.

\originalsection{Lane--Emden 方程的检验例}
从静水平衡可写 $\frac1{r^2}\frac{\dd}{\dd r}(\frac{r^2}{\rho}\frac{\dd P}{\dd r})=-4\pi G\rho$. 代 $P=K\rho_c^{1+1/n}\theta^{n+1}$、$\rho=\rho_c\theta^n$, 内部导数的 $\theta^n$ 与 $\rho$ 消去, 尺度 $\alpha$ 则消去剩余常数, 得式\eqref{eq:lane}. 此推导要求 $n>0$ 与正密度; $n=0$ 的不可压均匀密度球应直接取相应极限方程.

对 $n=0$, 直接积分 $\xi^2\theta'=-\xi^3/3$, 得 $\theta=1-\xi^2/6$, 第一零点 $\sqrt6$. 对 $n=1$, 令 $u=\xi\theta$, 方程化为 $u''+u=0$, 正则初值 $u(0)=0,u'(0)=1$, 故 $\theta=\sin\xi/\xi$, 第一零点 $\pi$, 质量参数 $-\pi^2\theta'(\pi)=\pi$. 两个解析例能检验中心处理与数值程序的归一化, 不能因它们可积就假定 $n=3/2,3$ 也有相同闭式.

无量纲曲线的作用是把中心密度与结构形状分开. 对固定 $n,K$, 改变 $\rho_c$ 只重缩放 $\alpha$ 和总质量, 无量纲 $\xi_1,w_n$ 不变. 因此质量半径尺度律是整族模型的结论, 不是某一次中心密度计算的巧合. 接近相对论简并极限时仍要注意, 真正的状态方程从 $5/3$ 连续过渡到 $4/3$, 并不是突然切换一个全星固定多方指数.

\addtocontents{toc}{\protect\Needspace{3\baselineskip}}
\originalchapter{星系的引力势与轨道}\label{ch:galaxy}
\originalsection{质量模型}
星系中恒星分布可用面亮度 $I(R)$ 描述, 转为质量密度需要质量光度比和投影模型. 指数盘 $I\propto\ee^{-R/R_s}$ 的面积积分给总光度 $2\pi I_0R_s^2$. 椭球星系常用 Sérsic 曲线 $I(R)=I_e\exp\{-b_n[(R/R_e)^{1/n}-1]\}$, $b_n$ 由 $R_e$ 包含半光度确定. 一个拟合曲线不能独自区分年龄、金属丰度、尘埃与恒星初始质量函数, 动力学测量能提供额外质量约束.

球对称势满足 $\Phi'(r)=GM(<r)/r^2$, 因而圆速度 $v_c^2=r\Phi'=GM(<r)/r$. 若 $v_c$ 长距离近似常数, 则 $M(<r)\propto r$, 密度近似 $r^{-2}$; 光度已经趋于有限时, 这种质量行为提示暗物质或引力模型需要解释. 不能仅由一条投影旋转曲线无条件唯一恢复三维密度.

常用势密度对包括 Plummer 模型
$\Phi=-GM/\sqrt{r^2+b^2}$、$\rho=3Mb^2/[4\pi(r^2+b^2)^{5/2}]$,
以及 NFW 密度 $\rho=\rho_s/[x(1+x)^2]$, $x=r/r_s$. 对 NFW 积分得 $M=4\pi\rho_sr_s^3[\ln(1+x)-x/(1+x)]$. 未截断的质量对数发散, 因而实际使用必须指定晕边界, 如平均密度达到背景或临界密度某倍数处.

\originalsection{有效势与本轮运动}
中心势中轨道平面固定, 单位质量角动量 $\ell=r^2\dot\phi$ 守恒. 径向运动满足 $\dot r^2/2+\Phi_{\rm eff}=E$, $\Phi_{\rm eff}=\Phi+\ell^2/(2r^2)$. 转向点是 $E=\Phi_{\rm eff}$, 圆轨道是有效势驻点. 对轴对称势, 以圆轨道半径 $R_0$ 为中心展开, 保持角动量而不是角速度固定, 得
\begin{equation}\label{eq:epicycle}
\kappa^2=\Phi''(R_0)+3\frac{\Phi'(R_0)}{R_0}
=R\frac{\dd\Omega^2}{\dd R}+4\Omega^2,
\qquad\Omega=\frac{v_c}R.
\end{equation}
$\kappa^2>0$ 时小径向扰动为谐振; $\kappa^2<0$ 时圆轨道不稳定. Kepler 势给 $\kappa=\Omega$, 平坦旋转曲线给 $\kappa=\sqrt2\Omega$, 后者的轨道一般不是闭合椭圆.

在太阳邻域把差分旋转一阶展开, Oort 常数为 $A=(v_c/R-v_c')/2$、$B=-(v_c/R+v_c')/2$, 因而 $\Omega=A-B$、$v_c'=-(A+B)$、$\kappa^2=-4B(A-B)$. 这些关系使用附近恒星、薄盘和近圆轨道假设, 远处的非圆运动不能都归入同一个局部线性场.

\originalsection{维里定理与弛豫}
令 $I=\sum m_ir_i^2$. 两次微分并代入 Newton 方程得 $\ddot I/2=2K+\sum\boldsymbol r_i\cdot\boldsymbol F_i$. 成对 $1/r$ 引力使后一项等于势能 $W$. 长期束缚稳态的时间平均 $\langle\ddot I\rangle=0$, 因而 $2\langle K\rangle+\langle W\rangle=0$. 非孤立系统需保留外场、表面压力等项, 突然坍缩过程中不能直接强加稳态等式.

速度为 $v$ 的粒子以冲量近似经过质量 $m$ 的散射体, 冲量 $\Delta v_\perp\simeq2Gm/(bv)$. 每时间的遭遇数是 $2\pi bn v\dd b$. 积分速度扩散率得
$\dd\langle(\Delta v)^2\rangle/\dd t\simeq8\pi G^2m^2n\ln\Lambda/v$,
$\ln\Lambda=\ln(b_{\max}/b_{\min})$.
定义 $t_{\rm rel}$ 为累计扰动达到 $v^2$ 的时间, 有 $t_{\rm rel}\sim v^3/(8\pi G^2m\rho\ln\Lambda)$. 星系的巨大 $N$ 通常使二体弛豫慢于年龄; 快速势变化引起的暴烈弛豫则不依靠逐次二体碰撞.

\addtocontents{toc}{\protect\Needspace{3\baselineskip}}
\originalchapter{分布函数、稳定性与星系演化}\label{ch:distribution}
\originalsection{从轨道到统计描述}
令 $f(\boldsymbol x,\boldsymbol v,t)\dd^3x\dd^3v$ 为相空间体元中的质量. Hamilton 流保持相空间体积, 没有碰撞时沿轨道 $\dd f/\dd t=0$, 即
\begin{equation}\label{eq:cbe}
\frac{\partial f}{\partial t}+\boldsymbol v\cdot\nabla_xf
-\nabla\Phi\cdot\nabla_vf=0,\qquad
\nabla^2\Phi=4\pi G\int f\dd^3v.
\end{equation}
势与分布必须自洽, 不能给任意 $f$ 再无条件套用另一个势. 对速度积分时要求 $f$ 在大速度处衰减, 边界项消失. 零阶矩给连续方程; 乘 $v_i$ 再积分得
\begin{equation}
\partial_t(\rho\bar v_i)+\partial_j(\rho\overline{v_iv_j})
=-\rho\partial_i\Phi.
\end{equation}
这是 Jeans 方程, 不是一个独立封闭状态方程. 稳态、球对称、无整体径向流时,
$\dd(\rho\sigma_r^2)/\dd r+2\beta\rho\sigma_r^2/r=-\rho GM(<r)/r^2$,
$\beta=1-(\sigma_\theta^2+\sigma_\phi^2)/(2\sigma_r^2)$. 不同 $\beta$ 可产生相近投影视向速度弥散, 形成质量与各向异性的退化.

\originalsection{盘的局部稳定性}
对薄气体盘取轴对称、波长远小于盘尺度的扰动, 局部径向惯性项给 $\kappa^2$, 压强给 $c_s^2k^2$. 薄片的扰动势满足 $\delta\Phi=-2\pi G\delta\Sigma/|k|$, 因而
\begin{equation}
\omega^2=\kappa^2+c_s^2k^2-2\pi G\Sigma|k|.
\end{equation}
对于 $k>0$, 二次式在 $k=\pi G\Sigma/c_s^2$ 最小. 所有局部轴对称模式稳定需 $Q_g=\kappa c_s/(\pi G\Sigma)\ge1$. 恒星盘的无碰撞响应不同, 不能把恒星速度弥散直接代入气体系数并声称同一个精确判据. 这也说明自引力、速度支持和旋转共同决定稳定性, 单凭盘半径无法在未指定模型时判断全局扰动.

\originalsection{恒星族与化学演化}
恒星族综合把初始质量函数 $\xi(m)$ 与单星年龄、组成依赖的光谱 $l_\nu(m,t,Z)$ 卷积, 再对过去的恒星形成率积分. 最亮星未必贡献多数质量, 尘埃还能改变光谱; 因而以光谱恢复历史是有退化的逆问题.

封闭箱模型假设没有流入流出、瞬时回收和充分混合. 令总质量 $M_T$ 固定, 气体质量 $M_g$, 锁入长寿命恒星的增量 $\dd M_s=-\dd M_g$. 金属气体质量的变化为 $\dd(ZM_g)=(y-Z)\dd M_s$, 其中 $y$ 是每锁入单位质量的新生成金属产额. 展开左侧并消去 $\dd M_g$ 得 $M_g\dd Z=y\dd M_s$, 初始 $Z=0$ 时积分为 $Z=-y\ln\mu$, $\mu=M_g/M_T$. 这条对数关系是模型推导, 不是任意星系的经验定律. G 型矮星问题、吸积和流出提示它的局限.

活动星系核的能量由黑洞吸积提供, 大致 $L=\eta\dot Mc^2$. 尘埃环的方向选择可使同一中心发动机在光谱上表现不同, 但所有 AGN 差异并不都由方向解释. 喷流需要磁场、角动量和相对论流动, 本卷使用能量与 Eddington 尺度, 不声称完成磁流体喷流理论.

\originalsection{局部盘色散关系的来源}
在半径 $R_0$ 的小区域, 取扰动径向波数 $k$ 与频率 $\omega$, 背景面密度近似常数. 线性连续式是 $-\ii\omega\delta\Sigma+\ii k\Sigma u_R=0$. 径向与切向 Euler 式为 $-\ii\omega u_R-2\Omega u_\phi=-\ii k(c_s^2\delta\Sigma/\Sigma+\delta\Phi)$, $-\ii\omega u_\phi+\kappa^2u_R/(2\Omega)=0$. 第二式所含的 $\kappa$ 来自差分旋转, 不是任意添上的振荡频率.

无限薄片的扰动 Poisson 方程为 $(\partial_z^2-k^2)\delta\Phi=4\pi G\delta\Sigma\delta(z)$. 两侧衰减解 $\delta\Phi=C\ee^{-|k||z|}$ 的导数跨片跳变 $-2|k|C$, 故 $C=-2\pi G\delta\Sigma/|k|$. 再由连续式消去 $\delta\Sigma$、切向式消去 $u_\phi$, 得正文色散关系. 若盘有限厚度, 引力响应因垂直分布被减弱; 若扰动非轴对称或波长接近盘尺度, 这个局部推导不能判断所有全局棒模或旋臂模.

Jeans 方程也可在一个自足模型中检验. 球对称、各向同性系统 $\rho=C/r^2$, 外部势给恒定圆速 $v_c$, 取常数 $\sigma_r^2$. Jeans 方程左侧为 $-2C\sigma_r^2/r^3$, 右侧为 $-Cv_c^2/r^3$, 因而 $\sigma_r^2=v_c^2/2$. 这是等温球尺度关系; 截断、各向异性或可见恒星作为示踪体时, 同一个比例不必保留. 通过检验例能看清密度梯度如何进入动力学质量估计.

\addtocontents{toc}{\protect\Needspace{3\baselineskip}}
\originalchapter{均匀膨胀宇宙}\label{ch:cosmo}
\originalsection{尺度因子与红移}
在足够大的尺度上采用均匀、各向同性的近似. 共动坐标把随 Hubble 流运动的星系位置固定, 其实际距离随尺度因子 $a(t)$ 改变. 取今日 $a_0=1$, 则 $H=\dot a/a$. $v=HD$ 是同一宇宙时刻相邻共动观察者的距离变化, 大距离的退行率可以超过 $c$, 它不等于局部狭义相对论中的相对速度.

沿平直 FLRW 径向光线有 $\dd\chi=c\dd t/a(t)$. 相邻两个波峰连接同一对共动观察者, 积分端点微小变化给 $\delta t_o/a_o=\delta t_e/a_e$, 所以 $1+z=\lambda_o/\lambda_e=a_o/a_e$. 红移由整个传播期间的几何变化给出, 不需要把遥远星系与地球的比较强行视为单次 Lorentz boost.

\originalsection{Friedmann 动力学}
以下调用广义相对论卷已经建立的 FLRW 方程. $\rho$ 为包含全部组分的质量等效密度, $p$ 为压强:
\begin{align}
H^2&=\frac{8\pi G}3\rho-\frac{kc^2}{a^2},\label{eq:friedmann}\\
\frac{\ddot a}a&=-\frac{4\pi G}3\left(\rho+\frac{3p}{c^2}\right),\qquad
\dot\rho+3H\left(\rho+\frac p{c^2}\right)=0.\label{eq:conservation}
\end{align}
曲率项不表示一种真实物质. 非相互作用组分分别守恒, 常数状态参数 $w=p/(\rho c^2)$ 给 $\rho\propto a^{-3(1+w)}$. 物质 $w=0$, 辐射 $w=1/3$, 真空 $w=-1$. 后者虽然密度恒定, 共动体积内能量随体积增加, 由负压强做功关系补偿, 不应套用无压尘埃的能量直觉.

平直单组分宇宙若 $w>-1$, 由 $\dot a\propto a^{-(1+3w)/2}$ 积分得 $a\propto t^{2/[3(1+w)]}$. 物质期 $a\propto t^{2/3}$, 辐射期 $a\propto t^{1/2}$; 真空期须单独积分, 得 $a\propto\ee^{Ht}$. 总密度临界值为 $\rho_c=3H^2/(8\pi G)$, 定义 $\Omega_i=\rho_i/\rho_c$, 则 $\Omega_k=-kc^2/(a^2H^2)$ 且 $\sum\Omega_i+\Omega_k=1$.

\originalsection{宇宙时间与视界}
多组分模型写作
\begin{equation}\label{eq:Hz}
E(z)^2\equiv\frac{H(z)^2}{H_0^2}
=\Omega_{r0}(1+z)^4+\Omega_{m0}(1+z)^3
+\Omega_{k0}(1+z)^2+\Omega_{\Lambda0}.
\end{equation}
回望时间 $t_0-t(z)=H_0^{-1}\int_0^z\dd z'/[(1+z')E(z')]$. 年龄是上限改为无穷的积分, 须给出早期组分以保证模型在积分全域有效. 平直物质与真空模型有 $t_0=2\operatorname{arsinh}\sqrt{\Omega_\Lambda/\Omega_m}/(3H_0\sqrt{\Omega_\Lambda})$.

粒子视界 $D_p=a(t)\int_{t_{\rm start}}^t c\dd t'/a(t')$ 描述到今天能收到的最远因果传播; 事件视界 $D_e=a(t)\int_t^\infty c\dd t'/a(t')$ 描述现在发出的信号能到达的共动范围. 后者仅在积分收敛且模型持续有效时存在. Hubble 半径 $c/H$ 是局部膨胀尺度, 通常不等于前两者. 例如 $a\propto t^p$、$0<p<1$ 给 $D_p=ct/(1-p)$, 而持续指数膨胀给 $D_e=c/H$.

\addtocontents{toc}{\protect\Needspace{3\baselineskip}}
\originalchapter{距离、吸收与观测推断}\label{ch:distance}
\originalsection{三种距离}
由径向光线得共动径向距离 $\chi(z)=cH_0^{-1}\int_0^z\dd z'/E(z')$. 曲空间横向距离为 $D_M=a_0S_k(\chi)$, 其中 $S_k$ 分别为球面、平直或双曲几何的正弦、恒等、双曲正弦形式, 并保留相应曲率半径. 以下数值题的默认平直情形有 $D_M=\chi$.

发射时大小为 $l$ 的物体张角 $\theta=l/D_A$, $D_A=D_M/(1+z)$. 每个光子能量下降 $1+z$, 到达间隔再增加 $1+z$, 所以 $F=L/[4\pi D_M^2(1+z)^2]$. 定义 $D_L$ 使 $F=L/(4\pi D_L^2)$, 得
\begin{equation}\label{eq:distances}
D_L=(1+z)D_M=(1+z)^2D_A.
\end{equation}
此关系需要光子数守恒和正常的度规光传播; 吸收、透镜选择效应或仪器选择会影响实际测量. 高红移物体的角大小不能用今日实际距离直接除, 宇宙学距离也不能统一称为一个「距离」而省掉定义.

\originalsection{超新星与参数退化}
Ia 型超新星通过光变宽度和颜色修正构造距离模数. 典型参数化为 $\mu_i=m_i^*-M+\alpha(s_i-1)-\beta c_i$, 与 $\mu_{\rm th}=5\log_{10}[D_L/(\mathrm{Mpc})]+25$ 比较. $M$ 与 $H_0$ 存在归一化退化, 不能只凭超新星相对距离各自精确确定两者. 若 $\sigma_i^2$ 依赖 $\alpha,\beta$, 简单只最小化 $\sum(\mu_i-\mu_{\rm th})^2/\sigma_i^2$ 会偏向放大分母; 需要完整统计模型或题面指定的迭代权重策略.

低红移展开为 $D_L=(cz/H_0)[1+(1-q_0)z/2+O(z^2)]$, $q_0=\Omega_{m0}/2-\Omega_{\Lambda0}$ 对尘埃与真空成立. 很低红移只确定第一项, 较高红移才能比较减速或加速历史. 星系间尘埃使超新星更暗, 可模拟更大距离; 颜色、红外、背景辐射和不同距离探针的比较为这种退化提供检验. 本卷不捏造没有输入数据的最佳拟合.

\originalsection{吸收统计与 Lyman--alpha 森林}
沿实际路径 $\dd l=c\dd z/[(1+z)H(z)]$, 共动密度恒定 $n_c$ 的吸收体有实际密度 $n_c(1+z)^3$. 对固定实际截面 $\sigma$, 遭遇概率给
\begin{equation}\label{eq:absorber}
\frac{\dd N}{\dd z}=n_c\sigma\frac{c(1+z)^2}{H(z)}.
\end{equation}
若指定固定共动截面, 则实际截面还按 $(1+z)^{-2}$ 变, 红移趋势因此改变. 题面中的截面约定必须先确定.

氢的 Lyman--alpha 共振发生在约 $1216$ 埃. 远处类星体光穿过多个中性氢区域, 不同位置的共振因红移落在不同观测波长, 形成森林. 在均匀 Hubble 流的窄线近似下, 由 $\tau_\nu=\int n_{\rm HI}\sigma_\nu\dd l$ 与频率沿程扫描的速率可得 Gunn--Peterson 光学深度
$\tau_{\rm GP}=\pi e^2f_\alpha\lambda_\alpha n_{\rm HI}/[m_ecH(z)]$, 此式使用高斯电磁单位. 即使中性分数很小也可能强烈吸收. 光电离平衡 $n_{\rm HI}\Gamma\simeq\alpha(T)n_en_p$ 使森林同时依赖密度、温度和电离背景, 因而吸收不是单一的密度计; 速度与热展宽也改变谱线.

\addtocontents{toc}{\protect\Needspace{3\baselineskip}}
\originalchapter{宇宙热史与轻元素}\label{ch:thermal}
\originalsection{相对论粒子与熵}
温度高于静止质量时, 粒子速度接近 $c$. 对自由玻色子与费米子, 占据数分别为 $(\ee^{E/T}-1)^{-1}$ 和 $(\ee^{E/T}+1)^{-1}$; 本节恢复 $\kb,\hbar,c$ 时写明单位. 动量积分给
\begin{equation}
u=\frac{\pi^2}{30}g_*\frac{(\kb T)^4}{(\hbar c)^3},\qquad
s=\frac{2\pi^2}{45}g_{*s}\kb\frac{(\kb T)^3}{(\hbar c)^3}.
\end{equation}
费米子每自由度对 $g_*$ 的权重为 $7/8$, 若不同粒子温度不同还乘 $(T_i/T)^4$; 熵权重则乘三次方. 辐射压强 $p=u/3$ 由各向同性动量流得到, $s=(u+p)/T$ 使用零化学势. 绝热膨胀时 $sa^3$ 守恒, 所以 $aT\propto g_{*s}^{-1/3}$. 粒子湮灭时 $g_{*s}$ 下降, 仍耦合的粒子升温得比简单 $a^{-1}$ 慢.

\originalsection{反应率与冻结}
是否保持平衡取决于每粒子反应率 $\Gamma=n\langle\sigma v\rangle$ 和膨胀率 $H$, 而不只是温度高低. $\Gamma\gg H$ 时体系能跟随平衡分布; $\Gamma\lesssim H$ 时反应来不及调整. 弱相互作用中微子粗略 $\Gamma\sim G_F^2T^5$, 辐射期 $H\sim T^2/\mpl$, 故宇宙冷却时前者下降更快, 在 MeV 尺度退耦. 此估计不包含精确输运和非瞬时退耦修正.

中微子先退耦后, 电子正电子湮灭只加热光子. 湮灭前与光子耦合的熵自由度为 $2+(7/8)4=11/2$, 后为 2, 因而 $T_\gamma/T_\nu=(11/4)^{1/3}$. 中微子温度继续按 $a^{-1}$ 下降. 暗物质若以热遗迹产生, 数密度满足 $\dot n+3Hn=-\langle\sigma v\rangle(n^2-n_{\rm eq}^2)$. 用 $Y=n/s$ 消去纯稀释后, 剩余丰度取决于冻结期间的反应积分, 不能只给 $n=n_{\rm eq}$ 并认为今天仍平衡.

\originalsection{核合成与复合}
弱平衡阶段的中子质子比近似 $n_n/n_p\simeq\exp[-(m_n-m_p)c^2/\kb T]$. 弱反应冻结后中子继续衰变, 直到氘能存活并迅速生成氦. 由于每重子有许多光子, 即使 $\kb T$ 已低于氘结合能, 高频尾仍能毁氘, 形成氘瓶颈. 当核反应展开, 大部分存活中子进入 ${}^4\mathrm{He}$, 故氦质量分数近似 $Y_p\simeq2r/(1+r)$, $r$ 是该时刻中子质子比. 精确 D、He、Li 丰度需要反应网络和速率数据; 此处不把估计标成精密 BBN 计算.

氢复合也发生在远低于 $13.6\,\mathrm{eV}$ 的温度, 原因同样与光子重子比有关. Saha 关系描述早期平衡, 但后期非平衡复合需要速率方程. 光子的 Thomson 散射率 $n_e\sigma_Tc$ 随自由电子减少而下降. 定义向今日的光学深度 $\tau$, 最后散射概率密度与 $\ee^{-\tau}|\dot\tau|$ 成正比. 复合、光子退耦与后来再电离是不同过程, 不能视为同一瞬间.

\addtocontents{toc}{\protect\Needspace{3\baselineskip}}
\originalchapter{扰动增长与结构形成}\label{ch:structure}
\originalsection{线性扰动}
写 $\rho=\bar\rho(1+\delta)$, 把物理速度分为 Hubble 流与特殊速度 $\boldsymbol u$. 在共动坐标中连续、Euler 与 Poisson 方程线性化为
\begin{align*}
\dot\delta+\frac1a\nabla\cdot\boldsymbol u&=0,\\
\dot{\boldsymbol u}+H\boldsymbol u&=-\frac1a\nabla\phi-\frac{c_s^2}a\nabla\delta,\\
\nabla^2\phi&=4\pi Ga^2\bar\rho\delta.
\end{align*}
对第一式求导, 消去速度与势, Fourier 模满足
\begin{equation}\label{eq:growth}
\ddot\delta_{\boldsymbol k}+2H\dot\delta_{\boldsymbol k}
+\left(\frac{c_s^2k^2}{a^2}-4\pi G\bar\rho\right)\delta_{\boldsymbol k}=0.
\end{equation}
此 Newton 近似用于远小于视界、非相对论组分的尺度, 超视界扰动需要相对论且必须指定规范. 在 Einstein--de Sitter 尘埃模型中 $H=2/(3t)$、$4\pi G\bar\rho=2/(3t^2)$, 代 $\delta\propto t^p$ 得 $p=2/3,-1$, 生长模 $\delta\propto a$. 辐射主导时物质扰动的增长受背景抑制, 暗能量主导后也不再近似 $D=a$.

\originalsection{统计描述}
取 $\delta(\boldsymbol x)=\int\dd^3k\,\delta(\boldsymbol k)\ee^{\ii\boldsymbol k\cdot\boldsymbol x}/(2\pi)^3$, 齐次场满足 $\langle\delta(\boldsymbol k)\delta^*(\boldsymbol k')\rangle=(2\pi)^3\delta_D(\boldsymbol k-\boldsymbol k')P(k)$. 相关函数 $\xi(r)$ 是 $P$ 的 Fourier 变换; $\Delta^2(k)=k^3P(k)/(2\pi^2)$ 是每对数波数的方差贡献. 原初谱经转移函数与生长因子传播, $P(k,t)=D(t)^2T(k)^2P_{\rm initial}(k)$ 只在适用的线性范围成立.

给定平滑窗口 $W(kR)$, $\sigma^2(R)=\int\Delta^2(k)|W(kR)|^2\dd\ln k$. 对实空间球顶帽 $W(x)=3(\sin x-x\cos x)/x^3$, 小 $x$ 极限为 1. 窗口及质量映射 $M=4\pi\bar\rho_0R^3/3$ 必须一起给出, 否则不同 $\sigma(M)$ 不能直接比较.

\originalsection{球对称坍缩}
均匀球顶帽在无压尘埃、无壳交叉条件下, 每层满足 $\ddot r=-GM/r^2$. 负能解写成 $r=A(1-\cos\eta)$、$t=B(\eta-\sin\eta)$、$A^3=GMB^2$. 在 $\eta=\pi$ 达最大半径, 在 $2\pi$ 形式上坍为一点. EdS 背景密度为 $1/(6\pi Gt^2)$, 代入球平均密度得到转向密度比 $9\pi^2/16$. 早期展开给线性外推 $\delta_L=(3/5)[3(\eta-\sin\eta)/4]^{2/3}$, 在形式坍缩时为 $\delta_c\simeq1.686$.

若能量守恒且最终满足维里关系, $W=-3GM^2/(5R)$ 导致 $R_{\rm vir}=R_{\rm ta}/2$. 因背景密度在 $t_{\rm coll}=2t_{\rm ta}$ 时下降四倍, 晕密度因半径减半提高八倍, 最终密度比为 $18\pi^2$. 这些数值是 EdS 理想化, 不可无条件用于任意红移的 $\Lambda$CDM.

\originalsection{质量函数与冷却}
Gaussian 平滑场超过阈值的体积分数是 $F(>M)=\frac12\operatorname{erfc}[\delta_c/(\sqrt2\sigma(M))]$. 在 $\sigma\to\infty$ 的小质量极限只到 $1/2$, 原始 Press--Schechter 计数漏掉嵌套区域. 锐波数滤波使 $\delta$ 随方差 $S$ 作 Markov 随机游走. 将阈值设为吸收壁, 镜像法从自由 Gaussian 核减去镜像核, 首次过壁分布为 $f(S)=\delta_c\exp[-\delta_c^2/(2S)]/(\sqrt{2\pi}S^{3/2})$, 积分为 1. 这解释因子 2, 但实空间顶帽游走并非 Markov, 椭球坍缩和数值晕质量函数还需更多修正.

气体能否形成恒星还取决于冷却. 光学薄气体损失率写作 $n_en_i\Lambda(T)$, 储能约为 $3n\kb T/2$, 因而 $t_{\rm cool}=3n\kb T/(2n_en_i\Lambda)$. 与 $t_{\rm dyn}\sim(G\rho)^{-1/2}$ 比较可估计哪些晕能快速形成冷气体. $\Lambda$ 的单位和电离、组成约定不能省略, 固定一常数冷却曲线只是数量级模型.

\originalsection{坍缩阈值的早期展开}
球顶帽的密度比可以直接用参数解而不用背诵常数. $\rho_{\rm halo}=3M/[4\pi A^3(1-\cos\eta)^3]$, EdS 背景 $\bar\rho=1/(6\pi GB^2(\eta-\sin\eta)^2)$, 结合 $A^3=GMB^2$ 得
\[
1+\delta=\frac92\frac{(\eta-\sin\eta)^2}{(1-\cos\eta)^3}.
\]
展开 $\eta-\sin\eta=\eta^3/6-\eta^5/120+\cdots$、$1-\cos\eta=\eta^2/2-\eta^4/24+\cdots$, 比值给 $\delta=3\eta^2/20+O(\eta^4)$. 同时 $t\simeq B\eta^3/6$, 线性增长 $t^{2/3}$ 因而唯一匹配为 $\delta_L=(3/20)[6(\eta-\sin\eta)]^{2/3}$, 与正文写法等价. 在 $\eta=2\pi$ 外推得 $\delta_c=(3/20)(12\pi)^{2/3}=1.68647$.

外推的 $\delta_L$ 不是坍缩时真实密度对比, 后者在理想无压模型发散. 两者用途不同: 前者把线性随机场与形成条件联系, 后者描述非线性形状. Press--Schechter 使用的是前者. 后续维里化并非由 $\eta=2\pi$ 的点状解自行产生, 而是额外加入能量守恒和稳态维里假设后用有限半径模型替换.

\addtocontents{toc}{\protect\Needspace{3\baselineskip}}
\originalchapter{微波背景与暴胀}\label{ch:cmb}
\originalsection{声学振荡与角谱}
复合前光子与重子强耦合. 光子提供压强, 重子增加惯性. 若 $R_b=3\rho_b/(4\rho_\gamma)$, 有效声速 $c_s=c/\sqrt{3(1+R_b)}$. 这是对混合物绝热扰动的 $\delta p/\delta\rho$ 计算, 不是重子的热声速. 引力势推动压缩, 压强推动回弹, 于是每个波数经历不同相位的声学振荡. 共动声视界 $r_s=\int_0^{t_*}c_s\dd t/a$ 与横向共动距离给角尺度 $\theta_s=r_s/D_M(z_*)$.

天空温度写为 $T(\boldsymbol n)=\bar T[1+\Theta(\boldsymbol n)]$, 展开 $\Theta=\sum a_{\ell m}Y_{\ell m}$. 统计各向同性使 $\langle a_{\ell m}a^*_{\ell'm'}\rangle=C_\ell\delta_{\ell\ell'}\delta_{mm'}$. 单一天空的估计 $\widehat C_\ell=(2\ell+1)^{-1}\sum_m|a_{\ell m}|^2$ 对 Gaussian 场的相对标准差为 $\sqrt{2/(2\ell+1)}$, 即宇宙方差. 小 $\ell$ 可用模式少, 提高仪器精度也不能消除这一统计限制.

在物质主导、大尺度、绝热扰动和忽略势随时间变化的近似下, 引力红移与本征温度共同给 $\Theta\simeq\Phi/(3c^2)$, 即普通 Sachs--Wolfe 项. 随时间变化的势还给积分项, 速度给 Doppler 项, 小尺度有扩散阻尼. 本卷解释来源与尺度, 不把完整 Einstein--Boltzmann 辐射输运求解标作已实现.

\originalsection{视界问题与加速膨胀}
在减速的辐射或物质期, 共动 Hubble 尺度 $(aH)^{-1}$ 增长. 今日遥远区域回溯到早期时, 可能没有足够时间通过信号交换达到相似温度. 暴胀的核心条件是 $\ddot a>0$, 等价于 $\dd(aH)^{-1}/\dd t<0$. 若一块原本有因果联系的区域在这一时期被拉伸, 以后可包含整个可观测宇宙. 这给视界问题一种机制, 并不自动证明某一个暴胀模型.

取 $c=\hbar=1$, 最小耦合标量场的均匀能量与压强为 $u=\dot\phi^2/2+V$、$p=\dot\phi^2/2-V$. 场方程 $\ddot\phi+3H\dot\phi+V'=0$ 与 $H^2=u/(3\bar M_{\rm Pl}^2)$ 相容, $\bar M_{\rm Pl}=(8\pi G)^{-1/2}$. 势能占优时负压强使膨胀加速. 慢滚近似忽略 $\ddot\phi$ 与动能, 得 $3H\dot\phi\simeq-V'$、$3\bar M_{\rm Pl}^2H^2\simeq V$. 其自洽条件为
\begin{equation}
\epsilon_V=\frac{\bar M_{\rm Pl}^2}2\left(\frac{V'}V\right)^2\ll1,
\qquad\left|\eta_V\right|=\left|\bar M_{\rm Pl}^2\frac{V''}V\right|\ll1.
\end{equation}
扩张倍数的对数 $N=\int H\dd t\simeq\bar M_{\rm Pl}^{-2}\int_{\phi_f}^{\phi_i}V\dd\phi/V'$; 所需 $N$ 随能标与再加热历史改变, 不能把固定「60」作为对所有模型的严格常数.

\originalsection{平直性与量子涨落的边界}
曲率偏离满足 $\Omega-1=kc^2/(a^2H^2)$. 减速期偏离放大, 加速期偏离缩小, 因而暴胀能把初始曲率压低. 暴胀结束后场能量必须转为粒子与辐射, 即再加热; 热大爆炸的初始条件并非与再加热无关.

自由轻标量在近 de Sitter 空间的量子涨落经拉伸可提供近尺度不变扰动. 在单场、规范真空、慢滚且扰动线性的假设下, 曲率扰动谱的领先阶结果为 $\mathcal P_{\mathcal R}=H^2/(8\pi^2\epsilon_H\bar M_{\rm Pl}^2)$, $\epsilon_H=-\dot H/H^2$. 这里将曲时空量子模的规范化作为高阶工具引用, 未完成规范不变量量子作用量、再加热微观过程、多场或永恒暴胀的完整论证. 这些边界不会影响前述背景动力学、热史和选定题解的自足推导.

\originalsection{声速和宇宙方差的推导}
紧耦合、绝热混合物中 $\delta\rho_b/\rho_b=3\delta\rho_\gamma/(4\rho_\gamma)$. 总压强扰动只有光子贡献 $\delta p=c^2\delta\rho_\gamma/3$, 总质量等效密度扰动 $\delta\rho=\delta\rho_\gamma[1+3\rho_b/(4\rho_\gamma)]$. 比值便是 $c_s^2=c^2/[3(1+R_b)]$. 若忽略宇宙膨胀和势缓慢变化, 对局部平面扰动的连续与 Euler 式求导消元得到 $\ddot\delta_\gamma+c_s^2k^2\delta_\gamma=\text{引力驱动项}$. 均匀解是余弦和正弦, 每个波数在最后散射时已有不同相位, 因而角谱有多峰而不是一条无结构温度曲线. 膨胀、重子加载及势变化会改变振幅和相位, 精密计算须恢复它们.

对于 Gaussian 各向同性天空, 在实球谐基中某个 $\ell$ 有 $2\ell+1$ 个独立实 Gaussian 模, 每个方差 $C_\ell$. 正规化的平方和为 $\chi^2_{2\ell+1}$, 其均值 $2\ell+1$、方差 $2(2\ell+1)$. 所以 $\operatorname{Var}(\widehat C_\ell)=2C_\ell^2/(2\ell+1)$. 这里假设完整天空、无噪声、Gaussian, 部分天空和噪声会进一步增大估计不确定性. 以复球谐直接计数时 $a_{\ell,-m}=(-1)^ma_{\ell m}^*$, 不可把 $m$ 与 $-m$ 又当作独立复变量并错误多算一倍.

\originalchapter{恒星观测、双星与潮汐演化习题}
本章对应8.901 PS1可公开读到题面的六个题组。PS1.1--4仅回指受限教材，未重构。解答全部为独立补写；原题要求的历史月食照片实际返回HTTP403，本章明确保留这一数据缺口。其余数值由随包代码实际求解，地月参数引用NASA事实表，题给的均匀球转动惯量仍按原要求使用。

\src{8.901 PS1.5，历史天文学，原题(a)--(c)，PDF第1--2页。}
\begin{exercise}
(a) 夏季某日Syene深井井底可被正午太阳直照，Alexandria同日正午太阳在天顶以南$7^\circ$；两地南北相距略少于5000 stadia，每stadium约160 m。把太阳光看成平行光，求地球半径，与6378 km比较。
(b) 使用原题指定Timothy Ferris月食合成照片，量暗本影圆直径及单个月亮直径；解释月心路径为何未经过本影中心。取太阳、月亮的角直径均$0.5^\circ$，由本影几何推月地半径比、误差及月地距离，以现代$6378/1738=3.67$比较。
(c) 月亮在上弦Q处满足地月日角$\angle EQS=90^\circ$。朔N到上弦Q比Q到望F少35 min，朔望周期29.53 d；用这一信息估日地距离／月地距离。
\end{exercise}
\begin{solution}
(a) 平行光使两地太阳天顶角差等于地心角，$R_E=800\ \mathrm{km}/(7\pi/180)=6548$ km，约比题给现代值大2.7\%。原题的“略少于”和stadium近似都影响精度，不应给到米级。

(b) 照片的\href{https://web.mit.edu/8.901/images/moon.png}{原链接}未取得，不能编造所量比值。设实际可量的比值为$q=d_{\rm umb}/d_M$，角直径$\alpha=0.5\pi/180$。太阳有有限角直径，本影半径随离地距离减小约$\alpha D/2$；月亮半径$R_M\simeq\alpha D/2$，所以
\[
 q=\frac{R_E-\alpha D/2}{R_M}=\frac{R_E}{R_M}-1,\quad
 \frac{R_E}{R_M}=q+1,\qquad \frac D{R_E}=\frac{2}{\alpha(q+1)}.
\]
忽略角直径误差时$\sigma_{R_E/R_M}=\sigma_q$、$\sigma_D/D=\sigma_q/(q+1)$。现代比值3.67对应$q\simeq2.67$和$D/R_E\simeq62.4$，只是\emph{由现代值反算的比较}，不是本照片测量。月轨道与黄道有倾角，月心可穿过本影圆的一条偏心弦，因而路径无需经过影心。

(c) 近匀角速近似下N到F耗时$T/2$。两段时间相差$\Delta t=35$ min，因此$t_{NQ}=T/4-\Delta t/2$。月亮由朔位置转过$\pi/2-\delta$，$\delta=\pi\Delta t/T$。直角三角形EQS给$D_M/D_S=\sin\delta$，故
\[
 \frac{D_S}{D_M}=\csc\left(\frac{\pi\times35}{29.53\times1440}\right)\simeq386.7.
\]
小角很难从肉眼相位时刻精确判断，且月轨道并非圆轨道。这是对题给理想信息的推算，不声称实际古代观测有这一精度。(b)的照片量测仍列为未取得依赖。
\end{solution}

\src{8.901 PS1.6，Kepler双曲运动，PDF第2页。}
\begin{exercise}
椭圆Kepler方程为$\omega t=\psi-e\sin\psi$。双曲轨道定义参数$H$使$r=a(e\cosh H-1)$，其中$e>1$、$a>0$，近心距离为$a(e-1)$。求从近心时刻起经过的时间$t(H)$。
\end{exercise}
\begin{solution}
以焦点为原点可取$X=a(e-\cosh H)$、$Y=a\sqrt{e^2-1}\sinh H$。距原点正是题给$r$。单位质量角动量为
\[
 \ell=X\dot Y-Y\dot X
 =a^2\sqrt{e^2-1}(e\cosh H-1)\dot H.
\]
Kepler轨道满足$\ell^2=GMa(e^2-1)$，所以$\dd t=\sqrt{a^3/(GM)}(e\cosh H-1)\dd H$。以$H=t=0$为近心，得
\[
 \sqrt{\frac{GM}{a^3}}\,t=e\sinh H-H.
\]
这里$M$是两体相对运动的总质量；$\sqrt{GM/a^3}$是时间尺度的倒数，双曲轨道没有绕行周期。$H<0$给接近近心的负时间，积分常数不能省略而又换任意起点。
\end{solution}

\src{8.901 PS1.7，地月潮汐，原题(a)--(i)，PDF第3页。}
\begin{exercise}
忽略月亮自转。地球自转频率$\omega$、月球圆轨道频率$\Omega$，质量$M_E,M_M$，地球惯量$I$、距离$a$。(a)写总能量与角动量；(b)以角动量消去$\omega$；(c)证明可写成$\epsilon=-1/s+\alpha(j-\sqrt s)^2$，其中$s=a/a_0$，能量单位$GM_EM_M/(2a_0)$，角动量单位$\mu a_0^2\Omega_0$，$\mu$为约化质量；(d)取$I=2M_ER_E^2/5$、$a_0=3.84\times10^5$ km，查质量求$\alpha,j$；(e)画能量曲线找两极值；(f)数值求根并证明两处$\omega=\Omega$，求同步周期；(g)求今日到未来同步态须失去的能量；(h)假设每天两次将海洋上层1 m水抬高1 m再放下，数个百分点机械能转成热，估耗散功率；(i)由能量与功率估同步所需时间。
\end{exercise}
\begin{solution}
(a)--(b) 二体质心系的圆轨道有$\Omega^2=G(M_E+M_M)/a^3$，故
\[
 E=-\frac{GM_EM_M}{2a}+\frac{I\omega^2}{2},\quad
 J=\mu a^2\Omega+I\omega,
\]
\[
 E=-\frac{GM_EM_M}{2a}
 +\frac{[J-\mu\sqrt{G(M_E+M_M)a}]^2}{2I}.
\]
潮汐把自转和轨道角动量交换；忽略外界力矩时$J$守恒而$E$可减小。若使用两体轨道能量却把全部月球相对速度动能都当质心动能，会多算地球平动部分。

(c)--(d) $\alpha=\mu a_0^2/I$、$j=J/(\mu a_0^2\Omega_0)$，代入即得所需无量纲式。取\href{https://nssdc.gsfc.nasa.gov/planetary/factsheet/earthfact.html}{NASA地球事实表}的$M_E=5.9722\times10^{24}$ kg、$R_E=6378.137$ km和恒星日23.9345 h，以及\href{https://nssdc.gsfc.nasa.gov/planetary/factsheet/moonfact.html}{月球事实表}的$M_M=7.346\times10^{22}$ kg，保持题给$I$，算得$\alpha=110.109$、$j=1.24809$。真实地球$I/(M_ER_E^2)$约0.3308，题设0.4是近似，不能悄悄换值而仍称原题结果。

(e)--(f) 导数为
\[
 \epsilon'(s)=\frac1{s^2}-\frac{\alpha(j-\sqrt s)}{\sqrt s},\qquad
 \frac\omega{\Omega_0}=\alpha(j-\sqrt s),\quad
 \frac\Omega{\Omega_0}=s^{-3/2}.
\]
导数为零恰好给$\omega=\Omega$。实际求根得$s_-=0.0423435$与$s_+=1.54595$；前者能量极大、不稳定，后者为极小。相应距离约16260和593644 km，周期分别0.23737 d与52.3637 d，两种钟同步。曲线由本式计算：
\begin{center}\begin{tikzpicture}\begin{axis}[width=.7\textwidth,height=4.6cm,xlabel={$s=a/a_0$},ylabel={$\epsilon$},xmin=0,xmax=3]
\addplot[structurecolor] table[x=s,y=epsilon,col sep=comma]{qa/stellar-tide.csv};
\draw[dashed] (axis cs:1,-70)--(axis cs:1,130);\node[anchor=north west] at (axis cs:1,120){今日};
\end{axis}\end{tikzpicture}\end{center}
系统今日在两极值之间，自转比轨道快，失去能量的外迁方向趋向外侧极小。内侧同步点不是今日系统会自行走向的稳定终点。

(g) 固定今日$J$比较两点，$\Delta E=[\epsilon(1)-\epsilon(s_+)]GM_EM_M/(2a_0)=2.45\times10^{29}$ J。(h) 原题没有给海洋面积和确切耗散百分比，示例取表面70\%为海洋、海水密度$1000\ \mathrm{kg\,m^{-3}}$、$g=9.81$和热比例$f=0.03$，则
\[
 P_{\rm tide}\sim\frac2{86400}
 [0.7(4\pi R_E^2)\rho_w(1\ \mathrm m)g(1\ \mathrm m)]f
 \simeq2.44\times10^{12}\ \mathrm W.
\]
(i) 恒功率假设下$\Delta E/P\simeq3.18\times10^9$ yr，结果正比$0.03/f$且反比假定海洋面积。真实海洋共振、潮汐响应、太阳力矩和地球结构随时间变，不能把这个练习当成太阳未来演化期间地月同步的确定预报。
\end{solution}

\src{8.901 PS1.8，食双星，PDF第3页。}
\begin{exercise}
两恒星在质心周围作圆轨道，中心相距$a$，半径$R_1,R_2$。视线与轨道角动量夹角$i\in[0,90^\circ]$。求恰能出现食的最小$i$。
\end{exercise}
\begin{solution}
合时中心连线的最小天空投影距离为$a\cos i$，两投影圆相切条件$a\cos i=R_1+R_2$。对互不重叠的球形星有
\[
 i_{\min}=\arccos\frac{R_1+R_2}{a}.
\]
若$R_1+R_2>a$，分离球模型已失效，不能继续对大于1的数取$\arccos$。偏心轨道还需用合时实际距离和近心方向，不能只把$a$替换成半长轴后沿用同式。
\end{solution}

\src{8.901 PS1.9，单线共生双星样本，原题(a)--(b)，PDF第3--4页。}
\begin{exercise}
单线光谱测到红巨星$M_1$的周期$P$和径向速度半振幅$K_1$，假定圆轨道。样本数值如下：
\begin{center}\begin{tabular}{lrrlrr}\toprule
对象&$P$(d)&$K_1$(km/s)&对象&$P$(d)&$K_1$(km/s)\\\midrule
EG And&482.6&7.3&Z And&759.0&6.7\\
T CrB&227.6&23.9&BF Cyg&757.2&6.7\\
V1329 Cyg&956.5&7.8&CI Cyg&853.8&6.7\\
AG Dra&548.7&5.9&V443 Her&599.4&2.5\\
BX Mon&1259&4.6&RS Oph&455.7&16.7\\
V2116 Oph&1042&16.0&AG Peg&818.2&5.4\\
AX Per&682.1&7.8&FG Ser&633.5&6.9\\
V343 Ser&450.5&2.7&&&\\\bottomrule
\end{tabular}\end{center}
(a) 从$f_1=(M_2\sin i)^3/(M_1+M_2)^2=4\pi^2(a_1\sin i)^3/(GP^2)$推$f_1(P,K_1)$并解释意义。
(b) 典型共生双星的红巨星$M_1=1.5\Ms$、白矮星$M_2=0.56\Ms$。随机轨道方向使$\cos i$均匀；画样本的严格小于型CDF及此模型期望CDF，判断是否相容，找异常对象，仅固定$M_1=1.5\Ms$时对伴星有何推断？
\end{exercise}
\begin{solution}
(a) $K_1=2\pi a_1\sin i/P$，故$f_1=PK_1^3/(2\pi G)$。它由可观测量决定，是质量与未知倾角的组合；$\sin i\leq1$且$M_1>0$给$f_1<M_2$，已知$M_1$时可求更强的最低$M_2$。

(b) 固定典型质量有$f_{\max}=0.56^3/(2.06)^2=0.0413837\Ms$。令$u=\cos i$在$[0,1]$均匀，$f=f_{\max}(1-u^2)^{3/2}$，所以
\[
 F(f)=\begin{cases}
 0,&f<0,\\
 1-\sqrt{1-(f/f_{\max})^{2/3}},&0\leq f\leq f_{\max},\\
 1,&f>f_{\max}.
 \end{cases}
\]
样本CDF在排序的第$j$个值取$(j-1)/15$，符合原题严格小于定义；右连续阶梯则在该点为$j/15$，两者跳点处不同但模型判断不变。实际计算与绘图：
\begin{center}\begin{tikzpicture}\begin{axis}[width=.7\textwidth,height=4.7cm,xlabel={$f_1/\Ms$},ylabel={CDF},xmin=0,xmax=.47,ymin=0,ymax=1.05,legend pos=south east]
\addplot[only marks,mark=*,structurecolor] table[x=f_solar,y=CDF_strict,col sep=comma]{qa/stellar-binary-cdf.csv};\addlegendentry{15个测量，严格小于}
\addplot[black] table[x=f_solar,y=CDF,col sep=comma]{qa/stellar-binary-expected.csv};\addlegendentry{固定典型质量，随机方向}
\end{axis}\end{tikzpicture}\end{center}
四个对象超出模型支持范围：T CrB、V1329 Cyg、RS Oph、V2116 Oph，$f_1/\Ms$分别0.32194、0.04703、0.21991、0.44222。仅保持$M_1=1.5\Ms$且$\sin i\leq1$，解$m^3/(1.5+m)^2=f_1/\Ms$，最低伴星质量依次1.391、0.590、1.158、1.630$\Ms$。V1329 Cyg只是略超典型0.56；其余三个显著偏离典型模型。这里没有原测量误差和样本选择函数，不能给精确总体显著性或从CDF唯一判定紧致伴星类型；原表数据事实保留Hinkle等(2003)来源，经原作业给出，不复制原论文版式。
\end{solution}

\src{8.901 PS1.10，光子计数和星等，原题(a)--(b)，PDF第4--5页。}
\begin{exercise}
Poisson计数的标准差为$\sqrt N$。(a) 测星等至$\pm0.02$需多少光子？一米口径、100\%总效率测$m_B=20$需多长曝光，忽略天光。(b) 改为$m_B=24$，天光22.5 mag/arcsec$^2$，星像分布于直径1角秒的圆；需多少源光子？一米及四米望远镜各需多久？
\end{exercise}
\begin{solution}
由$m=-2.5\log_{10}F+\mathrm{constant}$，小误差$\sigma_m=(2.5/\ln10)\sigma_F/F$。无天光时$N=[2.5/(0.02\ln10)]^2=2947$。记零星等B带光子通量$\Phi_{B0}$和收集面积$A=\pi D^2/4$，曝光
\[
 t_{20}=\frac{2947}{A\Phi_{B0}10^{-8}}.
\]
原题未给B带通、光谱和星等零点；100\%效率不决定零点，所以有秒数的答案须声明额外标定。作为可复现示例，取3950--4850埃的理想矩形带，AB零点$F_{\nu0}=3631$ Jy、恒定$F_\nu$，有
\[
 \Phi_0=\frac{F_{\nu0}}h\ln\frac{4850}{3950}
 =1.125\times10^{10}\ \mathrm{photons\,m^{-2}\,s^{-1}},\quad
 t_{20}(1\ \mathrm m)=33.36\ \mathrm s.
\]
这不是擅自声称真实Johnson B的Vega标定，只是补充标定下的时间示例；使用其他$\Phi_{B0}$按其反比缩放。

口径内天光／源光比$B/S=10^{0.4(24-22.5)}\pi(0.5)^2=3.1267$。背景均值若已准确知道，源净计数$S$的方差仍为$S+B$，故所需\emph{源}光子$S=2947(1+B/S)=12162$。总到达光子不是这个数，而是$S+B$。同一示例标定下，$t_{24}(1\ \mathrm m)=5481$ s，四米收集面积大16倍，$t_{24}(4\ \mathrm m)=342.5$ s。若背景由面积有限的另一个天空区估计，需再加背景估计方差；原题没有给该区面积，不能无条件忽略这一实际观测限制。
\end{solution}

\originalchapter{电离、不透明度与对流习题}
本章对应8.901 PS2.3--4的公开部分与PS3.1--2，答案为独立补写。PS2.4(a)、PS3.3--5的教材引用题面未公开，不以同主题自拟题替代。纯氦数值使用题目允许的密度$10^{-6}\ \mathrm{g\,cm^{-3}}$；分配函数限于基态，不能把结果推广到强非理想或简并等离子体。

\src{8.901 PS2.3，纯氦Saha方程，原题(a)--(d)。}
\begin{exercise}
固定密度、温度的纯氦气体有电离能$\chi_0=24.6$ eV、$\chi_1=54.4$ eV。所有原子、离子均在基态，数密度$n_e,n_0,n_1,n_2$分别为电子、中性氦、单电离氦、双电离氦；$n=n_0+n_1+n_2$，$x_e=n_e/n$、$x_i=n_i/n$，电中性。
(a) 确定统计权重和能量零点，写$n_1/n_0,n_2/n_1$；
(b) 消去其他变量，只留下$x_1,x_2$未知的两式，参数为$T,n$或$\rho=4n/N_A$；
(c) 在$4\times10^4\leq T\leq2\times10^5$ K密网格，选择$\rho=10^{-4},10^{-6},10^{-8}\ \mathrm{g\,cm^{-3}}$之一，用计算机同时求解并算$x_0,x_e$；
(d) 画$x_0,x_1,x_2$及电子比例，确定两电离阶段的半电离温度。
\end{exercise}
\begin{solution}
中性氦基态$1\,{}^1S_0$有$g_0=1$；类氢He$^+$的$1s$态总角动量$J=1/2$，$g_1=2$；裸核He$^{++}$按基态核自旋零取$g_2=1$。电子自旋另有因子2，不能再次混入离子权重。以每一步反应的游离电子动能零为参照，定义
\[
 b(T)=\left(\frac{2\pi m_e\kb T}{h^2}\right)^{3/2},\quad
 S_0=4b\ee^{-\chi_0/(\kb T)},\quad
 S_1=b\ee^{-\chi_1/(\kb T)}.
\]
于是$n_en_1/n_0=S_0$、$n_en_2/n_1=S_1$；第二步的能量是54.4 eV，不是把第一步和第二步相加的79 eV。

电中性给$x_e=x_1+2x_2$，数目守恒给$x_0=1-x_1-x_2$，故所求两式为
\[
 \frac{n(x_1+2x_2)x_1}{1-x_1-x_2}=S_0,\qquad
 \frac{n(x_1+2x_2)x_2}{x_1}=S_1.
\]
为数值稳定，也可固定$n_e$构造未归一权重$(1,S_0/n_e,S_0S_1/n_e^2)$，以对数求和归一，再用$n_e=n(x_1+2x_2)$求一个标量根。这同时保证各比例非负且总和为1，避免直接Newton迭代跳到无物理解。

实际对500个温度求解，$\rho=10^{-6}$、$n=\rho/(4m_u)$；全部点电荷约束残差小于$10^{-9}$。$T=10^5$ K时$(x_0,x_1,x_2)\simeq(3.7\times10^{-8},0.00217,0.99783)$，$x_e\simeq1.99783$。电子比例可到2，因为每个氦核释放两个电子，不是概率。
\begin{center}\begin{tikzpicture}\begin{axis}[width=.72\textwidth,height=4.8cm,xlabel={$T$（K）},ylabel={比例},xmin=40000,xmax=200000,ymin=0,ymax=2.05,legend pos=north east]
\addplot[structurecolor] table[x=T_K,y=He0,col sep=comma]{qa/stellar-helium-saha.csv};\addlegendentry{$x_0$}
\addplot[black,dashed] table[x=T_K,y=He1,col sep=comma]{qa/stellar-helium-saha.csv};\addlegendentry{$x_1$}
\addplot[gray] table[x=T_K,y=He2,col sep=comma]{qa/stellar-helium-saha.csv};\addlegendentry{$x_2$}
\addplot[structurecolor,dotted,thick] table[x=T_K,y=e_per_He,col sep=comma]{qa/stellar-helium-saha.csv};\addlegendentry{$x_e$}
\end{axis}\end{tikzpicture}\end{center}
第一阶段的半中性点$x_0=1/2$可能在题给绘图区间以下；必须另外扩大求根区间，不能把40\,000 K左端当成半电离温度。扩大区间后实际求得第一阶段约22047 K、第二阶段约53528 K，前者确在绘图区间以下；它们指该密度下的平衡比例，并非只由$\chi/\kb$决定。
\end{solution}

\src{8.901 PS2.4(b)，不透明度，(a)仅引用教材题面未取得。}
\begin{exercise}
假设地球空气具有太阳光球在5000埃处的质量不透明度$\kappa=0.264\ \mathrm{cm^2\,g^{-1}}$，空气密度$1.2\times10^{-3}\ \mathrm{g\,cm^{-3}}$。估能看穿多远。原题还括注光球密度为$2.5\times10^7\ \mathrm{g\,cm^{-3}}$。
\end{exercise}
\begin{solution}
光学深度$\tau=\kappa\rho\ell$，一个e折衰减长度为
\[
 \ell_{\tau=1}=\frac1{\kappa\rho}
 =3157\ \mathrm{cm}=31.57\ \mathrm m.
\]
这不是所有人眼识别阈值下同一“可视距离”，只是原题通常采用的平均自由程／光学深度尺度。括注光球密度明显不合通常光球数量级，可能缺负指数；它不参与已给$\kappa$与空气$\rho$的计算，故本卷保留原记录并注明问题，不自行改题给数据求另一答案。PS2.4(a)的教材题未重构。
\end{solution}

\src{8.901 PS3.1，辐射压力与Eddington极限，(a)--(c)。}
\begin{exercise}
(a) 光学薄云受吸收／动量转移力每单位质量为$\int\kappa_\nu F_\nu\dd\nu/c$；若$\kappa$不随频率变，证明被附近质量$M$、光度$L$的源辐射推出需$M/L<\kappa/(4\pi Gc)$。
(b) 云在距离$R$从静止出发，仅受引力与辐射，求终速并证明$v_\infty^2=(2GM/R)[\kappa L/(4\pi GMc)-1]$。
(c) 纯氢完全电离，$\sigma_T=6.65\times10^{-25}\ \mathrm{cm^2}$，用$\kappa=\sigma_T/m_H$估可不逐出氢的最大光度，并给$1.3\times10^{38}(M/\Ms)\ \mathrm{erg\,s^{-1}}$系数。
\end{exercise}
\begin{solution}
各向同性稳态源给$F=L/(4\pi r^2)$。每单位质量的径向净加速度是
\[
 \ddot r=\frac1{r^2}\left(\frac{\kappa L}{4\pi c}-GM\right)
 =\frac{GM}{r^2}(\Gamma-1),\quad
 \Gamma=\frac{\kappa L}{4\pi GMc}.
\]
(a) 静止云向外加速需要$\Gamma>1$，即所求比值。(b) 乘$\dot r$，由$R$到无穷积分：
$v_\infty^2/2=GM(\Gamma-1)\int_R^\infty\dd r/r^2=GM(\Gamma-1)/R$。$\Gamma<1$时不能把负数平方根当成实终速；静止起点的云向内落。

(c) 一个氢核约对应一个自由电子，$\kappa_T=\sigma_T/m_H\simeq0.398\ \mathrm{cm^2\,g^{-1}}$。于是
\[
 L_{\rm Edd}=\frac{4\pi GMcm_H}{\sigma_T}
 \simeq1.26\times10^{38}\frac M\Ms\ \mathrm{erg\,s^{-1}},
\]
按题面精度写1.3。中性气体、混合元素、谱线不透明度、几何准直和光学厚度都可改变有效条件；这一推导不是所有恒星与吸积流的硬光度上限。
\end{solution}

\src{8.901 PS3.2，对流稳定性，原题(a)--(b)。}
\begin{exercise}
已知稳定条件$\rho P'/(\gamma P)-\rho'>0$，其中撇号为对$r$求导，绝热气团满足$P\propto\rho^\gamma$。
(a) 用$P=\rho\kb T/(\mu m_p)$证明$T'>(1-1/\gamma)(T/P)P'$；再用辐射恒星结构式推出$L<(1-1/\gamma)16\pi a_{\rm rad}cT^4GM(r)/(3\kappa P)$。
(b) 理想单原子气体证明cgs形式$L<1.22\times10^{-18}\mu T^3M/(\kappa\rho)$。
\end{exercise}
\begin{solution}
固定成分时$\rho'/\rho=P'/P-T'/T$，代入稳定条件后直接得$T'/T>(1-1/\gamma)P'/P$。由于正常恒星$P'<0$，定义$\nabla=\dd\ln T/\dd\ln P$再除负数须反向，稳定条件才是$\nabla<\nabla_{\rm ad}=1-1/\gamma$。

流体静力$P'=-GM(r)\rho/r^2$，扩散输运给$T'=-3\kappa\rho L/(16\pi a_{\rm rad}cr^2T^3)$，所以
\[
 \nabla_{\rm rad}=\frac{3\kappa PL}{16\pi a_{\rm rad}cGMT^4}.
\]
要求$\nabla_{\rm rad}<\nabla_{\rm ad}$就得原题光度上限。单原子$\gamma=5/3$，以理想气体消$P$，系数为
\[
 \frac25\frac{16\pi a_{\rm rad}cGm_p}{3\kb}
 \simeq1.22\times10^{-18}\quad(\mathrm{cgs}).
\]
它是辐射梯度是否超过绝热梯度的局部检查，出现对流后实际梯度和输运需重算。部分电离会降低$\gamma$和$\nabla_{\rm ad}$，成分随$r$变化时还需分子量梯度，不能把上式当成无条件稳定性结论。
\end{solution}

\originalchapter{多方恒星与核反应习题}
本章完整对应8.901 PS4.2--8。PS4.1仅有教材引用而无公开题面，未重构。多方数值实际计算至首零点，核质量使用已核验NIST同位素原子质量，禁止以元素的平均原子量替代。所有解答为独立补写；计算结果是所声明模型的结果，不是完整恒星演化。

\src{8.901 PS4.2，多方恒星的性质，原题(a)--(d)。}
\begin{exercise}
设$P=K\rho^{1+1/n}$，$\rho=\rho_c\theta^n$，$r=\alpha\xi$，$\alpha^2=(n+1)K\rho_c^{1/n-1}/(4\pi G)$，Lane--Emden方程为$(\xi^2\theta')'=-\xi^2\theta^n$，中心$\theta(0)=1,\theta'(0)=0$，首零点$\xi_1$。
证明(a)理想气体有$T=T_c\theta$；(b)$M=-4\pi\alpha^3\rho_c\xi_1^2\theta'_1$；(c)$\bar\rho/\rho_c=-3\theta'_1/\xi_1$；(d)$P_c=GM^2/[4\pi(n+1)(\theta'_1)^2R^4]$。
\end{exercise}
\begin{solution}
(a) 成分固定的理想气体给$T=(\mu m_p/\kb)P/\rho=(\mu m_pK/\kb)\rho^{1/n}$，由$\rho/\rho_c=\theta^n$得$T/T_c=\theta$。简并压力占优时不能继续解释$\theta$为温度。

(b) 质量积分$M=4\pi\alpha^3\rho_c\int_0^{\xi_1}\xi^2\theta^n\dd\xi$，用微分方程把积分写成$-[\xi^2\theta']_0^{\xi_1}$。中心项因正则条件为零，得到所求质量。(c) $R=\alpha\xi_1$，$\bar\rho=3M/(4\pi R^3)$，立即给出所求比例。

(d) 中心压力$P_c=K\rho_c^{1+1/n}=4\pi G\alpha^2\rho_c^2/(n+1)$。从(b)有$M^2=16\pi^2\alpha^6\rho_c^2\xi_1^4(\theta'_1)^2$，除以$R^4=\alpha^4\xi_1^4$并整理，得题式。三条关系把一个无量纲首零点解转成实际$M,R,P_c$；还需指定$K,\rho_c$才能得到独一的恒星模型。
\end{solution}

\src{8.901 PS4.3，Lane--Emden数值解，原题(a)--(b)。}
\begin{exercise}
(a) 证明中心的正则幂级数仅含偶次项，并求$\theta=1-\xi^2/6+n\xi^4/120+\cdots$。
(b) 对$n=1,1.5,2,2.5,3,3.5$，化二阶式为$\theta'=u,\ u'=-2u/\xi-\theta^n$，用中心展开启动实际数值积分。把六条$\theta(\xi)$画在一张图，六条$\theta^n(\xi)$画在另一张图。
\end{exercise}
\begin{solution}
(a) 在中心设正则$\theta=\sum a_j\xi^j$，$\theta'(0)=0$给$a_1=0$。$\theta''+2\theta'/\xi$的$\xi^{j-2}$系数是$j(j+1)a_j$；右侧$-\theta^n$。常数项给$6a_2=-1$，下一偶次项给$20a_4=-na_2=n/6$，所以$a_2=-1/6,a_4=n/120$。由方程在$\xi\to-\xi$下不变且正则中心初值相同，局部解为偶函数；逐阶也可见所有奇次系数为零。

(b) 不能在$\xi=0$直接计算$2u/\xi$。本计算取$\xi_{\rm start}=10^{-5}$，$\theta=1-\xi^2/6+n\xi^4/120$、$u=-\xi/3+n\xi^3/30$，用DOP853、相对误差$2\times10^{-11}$、绝对误差$2\times10^{-13}$、最大步0.025积分；检测首零点并停止。非整数$n$不可把负$\theta$代入实数幂后继续走“外部恒星”。所得
\begin{center}\begin{tabular}{rrrr}\toprule
$n$&$\xi_1$&$-\xi_1^2\theta'_1$&$\rho_c/\bar\rho$\\\midrule
1.0&3.14159&3.14159&3.28987\\
1.5&3.65375&2.71406&5.99070\\
2.0&4.35287&2.41105&11.4025\\
2.5&5.35528&2.18720&23.4065\\
3.0&6.89685&2.01824&54.1825\\
3.5&9.53581&1.89056&152.884\\\bottomrule
\end{tabular}\end{center}
\begin{center}
\begin{tikzpicture}\begin{axis}[width=.47\textwidth,height=5cm,xlabel={$\xi$},ylabel={$\theta$},xmin=0,xmax=10,ymin=0,ymax=1.05,legend style={font=\footnotesize},legend pos=north east]
\addplot table[x=xi,y=theta,col sep=comma]{qa/stellar-lane-n1p0.csv};\addlegendentry{$n=1$}
\addplot table[x=xi,y=theta,col sep=comma]{qa/stellar-lane-n1p5.csv};\addlegendentry{$1.5$}
\addplot table[x=xi,y=theta,col sep=comma]{qa/stellar-lane-n2p0.csv};\addlegendentry{$2$}
\addplot table[x=xi,y=theta,col sep=comma]{qa/stellar-lane-n2p5.csv};\addlegendentry{$2.5$}
\addplot table[x=xi,y=theta,col sep=comma]{qa/stellar-lane-n3p0.csv};\addlegendentry{$3$}
\addplot table[x=xi,y=theta,col sep=comma]{qa/stellar-lane-n3p5.csv};\addlegendentry{$3.5$}
\end{axis}\end{tikzpicture}\hfill
\begin{tikzpicture}\begin{axis}[width=.47\textwidth,height=5cm,xlabel={$\xi$},ylabel={$\rho/\rho_c$},xmin=0,xmax=10,ymin=0,ymax=1.05,legend style={font=\footnotesize},legend pos=north east]
\addplot table[x=xi,y=rho_over_central,col sep=comma]{qa/stellar-lane-n1p0.csv};\addlegendentry{$n=1$}
\addplot table[x=xi,y=rho_over_central,col sep=comma]{qa/stellar-lane-n1p5.csv};\addlegendentry{$1.5$}
\addplot table[x=xi,y=rho_over_central,col sep=comma]{qa/stellar-lane-n2p0.csv};\addlegendentry{$2$}
\addplot table[x=xi,y=rho_over_central,col sep=comma]{qa/stellar-lane-n2p5.csv};\addlegendentry{$2.5$}
\addplot table[x=xi,y=rho_over_central,col sep=comma]{qa/stellar-lane-n3p0.csv};\addlegendentry{$3$}
\addplot table[x=xi,y=rho_over_central,col sep=comma]{qa/stellar-lane-n3p5.csv};\addlegendentry{$3.5$}
\end{axis}\end{tikzpicture}
\end{center}
独立解析检查使用$n=1$解$\theta=\sin\xi/\xi$，500点最大误差小于$10^{-9}$。其他$n$的首零点和质量积分仍是数值结果；更大的$n$趋向强中心集中，不是同样$M,R$下的六颗实际恒星。
\end{solution}

\src{8.901 PS4.4，经典Coulomb碰撞，原题(a)--(c)。}
\begin{exercise}
(a) 忽略量子作用，假设速度为Maxwell--Boltzmann均方根速度十倍的核可克服两质子Coulomb势垒，估所需温度并与太阳中心温度比较。
(b) 求十倍均方根速率与均方根速率处的粒子数比。
(c) 错误地把太阳当纯氢估核数，判断快质子是否足够解释太阳光度。
\end{exercise}
\begin{solution}
(a) 需先指定核接触距离$r_N$和十倍速率的参考分布。势垒是$E_C=e^2/(4\pi\epsilon_0r_N)=1.44\,(\mathrm{fm}/r_N)$ MeV。若使用\emph{相对速度}的Maxwell分布，约化质量$\mu_r=m_p/2$、$v_{\rm rms,rel}^2=3\kb T/\mu_r$，十倍速率的相对能量为$150\kb T$，于是
\[
 T\gtrsim1.11\times10^8\left(\frac{1\ \mathrm{fm}}{r_N}\right)\mathrm K.
\]
例如$r_N=2$ fm给$5.6\times10^7$ K，比太阳中心约$1.5\times10^7$ K高。若“十倍”指一个质子相对静止靶的速率，则相对能量为$75\kb T$，所需温度再大两倍；两快质子迎头相撞又是另一条件。原题未给核半径和这个速度约定，不能隐藏它们给一个无条件精确温度。

(b) 速率密度$f(v)\propto v^2\exp[-mv^2/(2\kb T)]$，$v_{\rm rms}^2=3\kb T/m$，相同小速率区间内
\[
 \frac{f(10v_{\rm rms})}{f(v_{\rm rms})}
 =100\ee^{-3(100-1)/2}=3.22\times10^{-63}.
\]
这是速率分布之比，不是三维速度密度之比；后者不含$v^2$。若要求整段高速尾概率，应另积$f$，结论仍极小。

(c) $N_p=\Ms/m_p\simeq1.19\times10^{57}$。即使粗略把全部核数乘上述比值，也只有约$4\times10^{-6}$个，实际高速尾同样不足一个核。经典热尾无法支持持续太阳核反应；量子隧穿允许能量远低于接触势垒的核发生反应，下一题给出它的关键指数。把速度尾比直接理解为一个核的分数是概率／期望值的解释，非真实存在可切开的质子。
\end{solution}

\src{8.901 PS4.5，Coulomb势垒穿透。}
\begin{exercise}
电荷$Z_1e,Z_2e$、质量数$A_1,A_2$的核迎头碰撞，相对速度$v$、远处质心动能$E$。用WKB证明穿透概率近似$\mathcal T=\exp[-2\pi Z_1Z_2e^2/(\hbar v)]$（Gaussian单位），并写成$\exp[-31.28Z_1Z_2A^{1/2}E_{\rm keV}^{-1/2}]$，$A=A_1A_2/(A_1+A_2)$。
\end{exercise}
\begin{solution}
令$C=Z_1Z_2e^2$、约化质量$\mu_r=Am_u$，外转折点$b=C/E$。忽略核接触半径相对$b$时，WKB指数为
\[
 \frac2\hbar\int_0^b\sqrt{2\mu_r(C/r-E)}\dd r.
\]
以$r=b\sin^2\theta$，$\sqrt{b/r-1}=\cot\theta$，积分变为
$2b\sqrt{2\mu_rE}\int_0^{\pi/2}\cos^2\theta\dd\theta
=\pi C\sqrt{2\mu_r/E}/2$。再乘$2/\hbar$，用$E=\mu_rv^2/2$，得$2\pi C/(\hbar v)$。

系数$2\pi\alpha_{\rm em}\sqrt{m_uc^2/(2\,\mathrm{keV})}\simeq31.29$，与题给31.28在常数取舍精度内相同。真实下限为$r_N>0$且还有核反应因子$S(E)$，这里主指数要求$r_N\ll b$；不能把Coulomb透射概率直接当成完整核反应截面。
\end{solution}

\src{8.901 PS4.6，核反应能量，原题(a)--(e)。}
\begin{exercise}
计算下列反应$Q$值，释放能量为正，说明吸热或放热：
(a)$^{12}\mathrm C+{}^{12}\mathrm C\to{}^{24}\mathrm{Mg}+\gamma$；
(b)$^{12}\mathrm C+{}^{12}\mathrm C\to{}^{16}\mathrm O+2\,{}^4\mathrm{He}$；
(c)$^{19}\mathrm F+{}^1\mathrm H\to{}^{16}\mathrm O+{}^4\mathrm{He}$；
(d)$p+p\to{}^2\mathrm H+e^++\nu$；
(e)$^{15}\mathrm N+{}^1\mathrm H\to{}^{12}\mathrm C+{}^4\mathrm{He}$。
\end{exercise}
\begin{solution}
采用\href{https://physics.nist.gov/cgi-bin/Compositions/stand_alone.pl?ele=H}{NIST同位素表}与相应He、Mg、F、N、O页面核验的中性原子质量，单位$u$：
\begin{center}\begin{tabular}{lr@{\qquad}lr}\toprule
同位素&原子质量&同位素&原子质量\\\midrule
$^1$H&1.00782503223&$^2$H&2.01410177812\\
$^4$He&4.00260325413&$^{12}$C&12（定义）\\
$^{16}$O&15.99491461957&$^{24}$Mg&23.985041697\\
$^{19}$F&18.99840316273&$^{15}$N&15.00010889888\\\bottomrule
\end{tabular}\end{center}
(a)、(b)、(c)、(e)两边总电荷相同，原子电子静质量相消，电子结合能在本题MeV精度下忽略，直接以原子质量差乘$m_uc^2\simeq931.49$ MeV。(d)初态两个质子，末态一个氘核加一个正电子；改用中性原子质量时须扣两份电子静质量：
\[
 Q_{pp}=[2M(^1\mathrm H)-M(^2\mathrm H)-2m_e]c^2.
\]
实际结果：
\begin{center}\begin{tabular}{lrl}\toprule
反应&$Q$（MeV）&性质\\\midrule
(a)&13.9336&放热\\
(b)&$-0.11283$&吸热\\
(c)&8.11361&放热\\
(d)&0.42022&放热\\
(e)&4.96549&放热\\\bottomrule
\end{tabular}\end{center}
随后正电子与环境电子湮灭可再提供1.022 MeV，但不属于已写反应(d)本身；中微子分走的能量也不全转成恒星热。$Q>0$只判断能量允许性，实际反应速率和选择规则还需核物理。
\end{solution}

\src{8.901 PS4.7，假想燃烧反应的温度比。}
\begin{exercise}
假想元素G、J、M的质量数分别$A,2A,4A$，电荷数为质量数的一半。两颗恒星分别进行$^A\mathrm G+{}^A\mathrm G\to{}^{2A}\mathrm J+\gamma$、$^{2A}\mathrm J+{}^{2A}\mathrm J\to{}^{4A}\mathrm M+\gamma$，中心反应率相同。忽略密度和$S(E)$等指数前因子，估中心温度比并评论合理性。
\end{exercise}
\begin{solution}
热率的主指数来自$\exp[-E/(\kb T)-\sqrt{E_G/E}]$，驻点满足$E_0=(E_G(\kb T)^2/4)^{1/3}$，峰指数为$3[E_G/(4\kb T)]^{1/3}$。$E_G\propto\mu_r(Z_1Z_2)^2$。第二反应每个质量数和电荷数都加倍，故$\mu_r$大2倍，$Z_1Z_2$大4倍，$E_G$大$2\times4^2=32$倍。保持主指数相同需
\[
 \frac{T_2}{T_1}=32.
\]
高电荷核要在更高温度燃烧符合物理方向；因指数对Coulomb排斥很敏感，变化不是简单两倍。这个模型故意忽略前因子和反应共振，不能把32当成真实所有相邻燃烧阶段的精确比。
\end{solution}

\src{8.901 PS4.8，两Maxwell分布的相对速度。}
\begin{exercise}
两群可区分粒子质量$m_1,m_2$，分别在同一温度$T$下有三维Maxwell分布
\[
 f_i(\mathbf v_i)=\left(\frac{m_i}{2\pi\kb T}\right)^{3/2}
 \exp\left(-\frac{m_iv_i^2}{2\kb T}\right).
\]
证明相对速度$\mathbf v=\mathbf v_1-\mathbf v_2$也服从Maxwell分布；可使用质心速度。
\end{exercise}
\begin{solution}
定义$M=m_1+m_2$、$\mu_r=m_1m_2/M$、$\mathbf V=(m_1\mathbf v_1+m_2\mathbf v_2)/M$。反变换$\mathbf v_1=\mathbf V+(m_2/M)\mathbf v$、$\mathbf v_2=\mathbf V-(m_1/M)\mathbf v$，每个分量的Jacobian绝对值为1，故六维体积不变。直接展开动能得
\[
 m_1v_1^2+m_2v_2^2=MV^2+\mu_rv^2.
\]
归一系数也因$m_1m_2=M\mu_r$分成两个Gaussian。将$\mathbf V$积分掉，得到
\[
 f_{\rm rel}(\mathbf v)=
 \left(\frac{\mu_r}{2\pi\kb T}\right)^{3/2}
 \exp\left(-\frac{\mu_rv^2}{2\kb T}\right).
\]
再积方向得速率密度$4\pi v^2f_{\rm rel}$，正是$\langle\sigma v\rangle$使用的分布。若两群有不同温度或相对整体漂移，结果的方差／均值会改变，不能直接沿用本题同温、零漂移式。
\end{solution}

\originalchapter{红巨星演化、氦闪与中微子习题}
本章对应8.901 PS5.2--4。原题给的经验函数与理想热模型均实际积分，解答为独立补写。它们是研究某个反馈或尺度的练习，不是完整恒星演化代码；额外参数和原式问题分别说明。

\src{8.901 PS5.2，红巨星支上的演化。}
\begin{exercise}
低质量恒星离开主序后，薄氢壳燃烧围绕简并氦核。题给$L=2\times10^5\Ls(m_c/\Ms)^6$、$R=3700\Rs(m_c/\Ms)^4$。用$L=0.007\dot m_cc^2$，从$m_c=0.1\Ms$积分到$0.45\Ms$，画$R,L,T_{\rm eff}$随时间，以及$\log L$对$\log T_{\rm eff}$的HR轨迹，温度向左增加。
\end{exercise}
\begin{solution}
令$x=m_c/\Ms$，$C=2\times10^5\Ls/(0.007\Ms c^2)$，则$\dot x=Cx^6$。积分给
\[
 x(t)=[x_i^{-5}-5Ct]^{-1/5},\quad
 t_f=\frac{x_i^{-5}-x_f^{-5}}{5C}\simeq1.035\times10^{10}\ \mathrm{yr}.
\]
把$x(t)$代回得到$L,R$；$T_{\rm eff}=T_\odot(L/\Ls)^{1/4}(R/\Rs)^{-1/2}\propto x^{-1/2}$。核增长时整星光度和半径迅速增加，表面反而冷却。取$T_\odot=5772$ K，初点$(L/\Ls,R/\Rs,T)=(0.2,0.37,6346)$，终点约$(1661,151.7,2991)$。
\begin{center}
\begin{tikzpicture}\begin{axis}[width=.47\textwidth,height=4.5cm,xlabel={$t$（yr）},ylabel={太阳单位},ymode=log,legend pos=north west]
\addplot[structurecolor] table[x=t_yr,y=L_solar,col sep=comma]{qa/stellar-giant-evolution.csv};\addlegendentry{$L/\Ls$}
\addplot[black,dashed] table[x=t_yr,y=R_solar,col sep=comma]{qa/stellar-giant-evolution.csv};\addlegendentry{$R/\Rs$}
\end{axis}\end{tikzpicture}\hfill
\begin{tikzpicture}\begin{axis}[width=.47\textwidth,height=4.5cm,xlabel={$t$（yr）},ylabel={$T_{\rm eff}$（K）}]
\addplot[structurecolor] table[x=t_yr,y=Teff_K,col sep=comma]{qa/stellar-giant-evolution.csv};
\end{axis}\end{tikzpicture}
\end{center}
\begin{center}\begin{tikzpicture}\begin{axis}[width=.55\textwidth,height=4.4cm,xmode=log,ymode=log,x dir=reverse,xlabel={$T_{\rm eff}$（K，左侧更热）},ylabel={$L/\Ls$}]
\addplot[structurecolor,thick] table[x=Teff_K,y=L_solar,col sep=comma]{qa/stellar-giant-evolution.csv};
\end{axis}\end{tikzpicture}\end{center}
极长时间主要来自$x=0.1$附近极低$L$的外推，而该初点连半径都不足太阳，不能说它是真实已膨胀红巨星。答案遵照题给函数，同时指出其低核质量端的适用性问题；不通过改起点缩短时间而仍声称解了原题。
\end{solution}

\src{8.901 PS5.3，氦闪，原题全部条件与数值要求。}
\begin{exercise}
一克纯氦初始$\rho=2\times10^5\ \mathrm{g\,cm^{-3}}$、$T_0=1.5\times10^8$ K。只考虑三氦反应，保持密度、氦质量分数$Y=1$且无热流，题给
\[
 \epsilon_{3\alpha}=5.1\times10^8\rho^2Y^3T_9^{-3}\ee^{-4.4027/T_9}
 \quad\mathrm{erg\,g^{-1}s^{-1}},\quad T_9=T/10^9\ \mathrm K.
\]
电子比热$C_{V,e}=1.35\times10^5T x\sqrt{1+x^2}/\rho$，$\rho=2\times10^6x^3$；离子比热按原式$C_{V,\mathrm{He}}=3\kb/(4m_p)$。实际求$T(t)$，到非相对论简并分界$\rho=10^{-8}T^{3/2}$停止，画时间以天为单位的温度曲线；给氦闪时间、数值精度检查，并说明结果的物理合理性。
\end{exercise}
\begin{solution}
固定$\rho$使$x=(0.1)^{1/3}=0.46416$。热方程$C_V(T)\dot T=\epsilon(T)$，其中$C_V=C_{V,e}+C_{V,\mathrm{He}}$。温度单调增加，可将时间微分方程改写成
\[
 t(T)=\int_{T_0}^T\frac{C_V(T')}{\epsilon(T')}\dd T'.
\]
这仍是实际数值计算，且无需用固定时间步跨过升温末段。分界$T_d=(\rho/10^{-8})^{2/3}=7.368\times10^8$ K。自适应求积相对误差$10^{-10}$所得$t(T_d)=667284.2$ s，即7.7232 d；温度加倍时667265.3 s，距分界约19 s。若定义“闪开始”为温度加倍，则为7.7230 d；模型没有天然不连续的开始时刻，必须给定义。
\begin{center}\begin{tikzpicture}\begin{axis}[width=.72\textwidth,height=4.8cm,xlabel={$t$（d）},ylabel={$T$（K）},legend pos=north west]
\addplot[structurecolor] table[x expr=\thisrow{t_printed_s}/86400,y=T_K,col sep=comma]{qa/stellar-helium-flash.csv};\addlegendentry{原题比热}
\addplot[black,dashed] table[x expr=\thisrow{t_ideal_ion_s}/86400,y=T_K,col sep=comma]{qa/stellar-helium-flash.csv};\addlegendentry{单原子离子比热修正}
\end{axis}\end{tikzpicture}\end{center}
纯氦离子的经典定容热容为$3\kb/2$，每克离子数$1/(4m_p)$，所以理想离子比热应为$3\kb/(8m_p)$，原式大两倍。不能静默修正后只给一个时间。另算修正式得到$t(T_d)=489384.3$ s，即5.6642 d。图中两条曲线分别保留题给计算和物理比热判断。

升温增强反应率，而简并电子压力暂时不强依赖温度，缺少普通热膨胀降温反馈，因而可有热失控。固定密度、无热流、固定$Y$在真实氦闪中会失效；停止边界也只是近似简并判据。应重算膨胀、燃料和输运，不能把这几天理解为对某颗红巨星实际爆发日期的预测。
\end{solution}

\src{8.901 PS5.4，中微子振荡，(a)--(b)。}
\begin{exercise}
两种味态$|\nu_e\rangle=\cos\theta|\nu_1\rangle+\sin\theta|\nu_2\rangle$、$|\nu_\mu\rangle=-\sin\theta|\nu_1\rangle+\cos\theta|\nu_2\rangle$，质量态能量$E_j=\sqrt{p^2c^2+m_j^2c^4}$。$t=0$在原点产生动量$p$的纯电子中微子，向$z$传播。
(a)用自由平面波、小质量展开求波函数，提出共同相位，取$E\simeq pc,z\simeq ct$；
(b)证明转成缪子味的概率$P=\sin^2(2\theta)\sin^2[c^3\Delta m^2z/(4\hbar E)]$和振荡长度。原作业提示太阳内部MSW效应，引用2004年$\sin^22\theta=0.82\pm0.07$、$\Delta m^2c^4=(8.2\pm0.6)\times10^{-5}\ \mathrm{eV^2}$。
\end{exercise}
\begin{solution}
小质量时$E_j\simeq pc+m_j^2c^3/(2p)$。共同平面相位$\exp[-\ii pc(t-z/c)/\hbar]$不影响味测量，移除后
\[
 |\psi(z)\rangle=\ee^{-\ii m_1^2c^3z/(2\hbar E)}\cos\theta|\nu_1\rangle
 +\ee^{-\ii m_2^2c^3z/(2\hbar E)}\sin\theta|\nu_2\rangle.
\]
投影到$\nu_\mu$得振幅$\sin\theta\cos\theta(\ee^{-\ii\phi_2}-\ee^{-\ii\phi_1})$；用$|\ee^{\ii a}-\ee^{\ii b}|^2=4\sin^2[(a-b)/2]$即得题式。概率周期使正弦自变量增加$\pi$，所以
\[
 L_{\rm osc}=\frac{4\pi\hbar E}{c^3|\Delta m^2|}.
\]
需要非零混合与质量差；绝对质量共同增量只改共同相位，振荡不单独确定最轻态质量。

太阳电子背景对电子味多一个前向散射势，改变传播本征态和有效混合；密度缓变可有MSW增强。因此太阳缺失量不必在所有能量都只由地日几何相位解释。两味平面波近似忽略波包分离、三味和介质非均匀性；2004年数值保留年代，不作当前最佳值。
\end{solution}

\originalchapter{白矮星冷却与恒星风习题}
本章完整对应8.901 PS6.1--5，均为独立补写。原作业中经验关系按其年代和单位使用；额外的包层模型、标准质量和初始质量分别声明，不把模型演算冒充当前精密观测。

\src{8.901 PS6.1，40 Eri B白矮星，(a)--(d)。}
\begin{exercise}
40 Eri B和C绕行周期247.9 yr、视差0.201角秒、相对轨道半长轴6.89角秒、$a_B/a_C=0.37$；A约在400 AU外，本题只处理B+C。
(a)求B、C质量；(b)B的绝对热星等9.6，求光度；
(c)$T_{\rm eff}=16900$ K，求半径、平均密度；
(d)与Sirius B的$R=0.008\Rs$、$\bar\rho=3.0\times10^6\ \mathrm{g\,cm^{-3}}$比较并说明密度差。
\end{exercise}
\begin{solution}
距离$D=1/0.201=4.975$ pc，半长轴$a=6.89D=34.279$ AU。在太阳质量、AU、yr单位中$M_B+M_C=a^3/P^2=0.6554\Ms$。质心关系$M_Ba_B=M_Ca_C$给$M_C/M_B=0.37$，所以$M_B=0.4784\Ms,M_C=0.1770\Ms$；位置比不能误当同方向质量比。

明确取热星等零点$M_{{\rm bol},\odot}=4.74$，$L_B/\Ls=10^{-0.4(9.6-4.74)}=0.01138$。有效温度定义给$R_B/\Rs=\sqrt{L_B/\Ls}(5772/16900)^2=0.01244$，平均密度$3M_B/(4\pi R_B^3)=3.50\times10^5\ \mathrm{g\,cm^{-3}}$；平均密度不是中心密度。

Sirius B平均密度约高8.6倍、半径更小。非相对论电子简并支持给大致$R\propto M^{-1/3}$，较大质量可有较小半径和更高密度；组成、温度及相对论修正亦影响。不能从“更亮”直接推“更密”。这里保持题给2006年轨道参数，不作为40 Eri当前更新质量。
\end{solution}

\src{8.901 PS6.2，结晶与潜热冷却，(a)--(d)。}
\begin{exercise}
纯碳白矮星$M=0.6\Ms,R=0.012\Rs$。相邻离子尺度定义$4\pi r_i^3/3=Am_p/\rho$，Coulomb参数$\Gamma=Z_i^2e^2/(r_i\kb T)$。
(a)求$r_i$；(b)以$\Gamma=160$估内部结晶温度；
(c)估此温度的光度，包层质量分数$X=0,Y=0.9,Z=0.1$；
(d)每核释放$\kb T$潜热时，可维持此光度多久？
\end{exercise}
\begin{solution}
(a) 平均密度$\bar\rho=4.895\times10^5\ \mathrm{g\,cm^{-3}}$，碳$A=12,Z_i=6$，得$r_i=2.139\times10^{-10}$ cm。
(b) Gaussian单位$e^2=\alpha\hbar c$，$T_c=36e^2/(160r_i\kb)=1.758\times10^6$ K。实际核心密度大于平均密度，结晶会分区发生；此单温度是题设平均模型估计。

(c) 题面没有给包层不透明度、边界或$L(T_c)$；成分本身不唯一决定光度。可明确补入Mestel型示例：理想包层$P=\rho\kb T/(\mu m_p)$，$\mu^{-1}=3Y/4+Z/2$，辐射Kramers不透明度$\kappa=\kappa_0\rho T^{-7/2}$，静力质量近$M$，表面$P,T$近零。将扩散式除以静力式，有
\[
 \frac{\dd T}{\dd P}=\frac{3\kappa_0L\mu m_p}{16\pi a_{\rm rad}cGM\kb}
 P T^{-15/2}.
\]
积分得
\[
 P^2=\frac{64\pi a_{\rm rad}cGM\kb}{51\kappa_0L\mu m_p}T^{17/2}.
\]
以理想气体压力等于非相对论电子压力$K_{\rm NR}\rho^{5/3}$作简化匹配，$\rho_b=[\kb T_c/(\mu m_pK_{\rm NR})]^{3/2}$，从而
\[
 L=\frac{64\pi a_{\rm rad}cGM}{51\kappa_0}
 \left(\frac{\mu m_p}{\kb}\right)^4K_{\rm NR}^3T_c^{7/2}.
\]
示例额外定义$\kappa_0=4.6712\times10^{23}$（相应cgs单位），$\mu=1.3793,\mu_e=2$，实际得到$L=6.81\times10^{-7}\Ls$。这不是题面给的系数；$L\propto\kappa_0^{-1}$，更透明包层、不同匹配或实际核密度都会显著改值。通式明确表达缺少哪些物理输入，不将示例说成原题唯一光度。

(d) 核数$N_i=M/(12m_p)$，潜热$E_{\rm lat}=N_i\kb T_c$，固定光度时间$\tau=E_{\rm lat}/L$。上述示例给$1.75\times10^{11}$ yr，巨大值提示极简不透明包层与平均密度组合不可作为真实冷却年龄：光度不会恒定，核心结晶结构不同，还存在普通热容量等能量。原题条件下可给$r_i,T_c$与条件化$L,\tau$；若需真实冷却预测必须补实际包层输运和结晶过程，不能复制未取得的教材冷却图补空。
\end{solution}

\src{8.901 PS6.3，白矮星计时与银河年龄。}
\begin{exercise}
本地白矮星光度函数截断$L\simeq10^{-4.4}\Ls$。初始主序质量$0.6<M/\Ms<10$的形成白矮星前时间满足
\[
 \log_{10}\tau_{\rm pn}(\mathrm{yr})=9.921-3.6648\ell+1.9697\ell^2-0.9369\ell^3,
 \quad\ell=\log_{10}(M/\Ms).
\]
随后冷却式为$\tau_{\rm wd}=8.8\times10^6(A/12)^{-1}(M_{\rm wd}/\Ms)^{5/7}(\mu/2)^{-2/7}(L/\Ls)^{-5/7}$ yr。按原题纯碳$A=\mu=12$及标准质量，估银河盘年龄；假设样本足够大，最暗端来自能成为白矮星的最大初始质量，并解释原因。晕的截断$T_{\rm eff}=4000$ K，取标准质量、半径再估年龄。
\end{exercise}
\begin{solution}
总年龄是前身演化$\tau_{\rm pn}$加冷却$\tau_{\rm wd}$。较高质量前身较早变成白矮星，有更长的冷却时间；足够大样本才可能采到这种暗端物体。最大前身质量未给，示例取$8\Ms$，代多项式得$3.37\times10^7$ yr；白矮星另取$0.6\Ms$，不能继续用$8\Ms$算冷却。

原题写$A=\mu=12$；这里忠实按这一给定冷却参数，盘冷却$5.088\times10^9$ yr，加前身时间约5.12 Gyr。晕若取$R=0.012\Rs$，则$L/\Ls=(0.012)^2(4000/5772)^4=3.321\times10^{-5}$，冷却约5.791 Gyr，总约5.82 Gyr。

原$\mu$称平均分子量，但充分电离纯碳气体的通常粒子分子量是12/7，不是12。若把冷却式$\mu$解释为这种气体分子量，时间比原计算大$7^{2/7}\simeq1.74$；不能悄悄混用。真实银河年龄还受初始—最终质量关系、形成历史、包层、结晶潜热和样本完备性影响。这里是给定时钟的条件化结果；题面已给截断事实，无需复制未取得的教材光度函数图。
\end{solution}

\src{8.901 PS6.4，热亮星的风，(a)--(b)。}
\begin{exercise}
题给$\log_{10}(\dot M v_\infty R^{1/2})=-1.37+2.07\log_{10}(L/10^6)$，$M,R,L$为太阳单位，$\dot M$为$\Ms$/yr，$v_\infty$为km/s。物理逃逸速度$v_{\rm esc}=[2(1-\Gamma_e)GM/R]^{1/2}$，在$T\gtrsim2.1\times10^4$ K原题取$v_\infty=v_{\rm esc}/2.6$。
(a)忽略辐射压，画$L=10^5,3\times10^5,10^6,2\times10^6\Ls$下$\log\dot M$对$M$；
(b)取$\Gamma_e=\kappa L/(4\pi cGM)$、$\kappa=0.3\ \mathrm{cm^2g^{-1}}$重画，说明影响。
\end{exercise}
\begin{solution}
记太阳单位$m,l,r$，逃逸速度$617.7\sqrt{(1-\Gamma_e)m/r}$ km/s。乘题式的$\sqrt r$后半径消掉，所以
\[
 \dot M=\frac{10^{-1.37}(l/10^6)^{2.07}}
 {(617.7/2.6)\sqrt m\sqrt{1-\Gamma_e}}\ \Ms/\mathrm{yr}.
\]
原题使用除2.6的约定，不因其他资料常见另一标定而无说明改成乘2.6。无辐射时$\dot M\propto m^{-1/2}$；同质量加入辐射使$\dot M$乘$(1-\Gamma_e)^{-1/2}$。实际四个光度的两组曲线：
\begin{center}\begin{tikzpicture}\begin{axis}[width=.76\textwidth,height=5.4cm,xmode=log,ymode=log,xlabel={$M/\Ms$},ylabel={$\dot M$（$\Ms$/yr）},legend style={font=\footnotesize},legend pos=north east,unbounded coords=jump]
\addplot[black,dashed] table[x=M_solar,y=Mdot_no_Gamma,col sep=comma]{qa/stellar-wind-L100000.csv};\addlegendentry{$10^5$，无辐射}
\addplot[black] table[x=M_solar,y=Mdot_Gamma,col sep=comma]{qa/stellar-wind-L100000.csv};\addlegendentry{$10^5$，含辐射}
\addplot[gray,dashed] table[x=M_solar,y=Mdot_no_Gamma,col sep=comma]{qa/stellar-wind-L300000.csv};\addlegendentry{$3\times10^5$，无辐射}
\addplot[gray] table[x=M_solar,y=Mdot_Gamma,col sep=comma]{qa/stellar-wind-L300000.csv};\addlegendentry{$3\times10^5$，含辐射}
\addplot[structurecolor,dashed] table[x=M_solar,y=Mdot_no_Gamma,col sep=comma]{qa/stellar-wind-L1000000.csv};\addlegendentry{$10^6$，无辐射}
\addplot[structurecolor] table[x=M_solar,y=Mdot_Gamma,col sep=comma]{qa/stellar-wind-L1000000.csv};\addlegendentry{$10^6$，含辐射}
\addplot[main,dotted] table[x=M_solar,y=Mdot_no_Gamma,col sep=comma]{qa/stellar-wind-L2000000.csv};\addlegendentry{$2\times10^6$，无辐射}
\addplot[main,thick] table[x=M_solar,y=Mdot_Gamma,col sep=comma]{qa/stellar-wind-L2000000.csv};\addlegendentry{$2\times10^6$，含辐射}
\end{axis}\end{tikzpicture}\end{center}
$\Gamma_e<1$才有实逃逸速度，含辐射曲线只画满足该域的点。$\Gamma_e\to1^-$时经验式形式发散，表示假设失效，不是恒星可维持无限质量流；$\Gamma_e\geq1$不画成复数速度曲线。
\end{solution}

\src{8.901 PS6.5，冷亮星的Reimers风，(a)--(d)。}
\begin{exercise}
题给$\dot M=4\times10^{-13}\eta L/(gR)\ \Ms$/yr，$L,R,g$均为太阳单位，$\eta$约为1。
(a)估$M=1\Ms,L=7000\Ls,T=3000$ K的AGB星失质量率；
(b)错误地假定$L,R,\eta$恒定，求$M(t)$；
(c)用$L=7000,R=310,M_0=1,\eta=1$画质量对$10^5$ yr单位时间；
(d)多久减至$0.6\Ms$碳氧核心？
\end{exercise}
\begin{solution}
太阳单位下$g=m/r^2$，给定流失量的正值为$4\times10^{-13}\eta lr/m$，恒星质量导数为负数。(a)光度和温度给$r=\sqrt{7000}(5772/3000)^2=309.7$，初始流失率约$8.67\times10^{-7}\Ms$/yr。

(b) $t$用yr、$m=M/\Ms$，有$m\dot m=-4\times10^{-13}\eta lr$，故
\[
 m(t)=\sqrt{m_0^2-8\times10^{-13}\eta lr\,t}.
\]
不能保持初始$g$不变再算线性失质量，因为$g$也随质量变。
(c)--(d) 按题指定$r=310$，到$m_f=0.6$需
\[
 t_f=\frac{1-0.6^2}{8\times10^{-13}\times7000\times310}
 =3.687\times10^5\ \mathrm{yr}.
\]
\begin{center}\begin{tikzpicture}\begin{axis}[width=.6\textwidth,height=4.3cm,xlabel={$t$（$10^5$ yr）},ylabel={$M/\Ms$},ymin=.55,ymax=1.02]
\addplot[structurecolor] table[x=t_1e5yr,y=M_solar,col sep=comma]{qa/stellar-reimers.csv};
\end{axis}\end{tikzpicture}\end{center}
在给定核质量停止。更晚平方根可失去实数解，是固定光度半径模型被过度外推，不能解释为真实核心质量为负或恒星突然全部消失。
\end{solution}

\addtocontents{toc}{\protect\Needspace{3\baselineskip}}
\originalchapter{恒星与气体习题}
本章均为 AI 独立推导. 题目采用原题参数; 未指定组成或燃料比例时, 解答列出所用假设. 原题的质量半径系数与物理常数近似另作说明.

\src{8.901, Spring 2006, PS2.1, (a)--(d).}
\begin{exercise}
给定 $B_\nu=2h\nu^3/[c^2(\ee^{h\nu/\kb T}-1)]$.
(a) 从频率谱推导波长谱 $B_\lambda$.
(b) 积分推导 $F=\sigma T^4$, 可用 $\int_0^\infty x^3/(\ee^x-1)\dd x=\pi^4/15$, 核查 $\sigma=5.67\times10^{-5}\,\mathrm{erg\,cm^{-2}s^{-1}K^{-4}}$.
(c) 求 $B_\lambda$ 峰值的 Wien 位移规律, 数值求解 $5(1-\ee^{-x})=x$.
(d) 求 $B_\nu$ 峰值与温度的关系; 两种峰对应的光子能量应否相同?
\end{exercise}
\begin{solution}
(a) 因 $\nu=c/\lambda$, 保留 Jacobian 得
$B_\lambda=B_\nu(c/\lambda)c/\lambda^2=2hc^2/[\lambda^5(\ee^{hc/\lambda\kb T}-1)]$.
(b) 半球角积分是 $\pi$, 换元 $x=h\nu/\kb T$ 得 $F=2\pi(\kb T)^4\int x^3/(\ee^x-1)\dd x/(h^3c^2)$, 所以 $\sigma=2\pi^5\kb^4/(15h^3c^2)$, 实算 $5.6703744\times10^{-5}$, 单位如题面.
(c) 对 $\lambda^{-5}/(\ee^{A/\lambda}-1)$ 求对数导数, 驻点满足 $5(1-\ee^{-x})=x$. 非零根 $x=4.9651142$, 故 $T\lambda_{\max}=hc/(\kb x)=0.289777\,\mathrm{cm\,K}$.
(d) 对频率谱同理得 $3(1-\ee^{-x})=x$, 根 $x=2.8214394$, $\nu_{\max}/T=5.87893\times10^{10}\,\mathrm{Hz\,K^{-1}}$. 两峰对应不同光子能量; 它们最大化的是两个不同横轴的密度, Jacobian 不是常数.
\end{solution}

\src{8.901 PS2.2, (a)--(c).}
\begin{exercise}
纯氢气体密度和温度固定, 电离能 $13.6\,\mathrm{eV}$, 原子与离子均取基态.
(a) 写出中性与电离氢的 Saha 方程并合理选配分函数.
(b) 用电中性 $n_e=n_1$ 与总核数 $n=n_0+n_1$ 消去密度, 以 $x=n_1/n$ 表示, 检查温度极限.
(c) 取 $n=\rho N_A$, $N_A=6.023\times10^{23}\,\mathrm{g^{-1}}$, 求 $x=1/2$ 的密度温度路径并在 $10^{-10}$--$10^{-2}\,\mathrm{g\,cm^{-3}}$ 作双对数图.
\end{exercise}
\begin{solution}
取氢基态电子简并 $U_0=2$, 裸质子内部权重 $U_1=1$, 核自旋共同因子约去, 定义\[S(T)=\left(\frac{2\pi m_e\kb T}{h^2}\right)^{3/2}\ee^{-\chi/\kb T}.\] 方程为 $n_en_1/n_0=S$. 代约束得 $x^2/(1-x)=S/n$. 令 $q=S/n$, 正根可稳定写为 $x=2q/[q+\sqrt{q^2+4q}]$. $T\to0$ 时 $q\to0$, $x\to0$; 模型形式上的 $T\to\infty$ 给 $x\to1$, 但极高温时非相对论假设需更新.

半电离时 $S=n/2$, 所以 $\rho=2S(T)/N_A$. 本图按题面 $N_A$ 作图; 数值脚本另使用 $n=\rho/m_p$ 的现代近似, 端点温度相差约千分之几, 不把两种常数定义混为同一次实算. 路径单调: 降低密度使电离所需温度降低.
\begin{center}\begin{tikzpicture}\begin{loglogaxis}[width=.75\linewidth,height=5cm,xlabel={$T\ (\mathrm K)$},ylabel={$\rho\ (\mathrm{g\,cm^{-3}})$},xmin=7000,xmax=80000,ymin=1e-10,ymax=1e-2]
\addplot[structurecolor,domain=8700:63000,samples=120] {2*2.41468e15*x^1.5*exp(-157821.4/x)/6.023e23};
\end{loglogaxis}\end{tikzpicture}\end{center}
\end{solution}

\src{8.901 PS5.1, (a)--(b).}
\begin{exercise}
主序低端 $M=0.085\Ms$, $\log_{10}T_e=3.438$, $\log_{10}(L/\Ls)=-3.297$; 高端 $M=90\Ms$, 对应为 $4.722,6.045$.
(a) 估计两者氢燃烧寿命. 低质量星完全对流, 可用全部氢; 高质量星只用内部约 $10\%$.
(b) 求两者半径.
\end{exercise}
\begin{solution}
题面未给氢丰度, 保留 $X$ 再以 $X=0.7$ 示例. 寿命 $0.007fXM c^2/L$ 在低端 $f=1$ 给 $1.22\times10^{13}\,\mathrm{yr}$, 高端 $f=0.1$ 给 $5.88\times10^5\,\mathrm{yr}$. 这些是把当前端点光度固定后的燃料估计, 不是完整演化年龄; 低端远大于宇宙年龄并不矛盾. 由 $L=4\pi R^2\sigma T_e^4$, $R/\Rs=\sqrt{L/\Ls}(5772\,\mathrm K/T_e)^2$, 依次为 $0.0996$ 与 $12.62$.
\end{solution}

\src{8.901 PS5.5, (a)--(c).}
\begin{exercise}
多方星 $P=K\rho^{1+1/n}$ 满足 Lane--Emden 方程, $R=\alpha\xi_1$, $M=4\pi\alpha^3\rho_c[-\xi_1^2\theta'(\xi_1)]$, $\alpha^2=(n+1)K\rho_c^{(1-n)/n}/(4\pi G)$.
(a) 消去 $\rho_c$ 推导质量半径关系.
(b) 对非相对论白矮星 $n=3/2$, 核查原题系数 $R=1.122\times10^4\,\mathrm{km}(\rho_c/10^6\,\mathrm{g\,cm^{-3}})^{-1/6}(\mu_e/2)^{-5/6}$、$M=0.4964\Ms(\rho_c/10^6)^{1/2}(\mu_e/2)^{-5/2}=0.7011\Ms(R/10^4\,\mathrm{km})^{-3}(\mu_e/2)^{-5}$.
(c) 对相对论白矮星 $n=3$, 核查 $R=3.347\times10^4\,\mathrm{km}(\rho_c/10^6)^{-1/3}(\mu_e/2)^{-2/3}$、$M=1.457\Ms(\mu_e/2)^{-2}$. 原题取 $n_e=\rho/(\mu_em_p)$、$\mu_e\simeq A/Z$.
\end{exercise}
\begin{solution}
令 $C=(n+1)K/(4\pi G)$、$w_n=-\xi_1^2\theta'(\xi_1)$. 当 $n\ne1$, 从 $R=\sqrt C\rho_c^{(1-n)/(2n)}\xi_1$ 解出 $\rho_c=(R/(\sqrt C\xi_1))^{2n/(1-n)}$. 代入质量得到
\[
M=4\pi C^{n/(n-1)}R^{(3-n)/(1-n)}\xi_1^{(n-3)/(1-n)}w_n.
\]
由动量积分得 $K_{\rm NR}=(3\pi^2)^{2/3}\hbar^2/[5m_e(\mu_em_b)^{5/3}]$、$K_{\rm R}=(3\pi^2)^{1/3}\hbar c/[4(\mu_em_b)^{4/3}]$. $m_b$ 是每核子质量近似. 数值积分给 $(\xi_1,w_n)=(3.653754,2.714055)$ 与 $(6.896849,2.018236)$. 于是分别有 $R\propto\rho_c^{-1/6}\mu_e^{-5/6}$、$M\propto\rho_c^{1/2}\mu_e^{-5/2}$, 和 $R\propto\rho_c^{-1/3}\mu_e^{-2/3}$、$M\propto\mu_e^{-2}$, 均符合题面尺度律.

原题印刷系数使用历史常数与每核子质量近似. 严格按现代 $m_b=m_p$ 实算, 非相对论组为 $11148\,\mathrm{km}$ 与 $0.4872\Ms$, 相对论组为 $33297\,\mathrm{km}$ 与 $1.4353\Ms$, 与原题差约百分之一至二. 若用 $m_b=m_u$ 并使 $\mu_e=A/Z$, 质量约为 $0.4961\Ms$、$1.4562\Ms$, 更接近原题. 不能同时要求现代质子质量的精确代入与印刷近似系数逐位相等. $n=3$ 的极限质量表示相对论简并支撑的尺度, 不是完整稳定性分析.
\end{solution}

\src{8.901 PS10.1, (a)--(c), 原作业不提交、不评分.}
\begin{exercise}
质量 $M_1,M_2$ 的圆轨道双星在共转坐标中的 Roche 势为 $\Phi=-GM_1/|\boldsymbol r-\boldsymbol r_1|-GM_2/|\boldsymbol r-\boldsymbol r_2|-|\boldsymbol\Omega\times\boldsymbol r|^2/2$. 星 2 经 $L_1=(x_L,0,0)$ 溢出.
(a) 用 $\nabla\Phi(L_1)=0$ 证明横向势增量数量级为 $\Omega^2(\Delta y)^2$.
(b) 用能量说明溢出斑块半径 $\Delta y\sim c_s/\Omega$.
(c) 推导 $-\dot M_2\sim\rho c_s^3P_{\rm orb}^2/(4\pi)$, 并说明深入恒星内部时流率为何增加.
\end{exercise}
\begin{solution}
应估计势\textbf{差}, 势零点可任意选择. 一阶 Taylor 项为零, 二阶横向系数为 $GM_1/d_1^3+GM_2/d_2^3-\Omega^2$, 与轨道 $\Omega^2$ 同阶, 故 $\Phi(x_L,\Delta y,0)-\Phi(L_1)\sim\Omega^2(\Delta y)^2$. 热能每单位质量约 $c_s^2$, 能越过势垒的横向范围满足 $\Omega^2(\Delta y)^2\lesssim c_s^2$. 穿流速度约 $c_s$, 面积约 $\pi(c_s/\Omega)^2$, 故 $-\dot M_2\sim\pi\rho c_s^3/\Omega^2=\rho c_s^3P_{\rm orb}^2/(4\pi)$. 系数仅为数量级. 深处 $\rho,T$ 增加, $c_s\propto T^{1/2}$, 流率随 $\rho T^{3/2}$ 增加.
\end{solution}

\src{8.901 PS10.2, (a)--(c).}
\begin{exercise}
中心星质量 $M$, 薄吸积盘在 $z=0$ 平面, 半径 $r$ 处半厚度 $H\ll r$, 盘自引力忽略.
(a) 写垂直静水平衡式.
(b) 证明 $H\simeq(c_s/v_K)r$, $c_s\simeq\sqrt{P/\rho}$.
(c) 解释薄盘条件与超声速轨道运动的等价性.
\end{exercise}
\begin{solution}
精确垂直重力为 $-GMz/(r^2+z^2)^{3/2}$, 在 $z\ll r$ 时 $\partial_zP=-\rho\Omega_K^2z$. 若等温, $P=c_s^2\rho$, 积分得 $\rho(z)=\rho(0)\exp[-z^2/(2H^2)]$, $H=c_s/\Omega_K=c_sr/v_K$. 非等温盘中这是尺度估计. 所以 $H/r=c_s/v_K$, 薄盘要求轨道速度远大于声速. 原题 (b) 的大写 $R$ 按上下文指当地半径 $r$.
\end{solution}

\src{8.901 PS10.3, (a)--(c).}
\begin{exercise}
均匀分子云 $T=150\,\mathrm K$、粒子数密度 $10^8\,\mathrm{cm^{-3}}$, 边缘压强近零. 径向运动方程 $\rho\ddot r=-GM_r\rho/r^2-\dd P/\dd r$.
(a) 用理想气体估计起始 $|\dd P/\dd r|\sim P_c/R_J$, 取 Jeans 质量与半径.
(b) 验证压强与引力在数量级上相当.
(c) 说明等温收缩后压强相对重力变弱.
\end{exercise}
\begin{solution}
假设分子为纯 ${\rm H}_2$, $\mu=2$; 题面未指定组成, 其他选择改变系数. 依第\ref{ch:compact}章的约定, $\rho=2m_pn$, $c_s^2=\kb T/(2m_p)$, 得 $R_J=986\,\mathrm{AU}$、$M_J=2.265\Ms$. $P_c/R_J\simeq1.40\times10^{-21}\,\mathrm{Pa\,m^{-1}}$, $GM_J\rho/R_J^2\simeq4.62\times10^{-21}\,\mathrm{N\,m^{-3}}$, 比值约 $0.30$, 在粗估精度内相当.

收缩后同一质量 $M$ 固定, $\rho\propto R^{-3}$. 等温时压强梯度 $P/R\propto R^{-4}$, 引力力密度 $GM\rho/R^2\propto R^{-5}$, 前者与后者比值正比于 $R$. 一旦引力胜过压力, 后续等温缩小不会单靠压力恢复平衡.
\end{solution}

\src{8.901 PS10.4.}
\begin{exercise}
对上一题的 Jeans 云估计引力势能每单位体积, 与 $B=10\,\mu\mathrm G$ 的磁能密度比较, 判断磁场是否显著影响该模型的坍缩.
\end{exercise}
\begin{solution}
均匀球 $|W|=3GM^2/(5R)$, 除体积得 $u_G=4.09\times10^{-7}\,\mathrm{J\,m^{-3}}$. $10\,\mu\mathrm G=10^{-9}\,\mathrm T$, SI 磁能密度 $u_B=B^2/(2\mu_0)=3.98\times10^{-13}\,\mathrm{J\,m^{-3}}$, 比值约 $9.7\times10^{-7}$. 对所给弱场和致密均匀云, 初始磁能很小. 结论不包括磁通冻结后的场增长、磁支撑的方向性和其他云参数.
\end{solution}

\src{8.901 PS10.5, (a)--(d).}
\begin{exercise}
(a) 写含固定轴旋转的云径向自由落体式. 初始径向速度零, $r_f\ll r_0$, 以角动量守恒证明停转向半径 $r_f\simeq\omega_0^2r_0^4/(2GM_r)$. 原题允许全程刚体转动这一不准确假设.
(b) $M=1\Ms,r_0=0.5\,\mathrm{pc},r_f=100\,\mathrm{AU}$, 求初始角速度与边缘线速度.
(c) 初始均匀球、末态均匀盘, 以总角动量守恒求末态边缘旋转速度.
(d) 求末态公转时间并与 Kepler 周期比较, 解释差异.
\end{exercise}
\begin{solution}
对无转矩质量元 $j=\omega_0r_0^2$, $\ddot r=-GM_r/r^2+j^2/r^3$. 乘 $\dot r$ 积分, 得 $\dot r^2/2-GM_r/r+j^2/(2r^2)=-GM_r/r_0+j^2/(2r_0^2)$. 若右侧相对转向点各项可忽略, 得 $r_f=j^2/(2GM_r)$. 这是偏心轨道的近日点估计, 不是圆化半径 $j^2/(GM_r)$.

代入给定量得 $\omega_0=2.6473\times10^{-16}\,\mathrm{s^{-1}}$, $v_0=408.4\,\mathrm{cm\,s^{-1}}$. (c) 使用总惯量 $I_0=2Mr_0^2/5$、$I_f=Mr_f^2/2$, 因而 $\omega_f=(4/5)\omega_0(r_0/r_f)^2$, $v_f=3.370\,\mathrm{km\,s^{-1}}$. 周期为 $884\,\mathrm{yr}$, Kepler 周期约 $1000\,\mathrm{yr}$. 末速是 $(4/5)\sqrt2$ 倍 Kepler 速度. (a) 保持单元 $j$, (c) 改用全体惯量并改变形状, 两者不构成严格自洽的单一刚体流体解; 也没有耗散圆化过程, 因此两周期不应精确相同. 这一差异不能被误写成稳定盘可以任意超 Kepler 转动.
\end{solution}

\originalchapter{真实星系图像与光谱的测量习题}
本章完整对应8.902 PS1的三个题组，均为独立补写解答。\href{https://ocw.mit.edu/courses/8-902-astrophysics-ii-fall-2023/mit8_902_f23_problem_set_1_additional_files.zip}{官方附加ZIP}已实际读取，四份FITS数组与独立读取器交叉核对一致。原始SDSS照片不收录；以下曲线为像素数值的环统计、谱线拟合与TeX重绘，保留SDSS署名和原文件例外许可。分析代码、输入哈希、拟合区间与测量结果见\texttt{scripts/analyze\_galaxy\_data.py}和\texttt{qa/ps1-measurements.json}。

两个星系的图像为\texttt{galaxy1.fit}和\texttt{galaxy2.fit}，题面给的光谱名后缀为\texttt{.fit}。

实际ZIP中光谱名为\texttt{galaxy1spec.fits}和\texttt{galaxy2spec.fits}。

图像大致中心分别为$(1831.5,636.5)$与$(1346.0,945.0)$。光谱为$4\times N$数组，第一行为每单位波长通量，波长为$\lambda_i=10^{w+0.0001i}$埃，$i=0,\ldots,N-1$，两者$(w,N)=(3.5823,3818),(3.5825,3815)$。实际文件第二行为扣连续谱的谱，第三行为误差，第四行为位掩码；不能把第三行误当成逆方差。

\src{8.902 PS1.1，原题(A)--(E)，官方PDF第1--2页。}
\begin{exercise}
(A) 以一像素宽的径向环画两星系的表面亮度曲线，纵轴为像素计数DN的对数，横轴为距中心角秒。估计并扣除非零天空背景，必要时调整中心；像素尺度为$0.394$角秒。
(B) 改以2角秒宽径向环画每环的中位表面亮度，以提高信噪比。
(C) 星等20的天体在星系1和2图像中分别贡献2328、2496个总DN，使用此校准把(B)纵轴变为每平方角秒的星等并重画。
(D) 分别拟合$I(r)=I_0\exp(-r/r_s)$和$I(r)=I_0\exp[-7.67(r/r_e)^{1/4}]$，估计指数尺度长$r_s$与de Vaucouleurs有效半径$r_e$，将模型叠在数据上；判断每个星系更接近哪一族。
(E) 单组分是否充分描述全部半径？若不充分，指出失效区域并说明物理原因。
\end{exercise}
\begin{solution}
图像FITS先使用\texttt{BSCALE,BZERO}还原DN。软件偏置约1000 DN已包含在原数组中，不能只减头部约90--96 DN的天空项。以核心椭圆Gaussian拟合定位，所得\emph{FITS一基}中心为$(1834.39,647.62)$和$(1346.60,945.67)$；星系1中心比题表明显偏移，若不调整会把尖核摊宽。在150--200像素背景环取中位数，得总底1096、1091 DN。下面不对原照片作任何再分发。

一像素环的平均DN扣底后取对数；2角秒环则直接对该环扣底像素取中位数$I$。像素面积$p^2=(0.394)^2$平方角秒，若校准总数为$C_{20}$，(C)给
\[
 \mu_g=20-2.5\log_{10}\left(\frac{I}{p^2C_{20}}\right).
\]
环面积影响总光度，却不再次乘进已经是每像素均值或中位数的$I$；漏$p^2$会错约2星等。非正扣底值不取对数，应作为低信噪限制处理。
\begin{center}\begin{tikzpicture}
\begin{axis}[width=.72\textwidth,height=4.7cm,xlabel={$r$（角秒）},ylabel={$\log_{10}(\mathrm{DN})$},xmin=0,xmax=45,legend pos=north east]
\addplot[structurecolor,thin] table[x=r_arcsec,y=log_dn,col sep=comma]{qa/ps1-galaxy1-onepixel.csv};\addlegendentry{星系1}
\addplot[black,dashed] table[x=r_arcsec,y=log_dn,col sep=comma]{qa/ps1-galaxy2-onepixel.csv};\addlegendentry{星系2}
\end{axis}\end{tikzpicture}\end{center}

在3--31角秒的两角秒环上，对$\ln I$做不加权最小二乘；指数模型对$r$线性，de Vaucouleurs模型对$r^{1/4}$线性。这个目标使多个亮度量级均有权重，不冒充依据完整像素协方差的最优统计拟合。结果如下，RMS以$\log_{10}I$表示。
\begin{center}\begin{tabular}{lrrrr}\toprule
对象与模型&尺度（角秒）&$I_0$（DN）&RMS（dex）&总模型星等\\\midrule
星系1，指数&7.406&199.2&0.1021&14.303\\
星系1，$r^{1/4}$&14.860&43579&0.0175&13.876\\
星系2，指数&9.129&228.4&0.0500&13.776\\
星系2，$r^{1/4}$&40.771&13888&0.1017&13.001\\\bottomrule
\end{tabular}\end{center}
\begin{center}
\begin{tikzpicture}\begin{axis}[width=.47\textwidth,height=4.7cm,xlabel={$r$（角秒）},ylabel={DN／像素},ymode=log,xmin=0,xmax=40,legend pos=north east]
\addplot[only marks,mark=*,mark size=1pt,structurecolor] table[x=r_arcsec,y=dn_median,col sep=comma]{qa/ps1-galaxy1-profile.csv};\addlegendentry{星系1}
\addplot[black,dashed] table[x=r_arcsec,y=fit_exp_dn,col sep=comma]{qa/ps1-galaxy1-profile.csv};\addlegendentry{指数}
\addplot[structurecolor] table[x=r_arcsec,y=fit_dev_dn,col sep=comma]{qa/ps1-galaxy1-profile.csv};\addlegendentry{$r^{1/4}$}
\end{axis}\end{tikzpicture}\hfill
\begin{tikzpicture}\begin{axis}[width=.47\textwidth,height=4.7cm,xlabel={$r$（角秒）},ymode=log,xmin=0,xmax=40,legend pos=north east]
\addplot[only marks,mark=*,mark size=1pt,structurecolor] table[x=r_arcsec,y=dn_median,col sep=comma]{qa/ps1-galaxy2-profile.csv};\addlegendentry{星系2}
\addplot[black,dashed] table[x=r_arcsec,y=fit_exp_dn,col sep=comma]{qa/ps1-galaxy2-profile.csv};\addlegendentry{指数}
\addplot[structurecolor] table[x=r_arcsec,y=fit_dev_dn,col sep=comma]{qa/ps1-galaxy2-profile.csv};\addlegendentry{$r^{1/4}$}
\end{axis}\end{tikzpicture}
\end{center}
星系1明显更接近$r^{1/4}$族，星系2更接近指数盘。$I_0$是模型外推的中心强度，不等于已受PSF卷积的实际中央像素；尤其de Vaucouleurs的巨大$I_0$不能解释成中央像素饱和测量。

\begin{center}\begin{tikzpicture}\begin{axis}[width=.72\textwidth,height=4.7cm,xlabel={$r$（角秒）},ylabel={$\mu_g$（星等／角秒$^2$）},y dir=reverse,xmin=0,xmax=45,legend pos=north west]
\addplot[only marks,mark=*,structurecolor] table[x=r_arcsec,y=mu_mag_arcsec2,col sep=comma]{qa/ps1-galaxy1-profile.csv};\addlegendentry{星系1}
\addplot[only marks,mark=square*,black] table[x=r_arcsec,y=mu_mag_arcsec2,col sep=comma]{qa/ps1-galaxy2-profile.csv};\addlegendentry{星系2}
\end{axis}\end{tikzpicture}\end{center}
单组分不是全径向精确模型：内约1--3角秒需PSF及核／核球结构，星系2在若干中间半径偏离平滑指数，与盘中非轴对称恒星结构及核球叠加相容；最外环又对天空底估计和邻近源敏感。环中位数降低邻近亮点影响，却也不等于真正椭圆等光度线拟合。以上判断来自实际数据残差及输入的私下检查，不能由RMS单独宣布某个物理组分已唯一识别。
\end{solution}

\src{8.902 PS1.2，原题(A)--(D)，官方PDF第2--3页。}
\begin{exercise}
(A) 把恒星光谱近似成黑体，用两星系的连续谱峰和Wien律估计“典型恒星”温度及粗略光谱型，不要求复杂恒星模型。
(B) 星系2在约6679埃有H$\alpha$发射峰，静止波长取6563埃；星系1在约8741埃有Ca II吸收，静止波长取8545埃。测线心并求退行速度。
(C) 取$H_0=71\ \mathrm{km\,s^{-1}\,Mpc^{-1}}$求距离，忽略宇宙学修正，用1(D)尺度换算发光物质的物理特征半径（kpc）。
(D) 积分1(D)模型并结合距离求总G带光度，以太阳为单位；题设太阳总光度90\%位于G带。
\end{exercise}
\begin{solution}
(A) 用实际第二行扣除谱恢复连续谱，并以200埃区间中位数找宽峰，星系1约7700埃，星系2约4700埃。经下列小红移修正到静止系后，Wien律$\lambda_{ \rm peak}T=2.89777\times10^7\ \mathrm{\AA K}$给约3850 K与6280 K，粗略是晚K／早M和F型温度。星系1的红色连续谱峰很宽，选不同合理平滑区间会得到数百K以上变化；复合恒星族群也不是真正单温黑体，不能把这个近似解释成测得了所有恒星的统一温度。

(B) 在Ca II的8705--8775埃和H$\alpha$的6668--6691埃内，用线性连续谱加吸收／发射Gaussian拟合，按实际第三行误差加权。历史SDSS掩码$2^{30}$表示发现发射线，\emph{不是坏像素}，否则会把所求H$\alpha$线删掉。所得线心为8741.15与6680.88埃；后者与题面约6679的粗定位有区别，因此直接用测量值而不伪装为同一个数。按题给静止波长，
\[
 z_1=0.0229554,\quad z_2=0.0179618,\qquad
 cz_1=6882,\quad cz_2=5385\ \mathrm{km\,s^{-1}}.
\]
光谱标为真空波长，题给的近似原子波长不足以作精密速度校准；这里按原题数值作低红移估计。H$\alpha$简单Gaussian的约化$\chi^2\simeq8.45$，模型存在可见系统残差；独立正通量一阶矩线心与拟合差约0.05埃，不能以形式上0.023埃的拟合误差宣称真实总误差如此小。

(C) $D=cz/H_0$给$D_1=96.93$、$D_2=75.84$ Mpc。一角秒对应$D/206265$，所以较优模型的星系1$r_e=6.98$ kpc，星系2$r_s=3.36$ kpc。指数盘的半光半径是$1.678r_s$，约5.63 kpc；不能将$r_s$直接称为全部星系的边界。两个模型均无有限外缘，“extent”在本题应明确是模型的特征半径。

(D) 若$I_0$单位为每像素DN、角半径用角秒，模型总数为
\[
 C_{ \rm exp}=\frac{2\pi I_0r_s^2}{p^2},\qquad
 C_{ \rm dev}=\frac{8\pi I_0r_e^2\Gamma(8)}{p^2\,7.67^8}.
\]
第二式来自变量$u=7.67(r/r_e)^{1/4}$，径向积分变成$u^7\ee^{-u}$；不能把1(D)的$I_0$当成$r_e$处强度。较优模型给$6.56\times10^5$、$7.70\times10^5$ DN，校准为$g=13.876,13.776$。

\textbf{总G带能量需要说明绝对波段约定}：DN与星等20只确定相对标定，不单独给能量积分的带宽。为做一个可复算数值例，采用历史SDSS的$g\simeq$AB零点3631 Jy、官方2001响应$R_g$归一至峰值1，并假定中心光谱形状代表整体。Photon加权AB均值为
\[
 \bar F_\nu=\frac{\int F_\lambda R_g\lambda\dd\lambda}
 {c\int R_g\dd\lambda/\lambda},\qquad
 \bar F_\nu=3631\ \mathrm{Jy}\,10^{-0.4g}.
\]
以实际谱形将它转换为\emph{明确约定的响应加权带通量}$F_G=\int F_\lambda R_g\dd\lambda$，再用$L_G=4\pi D^2F_G$、题设$L_{G,\odot}=0.9\Ls$，得约$5.16\times10^9$与$3.51\times10^9$个题设太阳G光度。响应文件、端点外推和忽略色梯度／消光／$K$修正的假设均在回执中，不能把它称为不依赖波段定义的唯一精密光度。若采用别的明确太阳G绝对星等或带通能量定义，应以$10^{-0.4(M_g-M_{g,\odot})}$或相应零通量重算。
\end{solution}

\src{8.902 PS1.3，原题(A)--(C)，官方PDF第3页。}
\begin{exercise}
(A) 对星系1约8740埃Ca II吸收拟合Gaussian，求标准差$\sigma_\lambda$，并用$\sigma_\lambda/\lambda=\sigma_v/c$给速度宽度。
(B) 同一仪器观测近邻相似恒星，含仪器与恒星内部原子运动的对照Gaussian宽度为95 km/s；按方差相加扣除它，估星系1恒星的轨道速度弥散。
(C) 假定该弥散代表1(D)特征半径处的恒星轨道速度，求内部质量，以太阳质量表示，并与2(D)所得的恒星数量估计比较。
\end{exercise}
\begin{solution}
(A) 实际加权拟合给$\sigma_\lambda=6.17$埃，形式统计误差约0.41埃。用\emph{同一观测系}线心8741.15埃换算得$\sigma_{ \rm obs}=212\ \mathrm{km\,s^{-1}}$。这是标准差，不是FWHM；后者为$2\sqrt{2\ln2}\sigma$。
\begin{center}\begin{tikzpicture}\begin{axis}[width=.70\textwidth,height=4.7cm,xlabel={观测波长（埃）},ylabel={连续谱归一通量},xmin=8705,xmax=8775]
\addplot[only marks,mark=*,mark size=1pt,black] table[x=wavelength_A,y=normalized_measurement,col sep=comma]{qa/ps1-galaxy1-line-measurements.csv};
\addplot[structurecolor,thick] table[x=wavelength_A,y=normalized_fit,col sep=comma]{qa/ps1-galaxy1-line-fit.csv};
\end{axis}\end{tikzpicture}\end{center}
(B) Gaussian卷积使方差相加，故$\sigma_*=[\sigma_{ \rm obs}^2-(95\ \mathrm{km\,s^{-1}})^2]^{1/2}\simeq189\ \mathrm{km\,s^{-1}}$。若观测宽度小于对照宽度，应报未分辨和上限，而不能取虚数弥散。本例没有此问题。

(C) 量纲估计为$M(<R)=C\sigma_{ \rm los}^2R/G$；$C$还取决于密度剖面、各向异性以及视向到三维速度的转换，题面未给它。按把弥散当作典型圆速的$C=1$约定、$R=r_e=6.98$ kpc，得$M\simeq5.81\times10^{10}\Ms$；若用维里估计的不同结构系数需相应乘$C$。

2(D)实际给\emph{光度}，不是唯一的恒星数。若粗略把每颗星都当成题设太阳，发光数量约$5.16\times10^9$颗，发光质量同量级；$C=1$的动力学质量约为其11倍。真实族群含不同单星光度、恒星遗迹和色梯度，且动力学系数未知，因此这一比较只说明简单纯太阳族群与动力学质量有差异，不能由它精确量出暗物质比例或证明真实恒星数已测定。
\end{solution}

\addtocontents{toc}{\protect\Needspace{3\baselineskip}}
\originalchapter{星系动力学与宇宙结构习题}
本章均为 AI 独立解答. 8.902 官方没有公开作业答案. 原题缺少唯一数值所需的参数时, 给出参数表达式, 再以明确假设算例; 算例不冒充源文件的数据测量.

\src{8.902 Fall 2023, PS2.1, 三倍权重.}
\begin{exercise}
球对称势 $\Phi=v_c^2\ln(r/r_{\rm ref})$ 中一恒星在远心点速度为 $v_c/2$. 数值积分并画轨道平面位置, 取绘图间隔约为径向周期的 $1/50$, 展示 50--100 个轨道周期, 同时给连续曲线与反映停留概率的点图, 简述结果. 数值步长可比绘图间隔小; 给出方程与代码.
\end{exercise}
\begin{solution}
取远心距 $r_a$ 为长度单位、$r_a/v_c$ 为时间单位, 平面方程为 $\ddot{\boldsymbol r}=-\boldsymbol r/r^2$, 初值 $(x,y,\dot x,\dot y)=(1,0,0,1/2)$. 能量 $E=1/8$, 角动量 $\ell=1/2$. 非平凡近日点满足 $\ln r+1/(8r^2)=1/8$, 得 $r_p=0.310885$. 径向周期积分为
\[
T_r=2\int_{r_p}^1\frac{\dd r}{\sqrt{2[1/8-\ln r-1/(8r^2)]}}=2.984886.
\]
代码使用 DOP853, 内步长最多 $T_r/150$, 相对容差 $10^{-11}$, 积分 75 个径向周期, 按 $T_r/50$ 采样 3750 点. 能量最大误差 $6.9\times10^{-13}$, 角动量误差 $4.9\times10^{-14}$; 已实际运行, 但有限误差检查不证明任意时间的数值稳定性.
\begin{center}\begin{tikzpicture}\begin{axis}[width=.48\linewidth,height=5.5cm,axis equal image,xlabel={$x/r_a$},ylabel={$y/r_a$},xmin=-1.1,xmax=1.1,ymin=-1.1,ymax=1.1]
\addplot[structurecolor,thin] table[x=x,y=y,col sep=comma]{qa/logarithmic-orbit.csv};
\end{axis}\end{tikzpicture}\hfill
\begin{tikzpicture}\begin{axis}[width=.48\linewidth,height=5.5cm,axis equal image,xlabel={$x/r_a$},ylabel={$y/r_a$},xmin=-1.1,xmax=1.1,ymin=-1.1,ymax=1.1]
\addplot[only marks,mark=*,mark size=.22pt,structurecolor] table[x=x,y=y,col sep=comma]{qa/logarithmic-orbit.csv};
\end{axis}\end{tikzpicture}\end{center}
轨道进动而形成玫瑰线; 等时间点在远心处较密, 因为径向和角运动都慢. 「75 周期」指径向周期, 不把不闭合轨道称为 75 个相同椭圆. 完整代码和采样点位于源码包 \texttt{scripts/} 和 \texttt{qa/}.
\end{solution}

\src{8.902 PS2.2, (A)--(C).}
\begin{exercise}
Mestel 盘势 $\Phi=v_c^2[\ln(r/r_{\rm ref})+\ln((1+|\cos\theta|)/2)]$.
(A) 证明盘面旋转曲线平坦.
(B) 证明质量密度在 $z\ne0$ 时为零, 盘面上呈无限薄集中.
(C) 对 $z$ 积分证明有限面密度 $\Sigma(R)=v_c^2/(2\pi GR)$.
\end{exercise}
\begin{solution}
用柱半径 $R$, 球半径 $r=\sqrt{R^2+z^2}$, 势写成 $v_c^2\ln[(r+|z|)/(2r_{\rm ref})]$. 在 $z=0$ 有 $R\partial_R\Phi=v_c^2$, 故圆速恒定. 对 $z>0$, $\partial_z\Phi=v_c^2/r$, $\partial_R\Phi=v_c^2R/[r(r+z)]$, 代柱坐标 Laplacian 可验证 $\nabla^2\Phi=0$; 下半空间由对称性相同. 跨盘面 $\partial_z\Phi$ 从 $-v_c^2/R$ 跳至 $+v_c^2/R$, 所以分布意义下 $\partial_z^2\Phi$ 有 $2v_c^2\delta(z)/R$ 项. 由 Poisson 方程, $\rho=\Sigma(R)\delta(z)$, $\Sigma=v_c^2/(2\pi GR)$. 「无穷密度」应理解为 delta 分布, 不是一个可当普通函数代入的无穷数. 未截断模型总质量发散, 实际盘须有限范围.
\end{solution}

\src{8.902 PS3.1, (A)--(E).}
\begin{exercise}
无限薄差分旋转盘圆速 $v_c(R)$, 附近恒星与太阳距离 $d\ll R_0$, 银经 $l$.
(A) 证明相对视向与横向速度 $v_r=Ad\sin2l$, $v_T=d(A\cos2l+B)$, $A=(v_c/R-v_c')/2$, $B=-(v_c/R+v_c')/2$.
取 $A=14.82\pm0.84,B=-12.37\pm0.64$ 单位 $\mathrm{km\,s^{-1}kpc^{-1}}$, $R_0=8.5\,\mathrm{kpc}$, 只用中心值:
(B) 求圆速与太阳公转周期; (C) 求局部剪切 $v_c'$; (D) 求本轮频率; (E) 求每次公转中的本轮次数.
\end{exercise}
\begin{solution}
用精确的局部观测几何式 $v_r=[R_0\Omega(R)-v_0]\sin l$, $v_T=R_0[\Omega(R)-\Omega_0]\cos l-\Omega(R)d$. 由 $R-R_0=-d\cos l+O(d^2/R_0)$ 及 $R_0\Omega'_0=v_c'-v_0/R_0=-2A$, 一阶给 $v_r=2Ad\cos l\sin l=Ad\sin2l$. 横向分量为 $2Ad\cos^2l-\Omega_0d=d[A\cos2l+(A-\Omega_0)]$. 使用 $A-\Omega_0=B$ 得题面式. 横向正向与银经增加方向一致; 改变正向会使整式变号.

$\Omega_0=A-B=27.19$, $v_0=R_0(A-B)=231.115\,\mathrm{km\,s^{-1}}$, $2\pi/\Omega_0=225.95\,\mathrm{Myr}$. 剪切 $v_c'=-(A+B)=-2.45\,\mathrm{km\,s^{-1}kpc^{-1}}$. 由式\eqref{eq:epicycle}, $\kappa=\sqrt{-4B(A-B)}=36.679\,\mathrm{km\,s^{-1}kpc^{-1}}$, 每个引导中心周期内 $\kappa/\Omega_0=1.349$ 次本轮振荡.
\end{solution}

\src{8.902 PS4.1, (A)--(C).}
\begin{exercise}
默认 $\Omega_m=0.3,\Omega_\Lambda=0.7,H_0=71\,\mathrm{km\,s^{-1}Mpc^{-1}}$. 两个 Hubble 流吸收星系的 Si$^{++}$ 线静止波长 $1206$ 埃, 观测波长 $4711,4836$ 埃.
(A) 求两星系共动视向距离; (B) 较近星系在其吸收光时刻看较远星系的红移是多少?
(C) 对小红移间隔, 用早期 $H(z)$ 建立 (A)、(B) 的近似关系并估计距离.
\end{exercise}
\begin{solution}
$z_1=2.906302,z_2=3.009950$, 共动间距 $\Delta\chi=(c/H_0)\int_{z_1}^{z_2}\dd z/E(z)=99.621\,\mathrm{Mpc}$. 同一段光线在近星系时波长相对远处发射波长增加 $a_1/a_2$, 故 $1+z_{21}=(1+z_2)/(1+z_1)$, $z_{21}=0.026534$. 小间隔下 $\Delta\chi\simeq c\Delta z/H(z_1)=c(1+z_1)z_{21}/H(z_1)$, $H(z)=H_0\sqrt{0.3(1+z)^3+0.7}$, 给约 $101.5\,\mathrm{Mpc}$. 若用中点 $H$ 则更接近精确积分. 对应近星系吸收时的实际间距须再除 $1+z_1$.
\end{solution}

\src{8.902 PS4.2, (A)--(B).}
\begin{exercise}
吸收体随机分布, 共动数密度 $n$ 恒定、吸收截面 $\sigma$ 恒定.
(A) 求平均视线每红移吸收体数, 判断随红移增减.
(B) 在 $z\sim1.5$--$2.5$, 观测 $\dd N/\dd z\sim2$, 典型亮星系 $n\sim0.001\,\mathrm{Mpc^{-3}}$, 求 $z\sim2$ 所需截面.
\end{exercise}
\begin{solution}
将恒定截面理解为\textbf{实际}面积, 式\eqref{eq:absorber}给 $\dd N/\dd z=nc\sigma(1+z)^2/H(z)$. 对题面的物质和真空模型, $\dd\ln[(1+z)^2/H]/\dd\ln(1+z)=2-3\Omega_m(z)/2>0$, 故增加. $z=2$ 给 $\sigma=0.1561\,\mathrm{Mpc^2}$, 圆形吸收半径 $\sqrt{\sigma/\pi}=0.223\,\mathrm{Mpc}$. 若截面指固定共动面积, 上式红移因子不同, 答案也改变; 本解没有把两种约定混用.
\end{solution}

\src{8.902 PS4.3, (A)--(D).}
\begin{exercise}
(A) 求复合面上包含一个星系质量的球区域张角.
(B) 求直径 $200$ 共动 Mpc 区域在复合面的张角.
(C) 求复合时粒子视界内总质量.
(D) 求宇宙刚进入物质主导时视界内总质量. 视界是自宇宙起始至该时刻的实际光传播距离.
\end{exercise}
\begin{solution}
原题未指定星系质量、复合红移与辐射密度, 因而不能唯一给四个数值. 令 $M_g,z_*,\Omega_{r0}$ 为参数. 共动球半径 $R_c=(3M_g/(4\pi\rho_{m0}))^{1/3}$, (A) 的直径张角为 $2R_c/\chi(z_*)$, (B) 为 $200\,\mathrm{Mpc}/\chi(z_*)$. 粒子视界共动半径
$\chi_p(a)=cH_0^{-1}\int_0^a\dd x/\sqrt{\Omega_{r0}+\Omega_{m0}x+\Omega_{\Lambda0}x^4}$,
实际半径为 $a\chi_p$. 物质质量是 $M_m=4\pi\rho_{m0}\chi_p^3/3$; 若「总质量」含辐射能的等效质量, 应为 $M_{\rm tot}=M_m[1+\Omega_{r0}/(\Omega_{m0}a)+\Omega_{\Lambda0}a^3/\Omega_{m0}]$. 真空项在本题早期极小.

示例明确取 $M_g=10^{12}\Ms,z_*=1100,\Omega_{r0}=8.5\times10^{-5}$. (A) 约 $0.914$ 角分, (B) 约 $0.8535^\circ$. (C) 物质质量 $3.565\times10^{18}\Ms$, 含辐射的总等效质量约 $4.68\times10^{18}\Ms$. (D) 等密度时 $a_{\rm eq}=\Omega_{r0}/\Omega_{m0}$, 物质质量 $2.184\times10^{17}\Ms$, 总等效质量约 $4.368\times10^{17}\Ms$. 仅含物质与真空且设 $\Omega_r=0$ 的模型没有辐射到物质的转折, 不能用于 (D).
\end{solution}

\src{8.902 PS5.1, (A)--(E).}
\begin{exercise}
银河与 M31 当前间距 $700\,\mathrm{kpc}$, 相对径向速度 $-121\,\mathrm{km\,s^{-1}}$, 宇宙年龄取 $13.7\times10^9\,\mathrm{yr}$. 理想径向两体模型 $r=A(1-\cos\eta),t=B(\eta-\sin\eta)$.
(A) 推导 $v=(r/t)\sin\eta(\eta-\sin\eta)/(1-\cos\eta)^2$.
(B) 用 Newton 方程证明 $GM_T=A^3/B^2=r^3(\eta-\sin\eta)^2/[t^2(1-\cos\eta)^3]$.
(C) 求当前 $\eta$; (D) 求总质量, 总光度 $3\times10^{10}\Ls$ 时的质量光度比.
(E) 若两星系等质量、银河圆速合理取值, 暗晕需延伸多远，两个暗晕分离还是重叠? 原题印有 $M(<R)=v_c^2R/(2G)$。
(F) 求历史最大间距、转向时宇宙年龄、距两星系碰撞还有多久；画从大爆炸到碰撞的$r(t)$并标示今日位置。
\end{exercise}
\begin{solution}
参数微分给 $v=(A/B)\sin\eta/(1-\cos\eta)$, 用 $A=r/(1-\cos\eta)$、$B=t/(\eta-\sin\eta)$ 得 (A). 再求加速度 $\ddot r=-A/[B^2(1-\cos\eta)^2]=-A^3/(B^2r^2)$, 与两体方程比较得 (B). 收缩支取 $\pi<\eta<2\pi$, 数值根 $\eta=4.297242$. 总质量 $3.994\times10^{12}\Ms$, 质量光度比 $133.1\,\Ms/\Ls$.

球对称圆轨道力平衡应为 $M(<R)=v_c^2R/G$, 原题分母的 2 与圆速定义不相容, 此处明示修正。取 $v_c=220\,\mathrm{km\,s^{-1}}$、银河为总质量的一半，$R=177.4\,\mathrm{kpc}$，两晕半径和约355 kpc小于间距700 kpc，按此尖锐截断估计不重叠。按印刷式会得单晕354.9 kpc、两晕几乎重叠，但它不满足同一圆速模型的力平衡。此timing模型忽略横向速度、质量增长、潮汐与真空力，Birkhoff定理也不能把任意有限星系对精确变成均匀闭宇宙，因此质量与晕边界均为理想化估计。

(F) 由当今$r,t,\eta$反求$A=498.82$ kpc、$B=2.6284$ Gyr。转向在$\eta=\pi$，故历史最大间距$2A=997.63$ kpc，年龄$\pi B=8.257$ Gyr；形式碰撞在$\eta=2\pi$，年龄$2\pi B=16.515$ Gyr，距今约2.815 Gyr。直接沿参数解采样得到下图，不将理想径向两体的零间距时刻当成已预测了真实复杂并合。
\begin{center}\begin{tikzpicture}\begin{axis}[width=.73\textwidth,height=4.6cm,xlabel={$t$（Gyr）},ylabel={$r$（kpc）},xmin=0,xmax=17,ymin=0,ymax=1100]
\addplot[structurecolor,thick] table[x=t_Gyr,y=r_kpc,col sep=comma]{qa/timing-complete.csv};
\addplot[only marks,mark=*,black] coordinates {(13.7,700)};
\node[anchor=south west] at (axis cs:13.7,700) {今日};
\end{axis}\end{tikzpicture}\end{center}
\end{solution}

\src{8.902 PS6.1, (A)--(D).}
\begin{exercise}
EdS 球顶帽在 $z=8$ 达到最大半径.
(A) 求动力学自由落体时间, 可引用转向密度比.
(B) 假设体系的全部引力势能在坍缩时用于加热重子，球内密度均匀，重子为理想气体，无任何冷却，求维里化后温度随$M$的关系，按$10^{12}\Ms$归一化。
(C) 重子质量占 $5\%$, 无冷却以外的机制, $\Lambda(T)=10^{-23}\,\mathrm{erg\,cm^3s^{-1}}$, 求冷却时间.
(D) 比较冷却和动力学时间, 估计能形成星系的最大特征质量. 原题未指定气体平均粒子质量或精确能量分配.
\end{exercise}
\begin{solution}
先列旧稿采用的$\Omega_{m0}=0.3,H_0=71$早期物质主导近似：这属于早期$\Lambda$CDM密度归一化，\textbf{不是严格EdS}。在此条件下， $\rho_{\rm ta}=(9\pi^2/16)\rho_{m0}(1+8)^3$. 由均匀球从静止坍缩时间 $t_{\rm ff}=\sqrt{3\pi/(32G\rho_{\rm ta})}$ 得 $621\,\mathrm{Myr}$. 维里半径 $R_v=R_{\rm ta}/2$, 晕密度 $8\rho_{\rm ta}$, $M=10^{12}\Ms$ 时 $R_v=56.0\,\mathrm{kpc}$.

题面的“全部势能加热重子”与普通示踪气体维里温度需要分开，因为重子只占$f_b=0.05$。令分给气体的热能$E_g=\alpha_E GM^2/R_v$，气体粒子数为$f_bM/(\mu m_p)$，于是
\[
 T=\frac{2\alpha_E\mu m_pGM}{3f_b\kb R_v}.
\]
普通维里估计把总动能$|W_v|/2=3GM^2/(10R_v)$按质量比例分给重子，$\alpha_E=3f_b/10$，给$T_{ \rm vir}=\mu m_pGM/(5\kb R_v)$。若将\emph{全部体系释放的}势能$|W_v|-|W_{ \rm ta}|=|W_v|/2$都给气体，则$\alpha_E=3/10$，温度为普通值的$1/f_b=20$倍；若字面上给整个$|W_v|$，为40倍，且还违反同一个封闭维里体系的总能量预算。不能仅给温度加倍而漏掉接收热量的质量是$f_bM$这一点。

取完全电离纯氢，$\mu=1/2$，普通维里解释得$T_{ \rm vir}=9.30\times10^5(M/10^{12}\Ms)^{2/3}$ K。冷却函数约定为体积损失$\mathcal L=n_en_p\Lambda$，$n_e=n_p=f_b\rho_v/m_p$，气体热能为$3n_p\kb T$，故
\[
 t_c=\frac{3\kb T}{n_p\Lambda}
 =444.1\ \mathrm{Myr}\,q\left(\frac{M}{10^{12}\Ms}\right)^{2/3},
 \quad q=T/T_{ \rm vir}.
\]
比较$t_c<t_{ \rm ff}=620.8$ Myr，给
\[
 M_{ \rm max}\simeq1.65\times10^{12}q^{-3/2}\Ms.
\]
严格按题面的EdS则应取$\Omega_{m0}=1$（这里明确仍取$H_0=71$）；重新计算得$t_{\rm ff}=340.04$ Myr、$R_v=37.493$ kpc、$T_{\rm vir}=1.3898\times10^6(M/10^{12}\Ms)^{2/3}$ K、$t_c=199.04q(M/10^{12}\Ms)^{2/3}$ Myr。因此严格EdS分支为$M_{\max}=2.233\times10^{12}q^{-3/2}\Ms$，$q=20$时$2.50\times10^{10}\Ms$；以下引用旧数值时仅指前面的非严格归一化，不能与此分支混用。两种归一化的独立实际计算见\texttt{qa/eds-cooling-normalization.json}。

普通气体维里温度$q=1$给约$10^{12}\Ms$的经典数量级；按全部释放能量分给气体的$q=20$给$1.84\times10^{10}\Ms$；把整个势能绝对值分给它的$q=40$形式上给$6.5\times10^9\Ms$，但不满足前述封闭模型预算。因此原题措辞不能同时支持唯一的标准维里温度与全部能量给5\%重子的假设。这里把(B)--(D)各解释的能量账和结果完整列出，而不把含混条件悄悄改成一项唯一答案。真实冷却曲线也不在所有这些温度上恒定；这里按题给常数计算。
\end{solution}

\src{8.902 PS6.2, (A)--(B).}
\begin{exercise}
(A) 积分 Press--Schechter 的原始体积分数, 证明只计入一半质量.
(B) 说明 excursion-set / extended Press--Schechter 的基本思想, 它怎样消除因子 2 问题. 原题引用 Bond 等人 1991 年工作作为阅读方向, 不要求复录论文.
\end{exercise}
\begin{solution}
以零均值 Gaussian 密度对比、方差 $S$ 为起点, $F=\int_{\delta_c}^\infty\exp[-\delta^2/(2S)]\dd\delta/\sqrt{2\pi S}=\frac12\operatorname{erfc}(\delta_c/\sqrt{2S})$. 从大质量 $S\to0$ 至小质量 $S\to\infty$, 总分数只能增加到 $1/2$. 漏计来源是某尺度上的低密度点仍可能属于更大尺度已经坍缩的区域, 即 cloud-in-cloud 问题.

锐波数滤波在独立 Gaussian 模假设下使增量独立, 分布满足 $\partial_SP=\frac12\partial_\delta^2P$. 在 $\delta_c$ 设吸收壁, 镜像法给 $P=[\ee^{-\delta^2/(2S)}-\ee^{-(2\delta_c-\delta)^2/(2S)}]/\sqrt{2\pi S}$, $\delta<\delta_c$. 首过壁概率流为 $f(S)=\delta_c\ee^{-\delta_c^2/(2S)}/(\sqrt{2\pi}S^{3/2})$. 换元 $u=\delta_c/\sqrt S$ 得 $\int_0^\infty f(S)\dd S=1$. 通过「第一次过壁」给每个质量元指定最大归属尺度, 不再重复或漏掉嵌套区域. 对非 Markov 滤波或非 Gaussian 场, 这一简单镜像结果不普遍成立.
\end{solution}

\src{8.902 期中题组 Q2, (A)--(B), 按作业计分.}
\begin{exercise}
封闭箱化学模型 $Z=-y\ln\mu$, 产额 $y$ 恒定, 当前气体金属丰度 $Z_0$.
(A) 推导归一化恒星累积分布 $M_s(<Z)=(1-\ee^{-Z/y})/(1-\ee^{-Z_0/y})$, 并用物理解释支持公式.
(B) 求 $\dd M_s/\dd\log_{10}Z$, 作示意并求峰位置.
\end{exercise}
\begin{solution}
充分混合使在某时刻形成的恒星继承当时气体丰度. 随气体分数由 1 降至 $\mu$, 锁入恒星的质量为 $M_T(1-\mu)$, 而 $\mu=\ee^{-Z/y}$. 以今天全部恒星质量 $M_T(1-\ee^{-Z_0/y})$ 归一化即得 (A). 它适用 $0\le Z\le Z_0$, 超出上界累积为 1.

微分得 $\dd M_s/\dd\log_{10}Z=(\ln10)Z\ee^{-Z/y}/[y(1-\ee^{-Z_0/y})]$. 无截断曲线在 $Z=y$ 达峰; 若 $Z_0<y$, 可观测区间峰在右端 $Z_0$. 图取示意 $Z_0=3y$, 横轴 $x=\log_{10}(Z/y)$:
\begin{center}\begin{tikzpicture}\begin{axis}[width=.65\linewidth,height=4.5cm,xlabel={$\log_{10}(Z/y)$},ylabel={每对数丰度区间的恒星分数},xmin=-2,xmax=.48]
\addplot[structurecolor,domain=-2:.477,samples=100]{ln(10)*10^x*exp(-10^x)/(1-exp(-3))};
\end{axis}\end{tikzpicture}\end{center}
\end{solution}

\src{8.902 期中题组 Q3.}
\begin{exercise}
只考虑二体弛豫, 总质量相同的星系若由各为 $10^6\Ms$ 的球状星团构成, 比由单个恒星构成的星系弛豫更快还是更慢? 给定量关系与直觉解释.
\end{exercise}
\begin{solution}
固定宏观密度与速度, $t_{\rm rel}\propto1/(m\ln\Lambda)$. 若单星质量取 $1\Ms$, 则 $t_{\rm cluster}/t_{\rm star}\simeq10^{-6}\ln\Lambda_{\rm star}/\ln\Lambda_{\rm cluster}$. $b_{\min}$ 和星团有限尺度会改变 Coulomb 对数, 所以不应把比例精确写死为 $10^{-6}$. 较重散射体造成的单次冲量大, 即使数目减少, $m^2n=m\rho$ 仍增大, 故弛豫较快. 这只用于星团可近似为离散散射中心的模型.
\end{solution}

\src{8.902 期中题组 Q4.}
\begin{exercise}
用维里定理推导银河中恒星的典型速度.
\end{exercise}
\begin{solution}
设势能为 $W=-C GM^2/R$, $K=M\langle v^2\rangle/2$, 稳态维里关系给 $\sqrt{\langle v^2\rangle}=\sqrt{CGM/R}$. 题面未给 $M,R$ 与结构系数, 只能先给此表达式. 取代表性内部质量 $M=10^{11}\Ms$, $R=10\,\mathrm{kpc}$、$C\sim1$, 得 $v\sim207\,\mathrm{km\,s^{-1}}$. 若均匀球 $C=3/5$ 则约 $161\,\mathrm{km\,s^{-1}}$. 视向一维弥散在各向同性下再除 $\sqrt3$, 不可将其与三维速度混称.
\end{solution}

\src{8.902 期中题组 Short Answer, (A)--(D), 原每项 5 分.}
\begin{exercise}
(A) 什么是 NFW 密度曲线?
(B) 什么是暴烈弛豫, 怎样发生, 为什么重要?
(C) 写无碰撞 Boltzmann 方程并说明如何导出 Jeans 方程.
(D) 银河中两恒星发生强遭遇的可能性如何?
\end{exercise}
\begin{solution}
(A) $\rho=\rho_s/[x(1+x)^2]$, 内部近 $r^{-1}$、外部近 $r^{-3}$, 为暗晕模拟常用拟合, 未截断质量对数发散.
(B) 体系快速形成或并合时 $\Phi(\boldsymbol x,t)$ 改变, 单星能量沿轨道满足 $\dd E/\dd t=\partial_t\Phi$, 轨道因而交换能量并发生相混合. 它在动力学时间内改变粗粒化分布, 不等于微观 $f$ 违反 Liouville 定理, 也不保证达到唯一 Maxwell 分布.
(C) 式\eqref{eq:cbe}沿轨道守恒. 乘 $v_i$ 积分, 用速度边界项消失和分部积分 $\int v_i\partial_{v_j}f\dd^3v=-\delta_{ij}\rho$, 得 $\partial_t(\rho\bar v_i)+\partial_j(\rho\overline{v_iv_j})=-\rho\partial_i\Phi$. 二阶矩未知, 仍需闭合或分布模型.
(D) 强偏转需要 $b\lesssim G(m_1+m_2)/v^2$. 取 $m_1=m_2=\Ms,v=200\,\mathrm{km\,s^{-1}}$, $b_{90}\sim0.044\,\mathrm{AU}$. 本地星数密度 $0.1\,\mathrm{pc^{-3}}$ 的示例遭遇率 $n\pi b_{90}^2v$ 约 $3\times10^{-18}\,\mathrm{yr^{-1}}$, 极低. 此数密度是解答的示例假设; 密集星团或银河中心须另估, 物理碰撞截面也不同于强引力偏转截面.
\end{solution}

\originalchapter{自由落体、盘稳定性与绘图补题}
本章补齐8.902中先前选讲未收录的完整公开题面。解答均为独立推导；原题给模型信息不足时明确补充条件，并区分其条件结果与一般定理。

\src{8.902 PS3.2，标为B\&T 3.4，但题面在官方作业PDF内完整给出。}
\begin{exercise}
密度为$\rho_0$的均匀无压强流体球由静止释放，证明其坍缩到一点所需时间为$t_{ \rm ff}=\sqrt{3\pi/(32G\rho_0)}$。
\end{exercise}
\begin{solution}
取初始半径$r_0$的一层壳，内质量$M=4\pi\rho_0r_0^3/3$。壳层未交叉时，壳层定理给$\ddot r=-GM/r^2$，初值$\dot r(r_0)=0$。乘$\dot r$并积分得
\[
 \dot r=-\sqrt{2GM\left(\frac1r-\frac1{r_0}\right)}.
\]
取负号是坍缩支。时间积分为
\[
 t_{ \rm ff}=\sqrt{\frac{r_0^3}{2GM}}\int_0^1\sqrt{\frac u{1-u}}\dd u.
\]
令$u=\sin^2\theta$，积分为$2\int_0^{\pi/2}\sin^2\theta\dd\theta=\pi/2$，代$M$即得所求。时间不依赖$r_0$，所以均匀球各壳同时到达中心，维持同调坍缩直到奇点。非均匀球、初始非零速度、压强、旋转或壳层交叉都会改变结果，不能把它当成所有引力体系的精确落体时间。
\end{solution}

\src{8.902期中替代题组Q1，25分，(A)--(D)。}
\begin{exercise}
有一个面密度$\Sigma$、半径$R$、角速度$\Omega$的气体盘，考察其边缘附近的试探粒子，并假定角动量守恒。
(A) 求近圆轨道径向本轮频率$\kappa$。
(B) 求$\kappa^2>0$所需的最小半径$R_{ \rm min}$。
(C) 若$R<R_{ \rm min}$会怎样？
(D) 为什么银河中更小的盘结构不经常发生此现象？
\end{exercise}
\begin{solution}
先明确一个问题：本轮频率由\emph{径向引力场的梯度}决定，仅有某个边缘的$\Sigma,R,\Omega$，没有盘的径向面密度函数、厚度、边界和外部势，一般不能唯一给出(A)、(B)。严格均匀、无限薄且硬截断的盘在边缘还存在引力梯度奇异，不能静默把它替换成一个平滑点质量势。

一般光滑轴对称势下，角动量$L$固定，径向有效势$\Phi_{ \rm eff}=\Phi+L^2/(2r^2)$。圆轨道要求$\Phi_R(R)=R\Omega^2$，其附近$r=R+\xi$满足
\[
 \ddot\xi+\kappa^2\xi=0,\qquad
 \kappa^2=\Phi_{RR}(R)+3\Omega^2=4\Omega^2+R\frac{\dd\Omega^2}{\dd R}.
\]
这才是(A)在给定势后可使用的完整式；(B)应由该势的$\kappa^2(R)>0$确定，未必存在一个唯一锐利下限。它与含压强、自引力的气体密度波Toomre判据不同。

若\emph{额外}作课堂常见的边缘点质量估计$M\simeq\pi R^2\Sigma$，并在试探粒子微小位移时固定该$M$，则$\Phi_{RR}\simeq-2GM/R^3$给
\[
 \kappa_{ \rm toy}^2\simeq3\Omega^2-\frac{2\pi G\Sigma}{R},\qquad
 R>\frac{2\pi G\Sigma}{3\Omega^2}.
\]
这是\emph{增加点质量假设后的条件表达式}，不是由原题全部信息唯一推得。还须注意：若同一个点质量势已经提供全部圆轨道平衡，则$GM/R^3=\Omega^2$，反而总有$\kappa^2=\Omega^2>0$；在该模型内不能再独立缩小$R$而保持$\Sigma,\Omega$固定，并声称原轨道仍是圆轨道。

(C) 对真正存在圆轨道的光滑模型，若其有效势二阶导数为负，径向偏离按$\exp(|\kappa|t)$增长，圆轨道不稳定，偏离方向决定向内或向外，不能一概说必然整体坍缩。(D) 银河中的气体不是孤立硬边薄盘，恒星与暗物质提供背景势，有限厚度、热压、湍流和角动量输运也影响集体扰动。应检验实际旋转曲线给出的$\kappa$和气体$Q$；小云确可引力坍缩形成恒星，但这不等于每一个小半径圆轨道都违反(A)的局部判据。没有公开官方答案，本解把信息缺口和常见近似一起说明，避免把不相容条件包装成唯一考试数值。
\end{solution}

\src{8.902 PS6.3，最后一次作业的5分绘图加分题。}
\begin{exercise}
画一个庆贺完成课程的星系。
\end{exercise}
\begin{solution}
以下为自绘TikZ示意，中央核球与旋臂仅用于这道绘图任务，没有把示意结构当成实际天体观测或密度拟合。
\begin{center}\begin{tikzpicture}[scale=.65]
\fill[structurecolor!15] (0,0) ellipse (2.5 and 1.5);
\foreach \a in {0,120,240}{
\draw[structurecolor,thick,rotate=\a] plot[domain=0:300,samples=80] ({.08*exp(0.009*\x)*cos(\x)},{.08*exp(0.009*\x)*sin(\x)});
}
\fill[structurecolor] (0,0) circle(.18);
\node[below] at (0,-2.2) {完成一程，继续探索};
\end{tikzpicture}\end{center}
它单列为创作任务，不将5分贺图虚计成一项物理证明。
\end{solution}

\originalchapter{以真实SNLS数据拟合超新星距离}
本章完整对应8.902 PS5.2，双倍权重。\href{https://ocw.mit.edu/courses/8-902-astrophysics-ii-fall-2023/mit8_902_f23_problem_set_5_additional_files.tar}{官方TAR}含Guy等人（2010）SNLS三年样本的CDS表9、10、11及ReadMe；本题实际只读取表11的SiFTO列。本卷不分发原始表或论文，提供独立分析脚本、派生数值曲线和结果。270行有完整数值且所报误差为正，7行占位／非正误差记录逐一列在回执；没有为改善拟合另删离群点。

\src{8.902 PS5.2，官方PDF第2--4页，(A)--(D)。}
\begin{exercise}
(A) 数值计算光度距离，画至$z=5$及放大到$z=1$的两张图。五种$(\Omega_m,\Omega_\Lambda)$为$(1,0)$、$(0,1)$、$(0.3,0.7)$、$(1.2,0)$、$(0.3,0)$，分别使用平直、正曲率、负曲率距离公式。取$H_0=71\ \mathrm{km\,s^{-1}\,Mpc^{-1}}$。

(B) 表10为光变观测，表11给SALT2与SiFTO的拟合参数；只用SiFTO的峰值星等$m_i^*$、形状$s_i$、颜色$c_i$及其方差与协方差。定义
\[
 \mu_i=m_i^*-M+\alpha s_i-\beta c_i,\quad
 \chi^2=\sum_i\frac{[\mu_i-\mu_{ \rm cosmo}(z_i)]^2}{\sigma_i^2+\sigma_{ \rm int}^2},
\]
忽略红移误差，且
\begin{align*}
 \sigma_i^2={}&\operatorname{Var}(m_i^*)+\alpha^2\operatorname{Var}(s_i)
 +\beta^2\operatorname{Var}(c_i)\\
 &+2\alpha\operatorname{Cov}(m_i^*,s_i)-2\beta\operatorname{Cov}(m_i^*,c_i)
 -2\alpha\beta\operatorname{Cov}(s_i,c_i).
\end{align*}
原题在$\sigma_{ \rm int}^2$旁给0.10，称作固有弥散。先假定平直，$\Omega_\Lambda=1-\Omega_m$，以$\alpha=\beta=1,M=-20$及一个初始$\Omega_m$开始，反复计算$\sigma_i$，固定它们后最小化$\chi^2$得到$(\alpha,\beta,M,\Omega_m)$，直至收敛。提交参数、各超新星$\mu_i$及最终误差的图，叠加最优曲线及保持平直而$\Omega_m$变化$\pm0.2$的曲线，并评论数据对$\Omega_m$的敏感度。

(C) 解释星系际尘埃消光怎样影响(A)及加速膨胀推断。
(D) 固定(B)所得$\alpha,\beta,M$，放松平直条件，扫描$\Omega_m,\Omega_\Lambda$。画包含最小值加2.3、6.2、11.8的$\chi^2$轮廓；讨论是否偏好平直、为何轮廓沿一个方向拉长，以及CMB等数据为何可帮助消除简并。
\end{exercise}
\begin{solution}
(A) 令$\Omega_K=1-\Omega_m-\Omega_\Lambda$，忽略辐射，有
\[
 I(z)=\int_0^z\frac{\dd z'}{\sqrt{\Omega_m(1+z')^3+\Omega_K(1+z')^2+\Omega_\Lambda}},
\]
\[
 D_L=\frac c{H_0}(1+z)\begin{cases}
 \sinh(\sqrt{\Omega_K}I)/\sqrt{\Omega_K},&\Omega_K>0,\\
 I,&\Omega_K=0,\\
 \sin(\sqrt{-\Omega_K}I)/\sqrt{-\Omega_K},&\Omega_K<0.
 \end{cases}
\]
用48点Gauss--Legendre积分逐点计算。96点独立加密复核的五组最大相对差小于$1.1\times10^{-14}$，物质平直和纯$\Lambda$还分别与$2(c/H_0)(1+z)[1-(1+z)^{-1/2}]$、$(c/H_0)z(1+z)$解析式核对。非平直不能只改变$E(z)$而漏掉横向正弦／双曲正弦。

\begin{center}\begin{tikzpicture}\begin{axis}[width=.80\textwidth,height=5.4cm,xlabel={$z$},ylabel={$D_L$（Gpc）},xmin=0,xmax=5,legend pos=north west]
\addplot[black] table[x=z,y expr=\thisrow{matter_flat}/1000,col sep=comma]{qa/ps5-distance-models.csv};\addlegendentry{$(1,0)$}
\addplot[structurecolor,dashed] table[x=z,y expr=\thisrow{Lambda_flat}/1000,col sep=comma]{qa/ps5-distance-models.csv};\addlegendentry{$(0,1)$}
\addplot[structurecolor,thick] table[x=z,y expr=\thisrow{model_0p3_0p7}/1000,col sep=comma]{qa/ps5-distance-models.csv};\addlegendentry{$(0.3,0.7)$}
\addplot[black,dotted] table[x=z,y expr=\thisrow{matter_closed}/1000,col sep=comma]{qa/ps5-distance-models.csv};\addlegendentry{$(1.2,0)$}
\addplot[structurecolor,dashdotted] table[x=z,y expr=\thisrow{matter_open}/1000,col sep=comma]{qa/ps5-distance-models.csv};\addlegendentry{$(0.3,0)$}
\end{axis}\end{tikzpicture}\end{center}
\begin{center}\begin{tikzpicture}\begin{axis}[width=.80\textwidth,height=5.1cm,xlabel={$z$},ylabel={$D_L$（Gpc）},xmin=0,xmax=1,ymin=0,ymax=9]
\addplot[black] table[x=z,y expr=\thisrow{matter_flat}/1000,col sep=comma]{qa/ps5-distance-models.csv};
\addplot[structurecolor,dashed] table[x=z,y expr=\thisrow{Lambda_flat}/1000,col sep=comma]{qa/ps5-distance-models.csv};
\addplot[structurecolor,thick] table[x=z,y expr=\thisrow{model_0p3_0p7}/1000,col sep=comma]{qa/ps5-distance-models.csv};
\addplot[black,dotted] table[x=z,y expr=\thisrow{matter_closed}/1000,col sep=comma]{qa/ps5-distance-models.csv};
\addplot[structurecolor,dashdotted] table[x=z,y expr=\thisrow{matter_open}/1000,col sep=comma]{qa/ps5-distance-models.csv};
\end{axis}\end{tikzpicture}\end{center}
第二图用同样的五条曲线和样式，显示低红移一阶斜率相近而较高红移才明显分离。

(B) 对固定$\Omega_m$与一次固定权重，$\alpha,\beta,M$都是线性参数。取设计矩阵的每行为$(s_i,-c_i,-1)$，目标为$\mu_{ \rm cosmo}-m_i^*$，加权线性最小二乘先精确消去这三个参数，再对$\Omega_m$作一维最小化。更新方差后重复，直到全部参数的最大变化小于$2\times10^{-8}$。这样执行的是原题的冻结权重迭代，没有让分母通过无约束放大$\alpha,\beta$降低残差。

表11的\texttt{e\_mb,e\_s,e\_c}是\emph{标准差}，进入方差时须平方；协方差直接按ReadMe列位读取。固定$H_0$还决定了$M$零点，超新星不能独立消除绝对光度与$H_0$的简并。

原题用$\sigma_{ \rm int}^2$称“弥散”有歧义。主结果按所写平方量\textbf{直接取方差0.10 mag$^2$}，另实际运行把“0.10”理解成标准差、方差0.01的对照；不悄悄切换约定。270颗有效SiFTO对象给
\begin{center}\begin{tabular}{lrrrrr}\toprule
固有方差&$\alpha$&$\beta$&$M$&$\Omega_m$&$\chi^2$/266\\\midrule
0.10&1.2390&2.9259&$-17.8463$&0.32332&61.66/266\\
0.01&1.2945&3.0289&$-17.8026$&0.31564&296.27/266\\\bottomrule
\end{tabular}\end{center}
分别6次与7次迭代收敛。主结果过小的约化$\chi^2$与把0.10当方差而得到约0.316 mag固有标准差一致；不能因为值小就称模型异常精确。题式用$s$而非通常的$s-1$，故$M$是平移后的截距；通常约定的绝对星等为$M_B=M-\alpha$，主结果约$-19.085$，并非$-17.846$的绝对峰值光度。

\begin{center}\begin{tikzpicture}\begin{axis}[width=.82\textwidth,height=6.0cm,xlabel={$z$},ylabel={$\mu$（mag）},xmin=.1,xmax=1.12,ymin=38,ymax=44.8,legend pos=south east]
\addplot[only marks,mark=*,mark size=.65pt,black,error bars/.cd,y dir=both,y explicit] table[x=z,y=mu,y error=err_mu,col sep=comma]{qa/snls-hubble-measurements.csv};\addlegendentry{实际SiFTO，误差$\sigma_i$}
\addplot[structurecolor,thick] table[x=z,y=best_mu,col sep=comma]{qa/snls-hubble-models.csv};\addlegendentry{最优平直}
\addplot[structurecolor,dashed] table[x=z,y=Omega_m_minus_0p2_mu,col sep=comma]{qa/snls-hubble-models.csv};\addlegendentry{$\Omega_m-0.2$}
\addplot[black,dotted] table[x=z,y=Omega_m_plus_0p2_mu,col sep=comma]{qa/snls-hubble-models.csv};\addlegendentry{$\Omega_m+0.2$}
\end{axis}\end{tikzpicture}\end{center}
误差棒按原要求画最终传播的$\sigma_i$，拟合分母另含固有方差。高红移较敏感，单颗散布仍大；增加样本可减统计误差，却不能取消颜色、标定和选择系统误差。实际主结果是这份课堂简化样本的拟合，不是重现原论文完整清选与系统协方差分析。

(C) 尘埃令$F$乘$\ee^{-\tau_d}$，附加星等$A=2.5\tau_d/\ln10\simeq1.086\tau_d$。若误以为超新星本征光度已知且无消光，便推断更大的$D_L$，可模仿某些暗能量模型的变暗趋势。普通尘埃还改变颜色，灰尘模型的波长依赖较弱而更难直接辨认；两者都需额外颜色、光谱和热辐射约束。不能仅凭这道课堂题就声称已排除所有尘埃模型。

(D) 固定主拟合的$\alpha,\beta,M$以及相应误差，在$\Omega_m\in[0,2.2]$、$\Omega_\Lambda\in[-0.6,2.6]$的$221\times257$网格计算完整曲率距离。所有路径要求$E^2>0$且$D_L>0$。网格最小值约在$(0.29,0.625)$，$\chi^2=61.640$；主平直点的$\Delta\chi^2\simeq0.019$，完全相容。下图由\emph{实际网格}的等值线顶点自绘；轮廓碰到$\Omega_m=0$物理边界，已截断，不能伪造闭合椭圆或宣布两侧均有测定误差。

\begin{center}\begin{tikzpicture}\begin{axis}[width=.76\textwidth,height=6.0cm,xlabel={$\Omega_m$},ylabel={$\Omega_\Lambda$},xmin=0,xmax=2.2,ymin=-.6,ymax=2.6,legend pos=north west,unbounded coords=jump]
\addplot[structurecolor,thick] table[x=Omega_m,y=Omega_lambda,col sep=comma]{qa/snls-contour-2p3.csv};\addlegendentry{$\Delta\chi^2=2.3$}
\addplot[structurecolor,dashed] table[x=Omega_m,y=Omega_lambda,col sep=comma]{qa/snls-contour-6p2.csv};\addlegendentry{6.2}
\addplot[black,dotted] table[x=Omega_m,y=Omega_lambda,col sep=comma]{qa/snls-contour-11p8.csv};\addlegendentry{11.8}
\addplot[black,domain=0:2.2]{1-x};\addlegendentry{平直线}
\addplot[only marks,mark=*,structurecolor] coordinates {(.29,.625)};
\end{axis}\end{tikzpicture}\end{center}
SN距离是$H(z)$和曲率映射的积分，不分别只依赖一种密度。低红移二阶项首先测$q_0=\Omega_m/2-\Omega_\Lambda$，故近似固定距离的方向有$\dd\Omega_\Lambda\sim\dd\Omega_m/2$，在参数平面呈正相关拉长；完整红移积分改变精确方向。CMB声视界与远距离几何约束不同组合，其他标准尺与增长数据亦可交叉限定。2.3、6.2、11.8仅是原题的两参数近似$\chi^2$阈值；固定烦扰参数、碰物理边界和忽略系统误差时，不将此图称为一般Bayes可信区域。
\end{solution}

\addtocontents{toc}{\protect\Needspace{3\baselineskip}}
\originalchapter{早期宇宙测验习题}
以下题面来自 8.286 Fall 2013 的正式测验, 不是复习题库. Quiz 1 第 3 题与 Quiz 3 第 2 题的官方答案仅回指作业答案, 未在该解答文件展开; 两题明确由 AI 独立补写. 其余六题依据官方解答核对并重新展开.

\src{8.286 Quiz 1, 2013年10月3日, Q2, (a)--(e), 原题 30 分; 官方解答标题误标 25 分.}
\begin{exercise}
平直宇宙 $a(t)=a_0\ee^{\chi t}$, 取 $\chi>0$. 两共动星系在 $t_1$ 的实际间距为 $\ell_p$, 一个在 $t_1$ 向另一个发光, 接收时刻 $t_2$.
(a) 求红移; (b) 仅以 $c,a_0,\chi,\ell_p,t_1$ 求 $t_2$; (c) 什么间距使光永不收到?
(d) 光在什么时候距两星系相等?
(e) 发射脉冲时长 $\Delta t$ 很短, 求收到的脉冲时长.
\end{exercise}
\begin{solution}
(a) $1+z=a(t_2)/a(t_1)=\ee^{\chi(t_2-t_1)}$.
(b) 共动间距 $\ell_c=\ell_p/a(t_1)$, 光传播给 $\ell_c=\int_{t_1}^{t_2}c\dd t/a(t)$. 因而 $\ell_p=(c/\chi)[1-\ee^{-\chi(t_2-t_1)}]$, 解得 $t_2=t_1-\chi^{-1}\ln(1-\chi\ell_p/c)$.
(c) 当 $\ell_p\ge c/\chi$ 时没有有限 $t_2$, 等号情形仅在无穷未来逼近.
(d) 同时的实际距离均含同一 $a(t)$, 中点因此也是共动中点. 积分路径改为 $\ell_c/2$, 得 $t_{\rm eq}=t_1-\chi^{-1}\ln[1-\chi\ell_p/(2c)]$. 中点能否到达的条件是 $\ell_p<2c/\chi$, 所以某些光永不到终点但可先过中点.
(e) 固定共动路径对端点微分给 $c\Delta t/a(t_1)=c\Delta t_o/a(t_2)$, 故 $\Delta t_o=(1+z)\Delta t$. 时长变长不是本地仪器走慢, 而是两次信号连接的宇宙时刻变化.
\end{solution}

\src{8.286 Quiz 1 Q3, (a)--(e), 25 分; 独立补写.}
\begin{exercise}
无限长均匀圆柱宇宙只沿 $x,y$ 膨胀, $z$ 方向静止. 初始密度 $\rho_i$, 初始速度 $\boldsymbol v_i=H_i\boldsymbol r$, 初始外半径有限. $\mu$ 为半径 $r$ 内每长度质量.
(a) Gauss 引力定律给 $\boldsymbol g=-A\mu\hat{\boldsymbol r}/r$, 求 $A$.
(b) 以初始半径 $r_i$ 写粒子 $\ddot r$.
(c) 证明 $u=r/r_i$ 与 $r_i$ 无关, 定义 $a(t)=u$.
(d) 求密度与尺度因子关系及 $\ddot a$.
(e) 找守恒量 $\mathcal E=\dot a^2/2+V(a)$, 判断永远膨胀还是再坍缩.
\end{exercise}
\begin{solution}
对单位长圆柱, $g_r2\pi r=-4\pi G\mu$, 故 $A=2G$. 未壳交叉时每层 $\mu=\pi\rho_ir_i^2$, $\ddot r=-2\pi G\rho_ir_i^2/r$. 代 $r=r_iu$ 得 $\ddot u=-2\pi G\rho_i/u$, 且对所有非中心层 $u(t_i)=1,\dot u(t_i)=H_i$. 正半轴的方程局部唯一, 相同初值因此给相同函数; 中心层由连续极限处理.

横截面积按 $a^2$ 变, $\rho=\rho_i/a^2$, 因而 $\ddot a=-2\pi G\rho a$. 乘 $\dot a$ 积分得 $V(a)=2\pi G\rho_i\ln a$、$\mathcal E=H_i^2/2$. 若初始膨胀, 转向点 $a_{\max}=\exp[H_i^2/(4\pi G\rho_i)]$ 有限; 对正密度和有限初始速度, 不存在一直向无穷半径逃逸的解. 对数势在大 $a$ 无界, 所以最终回缩; 这是无限长圆柱的性质, 不能套用到三维球宇宙.
\end{solution}

\src{8.286 Quiz 1 Q4, (a)--(b), 25 分.}
\begin{exercise}
今日宇宙时刻为 $t_0$, 星系当前实际距离 $\ell_p=a(t_0)\ell_c$. 视它为球, 发光时半径 $R_1$, 今日半径 $R_0$, 红移已知为 $z$, $\theta\ll1$.
(a) 不预设 $a(t)$ 的具体形式, 求张角.
(b) 发射总功率 $P$ 各向同性, 求今日地球收到的能量通量 $J$. 原题未写平直假设, 官方解答说明按当时课程的平直几何处理.
\end{exercise}
\begin{solution}
采用官方明确的平直假设. 星系发射时的共动半径是 $R_1/a_e$, 因此直径张角 $\theta=2R_1/(a_e\ell_c)=2R_1(1+z)/\ell_p$. 不应把 $R_0$ 代替 $R_1$, 引力束缚星系并不随 Hubble 流按统一比例膨胀.

今日波面面积 $4\pi\ell_p^2$, 每光子能量与到达率各降 $1+z$, 故 $J=P/[4\pi\ell_p^2(1+z)^2]$. 若空间弯曲, 应用横向距离 $D_M$ 替代简单 $\ell_p$, 题面原量不足以单独给此形式.
\end{solution}

\src{8.286 Quiz 2, 2013年11月7日, Q2, (a)--(d), 20 分.}
\begin{exercise}
平直宇宙由奇异物质主导, $\rho\propto a^{-5}$.
(a) 求尺度因子随时间的关系; (b) 求 $H(t)$; (c) 求实际粒子视界; (d) 求 $\rho(t)$.
\end{exercise}
\begin{solution}
从大爆炸起点 $t=0$ 计时, $\dot a/a\propto a^{-5/2}$, 故 $\dot a\propto a^{-3/2}$. 积分给 $a=Ct^{2/5}$, $H=2/(5t)$. 粒子视界 $D_p=a(t)c\int_0^t\dd t'/a(t')=5ct/3$. 由 Friedmann 式, $\rho=3H^2/(8\pi G)=3/(50\pi Gt^2)$. 能量守恒所对应状态参数为 $w=2/3$, 不等于通常辐射. 每项都适用此理想单组分模型, 不作为真实宇宙的历史.
\end{solution}

\src{8.286 Quiz 2 Q4, (a)--(c), 20 分.}
\begin{exercise}
假想物质能量密度随活塞体积 $V$ 按 $V^{-3/2}$ 下降, 初始为 $u_0$.
(a) 活塞拉大 $\Delta V$ 后, 求系统能量变化 $\Delta U$, 增加规定为正.
(b) 压强为 $p$, 求拉活塞的外界对系统做功 $\Delta W$.
(c) 用两者确定 $p$ 与 $u$ 的关系. 忽略二阶小量.
\end{exercise}
\begin{solution}
密度不是总能量. $U(V')=u_0(V/V')^{3/2}V'$, 在 $V'=V$ 求导给 $\Delta U=-(1/2)u_0\Delta V$. 外界对系统做功为 $\Delta W=-p\Delta V$. 没有额外热交换时两者相等, 故 $p=u/2$. 若把 $u\propto V^{-3/2}$ 直接当作 $U$ 的变化, 会错误得到 $3u/2$; 少掉的因子正是体积本身.
\end{solution}

\src{8.286 Quiz 3, 2013年12月5日, Q2, (a)--(c), 25 分; 独立补写.}
\begin{exercise}
缪子静止能 $106\,\mathrm{MeV}$, 电子 $0.511\,\mathrm{MeV}$, 两者均自旋 $1/2$, 缪子与反缪子电荷相反. 本题按理想阈值热辐射模型, 不考虑额外强子热贡献.
(a) 缪子与反缪子对 $u=(\pi^2/30)g(\kb T)^4/(\hbar c)^3$ 的 $g$ 贡献是多少?
(b) 当 $\kb T$ 刚高于 $106\,\mathrm{MeV}$ 时还有哪些轻粒子, 各贡献多少?
(c) 缪子消失时所有其他粒子仍保持共同热平衡, $aT$ 增加多少倍?
\end{exercise}
\begin{solution}
两电荷态各两自旋态, 费米权重 $7/8$, 所以 $g_\mu=4(7/8)=7/2$. 其余为光子两偏振 $g_\gamma=2$, 电子正电子 $g_e=7/2$, 三味中微子及反中微子共六个手征自由度 $g_\nu=6(7/8)=21/4$. 合计不含缪子 $g_{
\rm after}=43/4$, 含缪子 $g_{\rm before}=57/4$. 在全部共同温度且熵绝热守恒的模型中,
$[aT]_{\rm after}/[aT]_{\rm before}=(57/43)^{1/3}=1.0985$.
实际 $\kb T\sim106\,\mathrm{MeV}$ 并非严格的阶跃阈值, 接近质量阈值的粒子贡献连续变化, QCD 热史也更复杂. 原题的无中间质量叙述支持这一课堂自由粒子近似, 不把结果宣称为精确现代热史.
\end{solution}

\src{8.286 Quiz 3 Q3, (a)--(c), 25 分.}
\begin{exercise}
平直宇宙今日参数 $\Omega_{m0},\Omega_{r0},\Omega_{\Lambda0},H_0$, 三密度参数和为 1. 今日从原点发出球面光脉冲.
(a) 已知 $a(t)$ 时写永远能逼近的最大共动半径 $r_{\max}$.
(b) 换积分变量, 仅以今日参数和 $a_0$ 写此积分, 不要求求值.
(c) 对今天收到的光, 若源的共动位置恰在 $r_{\max}$, 写求源红移 $z_{\rm eh}$ 的显式积分方程.
\end{exercise}
\begin{solution}
零测地线给 $\dot r=c/a$. 最大半径为 $r_{\max}=c\int_{t_0}^\infty\dd t/a(t)$; 此积分需由未来模型保证收敛. 令 $x=a/a_0$, $H=H_0\sqrt{\Omega_{r0}x^{-4}+\Omega_{m0}x^{-3}+\Omega_{\Lambda0}}$, 则
\[
r_{\max}=\frac c{a_0H_0}\int_1^\infty
\frac{\dd x}{\sqrt{\Omega_{r0}+\Omega_{m0}x+\Omega_{\Lambda0}x^4}}.
\]
今天收到的光的过去传播距离也必须是这个共动数, 所以
\[
\int_0^{z_{\rm eh}}\frac{\dd z}{\sqrt{\Omega_{r0}(1+z)^4+\Omega_{m0}(1+z)^3+\Omega_{\Lambda0}}}
=\int_1^\infty\frac{\dd x}{\sqrt{\Omega_{r0}+\Omega_{m0}x+\Omega_{\Lambda0}x^4}}.
\]
「该源今天的位置在事件视界」不表示它的早期光永远不可见. 示例忽略辐射、取 $0.3,0.7,H_0=71$, 今日实际事件视界约 $4816\,\mathrm{Mpc}$, 对应今天收到的源红移约 $1.805$. 这只是给定未来真空不变模型的演示.
\end{solution}

\src{8.286 Quiz 3 Q4, (a)--(b), 15 分.}
\begin{exercise}
(a) 从 $H^2=8\pi G\rho/3-kc^2/a^2$ 推导 $(\Omega-1)/\Omega=A T^2/\rho$, 找常数 $A$.
(b) 对近乎平直、$\rho\propto a^{-5}$ 的假想宇宙, 求该偏离量随时间的变化. 温度作为光子温度, 假设 $aT=C$ 固定.
\end{exercise}
\begin{solution}
$\Omega=8\pi G\rho/(3H^2)$, 所以 $(\Omega-1)/\Omega=1-3H^2/(8\pi G\rho)=3kc^2/(8\pi G\rho a^2)$. 用 $aT=C$ 得 $A=3kc^2/(8\pi GC^2)$. 常数依尺度因子的规范与温度归一化而变, 整个偏离量不变. 近乎平直时上一章题解给 $a\propto t^{2/5}$, 所以 $T^2/\rho\propto a^3\propto t^{6/5}$. 原式在粒子熵自由度改变时 $aT$ 不恒定, 不能继续把 $A$ 视作同一个常数.
\end{solution}

\addtocontents{toc}{\protect\Needspace{3\baselineskip}}
\originalchapter{课程来源与使用范围}
本卷按主题重编, 原课程的单讲顺序不作为章节编号. 恒星与气体部分参照8.901的26次公开课历与47组公开题面；星系至Lyman--alpha主体参照8.902全部21份公开分讲，并保留独立数值项目的数据依赖；均匀膨胀、热史和暴胀参照8.286公开幻灯、十份作业、三次测验及三份复习题。分讲与合订本是重复呈现，不能叠加数目。章节内完整推导是本卷重构, 未取得的教材题不由想象填入.
\begin{thebibliography}{9}
\bibitem{a901} Deepto Chakrabarty. Astrophysics I, Spring 2006. MIT OpenCourseWare. \url{https://ocw.mit.edu/courses/8-901-astrophysics-i-spring-2006/}.
\bibitem{a902} Mark Vogelsberger, Stephanie O'Neil, David DePalma. Astrophysics II, Fall 2023. 讲义、作业与课程资料. \url{https://ocw.mit.edu/courses/8-902-astrophysics-ii-fall-2023/}.
\bibitem{a902full} 8.902 Fall 2023 讲义合订本, 122 PDF 页, 与分讲有重复. \url{https://ocw.mit.edu/courses/8-902-astrophysics-ii-fall-2023/mit8_902_f23_lec_full.pdf}.
\bibitem{a286} Alan Guth. The Early Universe, Fall 2013. MIT OpenCourseWare. \url{https://ocw.mit.edu/courses/8-286-the-early-universe-fall-2013/}.
\bibitem{a286q1review} Alan Guth. 8.286 Fall 2013 Quiz 1 Review Problems and official solutions, 68 PDF 页. \url{https://ocw.mit.edu/courses/8-286-the-early-universe-fall-2013/16163bd7e7d317ede49c779b99ca9f23_MIT8_286F13_q1review.pdf}.
\bibitem{a286q2review} Alan Guth. 8.286 Fall 2013 Quiz 2 Review Problems and official solutions, 70 PDF 页. \url{https://ocw.mit.edu/courses/8-286-the-early-universe-fall-2013/fda6167ffbceff35ad346026f920ae77_MIT8_286F13_q2review.pdf}.
\bibitem{a902project} MIT OpenCourseWare. 8.902 Fall 2023, Hands-on Numerical Project, 5页. \url{https://ocw.mit.edu/courses/8-902-astrophysics-ii-fall-2023/mit8_902_f23_numerical_project.pdf}. Canvas配置和快照未在此公开文件中取得.
\bibitem{sc1942} M. Schönberg and S. Chandrasekhar. On the Evolution of the Main-Sequence Stars, 1942, \emph{Astrophysical Journal} 96, 161--172. \url{https://articles.adsabs.harvard.edu/pdf/1942ApJ....96..161S}. 本卷独立解释匹配质量限制，不转载原图或原文.
\bibitem{license} MIT OpenCourseWare. 使用条款与逐件权利例外. \url{https://ocw.mit.edu/pages/privacy-and-terms-of-use/}.
\end{thebibliography}

\end{document}
